% Generated by scripts/build_readers.py; change the reviewed sources and rebuild.
\def\EditionMode{1}
\def\EditionLabel{Annotated master}
\def\EditionStatus{WORKING EDITION: 94 of 706 source pages reviewed and included}
\documentclass[10pt,oneside]{book}
\usepackage{fontspec}
\setmainfont{Latin Modern Roman}
\setsansfont{Latin Modern Sans}
\usepackage{amsmath,amssymb,mathtools}
\usepackage{unicode-math}
\setmathfont{Latin Modern Math}
\usepackage{xcolor,tikz-cd,lua-ul,marginnote,pdflscape}
\usetikzlibrary{arrows.meta,decorations.pathmorphing,calc,shapes.geometric}
\tikzcdset{ampersand replacement=\&}
\tikzcdset{every label/.append style={text=black}}
\usepackage[a4paper,left=18mm,right=45mm,top=18mm,bottom=22mm,marginparwidth=28mm,marginparsep=7mm]{geometry}
\usepackage[unicode,colorlinks=true,urlcolor=black,linkcolor=black]{hyperref}
\usepackage{bookmark}
\setlength{\parindent}{0pt}
\setlength{\parskip}{4pt plus 1pt minus 1pt}
\setlength{\emergencystretch}{1.5em}
\raggedbottom
\definecolor{PSHand}{RGB}{25,58,132}
\definecolor{PSDelete}{RGB}{146,37,37}
\definecolor{PSTypedOld}{RGB}{95,95,95}
\ifnum\EditionMode=2\colorlet{blue}{black}\else\colorlet{blue}{PSHand}\fi
\RenewDocumentCommand{\textcircled}{m}{\tikz[baseline=(PSCircle.base)]\node[draw,circle,inner sep=.5pt,line width=.4pt](PSCircle){#1};}
\newif\ifPSInApparatus
\newcounter{psnote}
\newcounter{psvariant}
\newcounter{pslinkmark}
\newif\ifPSInLink
\newsavebox{\PSNoteBox}
\newdimen\PSNoteLimit
{\footnotesize\global\PSNoteLimit=3\baselineskip}
\newcommand{\CurrentScanURL}{}
\newcommand{\CurrentKey}{}
\newcommand{\CurrentScan}{}
\newcommand{\pssource}[3]{\psflushmargins\clearpage\gdef\CurrentKey{#1}\gdef\CurrentScan{#2}\gdef\CurrentScanURL{#3}\setcounter{psvariant}{0}\typeout{PS-SOURCE: #1}}
\newcommand{\pslinkmarker}[2]{\stepcounter{pslinkmark}\latelua{texio.write_nl('log','PS-LINK-SHIP|#1|\CurrentKey|\arabic{pslinkmark}|#2')}}
\newcommand{\pslink}[2]{\ifPSInLink\pslinkmarker{#1}{nested}#2\else\begingroup\PSInLinktrue\href{\CurrentScanURL}{\pslinkmarker{#1}{outer}#2}\endgroup\fi}
\newcommand{\scanpage}[1]{\ifnum\EditionMode=2\marginnote{\scriptsize\sffamily\pslink{scanpage}{scan #1}}\else\par\noindent{\small\sffamily\pslink{scanpage}{[scan #1]}}\par\medskip\fi}
\newcommand{\mspage}[1]{\ifnum\EditionMode=2\marginnote{\scriptsize\sffamily p.\,#1}\else{\small\sffamily[p.\,#1]}\fi}
\newcommand{\hand}[1]{{\ifnum\EditionMode=2\color{black}\else\color{PSHand}\fi #1}}
\let\over\relax
\NewDocumentCommand{\over}{O{} m m}{{\textcolor{PSTypedOld}{#2}\ensuremath{^{\!\nearrow}}\if\relax\detokenize{#1}\relax\hand{#3}\else #3\fi}}
\NewDocumentCommand{\inserted}{O{} m}{{\if\relax\detokenize{#1}\relax\hand{\ensuremath{\langle}#2\ensuremath{\rangle}}\else\ensuremath{\langle}#2\ensuremath{\rangle}\fi}}
\newcommand{\deleted}[1]{{\color{PSDelete}\ifmmode\text{[del.: }#1\text{]}\else[\textsf{del.:} #1]\fi}}
\newcommand{\xout}[1]{{\color{PSTypedOld}\ifmmode\text{[x: }#1\text{]}\else[\textsf{x:} #1]\fi}}
\let\PSMathUnderline\underline
\RenewDocumentCommand{\underline}{+m}{\ifmmode\PSMathUnderline{#1}\else\underLine{#1}\fi}
\NewDocumentCommand{\uline}{+m}{\underline{#1}}
\NewDocumentCommand{\uuline}{+m}{\ifmmode\PSMathUnderline{\PSMathUnderline{#1}}\else\underLine{\underLine{#1}}\fi}
\newunderlinetype\PSWaveLine{\leaders\hbox{\lower.7ex\hbox{\tikz[x=1pt,y=1pt,baseline=0pt]{\draw[line width=.3pt] (0,0) .. controls (1,1) and (2,1) .. (3,0) .. controls (4,-1) and (5,-1) .. (6,0);}}}}
\NewDocumentCommand{\uwave}{+m}{{\PSWaveLine #1}}
\newcommand{\unsure}[1]{\pslink{unsure}{[#1\ifmmode\,?\else\,?\fi]}}
\newcommand{\illegible}[1]{\ifmmode\text{\pslink{illegible}{[illeg.\ #1]}}\else\pslink{illegible}{[illeg.\ #1]}\fi}
\NewDocumentCommand{\diagram}{+m}{\ifmmode\vcenter{\hbox{#1}}\else\par\medskip{\centering #1\par}\medskip\psflushmargins\fi}
\newenvironment{diagramblock}{\par\begin{center}}{\end{center}\par}
\newcommand{\PSMathNotes}{}
\makeatletter
\NewDocumentCommand{\psplacednote}{m +m}{%
      \sbox{\PSNoteBox}{\parbox{\marginparwidth}{\footnotesize\raggedright\hand{\textsf{m#1.} #2}}}%
      \ifdim\dimexpr\ht\PSNoteBox+\dp\PSNoteBox\relax<\PSNoteLimit
        \leavevmode\marginpar{\usebox{\PSNoteBox}}%
      \else
        \par\nobreak\smallskip\hrule\nobreak\smallskip
        {\footnotesize\textsf{Marginal note m#1.}\quad\hand{#2}\par}
        \nobreak\smallskip\hrule\smallskip
      \fi}
\newcommand{\psflushmargins}{\ifmmode\else\ifinner\else\ifx\PSMathNotes\@empty\else\let\PSTempNotes\PSMathNotes\global\let\PSMathNotes\@empty\PSTempNotes\fi\fi\fi}
\NewDocumentCommand{\margin}{+m}{%
  \ifPSInApparatus\ifmmode\hand{\text{[marg.: }#1\text{]}}\else\hand{\textsf{[marg.: }#1\textsf{]}}\fi
  \else\stepcounter{psnote}%
    \ifmmode
      \text{\textsuperscript{m\arabic{psnote}}}%
      \begingroup\edef\PSQueue{\endgroup\noexpand\g@addto@macro\noexpand\PSMathNotes{\noexpand\psplacednote{\arabic{psnote}}{\unexpanded{$#1$}}}}\PSQueue
    \else
      \textsuperscript{m\arabic{psnote}}\psplacednote{\arabic{psnote}}{#1}%
    \fi
  \fi}
\AtBeginDocument{\everymath\expandafter{\the\everymath\aftergroup\psflushmargins}\everydisplay\expandafter{\the\everydisplay\aftergroup\psflushmargins}}
\makeatother
\newcommand{\psnumber}[2]{\expandafter\gdef\csname psnum@#1\endcsname{#2}}
\NewDocumentCommand{\var}{+m +m m m}{\ifnum\EditionMode=0 #1\else\ifnum\EditionMode=1\ifmmode\text{\textsuperscript{\tiny\csname psnum@#4\endcsname}}\else\textsuperscript{\tiny\csname psnum@#4\endcsname}\fi\fi #2\fi}
\newcommand{\psapparatusstart}{\par\bigskip\hrule\smallskip\PSInApparatustrue\begingroup\footnotesize\textsf{Apparatus for scan \CurrentScan}\par}
\newcommand{\psapparatusend}{\par\endgroup\PSInApparatusfalse}
\NewDocumentCommand{\psentry}{m m m +m +m m}{\par\noindent\textbf{#2. #3.}\quad\textit{Typescript:} #4\quad\textit{Corrected:} \if\relax\detokenize{#5}\relax\textnormal{[omitted]}\else #5\fi\quad #6\par}
\newcommand{\psfrontmatter}{%
\hypersetup{pageanchor=false}\begin{titlepage}\raggedright
{\Large Alexander Grothendieck\par}\vspace{12mm}
{\Huge Pursuing Stacks\par}\medskip{\LARGE 1983\par}\vspace{12mm}
{\Large\EditionLabel\par}\vspace{8mm}
{\large\EditionStatus\par}\vfill
Université de Montpellier, Grothendieck archive, cotes 134-2--134-8.\par
Including the 26 source pages of R. Brown's translation at cote 134-2, sites 115--143.\par
\end{titlepage}\hypersetup{pageanchor=true}
\chapter*{Editorial statement}
AI-transcribed by Claude Opus 5.5 (Anthropic, effort high); checked against the scans, and this edition built, by GPT-6 Astra (OpenAI, ChatGPT Pro mode); prompts and records by Claude (Anthropic).

The witness consulted and not reproduced is the LaTeX edition of Mateo Carmona with the collaboration of Ulrik Buchholtz (arXiv 2111.01000 v2, 2021), extended from the 2015 edition by ``the scrivener'', CC0 1.0.

The translation in the appendix is R. Brown's, prepared with assistance from L. Breen and J-L. Loday, as his note says; the rights in it belong to its translator or his heirs.

The rights in the original work belong to its author or the author's representatives or heirs. Where we hold rights, in the AI-produced transcription, translation, corrections and typesetting, they are dedicated to the public domain under CC0 1.0. This edition is published for educational, non-profit purposes, and we claim no rights we do not hold. If you hold rights in this work and object to its publication here, please open a GitHub issue and we will take it down.

\section*{How to read this form}
\ifnum\EditionMode=0
The diplomatic form preserves typed and handwritten wording, cancellations, insertions, overwritten readings, underlining, doubts, and original page numbers. Blue marks handwriting; red labels handwritten cancellations; gray labels typewriter cancellations and underlying overwritten text. Angle brackets mark insertions; a superscript arrow separates an overwritten reading from its replacement. Typed insertions and typed overwriting remain plain. The disposition of long marginal notes is identified explicitly.
\else
The corrected text applies the author's insertions, replacements, cancellations, and recorded transpositions. It retains his English and notation, including consistent forms such as ``standart'', ``developped'', ``it's'' for ``its'', and ``A la''. An isolated slip is corrected only when specifically recorded and justified. Mathematical anomalies are preserved and discussed in the guide.
\ifnum\EditionMode=1
Blue continues to identify handwriting. A numbered call identifies each maximal editorial operation, with its full diplomatic and corrected readings and reason in the apparatus following the source page. Internal layers of a wholly replaced or deleted passage are preserved within that entry rather than duplicated as separate entries. These entries are generated from the same variants that govern the running text.
\else
This reading form omits the apparatus, cancelled material, superseded readings, insertion brackets, and color distinctions. It retains the resulting text, marginal notes, underlining, and unresolved doubts. Scan and source-page references appear in the margin.
\fi\fi

Every scan mark and remaining doubt links to its source page in the archive viewer. Doubts are editorial reading marks, not emendations. The source-page breaks identify the archival sequence; a source page may occupy more than one printed page. A partial working edition is explicitly labelled and contains only reviewed source pages.
\clearpage}

\begin{document}
\psfrontmatter
\pssource{ps_134-7_001}{134-7/1}{https://grothendieck.umontpellier.fr/web/viewer.html?file=../134-7.pdf\#page=2}
\psnumber{ps_134-7_001-v0001}{1}
\psnumber{ps_134-7_001-v0002}{2}
\psnumber{ps_134-7_001-v0003}{3}
\psnumber{ps_134-7_001-v0004}{4}
\psnumber{ps_134-7_001-v0005}{5}
\psnumber{ps_134-7_001-v0006}{6}
\psnumber{ps_134-7_001-v0007}{7}
\psnumber{ps_134-7_001-v0008}{8}
\psnumber{ps_134-7_001-v0009}{9}
\psnumber{ps_134-7_001-v0010}{10}
%% BEGIN SOURCE ps_134-7_001
% Pursuing Stacks (cotes 134-2 to 134-8): archive cote 134-7, site page 1 (PDF page 2), typescript page 443.
% typescript. Draft compared: ps6.tex:1-58.
% AI-transcribed by Claude Opus 5.5 (claude-opus-5-5, adaptive thinking, effort high), Message Batches API, batch ps3 (msgbatch_01XdFmuVVBAcgc7EFxdtCqsj), 2026-11-12; reviewed against both supplied scan halves; see CHECK_DECISIONS.tsv.
\scanpage{134-7/1}
\margin{\uline{VI Schematiza-tion}}

\mspage{443}

24.8.

110.\quad I pondered some more about homotopy types over a grou\unsure{nd r}ing $k$,
% NOTE: a rust stain runs down the right side of the page and hides letters at many line ends
just enough to become familiar again with the idea\hand{,}and more or less con-%
vinced that there should exist such a thing, which \var{\xout{amounts}}{}{typewriter deletion}{ps_134-7_001-v0001} shoul\unsure{d} amount\hand{,}
kind of\hand{,} to putting a ``continuous'' structure \var{\over{\unsure{s}}{(}}{\hand{(}}{writing over}{ps_134-7_001-v0002}namely the very rich struc-%
ture of a s\unsure{\hand{c}}heme) upon something usually visualized as something ``discrete''
- namely a homotopy type. The basic analogy here is free $\underline{Z}$-modules $M$
of finite type - a typical case of a ``\uline{discrete}'' structure. It gives
rise though to a \uline{vector bundle} $W(M)$ over the absolute base $S_o=\mathrm{Spec}(\underline{Z})$,
whose $\underline{Z}$-module of sections is $M$ , and the functor
\[ M \hand{\mapsto} W(M) \]
from free $\underline{Z}$-modules of finite type to vector bundles over $S_o$ is fully
faithful. When $M$ is an arbitrary $\underline{Z}$-module i.e. an abelian group, $W(M)$
still makes a sense, namely as a functor
\[ k \hand{\mapsto} M\otimes_{\underline{Z}}k \]
on the category of all commutative $\underline{Z}$-algebras (i.e. just commutative
rings); it is no longer representable by a scheme over $S_o$ (except \var{\deleted{when}}{}{deletion}{ps_134-7_001-v0003}
precisely \var{\inserted{when}}{\hand{when}}{insertion}{ps_134-7_001-v0004} $M$ is free of finite type), but it is very close \var{\xout{\illegible{1}}}{}{typewriter deletion}{ps_134-7_001-v0005}\var{\inserted[typed]{still}}{still}{typed insertion}{ps_134-7_001-v0006}, intui-%
tively and technically too, to a usual vector bundle (the ``vect\unsure{or}'' struc-%
ture coming from the $k$-module structure on $W(M)(k) = M\otimes_{\underline{Z}}k$ ). Again,
the functor $M\hand{\mapsto} W(M)$ from Ab to the category of ``generalized \unsure{ve}ctor
bundles'' over $S_o$ is fully faithful. Working with semisimplic\unsure{ia}l
$\underline{Z}$-modules (say) rather than \unsure{\var{\over{j}{j}}{\hand{j}}{writing over}{ps_134-7_001-v0007}}ust $\underline{Z}$-modules, and more specifical\unsure{ly}
with those corresponding to $K(\hand{\pi},n)$ types, and using Postnikov ``\unsure{dé}vissa-%
ge'' of a general homotopy type, one may hope to ``represent'' th\unsure{is} type,
in a more or less canonical way in terms of the successive semisi\unsure{m}plicial
P\hand{\u{o}}stnikov fibrations, by a semi-simplicial object in the topo\unsure{s} (\unsure{s}ay) of
all functors from $\underline{Z}$-algebras to sets which are ``sheaves'' for a s\unsure{u}itable
site structure on the dual category (namely the category of affi\unsure{ne}
sch\unsure{em|mm}es over $\underline{Z}$ ) - the so-called ``flat'' topology seems \var{\xout{to}}{}{typewriter deletion}{ps_134-7_001-v0008} OK. (NB To
eliminate logical difficulties, we may have to restrict somewhat the
rings $k$ used as arguments, for instance take them to be of fini\unsure{te} type
over $\underline{Z}$ - never mind such technicalities !). This approach \var{\xout{is}}{}{typewriter deletion}{ps_134-7_001-v0009} ma\unsure{y}
possibly work, when restricting to \var{\xout{simply}}{}{typewriter deletion}{ps_134-7_001-v0010} 1-connected homotopy \unsure{t}ypes,
or at any rate to the case when the fundamental groups of the con\unsure{n}ected
\margin{*} components are abelian. If we wish a ``schemati\unsure{z}ation'' of arbitrary
homotopy types, we may think of going about it by keeping the f\unsure{u}ndamen-%
tal group or groupoid ``discrete'', and ``schematize'' the P\hand{\u{o}}stni\unsure{ko}v

\margin{1}

% Visual collation: both supplied image halves inspected; see CHECK_DECISIONS_134_7_8.tsv.

%% END SOURCE ps_134-7_001
\ifnum\EditionMode=1
\psapparatusstart
\psentry{ps_134-7_001-v0001}{1}{typewriter deletion}{\xout{amounts}}{}{The typewriter strikes out this material; it leaves the corrected running text.}
\psentry{ps_134-7_001-v0002}{2}{writing over}{\over{\unsure{s}}{(}}{\hand{(}}{The superposed handwritten reading replaces the underlying reading.}
\psentry{ps_134-7_001-v0003}{3}{deletion}{\deleted{when}}{}{The author cancels this material by hand; it leaves the corrected running text.}
\psentry{ps_134-7_001-v0004}{4}{insertion}{\inserted{when}}{\hand{when}}{The handwritten insertion is incorporated at its indicated position; insertion brackets leave the running text.}
\psentry{ps_134-7_001-v0005}{5}{typewriter deletion}{\xout{\illegible{1}}}{}{The typewriter strikes out this material; it leaves the corrected running text.}
\psentry{ps_134-7_001-v0006}{6}{typed insertion}{\inserted[typed]{still}}{still}{The typed insertion is incorporated at its indicated position; insertion brackets leave the running text.}
\psentry{ps_134-7_001-v0007}{7}{writing over}{\over{j}{j}}{\hand{j}}{The superposed handwritten reading replaces the underlying reading.}
\psentry{ps_134-7_001-v0008}{8}{typewriter deletion}{\xout{to}}{}{The typewriter strikes out this material; it leaves the corrected running text.}
\psentry{ps_134-7_001-v0009}{9}{typewriter deletion}{\xout{is}}{}{The typewriter strikes out this material; it leaves the corrected running text.}
\psentry{ps_134-7_001-v0010}{10}{typewriter deletion}{\xout{simply}}{}{The typewriter strikes out this material; it leaves the corrected running text.}
\psapparatusend
\fi
\pssource{ps_134-7_002}{134-7/2}{https://grothendieck.umontpellier.fr/web/viewer.html?file=../134-7.pdf\#page=3}
\psnumber{ps_134-7_002-v0001}{1}
\psnumber{ps_134-7_002-v0002}{2}
\psnumber{ps_134-7_002-v0003}{3}
\psnumber{ps_134-7_002-v0004}{4}
\psnumber{ps_134-7_002-v0005}{5}
\psnumber{ps_134-7_002-v0006}{6}
\psnumber{ps_134-7_002-v0007}{7}
\psnumber{ps_134-7_002-v0008}{8}
\psnumber{ps_134-7_002-v0009}{9}
\psnumber{ps_134-7_002-v0010}{10}
\psnumber{ps_134-7_002-v0011}{11}
\psnumber{ps_134-7_002-v0012}{12}
\psnumber{ps_134-7_002-v0013}{13}
\psnumber{ps_134-7_002-v0014}{14}
\psnumber{ps_134-7_002-v0015}{15}
\psnumber{ps_134-7_002-v0016}{16}
\psnumber{ps_134-7_002-v0017}{17}
\psnumber{ps_134-7_002-v0018}{18}
%% BEGIN SOURCE ps_134-7_002
% Pursuing Stacks (cotes 134-2 to 134-8): archive cote 134-7, site page 2 (PDF page 3), typescript page 444.
% typescript. Draft compared: ps6.tex:46-100.
% AI-transcribed by Claude Opus 5.5 (claude-opus-5-5, adaptive thinking, effort high), Message Batches API, batch ps3 (msgbatch_01XdFmuVVBAcgc7EFxdtCqsj), 2026-11-12; reviewed against both supplied scan halves; see CHECK_DECISIONS.tsv.
\scanpage{134-7/2}
\mspage{444}

truncation involving only the higher homotopy groups ($\hand{\pi}_i$ with $i\geqslant 2$).
This suggests that for any inte\var{\xout{\unsure{x}}}{}{typewriter deletion}{ps_134-7_002-v0001}ger $n\geqslant 1$, there may be a ``schematization
% NOTE: the right edge of the scan cuts this line after "schematization"; the closing quotation marks are not visible
above level n'' for a given homotopy type, leaving the Cartan-Serre
truncation of \var{\xout{dimension}}{}{typewriter deletion}{ps_134-7_002-v0002} level n discrete, and ``schematizing'' the
corresponding \var{\inserted{total}}{\hand{total}}{insertion}{ps_134-7_002-v0003} Postnikov fiber (involving homotopy groups $\hand{\pi}_i$ for $i\geqslant n$).
Maybe such a schematization can be constructed equally for level 0 ,
even without assuming the fundamental groupoid to be abelian, \var{\xout{by replac}}{}{typewriter deletion}{ps_134-7_002-v0004}
only nilpotent - but then we may have to change from ground ring $\underline{Z}$ to
the considerably coarser one $\underline{Q}$ (compare comments at the end of section
94). In any case, the key step in this approach would consist in \var{\xout{the}}{}{typewriter deletion}{ps_134-7_002-v0005}
checking that, after schematization has been carried through succesful-%
ly up to a certain level in the successive elementary Postnikov fibrati-%
ons, the next elementary fibration (described by a cohomology class in
$H^{n+2}(X_n,\hand{\pi}_{n+1})$ ) comes from a``schematic'' one, and that the latter is
essentially unique; in other words, \uline{that the canonical map from the}
\uline{``schematic''} $H^{n+2}$ (\uline{with``quasi-coherent''coefficients}) \uline{to the usual}
\uline{``discrete''} \uline{one is bijective}. Maybe this hope is wholly unrealistic\hand{\unsure{,}}
though. One fact which calls for some scepticism about this approach,
comes in when looking at the case of an ``abelian'' homotopy type, descri-%
% NOTE: the right edge of the scan cuts this line after "descri"
bed by a semisimplicial abelian group $X_{\hand{\ast}}$ , in which case we expect that
% NOTE: the right edge of the scan cuts the last letters of "that"
base change $\underline{Z}\hand{\rightarrow} k$ should be just the usual base change
\[ X_{\hand{\ast}} \hand{\rightarrow} X_{\hand{\ast}}\otimes_{\underline{Z}} k \]
(if X. has torsion-free \var{\xout{components}}{}{typewriter deletion}{ps_134-7_002-v0006} i.e. flat components,at any rate).
But when X. has homology torsion, the universal coefficients formula
shows us that the homology \var{\deleted{(=}}{}{deletion}{ps_134-7_002-v0007}\var{\inserted{(or rather,}}{\hand{(or rather,}}{insertion}{ps_134-7_002-v0008} homotopy) groups of $X_{\hand{\ast}}\otimes_{\underline{Z}}k$ are \uline{not} just
the groups $\unsure{H}_{\unsure{i}}(X_{\hand{\ast}})\otimes_{\underline{Z}}k = \hand{\pi}_i(X_{\hand{\ast}})\otimes_{\underline{Z}}k$ , as we implicitly were assuming it
seems in the approach sketched above, when schematizing the homotopy
groups $\hand{\pi}_i$ \var{\xout{by}}{}{typewriter deletion}{ps_134-7_002-v0009} one by one via $W(\hand{\pi}_i)$. Thus \var{\inserted[typed]{maybe}}{maybe}{typed insertion}{ps_134-7_002-v0010} Postnikov dévissage \var{\inserted[typed]{is'nt}}{is'nt}{typed insertion}{ps_134-7_002-v0011} \var{\xout{may not be}}{}{typewriter deletion}{ps_134-7_002-v0012}
a possible approach towards schematization of homotopy types, and one
will have to work out rather a comprehensive yoga of reconstructing
a homotopy type from \var{\xout{its abelianization}}{}{typewriter deletion}{ps_134-7_002-v0013} one kind or other of``abelianiza-%
tion'' or ``linearization'' \var{\inserted{of homotopy types,}}{\hand{of homotopy types,}}{insertion}{ps_134-7_002-v0014} endowed with suitable extra structure embody-%
ing ``multiplicative''\var{\xout{\unsure{f}}}{}{typewriter deletion}{ps_134-7_002-v0015} features of the homology and cohomology structure.
At any rate, I did not hit upon any ``simplistic way'' to define homotopy
types over\var{\deleted{\illegible{1}}}{}{deletion}{ps_134-7_002-v0016} any ground ring, and I have some doubts there is any, in
terms of the general non-sense \var{\xout{of}}{}{typewriter deletion}{ps_134-7_002-v0017} we did so far. \var{\xout{at any rate}}{}{typewriter deletion}{ps_134-7_002-v0018}

\margin{2}

% Visual collation: both supplied image halves inspected; see CHECK_DECISIONS_134_7_8.tsv.

%% END SOURCE ps_134-7_002
\ifnum\EditionMode=1
\psapparatusstart
\psentry{ps_134-7_002-v0001}{1}{typewriter deletion}{\xout{\unsure{x}}}{}{The typewriter strikes out this material; it leaves the corrected running text.}
\psentry{ps_134-7_002-v0002}{2}{typewriter deletion}{\xout{dimension}}{}{The typewriter strikes out this material; it leaves the corrected running text.}
\psentry{ps_134-7_002-v0003}{3}{insertion}{\inserted{total}}{\hand{total}}{The handwritten insertion is incorporated at its indicated position; insertion brackets leave the running text.}
\psentry{ps_134-7_002-v0004}{4}{typewriter deletion}{\xout{by replac}}{}{The typewriter strikes out this material; it leaves the corrected running text.}
\psentry{ps_134-7_002-v0005}{5}{typewriter deletion}{\xout{the}}{}{The typewriter strikes out this material; it leaves the corrected running text.}
\psentry{ps_134-7_002-v0006}{6}{typewriter deletion}{\xout{components}}{}{The typewriter strikes out this material; it leaves the corrected running text.}
\psentry{ps_134-7_002-v0007}{7}{deletion}{\deleted{(=}}{}{The author cancels this material by hand; it leaves the corrected running text.}
\psentry{ps_134-7_002-v0008}{8}{insertion}{\inserted{(or rather,}}{\hand{(or rather,}}{The handwritten insertion is incorporated at its indicated position; insertion brackets leave the running text.}
\psentry{ps_134-7_002-v0009}{9}{typewriter deletion}{\xout{by}}{}{The typewriter strikes out this material; it leaves the corrected running text.}
\psentry{ps_134-7_002-v0010}{10}{typed insertion}{\inserted[typed]{maybe}}{maybe}{The typed insertion is incorporated at its indicated position; insertion brackets leave the running text.}
\psentry{ps_134-7_002-v0011}{11}{typed insertion}{\inserted[typed]{is'nt}}{is'nt}{The typed insertion is incorporated at its indicated position; insertion brackets leave the running text.}
\psentry{ps_134-7_002-v0012}{12}{typewriter deletion}{\xout{may not be}}{}{The typewriter strikes out this material; it leaves the corrected running text.}
\psentry{ps_134-7_002-v0013}{13}{typewriter deletion}{\xout{its abelianization}}{}{The typewriter strikes out this material; it leaves the corrected running text.}
\psentry{ps_134-7_002-v0014}{14}{insertion}{\inserted{of homotopy types,}}{\hand{of homotopy types,}}{The handwritten insertion is incorporated at its indicated position; insertion brackets leave the running text.}
\psentry{ps_134-7_002-v0015}{15}{typewriter deletion}{\xout{\unsure{f}}}{}{The typewriter strikes out this material; it leaves the corrected running text.}
\psentry{ps_134-7_002-v0016}{16}{deletion}{\deleted{\illegible{1}}}{}{The author cancels this material by hand; it leaves the corrected running text.}
\psentry{ps_134-7_002-v0017}{17}{typewriter deletion}{\xout{of}}{}{The typewriter strikes out this material; it leaves the corrected running text.}
\psentry{ps_134-7_002-v0018}{18}{typewriter deletion}{\xout{at any rate}}{}{The typewriter strikes out this material; it leaves the corrected running text.}
\psapparatusend
\fi
\pssource{ps_134-7_003}{134-7/3}{https://grothendieck.umontpellier.fr/web/viewer.html?file=../134-7.pdf\#page=4}
\psnumber{ps_134-7_003-v0001}{1}
\psnumber{ps_134-7_003-v0002}{2}
\psnumber{ps_134-7_003-v0003}{3}
\psnumber{ps_134-7_003-v0004}{4}
\psnumber{ps_134-7_003-v0005}{5}
\psnumber{ps_134-7_003-v0006}{6}
\psnumber{ps_134-7_003-v0007}{7}
\psnumber{ps_134-7_003-v0008}{8}
\psnumber{ps_134-7_003-v0009}{9}
\psnumber{ps_134-7_003-v0010}{10}
\psnumber{ps_134-7_003-v0011}{11}
\psnumber{ps_134-7_003-v0012}{12}
\psnumber{ps_134-7_003-v0013}{13}
\psnumber{ps_134-7_003-v0014}{14}
\psnumber{ps_134-7_003-v0015}{15}
\psnumber{ps_134-7_003-v0016}{16}
\psnumber{ps_134-7_003-v0017}{17}
\psnumber{ps_134-7_003-v0018}{18}
\psnumber{ps_134-7_003-v0019}{19}
%% BEGIN SOURCE ps_134-7_003
% Pursuing Stacks (cotes 134-2 to 134-8): archive cote 134-7, site page 3 (PDF page 4), typescript page 445.
% typescript. Draft compared: ps6.tex:86-143.
% AI-transcribed by Claude Opus 5.5 (claude-opus-5-5, adaptive thinking, effort high), Message Batches API, batch ps3 (msgbatch_01XdFmuVVBAcgc7EFxdtCqsj), 2026-11-12; reviewed against both supplied scan halves; see CHECK_DECISIONS.tsv.
\scanpage{134-7/3}
\mspage{445}

Besides this,  I spent hours to try and put some order into the mess
of \var{\inserted[typed]{all}}{all}{typed insertion}{ps_134-7_003-v0001} Hom and tensor product type operations between categories $A^{\wedge}_{k}$ , $A^{\wedge}_{\hand{\mathcal{M}}}$
$B^{\wedge}_{k}$ , $B^{\wedge}_{\hand{\mathcal{M}}}$ (\var{\xout{and}}{}{typewriter deletion}{ps_134-7_003-v0002} \var{\inserted[typed]{or}}{or}{typed insertion}{ps_134-7_003-v0003} their \&-style generalizations), and the duality and Cartan-%
type isomorphisms between these.  There are a few more still than the
fair bunch met with in these notes so far  -  I finally renounced to get
really through and work out a wholly satisfactory set of notations,
taking into account all symmetries in the situation.  I realized this
\var{\xout{would take}}{}{typewriter deletion}{ps_134-7_003-v0004} \var{\inserted[typed]{might}}{might}{typed insertion}{ps_134-7_003-v0005} well take days of work,  while at present there is no real
need yet for it.  I sometimes find it difficult to find a proper balance
in these notes between the need of working on reasonably firm ground,
and \var{\xout{developping}}{}{typewriter deletion}{ps_134-7_003-v0006} working out suggestive terminology and notations for
what is coming up,  and on the other hand my resolve not to get caught
\var{\xout{by the}}{}{typewriter deletion}{ps_134-7_003-v0007} again by the ``El\hand{é}ments de G\hand{é}om\hand{é}trie Algébrique'' style\var{\xout{x}}{}{typewriter deletion}{ps_134-7_003-v0008} of work,
when it was understood that everything had to be worked out in complete
detail and in greatest generality,  for the benefit of generations of
``usagers'' (besides my own till the end of my life !..).  As a matter of
fact,  this whole ``abelianization story''\hand{,} going on now for well over \hand{a}
hundred pages and nothing really startling coming out  -  just things I
feel I should have known for ages\hand{,} \var{\deleted{-}}{}{deletion}{ps_134-7_003-v0009} has been won (so to say) over an
inner reluctance against these ``digressions'' in the main line of thought\hand{,}
the reluctance of one who is in a hurry to get through.  I know well
this old reluctance,feeling silly when\var{\inserted{ever}}{\hand{ever}}{insertion}{ps_134-7_003-v0010} working out ``trivial\var{\deleted{\unsure{''}}}{}{deletion}{ps_134-7_003-v0011} details'' \var{\deleted{\illegible{1}}}{}{deletion}{ps_134-7_003-v0012}
with utmost care;  as I know too that through this work only would I
come to a thorough understanding of what \var{\over{\unsure{was}}{is}}{\hand{is}}{writing over}{ps_134-7_003-v0013} going on,  and \var{\xout{unexpected}}{}{typewriter deletion}{ps_134-7_003-v0014} \var{\inserted[typed]{new}}{new}{typed insertion}{ps_134-7_003-v0015}
intuitions or relationships would flash up sometimes and open up
unexpected landscapes  and provide fresh impetus.  The same has happened
innumerable times too within the last seven years,  when ``meditating''
on personal matters  -  constantly ``the-one-in-a-hurry'' has turned out
to be just the servant of the inner resistances against renewal,
against a fresh, innocent look upon things \hand{\unsure{,}}familiar\hand{,} and \var{\xout{long}}{}{typewriter deletion}{ps_134-7_003-v0016} consistent-%
ly ignored as ``irrelevant''.  It does'nt seem the-one-in-a-hurry gets
at all discouraged for not getting his way many times  -  he seems to be
just as stubborn as the one who likes to take \var{\xout{t}}{}{typewriter deletion}{ps_134-7_003-v0017}his time and \var{\inserted[typed]{look up}}{look up}{typed insertion}{ps_134-7_003-v0018} \var{\deleted{do}}{}{deletion}{ps_134-7_003-v0019} things
thoroughly !

\margin{3}

% Visual collation: both supplied image halves inspected; see CHECK_DECISIONS_134_7_8.tsv.

%% END SOURCE ps_134-7_003
\ifnum\EditionMode=1
\psapparatusstart
\psentry{ps_134-7_003-v0001}{1}{typed insertion}{\inserted[typed]{all}}{all}{The typed insertion is incorporated at its indicated position; insertion brackets leave the running text.}
\psentry{ps_134-7_003-v0002}{2}{typewriter deletion}{\xout{and}}{}{The typewriter strikes out this material; it leaves the corrected running text.}
\psentry{ps_134-7_003-v0003}{3}{typed insertion}{\inserted[typed]{or}}{or}{The typed insertion is incorporated at its indicated position; insertion brackets leave the running text.}
\psentry{ps_134-7_003-v0004}{4}{typewriter deletion}{\xout{would take}}{}{The typewriter strikes out this material; it leaves the corrected running text.}
\psentry{ps_134-7_003-v0005}{5}{typed insertion}{\inserted[typed]{might}}{might}{The typed insertion is incorporated at its indicated position; insertion brackets leave the running text.}
\psentry{ps_134-7_003-v0006}{6}{typewriter deletion}{\xout{developping}}{}{The typewriter strikes out this material; it leaves the corrected running text.}
\psentry{ps_134-7_003-v0007}{7}{typewriter deletion}{\xout{by the}}{}{The typewriter strikes out this material; it leaves the corrected running text.}
\psentry{ps_134-7_003-v0008}{8}{typewriter deletion}{\xout{x}}{}{The typewriter strikes out this material; it leaves the corrected running text.}
\psentry{ps_134-7_003-v0009}{9}{deletion}{\deleted{-}}{}{The author cancels this material by hand; it leaves the corrected running text.}
\psentry{ps_134-7_003-v0010}{10}{insertion}{\inserted{ever}}{\hand{ever}}{The handwritten insertion is incorporated at its indicated position; insertion brackets leave the running text.}
\psentry{ps_134-7_003-v0011}{11}{deletion}{\deleted{\unsure{''}}}{}{The author cancels this material by hand; it leaves the corrected running text.}
\psentry{ps_134-7_003-v0012}{12}{deletion}{\deleted{\illegible{1}}}{}{The author cancels this material by hand; it leaves the corrected running text.}
\psentry{ps_134-7_003-v0013}{13}{writing over}{\over{\unsure{was}}{is}}{\hand{is}}{The superposed handwritten reading replaces the underlying reading.}
\psentry{ps_134-7_003-v0014}{14}{typewriter deletion}{\xout{unexpected}}{}{The typewriter strikes out this material; it leaves the corrected running text.}
\psentry{ps_134-7_003-v0015}{15}{typed insertion}{\inserted[typed]{new}}{new}{The typed insertion is incorporated at its indicated position; insertion brackets leave the running text.}
\psentry{ps_134-7_003-v0016}{16}{typewriter deletion}{\xout{long}}{}{The typewriter strikes out this material; it leaves the corrected running text.}
\psentry{ps_134-7_003-v0017}{17}{typewriter deletion}{\xout{t}}{}{The typewriter strikes out this material; it leaves the corrected running text.}
\psentry{ps_134-7_003-v0018}{18}{typed insertion}{\inserted[typed]{look up}}{look up}{The typed insertion is incorporated at its indicated position; insertion brackets leave the running text.}
\psentry{ps_134-7_003-v0019}{19}{deletion}{\deleted{do}}{}{The author cancels this material by hand; it leaves the corrected running text.}
\psapparatusend
\fi
\pssource{ps_134-7_004}{134-7/4}{https://grothendieck.umontpellier.fr/web/viewer.html?file=../134-7.pdf\#page=5}
\psnumber{ps_134-7_004-v0001}{1}
\psnumber{ps_134-7_004-v0002}{2}
\psnumber{ps_134-7_004-v0003}{3}
\psnumber{ps_134-7_004-v0004}{4}
\psnumber{ps_134-7_004-v0005}{5}
\psnumber{ps_134-7_004-v0006}{6}
%% BEGIN SOURCE ps_134-7_004
% Pursuing Stacks (cotes 134-2 to 134-8): archive cote 134-7, site page 4 (PDF page 5), typescript page 446.
% typescript. Draft compared: ps6.tex:125-206.
% AI-transcribed by Claude Opus 5.5 (claude-opus-5-5, adaptive thinking, effort high), Message Batches API, batch ps3 (msgbatch_01XdFmuVVBAcgc7EFxdtCqsj), 2026-11-12; reviewed against both supplied scan halves; see CHECK_DECISIONS.tsv.
\scanpage{134-7/4}
\mspage{446}

25.8.

\margin{111.}\var{\xout{\unsure{141}}}{}{typewriter deletion}{ps_134-7_004-v0001}. Still about "schematization" of homotopy types ! Here is a ten-%
tative approach,  without any explicit use of Postnikov fibrations nor
abelianization,  although both are involved implicitly.  If n is any natu-%
ral integer,  I'll denote by
\[ \mathrm{Hot}_n \]
the full subcategory of the \var{\inserted[typed]{pointed}}{pointed}{typed insertion}{ps_134-7_004-v0002} homotopy category  $\mathrm{Hot}^{\bullet}$,  made up with n-%
connected homotopy types,  with the extra assumption for n=0 that the
fundamental group be abelian.  For any (commutative) ring k,  I want to
define a category
\[ \hand{(a)} \qquad \mathrm{Hot}_n(k) \quad , \]
depending covariantly on k,  in such a way that we have an equivalence
\[ \hand{(b)} \qquad \mathrm{Hot}_n(\underline{Z}) \hand{\xrightarrow{\approx}} \mathrm{Hot}_n \quad , \]
which should come from a canonical functor
\[ \hand{(c)} \qquad \mathrm{Hot}_n(k) \hand{\longrightarrow} \mathrm{Hot}_n \]
defined \var{\xout{frxm}}{}{typewriter deletion}{ps_134-7_004-v0003} for any \hand{k,}\var{\xout{\unsure{k}}}{}{typewriter deletion}{ps_134-7_004-v0004}and which should be viewed as ground-ring restri-%
tion from k to $\underline{Z}$  -  more generally,  for any ring homomorphism
\[ k' \hand{\longrightarrow} k \]
we expect a restriction functor,  besides the ring extension functor
\[ \hand{(d)} \qquad \mathrm{Hot}_n(k') \hand{\longrightarrow} \mathrm{Hot}_n(k) \quad \text{and} \quad \mathrm{Hot}_n(k) \hand{\longrightarrow} \mathrm{Hot}_n(k') \quad . \]
Among other important features to expect,  is that for any object X in
$\mathrm{Hot}_n(k)$ ,  the homotopy groups $\hand{\pi}_i(X)$  (defined via \hand{(c)}) should be natu-%
rally endowed with structures of k-modules,  and the ring extension and
restriction functors\hand{(d)} should be compatible with these.

Here is an idea for getting such a theory.  For given k,  we first
define an auxiliary category
\[ \hand{(e)} \qquad U(k) \quad , \]
whose objects may be called "unipotent bundles over k".  \var{\xout{One possible}}{}{typewriter deletion}{ps_134-7_004-v0005}
These \var{\xout{\illegible{1}}}{}{typewriter deletion}{ps_134-7_004-v0006}"bundles" will not be quite schemes over k,  they will be defined
as functors
\[ \hand{(f)} \qquad \mathrm{Alg}_{/k} \hand{\longrightarrow} (\mathrm{Sets}) \]
where $\mathrm{Alg}_{/k}$  is the category of \hand{(}commutative) k-algebras (in the basic
universe $\underline{U}$ ).  The opposite category may be identified \hand{with} the category
\[ \mathrm{Aff}_{/k} \]

\margin{4}

% Visual collation: both supplied image halves inspected; see CHECK_DECISIONS_134_7_8.tsv.

%% END SOURCE ps_134-7_004
\ifnum\EditionMode=1
\psapparatusstart
\psentry{ps_134-7_004-v0001}{1}{typewriter deletion}{\xout{\unsure{141}}}{}{The typewriter strikes out this material; it leaves the corrected running text.}
\psentry{ps_134-7_004-v0002}{2}{typed insertion}{\inserted[typed]{pointed}}{pointed}{The typed insertion is incorporated at its indicated position; insertion brackets leave the running text.}
\psentry{ps_134-7_004-v0003}{3}{typewriter deletion}{\xout{frxm}}{}{The typewriter strikes out this material; it leaves the corrected running text.}
\psentry{ps_134-7_004-v0004}{4}{typewriter deletion}{\xout{\unsure{k}}}{}{The typewriter strikes out this material; it leaves the corrected running text.}
\psentry{ps_134-7_004-v0005}{5}{typewriter deletion}{\xout{One possible}}{}{The typewriter strikes out this material; it leaves the corrected running text.}
\psentry{ps_134-7_004-v0006}{6}{typewriter deletion}{\xout{\illegible{1}}}{}{The typewriter strikes out this material; it leaves the corrected running text.}
\psapparatusend
\fi
\pssource{ps_134-7_005}{134-7/5}{https://grothendieck.umontpellier.fr/web/viewer.html?file=../134-7.pdf\#page=6}
\psnumber{ps_134-7_005-v0001}{1}
\psnumber{ps_134-7_005-v0002}{2}
\psnumber{ps_134-7_005-v0003}{3}
\psnumber{ps_134-7_005-v0004}{4}
\psnumber{ps_134-7_005-v0005}{5}
\psnumber{ps_134-7_005-v0006}{6}
%% BEGIN SOURCE ps_134-7_005
% Pursuing Stacks (cotes 134-2 to 134-8): archive cote 134-7, site page 5 (PDF page 6), typescript page 447.
% typescript. Draft compared: ps6.tex:185-244.
% AI-transcribed by Claude Opus 5.5 (claude-opus-5-5, adaptive thinking, effort high), Message Batches API, batch ps3 (msgbatch_01XdFmuVVBAcgc7EFxdtCqsj), 2026-11-12; reviewed against both supplied scan halves; see CHECK_DECISIONS.tsv.
\scanpage{134-7/5}
\mspage{447}

of affine schemes over k , thus we'll be working in the category of
functors (or presheaves over $\mathrm{Aff}_{/k}$ )
\[ \text{(\var{\over{\unsure{g}}{f}}{\hand{f}}{writing over}{ps_134-7_005-v0001}')} \qquad (\mathrm{Aff}_{/k})^{o} \hand{\longrightarrow} (\mathrm{Sets}) \quad , \]
more specifically U(k) will be a full subcategory of this category of
functors. We'll endow $\mathrm{Aff}_{/k}$ with one of the standard site structures,
the most convenient one here is the fpqc topology (faithfully flat qua-%
si compact topology), and work in the category of \uline{sheaves} in the
latter. In terms of the interpretation \var{\over{\unsure{g}}{f}}{\hand{f}}{writing over}{ps_134-7_005-v0002}) as covariant functors\var{\inserted{$X$}}{\hand{$X$}}{insertion}{ps_134-7_005-v0003} on
$\mathrm{Alg}_{/k}$ , this just means that we are restricting to functors which
1) commute to finite products and 2) are``compatible with faithfully
flat descent'' i.e. for any map
\[ k' \hand{\longrightarrow} k'' \]
in $\mathrm{Alg}_{/k}$ such that k'' becomes a faithfully flat algebra over k' ,
the following diagram in (Sets)
\[ X(k') \hand{\longrightarrow} X(k'') \hand{\rightrightarrows} X(k''\otimes_{k'}k'') \]
is exact. Thus U(k) will be defined as a full subcategory of the\var{\xout{x}}{}{typewriter deletion}{ps_134-7_005-v0004}
category of such functors, or ``sheaves''.

One way for defining U(k), is to ppesent it as the union of a
sequence of subcategories $U_m(k)$ (m a natural integer). We'll take
$U_o(k)$ to be just reduced to the final functors (i.e. X(k') is a one
point set for any k' in $\mathrm{Alg}_{/k}$ ), and define inductively $U_{m+1}(k)$ in
terms of $U_m(k)$ as follows. For any k-module M, let W(M) be the corres-%
ponding ``vector bundle'' , defined by
\[ \text{\hand{(g)}} \qquad W(M)(k') = M\otimes_k k' \quad , \]
then an object is in $U_{m+1}(k)$ iff there exists an object $X_m$ in $U_m(k)$,
and a k-module $M_{\var{\xout{\unsure{m}}}{}{typewriter deletion}{ps_134-7_005-v0005}} = M_{m+1}$ , in such a way that $X=X_{m+1}$ should be iso-%
morphic to a ``torsor'' over $X_{m\var{\xout{\unsure{+1}}}{}{typewriter deletion}{ps_134-7_005-v0006}}$ , with group W(M). I am not too sure
here whether we \unsure{shoubd|should} view the objects of $U_m(k)$ as endowed with the
extra structure consisting in giving the modules $M_1$, ... , $M_m$ used in
the inductive construction, and mo\unsure{re}over the successive fibrations -
if so, then of course the categories $U_m(k)$, and their union U(k),
will no longer be interpreted as a mere subcategory of the category of

\margin{5}

% Visual collation: both supplied image halves inspected; see CHECK_DECISIONS_134_7_8.tsv.

%% END SOURCE ps_134-7_005
\ifnum\EditionMode=1
\psapparatusstart
\psentry{ps_134-7_005-v0001}{1}{writing over}{\over{\unsure{g}}{f}}{\hand{f}}{The superposed handwritten reading replaces the underlying reading.}
\psentry{ps_134-7_005-v0002}{2}{writing over}{\over{\unsure{g}}{f}}{\hand{f}}{The superposed handwritten reading replaces the underlying reading.}
\psentry{ps_134-7_005-v0003}{3}{insertion}{\inserted{$X$}}{\hand{$X$}}{The handwritten insertion is incorporated at its indicated position; insertion brackets leave the running text.}
\psentry{ps_134-7_005-v0004}{4}{typewriter deletion}{\xout{x}}{}{The typewriter strikes out this material; it leaves the corrected running text.}
\psentry{ps_134-7_005-v0005}{5}{typewriter deletion}{$\xout{\unsure{m}}$}{}{The typewriter strikes out this material; it leaves the corrected running text.}
\psentry{ps_134-7_005-v0006}{6}{typewriter deletion}{$\xout{\unsure{+1}}$}{}{The typewriter strikes out this material; it leaves the corrected running text.}
\psapparatusend
\fi
\pssource{ps_134-7_006}{134-7/6}{https://grothendieck.umontpellier.fr/web/viewer.html?file=../134-7.pdf\#page=7}
\psnumber{ps_134-7_006-v0001}{1}
\psnumber{ps_134-7_006-v0002}{2}
\psnumber{ps_134-7_006-v0003}{3}
\psnumber{ps_134-7_006-v0004}{4}
\psnumber{ps_134-7_006-v0005}{5}
\psnumber{ps_134-7_006-v0006}{6}
\psnumber{ps_134-7_006-v0007}{7}
%% BEGIN SOURCE ps_134-7_006
% Pursuing Stacks (cotes 134-2 to 134-8): archive cote 134-7, site page 6 (PDF page 7), typescript page 448.
% typescript. Draft compared: ps6.tex:232-295.
% AI-transcribed by Claude Opus 5.5 (claude-opus-5-5, adaptive thinking, effort high), Message Batches API, batch ps3 (msgbatch_01XdFmuVVBAcgc7EFxdtCqsj), 2026-11-12; reviewed against both supplied scan halves; see CHECK_DECISIONS.tsv.
\scanpage{134-7/6}
\mspage{448}

\var{\deleted{\illegible{1}}}{}{deletion}{ps_134-7_006-v0001}
sheaves \var{\xout{of sets}}{}{typewriter deletion}{ps_134-7_006-v0002} (of sets) just described.  Possibly both approaches are
of interest,  and will yield non equivalent notions of schematization.
On the other hand,  although definitely X is not representable by a
usual \var{\xout{affine}}{}{typewriter deletion}{ps_134-7_006-v0003} scheme over k unless the k-modules $M_i$ are projective of
finite type (in which case X will be an affine scheme,  and even isomor-%
phic,  at least locally over Spec(k),  to standard affine space $E_k^d$ for
suitable  d),  it is felt that X,  as far as cohomology properties go,
should be very close to being an affine scheme,  and that presumably
it's cohomology groups $H^i$  with coefficients in ``quasi-coherent'' \var{\xout{\illegible{1}}}{}{typewriter deletion}{ps_134-7_006-v0004}
sheaves such as W(M),  should vanish for $i>0$;  consequently,  presumably
the torsors used for the inductive construction of X are trivial,  which
means that X is in fact isomorphic to the product of all $W(M_i)$'s .  In
the case when we disregard the successive fibration structure,  this
means that the objects of U(k) are just sheaves of sets which are isomor-%
phic to some W(M)  (where morally M is the direct sum of the modules $M_i$
which\var{\xout{\illegible{1}}}{}{typewriter deletion}{ps_134-7_006-v0005} have been used in our inductive definition).  This gives then a
rather trivial description of the objects of U(k) (and all $U_m(k)$'s are
already equal to U(k)\hand{,} for $m \var{\over{>}{\geq}}{\hand{\geq}}{writing over}{ps_134-7_006-v0006} 1$ ! ),  it should be remembered however \var{\xout{\unsure{that}}}{}{typewriter deletion}{ps_134-7_006-v0007}
that maps in U(k)  from a W(M) to a W(M')  are a lot more general than
just k-linear maps  $M \hand{\to} M'$  (they may be viewed as ``polynomial maps
from M to M'\,'').

The category U(k) will be endowed with the \uline{sections} \uline{functor}
\[ \text{(h)} \qquad X \hand{\mapsto} X(k) \ : \ U(k) \hand{\longrightarrow} \text{(Sets)} \quad . \]
Now let A be any test category,  for instance A = \hand{$\Delta$},  and consider
the functor
\[ \hand{(h')} \qquad \underline{\mathrm{Hom}}(A^{o},U(k)) \hand{\longrightarrow} \underline{\mathrm{Hom}}(A^{o},\text{(Sets)}) \]
induced by h).  The second member modelizes homotopy types,  which there-%
fore allows us to defined homotopy invariants for objects in the first
member,  and hence to define the property of n-connectedness,  and (if
n=0)  of 0-connectedness with abelian fundamental group.  As a matter of
fact,  we would like to define a subcategory $\hand{\mathcal{M}}_m(k)$   of the first mem-%
ber,  so that it should become clear that for an object $X_{\hand{*}}$  in it, it's
homotopy groups are endowed with structures of k-modules,  so that for
n=0  the abelian restriction on $\hand{\pi}_1$ is superfluous,  because automatic.
To be specific,  we better restrict maybe to A= \hand{$\Delta$},  more familiar to
us,  so that $X_{\hand{*}}$ may be viewed as a ``complex''  ($n \hand{\mapsto} X_n$)  with components
in U(k).  The kind of restriction I am thinking of for defining the

\margin{6}

% Visual collation: both supplied image halves inspected; see CHECK_DECISIONS_134_7_8.tsv.

%% END SOURCE ps_134-7_006
\ifnum\EditionMode=1
\psapparatusstart
\psentry{ps_134-7_006-v0001}{1}{deletion}{\deleted{\illegible{1}}}{}{The author cancels this material by hand; it leaves the corrected running text.}
\psentry{ps_134-7_006-v0002}{2}{typewriter deletion}{\xout{of sets}}{}{The typewriter strikes out this material; it leaves the corrected running text.}
\psentry{ps_134-7_006-v0003}{3}{typewriter deletion}{\xout{affine}}{}{The typewriter strikes out this material; it leaves the corrected running text.}
\psentry{ps_134-7_006-v0004}{4}{typewriter deletion}{\xout{\illegible{1}}}{}{The typewriter strikes out this material; it leaves the corrected running text.}
\psentry{ps_134-7_006-v0005}{5}{typewriter deletion}{\xout{\illegible{1}}}{}{The typewriter strikes out this material; it leaves the corrected running text.}
\psentry{ps_134-7_006-v0006}{6}{writing over}{$\over{>}{\geq}$}{$\hand{\geq}$}{The superposed handwritten reading replaces the underlying reading.}
\psentry{ps_134-7_006-v0007}{7}{typewriter deletion}{\xout{\unsure{that}}}{}{The typewriter strikes out this material; it leaves the corrected running text.}
\psapparatusend
\fi
\pssource{ps_134-7_007}{134-7/7}{https://grothendieck.umontpellier.fr/web/viewer.html?file=../134-7.pdf\#page=8}
\psnumber{ps_134-7_007-v0001}{1}
\psnumber{ps_134-7_007-v0002}{2}
\psnumber{ps_134-7_007-v0003}{3}
\psnumber{ps_134-7_007-v0004}{4}
\psnumber{ps_134-7_007-v0005}{5}
\psnumber{ps_134-7_007-v0006}{6}
\psnumber{ps_134-7_007-v0007}{7}
\psnumber{ps_134-7_007-v0008}{8}
\psnumber{ps_134-7_007-v0009}{9}
\psnumber{ps_134-7_007-v0010}{10}
\psnumber{ps_134-7_007-v0011}{11}
\psnumber{ps_134-7_007-v0012}{12}
%% BEGIN SOURCE ps_134-7_007
% Pursuing Stacks (cotes 134-2 to 134-8): archive cote 134-7, site page 7 (PDF page 8), typescript page 449.
% typescript. Draft compared: ps6.tex:279-349.
% AI-transcribed by Claude Opus 5.5 (claude-opus-5-5, adaptive thinking, effort high), Message Batches API, batch ps3 (msgbatch_01XdFmuVVBAcgc7EFxdtCqsj), 2026-11-12; reviewed against both supplied scan halves; see CHECK_DECISIONS.tsv.
\scanpage{134-7/7}
\mspage{449}

model category \hand{$\mathcal{M}$}$_n(k)$ is :

(i) \qquad $X_i = e$ \quad for \quad $i \hand{\leq} n$ \quad ,

where e is the final object of U(k), and also maybe, if we adopt the
more refined version of U(k)\var{\deleted{,}}{}{deletion}{ps_134-7_007-v0001}as a strictly increasing union of subcate-%
gories $U_m(k)$\hand{,}

\hand{(}i') \qquad $X_i$ is in $U_{\var{\over{n}{i}}{\hand{i}}{writing over}{ps_134-7_007-v0002}-n}(k)$ , for any $i \hand{\geq} n$ \quad .

One may have to play around some more to get ``the correct'' description
of the model category, \var{\xout{xxx(k)}}{}{typewriter deletion}{ps_134-7_007-v0003} which I tentatively propose to define \var{\inserted[typed]{simply}}{simply}{typed insertion}{ps_134-7_007-v0004} as
a \var{\inserted[typed]{suitable}}{suitable}{typed insertion}{ps_134-7_007-v0005} full subcategory

(j) \qquad $\hand{\mathcal{M}}_n(k) \hand{\subset} \underline{\mathrm{Hom}}(A^{\mathrm{o}}, U(k))$ \quad .

The functor\hand{(}\var{\over{\unsure{i}}{h}}{\hand{h}}{writing over}{ps_134-7_007-v0006}') allows to define a notion of weak equivalence in
\hand{$\mathcal{M}$}$_n(k)$, hence a localized category Hot$_n(k)$, and a functor\hand{(}\var{\over{\unsure{b}}{c}}{\hand{c}}{writing over}{ps_134-7_007-v0007}) from
this category to Hot$_n$ . The ring extension and restriction functors
\hand{(}d) are equally defined in an evident way, via corresponding functors
on the model categories (with the task however to check that these
are compatible with weak equivalences). The \var{\xout{main}}{}{typewriter deletion}{ps_134-7_007-v0008}\var{\inserted[typed]{key}}{key}{typed insertion}{ps_134-7_007-v0009} point here is to
check that for k=$\underline{Z}$ , the functor\hand{(}c) (namely\hand{(}b)) is indeed an equiva-%
lence of categories. Thus the main task seems to cut out carefully a
description of a model category \hand{$\mathcal{M}$}$_n(k)$, in terms of semisimplicial
objects say in a category such as U(k) , in such a way as to give rise
to an equivalence of categories\hand{(}b).

One point which is still somewhat misty in this (admittedly overall
misty !) picture, is how to get, for an object X in Hot$_n(k)$, the
promised operation of k on the homotopy groups $\hand{\pi}_i(X)$. I was thinking
about this when suggesting the conditions\hand{(}i) and\hand{(}i') above on k-models
for homotopy types - but I really doubt these are enough. On the
other hand, it seems hard to imagine \var{\xout{a notion}}{}{typewriter deletion}{ps_134-7_007-v0010} there be a good notion
of homotopy types over k, without the homotopy groups to be k-modules
over k, not just abelian groups. Even more still, there should be
moreover a\var{\xout{n}}{}{typewriter deletion}{ps_134-7_007-v0011} ``linearization functor''

\hand{(}k) \qquad Hot$_n(k) \hand{\longrightarrow}$ D.(Ab$_k$) \quad ,

with values in the derived category of the category\var{\inserted{of chain complexes in}}{\hand{of chain complexes in}}{insertion}{ps_134-7_007-v0012} Ab$_k$ = k-Mod of
k-modules, presumably coming by localization from a functor

\hand{(k')} \qquad \hand{$\mathcal{M}$}$_n(k) \hand{\longrightarrow}$ Ch.(Ab$_k$ ) \quad ,

and giving rise to a commutative diagram \margin{only for \unsure{$k=\mathbb{Z}$|$n=\mathbb{Z}$}}

\margin{7}

% Visual collation: both supplied image halves inspected; see CHECK_DECISIONS_134_7_8.tsv.

%% END SOURCE ps_134-7_007
\ifnum\EditionMode=1
\psapparatusstart
\psentry{ps_134-7_007-v0001}{1}{deletion}{\deleted{,}}{}{The author cancels this material by hand; it leaves the corrected running text.}
\psentry{ps_134-7_007-v0002}{2}{writing over}{$\over{n}{i}$}{$\hand{i}$}{The superposed handwritten reading replaces the underlying reading.}
\psentry{ps_134-7_007-v0003}{3}{typewriter deletion}{\xout{xxx(k)}}{}{The typewriter strikes out this material; it leaves the corrected running text.}
\psentry{ps_134-7_007-v0004}{4}{typed insertion}{\inserted[typed]{simply}}{simply}{The typed insertion is incorporated at its indicated position; insertion brackets leave the running text.}
\psentry{ps_134-7_007-v0005}{5}{typed insertion}{\inserted[typed]{suitable}}{suitable}{The typed insertion is incorporated at its indicated position; insertion brackets leave the running text.}
\psentry{ps_134-7_007-v0006}{6}{writing over}{\over{\unsure{i}}{h}}{\hand{h}}{The superposed handwritten reading replaces the underlying reading.}
\psentry{ps_134-7_007-v0007}{7}{writing over}{\over{\unsure{b}}{c}}{\hand{c}}{The superposed handwritten reading replaces the underlying reading.}
\psentry{ps_134-7_007-v0008}{8}{typewriter deletion}{\xout{main}}{}{The typewriter strikes out this material; it leaves the corrected running text.}
\psentry{ps_134-7_007-v0009}{9}{typed insertion}{\inserted[typed]{key}}{key}{The typed insertion is incorporated at its indicated position; insertion brackets leave the running text.}
\psentry{ps_134-7_007-v0010}{10}{typewriter deletion}{\xout{a notion}}{}{The typewriter strikes out this material; it leaves the corrected running text.}
\psentry{ps_134-7_007-v0011}{11}{typewriter deletion}{\xout{n}}{}{The typewriter strikes out this material; it leaves the corrected running text.}
\psentry{ps_134-7_007-v0012}{12}{insertion}{\inserted{of chain complexes in}}{\hand{of chain complexes in}}{The handwritten insertion is incorporated at its indicated position; insertion brackets leave the running text.}
\psapparatusend
\fi
\pssource{ps_134-7_008}{134-7/8}{https://grothendieck.umontpellier.fr/web/viewer.html?file=../134-7.pdf\#page=9}
\psnumber{ps_134-7_008-v0001}{1}
\psnumber{ps_134-7_008-v0002}{2}
\psnumber{ps_134-7_008-v0003}{3}
\psnumber{ps_134-7_008-v0004}{4}
\psnumber{ps_134-7_008-v0005}{5}
\psnumber{ps_134-7_008-v0006}{6}
%% BEGIN SOURCE ps_134-7_008
% Pursuing Stacks (cotes 134-2 to 134-8): archive cote 134-7, site page 8 (PDF page 9), typescript page 450.
% typescript. Draft compared: ps6.tex:329-399.
% AI-transcribed by Claude Opus 5.5 (claude-opus-5-5, adaptive thinking, effort high), Message Batches API, batch ps3 (msgbatch_01XdFmuVVBAcgc7EFxdtCqsj), 2026-11-12; reviewed against both supplied scan halves; see CHECK_DECISIONS.tsv.
\scanpage{134-7/8}
\mspage{450}

\diagram{\begin{tikzcd}
Hot_n(k) \arrow[color=blue,r] \arrow[color=blue,d] \& D.(Ab_k) \arrow[color=blue,d] \\
Hot_n \arrow[color=blue,r] \& D\hand{.}(Ab)
\end{tikzcd}}
\margin{(l)}
,

where the second horizontal arrow is the usual abelianization functor
for homotopy types,  and the second vertical one comes from ring restric-%
tion functor  $Ab_k \hand{\longrightarrow} Ab = Ab_{\underline{Z}}$ .  In \hand{(}k'\hand{),} Ch.  denotes the category
of chain complexes,  and it looks rather mysterious again how to get
such a functor\hand{(}k').  \var{\xout{But}}{}{typewriter deletion}{ps_134-7_008-v0001} We may of course think of the trivial abeliani-%
zation
\[ X. \hand{\mapsto} k^{(X.)} \stackrel{def}{=} (n \hand{\mapsto} k^{(X_n)} ) \quad , \]
where for an object X in U(k),  or more generally any sheaf on $Aff_{/k}$ ,
$k^{(X)}$ defines trivial k-linearization of this sheaf,  in the sense of the
topos of all such sheaves.  Anyhow $k^{(X)}$ is a sheaf,  not a k-module,
so we should still take sections to get what we want  -  but this \var{\inserted{functor}}{\hand{functor}}{insertion}{ps_134-7_008-v0002} looks
not only prohibitively large and unaccessible,  but just silly\hand{!}  A much
better choice for k-linearizing objects of U(k) specifically seems the
following.  Disregarding the fibration structures,  such an object X is
isomorphic to an object W(M),  M some k-module.  We look for a k-lineari-%
zation
\[ \text{(m)} \qquad X \hand{(}\var{\over{=}{\simeq}}{\hand{\simeq}}{writing over}{ps_134-7_008-v0003}W(M)) \hand{\longrightarrow} W(L(\hand{X})) \quad , \]
where L(X) is a suitable k-module.  Now,  among all maps
\[ X \hand{\longrightarrow} W(N) \]
\margin{X} of X into sheaves \uline{of the type} W(N),  there is a universal one,  which
in terms of M can be described as
\[ N = \hand{\Gamma}_k(M) \quad , \]
where $\hand{\Gamma}_k$  denotes the ``algebra with devided powers\var{\deleted{''}}{}{deletion}{ps_134-7_008-v0004} generated by M\hand{''},
the canonical map  $M \hand{\longrightarrow} \hand{\Gamma}_k(M)$  or rather
\[ \hand{(}\text{n}) \qquad W(M) \hand{\longrightarrow} W(\hand{\Gamma}_k(M)) \quad , \quad x \hand{\mapsto} exp(x) = \hand{\sum_{i\geq 0}} x^{(i)} \]
being \uline{the} ``\uline{universal}\var{\deleted{''}}{}{deletion}{ps_134-7_008-v0005} polynomial map\hand{''}  of M with values in a module \var{\over{M}{N}}{\hand{N}}{writing over}{ps_134-7_008-v0006}
(or rather,  of W(M) into W(N)).   Here  $x^{(i)}$ denotes the i.th devided
power of x ,  which is an element of  $\hand{\Gamma}_k^{\,i}(M)$ .  It just occurs to me
that this expression of exp(x),  the universal map,  is infinite,  thus
it does'nt take its values in W(N)  actually,  but in a suitable comple-%
tion of it - this does'nt seem too serious a drawback though !  The
point I wish to make here,  is that for given X in U(k),  defining

\margin{8}

% Visual collation: both supplied image halves inspected; see CHECK_DECISIONS_134_7_8.tsv.

%% END SOURCE ps_134-7_008
\ifnum\EditionMode=1
\psapparatusstart
\psentry{ps_134-7_008-v0001}{1}{typewriter deletion}{\xout{But}}{}{The typewriter strikes out this material; it leaves the corrected running text.}
\psentry{ps_134-7_008-v0002}{2}{insertion}{\inserted{functor}}{\hand{functor}}{The handwritten insertion is incorporated at its indicated position; insertion brackets leave the running text.}
\psentry{ps_134-7_008-v0003}{3}{writing over}{$\over{=}{\simeq}$}{$\hand{\simeq}$}{The superposed handwritten reading replaces the underlying reading.}
\psentry{ps_134-7_008-v0004}{4}{deletion}{\deleted{''}}{}{The author cancels this material by hand; it leaves the corrected running text.}
\psentry{ps_134-7_008-v0005}{5}{deletion}{\deleted{''}}{}{The author cancels this material by hand; it leaves the corrected running text.}
\psentry{ps_134-7_008-v0006}{6}{writing over}{\over{M}{N}}{\hand{N}}{The superposed handwritten reading replaces the underlying reading.}
\psapparatusend
\fi
\pssource{ps_134-7_009}{134-7/9}{https://grothendieck.umontpellier.fr/web/viewer.html?file=../134-7.pdf\#page=10}
\psnumber{ps_134-7_009-v0001}{1}
\psnumber{ps_134-7_009-v0002}{2}
\psnumber{ps_134-7_009-v0003}{3}
\psnumber{ps_134-7_009-v0004}{4}
\psnumber{ps_134-7_009-v0005}{5}
\psnumber{ps_134-7_009-v0006}{6}
\psnumber{ps_134-7_009-v0007}{7}
\psnumber{ps_134-7_009-v0008}{8}
\psnumber{ps_134-7_009-v0009}{9}
\psnumber{ps_134-7_009-v0010}{10}
\psnumber{ps_134-7_009-v0011}{11}
\psnumber{ps_134-7_009-v0012}{12}
\psnumber{ps_134-7_009-v0013}{13}
\psnumber{ps_134-7_009-v0014}{14}
\psnumber{ps_134-7_009-v0015}{15}
\psnumber{ps_134-7_009-v0016}{16}
\psnumber{ps_134-7_009-v0017}{17}
\psnumber{ps_134-7_009-v0018}{18}
\psnumber{ps_134-7_009-v0019}{19}
\psnumber{ps_134-7_009-v0020}{20}
\psnumber{ps_134-7_009-v0021}{21}
\psnumber{ps_134-7_009-v0022}{22}
%% BEGIN SOURCE ps_134-7_009
% Pursuing Stacks (cotes 134-2 to 134-8): archive cote 134-7, site page 9 (PDF page 10), typescript page 451.
% typescript. Draft compared: ps6.tex:383-449.
% AI-transcribed by Claude Opus 5.5 (claude-opus-5-5, adaptive thinking, effort high), Message Batches API, batch ps3 (msgbatch_01XdFmuVVBAcgc7EFxdtCqsj), 2026-11-12; reviewed against both supplied scan halves; see CHECK_DECISIONS.tsv.
\scanpage{134-7/9}
\mspage{451}

\var{\xout{xx}}{}{typewriter deletion}{ps_134-7_009-v0001})
(o) \qquad $L_k(X) = \var{\deleted{\illegible{1}}}{}{deletion}{ps_134-7_009-v0002}\,\hand{\Gamma}_k^{\hand{\wedge}}(M)\var{\deleted{)}}{}{deletion}{ps_134-7_009-v0003}\,\var{\inserted{\overset{\text{def}}{=}\prod_{i\geq 0}\Gamma^{i}(M)\deleted{)}}}{\hand{\overset{\text{def}}{=}\prod_{i\geq 0}\Gamma^{i}(M)}}{insertion}{ps_134-7_009-v0004}$ \hand{(}(k-linearization of X) ,
where M is any k-module endowed with an isomorphism
\[ u:\ X \hand{\simeq} W(M) \quad , \]
the \var{\inserted[typed]{k-module}}{k-module}{typed insertion}{ps_134-7_009-v0005} \var{\xout{object}}{}{typewriter deletion}{ps_134-7_009-v0006} $L_k(X)$ \uline{does not depend}, \uline{up to unique isomorphism}, \uline{on the
choice of a pair} (M,u), because \var{\xout{anyxx}}{}{typewriter deletion}{ps_134-7_009-v0007} for two k-modules M,M', any
morphism of \uline{sheaves of sets}
\[ v:\ W(M) \hand{\longrightarrow} W(M') \]
induces a\var{\xout{n}}{}{typewriter deletion}{ps_134-7_009-v0008} \var{\xout{isomorphism}}{}{typewriter deletion}{ps_134-7_009-v0009} homomorphism of k-modules
\[ \hand{\Gamma}_k^{\hand{\wedge}}(v):\quad \hand{\Gamma}_k^{\hand{\wedge}}(M) \hand{\longrightarrow} \hand{\Gamma}_k^{\hand{\wedge}}(M') \]
(compatible not with multiplications, but with diagonal maps ...), \var{\xout{whi}}{}{typewriter deletion}{ps_134-7_009-v0010}
which will be an isomorphism if v is.

Thus we do have, it seems to me, a good candidate for k-lineariza-%
tion. To check it is suitable indeed, the main point seems to check
that the corresponding diagram \hand{(l)} commutes up to canonical isomorphism,
the crucial case of course being $k=\underline{Z}$ . This now looks like a rather down
to earth question, which seems to me a pretty good test, whether the
intuition of schematization of homotopy types is a sound one. Let's
rephrase it here\var{\over{,}{.}}{\hand{.}}{writing over}{ps_134-7_009-v0011} \var{\over{f}{F}}{\hand{F}}{writing over}{ps_134-7_009-v0012}or this, let's first restate the \var{\xout{definitionxx}}{}{typewriter deletion}{ps_134-7_009-v0013}
description of the \var{\xout{\illegible{1}}}{}{typewriter deletion}{ps_134-7_009-v0014} category U(k) (coarse version) in the more
\var{\xout{\illegible{1}}}{}{typewriter deletion}{ps_134-7_009-v0015}
down-to-earth terms of linear algebra. Objects may be viewed as just
k-modules M, whereas \var{\xout{maps}}{}{typewriter deletion}{ps_134-7_009-v0016} (non additive\hand{)} ``maps''from M to M' (defined
previously as maps $W(M)\hand{\longrightarrow} W(M')$ of sheaves of sets) are described as
just continuous k-linear maps
\[ \hand{(p)} \qquad \hand{f}:\ \hand{\Gamma}_k^{\hand{\wedge}}(M) \hand{\rightarrow} \hand{\Gamma}_k^{\hand{\wedge}}(M') \quad , \]
which are mor\hand{e}over compatible with the natural augmentations to k,
and with the natural diagonal maps \hand{:}
\[ \hand{\varepsilon:\Gamma_k^{\wedge}(M)\rightarrow k,}\quad \hand{\Delta}:\ \hand{\Gamma}_k^{\hand{\wedge}}(M)\hand{\longrightarrow} \hand{\Gamma}_k^{\hand{\wedge}}(M)\ \hand{\hat{\ }}\!\otimes_k\ \hand{\Gamma}_k^{\hand{\wedge}}(M) \quad (\hand{\simeq}\ \hand{\Gamma}_k^{\hand{\wedge}}(MxM)\ )\hand{,} \]
\var{\inserted{(the latter,}}{\hand{(the latter,}}{insertion}{ps_134-7_009-v0017} \var{\deleted{(}}{}{deletion}{ps_134-7_009-v0018}deduced from the usual linear diagonal map $M\hand{\rightarrow} MxM$). When M is
looked \var{\inserted{at}}{\hand{at}}{insertion}{ps_134-7_009-v0019} as being embedded in $\hand{\Gamma}_k^{\hand{\wedge}}(M)$ by the exponenti\var{\xout{t}}{}{typewriter deletion}{ps_134-7_009-v0020}al map \var{\xout{(xx}}{}{typewriter deletion}{ps_134-7_009-v0021}\hand{(}n),
it is identified (if I remember it right) to the set of elements \var{\inserted{$\xi$}}{\hand{$\xi$}}{insertion}{ps_134-7_009-v0022} in
$\hand{\Gamma}_k^{\hand{\wedge}}(M)$ satisfying the relations
\[ \text{(q)} \qquad \hand{\varepsilon}(\hand{\xi}) = 1 \ ,\ \hand{\Delta}(\hand{\xi}) = \hand{\xi}\otimes\hand{\xi} \quad , \]
where \hand{$\varepsilon$} is the augmentation and \hand{$\Delta$} the diagonal map, hence p) induces
a map (in general not additive)

\margin{9}

% Visual collation: both supplied image halves inspected; see CHECK_DECISIONS_134_7_8.tsv.

%% END SOURCE ps_134-7_009
\ifnum\EditionMode=1
\psapparatusstart
\psentry{ps_134-7_009-v0001}{1}{typewriter deletion}{\xout{xx}}{}{The typewriter strikes out this material; it leaves the corrected running text.}
\psentry{ps_134-7_009-v0002}{2}{deletion}{$\deleted{\illegible{1}}$}{}{The author cancels this material by hand; it leaves the corrected running text.}
\psentry{ps_134-7_009-v0003}{3}{deletion}{$\deleted{)}$}{}{The author cancels this material by hand; it leaves the corrected running text.}
\psentry{ps_134-7_009-v0004}{4}{insertion}{$\inserted{\overset{\text{def}}{=}\prod_{i\geq 0}\Gamma^{i}(M)\deleted{)}}$}{$\hand{\overset{\text{def}}{=}\prod_{i\geq 0}\Gamma^{i}(M)}$}{The handwritten insertion is incorporated at its indicated position; insertion brackets leave the running text.}
\psentry{ps_134-7_009-v0005}{5}{typed insertion}{\inserted[typed]{k-module}}{k-module}{The typed insertion is incorporated at its indicated position; insertion brackets leave the running text.}
\psentry{ps_134-7_009-v0006}{6}{typewriter deletion}{\xout{object}}{}{The typewriter strikes out this material; it leaves the corrected running text.}
\psentry{ps_134-7_009-v0007}{7}{typewriter deletion}{\xout{anyxx}}{}{The typewriter strikes out this material; it leaves the corrected running text.}
\psentry{ps_134-7_009-v0008}{8}{typewriter deletion}{\xout{n}}{}{The typewriter strikes out this material; it leaves the corrected running text.}
\psentry{ps_134-7_009-v0009}{9}{typewriter deletion}{\xout{isomorphism}}{}{The typewriter strikes out this material; it leaves the corrected running text.}
\psentry{ps_134-7_009-v0010}{10}{typewriter deletion}{\xout{whi}}{}{The typewriter strikes out this material; it leaves the corrected running text.}
\psentry{ps_134-7_009-v0011}{11}{writing over}{\over{,}{.}}{\hand{.}}{The superposed handwritten reading replaces the underlying reading.}
\psentry{ps_134-7_009-v0012}{12}{writing over}{\over{f}{F}}{\hand{F}}{The superposed handwritten reading replaces the underlying reading.}
\psentry{ps_134-7_009-v0013}{13}{typewriter deletion}{\xout{definitionxx}}{}{The typewriter strikes out this material; it leaves the corrected running text.}
\psentry{ps_134-7_009-v0014}{14}{typewriter deletion}{\xout{\illegible{1}}}{}{The typewriter strikes out this material; it leaves the corrected running text.}
\psentry{ps_134-7_009-v0015}{15}{typewriter deletion}{\xout{\illegible{1}}}{}{The typewriter strikes out this material; it leaves the corrected running text.}
\psentry{ps_134-7_009-v0016}{16}{typewriter deletion}{\xout{maps}}{}{The typewriter strikes out this material; it leaves the corrected running text.}
\psentry{ps_134-7_009-v0017}{17}{insertion}{\inserted{(the latter,}}{\hand{(the latter,}}{The handwritten insertion is incorporated at its indicated position; insertion brackets leave the running text.}
\psentry{ps_134-7_009-v0018}{18}{deletion}{\deleted{(}}{}{The author cancels this material by hand; it leaves the corrected running text.}
\psentry{ps_134-7_009-v0019}{19}{insertion}{\inserted{at}}{\hand{at}}{The handwritten insertion is incorporated at its indicated position; insertion brackets leave the running text.}
\psentry{ps_134-7_009-v0020}{20}{typewriter deletion}{\xout{t}}{}{The typewriter strikes out this material; it leaves the corrected running text.}
\psentry{ps_134-7_009-v0021}{21}{typewriter deletion}{\xout{(xx}}{}{The typewriter strikes out this material; it leaves the corrected running text.}
\psentry{ps_134-7_009-v0022}{22}{insertion}{\inserted{$\xi$}}{\hand{$\xi$}}{The handwritten insertion is incorporated at its indicated position; insertion brackets leave the running text.}
\psapparatusend
\fi
\pssource{ps_134-7_010}{134-7/10}{https://grothendieck.umontpellier.fr/web/viewer.html?file=../134-7.pdf\#page=11}
\psnumber{ps_134-7_010-v0001}{1}
\psnumber{ps_134-7_010-v0002}{2}
\psnumber{ps_134-7_010-v0003}{3}
\psnumber{ps_134-7_010-v0004}{4}
\psnumber{ps_134-7_010-v0005}{5}
\psnumber{ps_134-7_010-v0006}{6}
\psnumber{ps_134-7_010-v0007}{7}
%% BEGIN SOURCE ps_134-7_010
% Pursuing Stacks (cotes 134-2 to 134-8): archive cote 134-7, site page 10 (PDF page 11), typescript page 452.
% typescript. Draft compared: ps6.tex:430-509.
% AI-transcribed by Claude Opus 5.5 (claude-opus-5-5, adaptive thinking, effort high), Message Batches API, batch ps3 (msgbatch_01XdFmuVVBAcgc7EFxdtCqsj), 2026-11-12; reviewed against both supplied scan halves; see CHECK_DECISIONS.tsv.
\scanpage{134-7/10}
\mspage{452}

\hand{(}p') \qquad $\hand{\Gamma}(f) \ : \ M \hand{\to} M'$

(corresponding to the action of f, viewed as a map $W(M)\hand{\to} W(M')$, on
\uline{sections} of $W(M)$) - and likewise after any ring extension $k\hand{\to} k'$ ,
defining a map
\[ \hand{\Gamma}(f)_{k'} \ : \ M\otimes_k k' \hand{\to} M'\otimes_k k' \]
from $W(M)$ to $W(M')$ - which is the description of the map of sheaves
$W(M)\hand{\longrightarrow} W(M')$ associated to a map \hand{(}p). We have thus a description, in
terms of linear algebra, of a category $U(k)$, and of a \var{\inserted{``sections''}}{\hand{``sections''}}{insertion}{ps_134-7_010-v0001} functor

\var{\xout{\unsure{hx}}}{}{typewriter deletion}{ps_134-7_010-v0002}(r) \qquad $\hand{\Gamma} \ : \ U(k) \hand{\longrightarrow} (Sets) \quad , \qquad X \hand{\mapsto} X(k)$ ,

which is essentially the functor \hand{(}h) above, viewed in a different light.
Now to our

\uline{Question}: \margin{\unsure{$\#$}}The category $U(k)$ and the functor r) being defined as above
in terms of linear algebra over the ground ring k , let $X_{\hand{*}} = (n\hand{\mapsto} X_n)$
be a semisimplicial object in $U(k)$, consider the corresponding semi-sim-%
plicial set $X_{\hand{*}}(k) = (n\hand{\mapsto} X_n(k))$, and the semi-sipplicial k-module
\margin{* One has to take $k=\underline{Z}$, and components $X_n$ projective.}
$L(X_{\hand{*}}) = (n\hand{\mapsto} L(X_n))$ , where for an object M of $U(k)$, $L(M)$ is defined as
\[ \text{(s)} \qquad L(M) = \hand{\Gamma}_k{}^{\hand{\wedge}}(M) \quad , \]
which depends functorially on M in $U(k)$. Then in the derived category
of (Ab), is there a canonical isomorphism between $L(X_{\hand{*}})$ (with ground
ring restricted from k to $\underline{Z}$ ) and the abelianization $\underline{Z}^{(X_{\hand{*}}(k))}$ of $X_{\hand{*}}(k)$?

We may have to throw in some extra assumption on $X_{\hand{*}}$ , at any rate
\[ \hand{(}\text{t)} \qquad X_o = e \var{\deleted{,}}{}{deletion}{ps_134-7_010-v0003} \hand{(final object of U(k)),} \]
\var{\deleted{(the final object of U(k)),}}{}{deletion}{ps_134-7_010-v0004} giving rise to \var{\xout{the}}{}{typewriter deletion}{ps_134-7_010-v0005} $L(X_o)\hand{\simeq} k$ . Also we may
have to restrict to $k=\underline{Z}$ , or otherwise correct the obvious drawback that
the two chain complexes do'nt have isomorphic $H_o$ (one is k I guess,
a k-module in any case, the other is $\underline{Z}$ ), by truncating accordingly
\margin{X}the two chain complexes ("killing" their $H_o$). The\var{\over{\unsure{re}}{re}}{\hand{re}}{writing over}{ps_134-7_010-v0006} is a natural can-%
didate \var{\over{\unsure{for}}{for}}{\hand{for}}{writing over}{ps_134-7_010-v0007} a map
\[ \hand{(}\text{u)} \qquad \underline{Z}^{(X_{\hand{*}}(k))} \hand{\longrightarrow} L(X_{\hand{*}}) \quad , \]
by using the functorial map
\[ \hand{(}\text{u')} \qquad \underline{Z}^{(M)} \hand{\longrightarrow} L(M) = \hand{\Gamma}_k{}^{\hand{\wedge}}(M) \quad , \]
deduced from the inclusion \hand{(n)}
\[ M \hand{\hookrightarrow} L(M) \quad , \]
and we may still specify the question above, by asking \uline{whether} \hand{(}u) \uline{induces}
\uline{an isomorphism for homology groups} in dimension $i>0$.

\margin{10}

% Visual collation: both supplied image halves inspected; see CHECK_DECISIONS_134_7_8.tsv.

%% END SOURCE ps_134-7_010
\ifnum\EditionMode=1
\psapparatusstart
\psentry{ps_134-7_010-v0001}{1}{insertion}{\inserted{``sections''}}{\hand{``sections''}}{The handwritten insertion is incorporated at its indicated position; insertion brackets leave the running text.}
\psentry{ps_134-7_010-v0002}{2}{typewriter deletion}{\xout{\unsure{hx}}}{}{The typewriter strikes out this material; it leaves the corrected running text.}
\psentry{ps_134-7_010-v0003}{3}{deletion}{$\deleted{,}$}{}{The author cancels this material by hand; it leaves the corrected running text.}
\psentry{ps_134-7_010-v0004}{4}{deletion}{\deleted{(the final object of U(k)),}}{}{The author cancels this material by hand; it leaves the corrected running text.}
\psentry{ps_134-7_010-v0005}{5}{typewriter deletion}{\xout{the}}{}{The typewriter strikes out this material; it leaves the corrected running text.}
\psentry{ps_134-7_010-v0006}{6}{writing over}{\over{\unsure{re}}{re}}{\hand{re}}{The superposed handwritten reading replaces the underlying reading.}
\psentry{ps_134-7_010-v0007}{7}{writing over}{\over{\unsure{for}}{for}}{\hand{for}}{The superposed handwritten reading replaces the underlying reading.}
\psapparatusend
\fi
\pssource{ps_134-7_011}{134-7/11}{https://grothendieck.umontpellier.fr/web/viewer.html?file=../134-7.pdf\#page=12}
\psnumber{ps_134-7_011-v0001}{1}
\psnumber{ps_134-7_011-v0002}{2}
\psnumber{ps_134-7_011-v0003}{3}
\psnumber{ps_134-7_011-v0004}{4}
\psnumber{ps_134-7_011-v0005}{5}
\psnumber{ps_134-7_011-v0006}{6}
\psnumber{ps_134-7_011-v0007}{7}
\psnumber{ps_134-7_011-v0008}{8}
\psnumber{ps_134-7_011-v0009}{9}
\psnumber{ps_134-7_011-v0010}{10}
\psnumber{ps_134-7_011-v0011}{11}
\psnumber{ps_134-7_011-v0012}{12}
%% BEGIN SOURCE ps_134-7_011
% Pursuing Stacks (cotes 134-2 to 134-8): archive cote 134-7, site page 11 (PDF page 12), typescript page 453.
% typescript. Draft compared: ps6.tex:491-560.
% AI-transcribed by Claude Opus 5.5 (claude-opus-5-5, adaptive thinking, effort high), Message Batches API, batch ps3 (msgbatch_01XdFmuVVBAcgc7EFxdtCqsj), 2026-11-12; reviewed against both supplied scan halves; see CHECK_DECISIONS.tsv.
\scanpage{134-7/11}
\mspage{453}

I am not too sure whether all this is'nt just complete nonsense -
it is worth getting it clear whether it is or not, at any rate ! There
is one case of special interest, the ``simplest'' one in a way, namely
when \var{\xout{Xx}}{}{typewriter deletion}{ps_134-7_011-v0001} the simplicial maps between the $X_n$'s , each represented by a
k-module $M_n$ , are in fact k-linear, in other words, when $X_{\hand{\ast}}$ comes
from a semisimplicial k-module $M_{\hand{\ast}}$ - more specifically still, when
this is a $K(\hand{\pi},n)$ type, say the nicest semisimplicial model of this,
using the Kan-Dold-Puppe functor for the chain complex of k-modules,
having $\hand{\pi}$ in degree n and zero elsewhere. Then the left hand side of u)
gives rise to the Eilenberg-Maclane homology groups
\[ \hand{(v)} \qquad H_i(\hand{\pi},n;\underline{Z}) \quad , \]
\margin{$\times$} which I guess should be k-modules a priori, because of the operations
of k upon $\hand{\pi}$ , and the question then arises whether these can be compu-%
ted using the right hand complex $L(X_{\hand{\ast}})$ . Maybe such a thing is even a
familiar fact for people in the know ? If it turned out to be false
even for $k=\underline{Z}$ , my faith in schematization of homotopy types would be
seriously shaken I confess ...

\hand{* * *}

After a little break for dinner, just one more afterthought. Wor-%
king with the completions $\hand{\Gamma}_k^{\hand{\wedge}}$ may seem a little forbidding, and all
the more so if it is used for computing \uline{homology} invariants, not cohomo-%
logy (the latter more likely to involve \var{\inserted[typed]{infinite}}{infinite}{typed insertion}{ps_134-7_011-v0002} products in \var{\xout{the}}{}{typewriter deletion}{ps_134-7_011-v0003} the correspon-%
ding (cochain) complexes ...). On the other hand, as was explicitly
stated from the beginning, the natural context here seems to be
\uline{pointed} homotopy types, and hence ``pointed'' algebraic para\var{\xout{\unsure{g}}}{}{typewriter deletion}{ps_134-7_011-v0004}digms for
these - an aspect which we lost sight of, when looking for a suitable
description of some category U(k) of ``unipotent bundles'' over k. It
would seem that in the ``question'' above, we should therefore insist
that $X_{\hand{\ast}}$ should be a semi-simplicial object of the category $U(k)^{\bullet}$ of
``pointed'' objects of U(k), namely objects \var{\inserted[typed]{X}}{X}{typed insertion}{ps_134-7_011-v0005} endowed with a section over
the final object e (i.e. with an \var{\deleted{object of}}{}{deletion}{ps_134-7_011-v0006} \var{\inserted{element in}}{\hand{element in}}{insertion}{ps_134-7_011-v0007} X(k)). This will be automa-%
tic at any rate in terms of the condition \hand{(}t\hand{)} , $X_o = e$ . The point I
wish to make is that the category of pointed objects of U(k) admits a
somewhat simpler description (by cho\var{\inserted{o}}{\hand{o}}{insertion}{ps_134-7_011-v0008}sing the marked point as the ``origin''
for \var{\deleted{a given}}{}{deletion}{ps_134-7_011-v0009} parametrisation of \hand{the} \var{\inserted{given}}{\hand{given}}{insertion}{ps_134-7_011-v0010} object X of U(k) by a k-module M),
by model-objects which are still arbitrary k-modules M, but the ``maps''
now being k-\uline{linear} \var{\deleted{maps}}{}{deletion}{ps_134-7_011-v0011} \var{\inserted{\uline{continuous maps}}}{\hand{\uline{continuous maps}}}{insertion}{ps_134-7_011-v0012}
\[ \hand{(}p') \qquad f : \hand{\Gamma}_k(M) \hand{\longrightarrow} \hand{\Gamma}_k(M') \]

\margin{11}

% Visual collation: both supplied image halves inspected; see CHECK_DECISIONS_134_7_8.tsv.

%% END SOURCE ps_134-7_011
\ifnum\EditionMode=1
\psapparatusstart
\psentry{ps_134-7_011-v0001}{1}{typewriter deletion}{\xout{Xx}}{}{The typewriter strikes out this material; it leaves the corrected running text.}
\psentry{ps_134-7_011-v0002}{2}{typed insertion}{\inserted[typed]{infinite}}{infinite}{The typed insertion is incorporated at its indicated position; insertion brackets leave the running text.}
\psentry{ps_134-7_011-v0003}{3}{typewriter deletion}{\xout{the}}{}{The typewriter strikes out this material; it leaves the corrected running text.}
\psentry{ps_134-7_011-v0004}{4}{typewriter deletion}{\xout{\unsure{g}}}{}{The typewriter strikes out this material; it leaves the corrected running text.}
\psentry{ps_134-7_011-v0005}{5}{typed insertion}{\inserted[typed]{X}}{X}{The typed insertion is incorporated at its indicated position; insertion brackets leave the running text.}
\psentry{ps_134-7_011-v0006}{6}{deletion}{\deleted{object of}}{}{The author cancels this material by hand; it leaves the corrected running text.}
\psentry{ps_134-7_011-v0007}{7}{insertion}{\inserted{element in}}{\hand{element in}}{The handwritten insertion is incorporated at its indicated position; insertion brackets leave the running text.}
\psentry{ps_134-7_011-v0008}{8}{insertion}{\inserted{o}}{\hand{o}}{The handwritten insertion is incorporated at its indicated position; insertion brackets leave the running text.}
\psentry{ps_134-7_011-v0009}{9}{deletion}{\deleted{a given}}{}{The author cancels this material by hand; it leaves the corrected running text.}
\psentry{ps_134-7_011-v0010}{10}{insertion}{\inserted{given}}{\hand{given}}{The handwritten insertion is incorporated at its indicated position; insertion brackets leave the running text.}
\psentry{ps_134-7_011-v0011}{11}{deletion}{\deleted{maps}}{}{The author cancels this material by hand; it leaves the corrected running text.}
\psentry{ps_134-7_011-v0012}{12}{insertion}{\inserted{\uline{continuous maps}}}{\hand{\uline{continuous maps}}}{The handwritten insertion is incorporated at its indicated position; insertion brackets leave the running text.}
\psapparatusend
\fi
\pssource{ps_134-7_012}{134-7/12}{https://grothendieck.umontpellier.fr/web/viewer.html?file=../134-7.pdf\#page=13}
\psnumber{ps_134-7_012-v0001}{1}
\psnumber{ps_134-7_012-v0002}{2}
\psnumber{ps_134-7_012-v0003}{3}
\psnumber{ps_134-7_012-v0004}{4}
\psnumber{ps_134-7_012-v0005}{5}
\psnumber{ps_134-7_012-v0006}{6}
\psnumber{ps_134-7_012-v0007}{7}
\psnumber{ps_134-7_012-v0008}{8}
\psnumber{ps_134-7_012-v0009}{9}
\psnumber{ps_134-7_012-v0010}{10}
%% BEGIN SOURCE ps_134-7_012
% Pursuing Stacks (cotes 134-2 to 134-8): archive cote 134-7, site page 12 (PDF page 13), typescript page 454.
% typescript. Draft compared: ps6.tex:544-612.
% AI-transcribed by Claude Opus 5.5 (claude-opus-5-5, adaptive thinking, effort high), Message Batches API, batch ps3 (msgbatch_01XdFmuVVBAcgc7EFxdtCqsj), 2026-11-12; reviewed against both supplied scan halves; see CHECK_DECISIONS.tsv.
\scanpage{134-7/12}
\mspage{454}

between the k-modules $\hand{\Gamma}_k$ ,  without having to pass to completions,
satisfying compatibility wi th augmentations and diagonal maps,  and
the extra condition  (expressing that $\hand{\Gamma}(f)(0) = 0$ \var{\xout{)}}{}{typewriter deletion}{ps_134-7_012-v0001} in exponential
notation) :
\[ f(1) = 1 \]
i.e. f reduced to component \var{\xout{of degree}}{}{typewriter deletion}{ps_134-7_012-v0002} $\hand{\Gamma}_k^{0}(M)\hand{\simeq} k$  of degree zero is just
the identity of k with  $k\hand{\simeq}\hand{\Gamma}_k^{0}(M')$ .  Accordingly,  we have a less awk-%
ward \var{\inserted[typed]{k-linearization}}{k-linearization}{typed insertion}{ps_134-7_012-v0003} functor than L in \hand{(}s) ,  namely ``pointed linearization'' $L_{pt}$ :
\[ \text{(s')} \qquad M\hand{\longmapsto} L_{pt}(M)\overset{def}{=}\hand{\Gamma}_k(M) \quad : \quad U(k)^{\bullet}\hand{\to} Ab_k \unsure{\hand{)}} \qquad , \]
\var{\xout{and}}{}{typewriter deletion}{ps_134-7_012-v0004} which seems to me the better candidate for describing linearization.
Thus we better rephrase now the ``question'' above in terms of \hand{(}s') rather
than \hand{(}s).  One trouble however is that the comparison map \hand{(}u)  takes \var{\xout{value}}{}{typewriter deletion}{ps_134-7_012-v0005}
values in $L(X_{\hand{*}})$ ,  not $L_{pt}(X_{\hand{*}})$,  therefore we may still have to use
the ``prohibitive'' $L(X_{\hand{*}})$  as an intermediary for comparing  the complexes
$\underline{Z}^{(X_{\hand{*}}(k))}$  and $L_{pt}(X_{\hand{*}})$.  It may be noted now that,  while the first chain
\margin{(v)} complex embodies  Eilenberg-Maclane homology \hand{(v)}  (in the special case
considered above\hand{)}, the second one  $L_{pt}(X_{\hand{*}})$  (in that same case) describes
the value of the total derived functor of the familiar $\hand{\Gamma}_k$ functor,  on
the ``argument'' $\hand{\pi}$ placed in degree n ,  and the statement that the two
\var{\inserted[typed]{are}}{are}{typed insertion}{ps_134-7_012-v0006}\var{\xout{should be}}{}{typewriter deletion}{ps_134-7_012-v0007}``the same''does sound like some standard Dold-Puppe type result
which everybody is supposed to know from the craddle  - sorry !

\begin{center}
* \qquad *\\
*
\end{center}

After another break (visit, tentative nap), still another afterthought\unsure{.}
The final shape we arrived at for the ``question'' above,  when working in
the ``pointed'' category  $U(k)^{\bullet}$  of ``pointed$^{*}$ unipotent bundles over k '' ,
was whether for any semisimplicial object $X_{*}$  in $U(k)^{\bullet}$ satisfying\hand{(}t) above
i.e.  $X_o = e$ (final object of $U(k)^{\bullet}$ ),  the two canonical maps of chain
complexes (in fact,  semisimplicial abelian groups)
\[ \text{(w)} \qquad \diagram{\begin{tikzcd} \underline{Z}^{(X_{\hand{*}}(k))} \arrow[color=blue,r, "\hand{}"] \& L\var{\deleted{\illegible{1}}}{}{deletion}{ps_134-7_012-v0008}(X_{*}) \& L_{pt}(X_{*}) \arrow[color=blue,l, hook', "\hand{}"'] \end{tikzcd}} \]
are quasi-isomorphisms\var{\inserted[typed]{(for H$_i$ with $i>0$).}}{(for H$_i$ with $i>0$).}{typed insertion}{ps_134-7_012-v0009}\unsure{\hand{$*$}} I am not too sure yet if some extra conditions
on $X_{*}$  are not required for this to be reasonable  -  I want to review
two \var{\over{\unsure{th}}{th}}{\hand{th}}{writing over}{ps_134-7_012-v0010}at came into my mind.

As the maps \hand{(}w) are functorial for varying $X_{\hand{*}}$ ,  it would follow

\margin{12}

% Visual collation: both supplied image halves inspected; see CHECK_DECISIONS_134_7_8.tsv.

%% END SOURCE ps_134-7_012
\ifnum\EditionMode=1
\psapparatusstart
\psentry{ps_134-7_012-v0001}{1}{typewriter deletion}{\xout{)}}{}{The typewriter strikes out this material; it leaves the corrected running text.}
\psentry{ps_134-7_012-v0002}{2}{typewriter deletion}{\xout{of degree}}{}{The typewriter strikes out this material; it leaves the corrected running text.}
\psentry{ps_134-7_012-v0003}{3}{typed insertion}{\inserted[typed]{k-linearization}}{k-linearization}{The typed insertion is incorporated at its indicated position; insertion brackets leave the running text.}
\psentry{ps_134-7_012-v0004}{4}{typewriter deletion}{\xout{and}}{}{The typewriter strikes out this material; it leaves the corrected running text.}
\psentry{ps_134-7_012-v0005}{5}{typewriter deletion}{\xout{value}}{}{The typewriter strikes out this material; it leaves the corrected running text.}
\psentry{ps_134-7_012-v0006}{6}{typed insertion}{\inserted[typed]{are}}{are}{The typed insertion is incorporated at its indicated position; insertion brackets leave the running text.}
\psentry{ps_134-7_012-v0007}{7}{typewriter deletion}{\xout{should be}}{}{The typewriter strikes out this material; it leaves the corrected running text.}
\psentry{ps_134-7_012-v0008}{8}{deletion}{$\deleted{\illegible{1}}$}{}{The author cancels this material by hand; it leaves the corrected running text.}
\psentry{ps_134-7_012-v0009}{9}{typed insertion}{\inserted[typed]{(for H$_i$ with $i>0$).}}{(for H$_i$ with $i>0$).}{The typed insertion is incorporated at its indicated position; insertion brackets leave the running text.}
\psentry{ps_134-7_012-v0010}{10}{writing over}{\over{\unsure{th}}{th}}{\hand{th}}{The superposed handwritten reading replaces the underlying reading.}
\psapparatusend
\fi
\pssource{ps_134-7_013}{134-7/13}{https://grothendieck.umontpellier.fr/web/viewer.html?file=../134-7.pdf\#page=14}
\psnumber{ps_134-7_013-v0001}{1}
\psnumber{ps_134-7_013-v0002}{2}
\psnumber{ps_134-7_013-v0003}{3}
\psnumber{ps_134-7_013-v0004}{4}
\psnumber{ps_134-7_013-v0005}{5}
\psnumber{ps_134-7_013-v0006}{6}
\psnumber{ps_134-7_013-v0007}{7}
%% BEGIN SOURCE ps_134-7_013
% Pursuing Stacks (cotes 134-2 to 134-8): archive cote 134-7, site page 13 (PDF page 14), typescript page 455.
% typescript. Draft compared: ps6.tex:590-663.
% AI-transcribed by Claude Opus 5.5 (claude-opus-5-5, adaptive thinking, effort high), Message Batches API, batch ps3 (msgbatch_01XdFmuVVBAcgc7EFxdtCqsj), 2026-11-12; reviewed against both supplied scan halves; see CHECK_DECISIONS.tsv.
\scanpage{134-7/13}
\mspage{455}

from a positive answer that whenever
\[ \hand{(x)} \qquad X_* \hand{\longrightarrow} X'_* \]
is a map of semisimplicial objects in $U(k)^\bullet$ satisfying condition \hand{(t)},
and such that the corresponding map
\[ \hand{(x')} \qquad X_*(k) \hand{\to} X'_*(k) \]
is a\var{\xout{n}}{}{typewriter deletion}{ps_134-7_013-v0001} weak equivalence, and \var{\xout{also}}{}{typewriter deletion}{ps_134-7_013-v0002} hence the map between the $\underline{Z}$-abelianiza-%
tions is a weak equivalence too i.e. a quasi-isomorphism, that the same
holds for the corresponding map
\[ \hand{(y)} \qquad L_{pt}(X_*) \hand{\to} L_{pt}(X'_*) \]
for the "schematic" k-linearizations. Now this is far from being an
evident fact by itself, except of course in the case when the map \hand{(x)}
above is a homotopism. Take for instance the case when we start with
a map of chain complexes in $Ab_k$
\[ \hand{(x'')} \qquad M. \hand{\to} \text{\var{\xout{\illegible{1}}}{}{typewriter deletion}{ps_134-7_013-v0003}}\ M.' \quad , \]
hence a map between the associated semisimplicial k-modules
\[ \hand{(x''')} \qquad M_* \hand{\longrightarrow} M'_* \quad , \quad \unsure{.} \]
which may be viewed as giving rise to a (componentwise linear) map bet-%
ween the associated semisimplicial objects $X_*$ , $X'_*$ in $U(k)^\bullet$ via the
canonical functor
\[ Ab_k \hand{\longrightarrow} U(k)^\bullet \ ; \]
the corresponding map \hand{(y)} is then just the componentwise extension of
\hand{(x''')} to the enveloping algebras with devided powers
\[ \hand{(y')} \qquad \hand{\Gamma}_k(M_*) \hand{\longrightarrow} \hand{\Gamma}_k(M'_*) \quad . \]
The map \hand{(x')} can now be identified with \hand{(x''')} , hence it is a weak equiva-%
lence iff \hand{(x''')} is, i.e. iff \hand{(x'')} is a quasi-isomorphism. If we assume
moreover the components of $M.$ , \var{\xout{\illegible{1}}}{}{typewriter deletion}{ps_134-7_013-v0004} $M.'$ to be projective objects in $Ab_k$,
then from the assumption that \hand{(x'')} is a quasi-isomorphism it does follow
that it is a chain homotopism, hence by Kan-Dold-Puppe the map \hand{(x''')}
is a semisimplicial homotopism, and hence the same holds for \hand{(y')}\hand{,} and
\hand{(y')} therefore is indeed a quasi-equivalence. But without the assumpti-%
on that components are projective, it is surely false that the mere
fact that \hand{(x'')} is a quasi-isomorphism, implies that \hand{(y')} is - otherwise
this wou\var{\over{\unsure{d}}{l}}{\hand{l}}{writing over}{ps_134-7_013-v0005}d mean that in order to compute the \var{\inserted{left}}{\hand{left}}{insertion}{ps_134-7_013-v0006} derived functor of the
non-additive functor \var{\xout{$\hand{\Gamma}_k$\illegible{1}}}{}{typewriter deletion}{ps_134-7_013-v0007}

\margin{13}

% Visual collation: both supplied image halves inspected; see CHECK_DECISIONS_134_7_8.tsv.

%% END SOURCE ps_134-7_013
\ifnum\EditionMode=1
\psapparatusstart
\psentry{ps_134-7_013-v0001}{1}{typewriter deletion}{\xout{n}}{}{The typewriter strikes out this material; it leaves the corrected running text.}
\psentry{ps_134-7_013-v0002}{2}{typewriter deletion}{\xout{also}}{}{The typewriter strikes out this material; it leaves the corrected running text.}
\psentry{ps_134-7_013-v0003}{3}{typewriter deletion}{\xout{\illegible{1}}}{}{The typewriter strikes out this material; it leaves the corrected running text.}
\psentry{ps_134-7_013-v0004}{4}{typewriter deletion}{\xout{\illegible{1}}}{}{The typewriter strikes out this material; it leaves the corrected running text.}
\psentry{ps_134-7_013-v0005}{5}{writing over}{\over{\unsure{d}}{l}}{\hand{l}}{The superposed handwritten reading replaces the underlying reading.}
\psentry{ps_134-7_013-v0006}{6}{insertion}{\inserted{left}}{\hand{left}}{The handwritten insertion is incorporated at its indicated position; insertion brackets leave the running text.}
\psentry{ps_134-7_013-v0007}{7}{typewriter deletion}{\xout{$\hand{\Gamma}_k$\illegible{1}}}{}{The typewriter strikes out this material; it leaves the corrected running text.}
\psapparatusend
\fi
\pssource{ps_134-7_014}{134-7/14}{https://grothendieck.umontpellier.fr/web/viewer.html?file=../134-7.pdf\#page=15}
\psnumber{ps_134-7_014-v0001}{1}
\psnumber{ps_134-7_014-v0002}{2}
\psnumber{ps_134-7_014-v0003}{3}
\psnumber{ps_134-7_014-v0004}{4}
\psnumber{ps_134-7_014-v0005}{5}
\psnumber{ps_134-7_014-v0006}{6}
\psnumber{ps_134-7_014-v0007}{7}
\psnumber{ps_134-7_014-v0008}{8}
\psnumber{ps_134-7_014-v0009}{9}
\psnumber{ps_134-7_014-v0010}{10}
\psnumber{ps_134-7_014-v0011}{11}
\psnumber{ps_134-7_014-v0012}{12}
%% BEGIN SOURCE ps_134-7_014
% Pursuing Stacks (cotes 134-2 to 134-8): archive cote 134-7, site page 14 (PDF page 15), typescript page 456.
% typescript. Draft compared: ps6.tex:650-704.
% AI-transcribed by Claude Opus 5.5 (claude-opus-5-5, adaptive thinking, effort high), Message Batches API, batch ps3 (msgbatch_01XdFmuVVBAcgc7EFxdtCqsj), 2026-11-12; reviewed against both supplied scan halves; see CHECK_DECISIONS.tsv.
\scanpage{134-7/14}
\mspage{456}

\[ \hand{\Gamma}_k : Ab_k \hand{\to} Ab_k \qquad , \]
it is enough\var{\deleted{\illegible{1}}}{}{deletion}{ps_134-7_014-v0001}\hand{,} \hand{f}or get\var{\inserted{ting}}{\hand{ting}}{insertion}{ps_134-7_014-v0002} it's value on a chain complex \var{\xout{MxxxMxxxx}}{}{typewriter deletion}{ps_134-7_014-v0003} $M.'$ say ,
to replace \var{\xout{Mx}}{}{typewriter deletion}{ps_134-7_014-v0004} $M.'$ by $M_*'$ and apply the functor $\hand{\Gamma}_k$ componentwise,
without first having to take a projective resolution $M.$ of $M.'$ - some-%
thing rather absurd indeed ! Thus the statement made on page 454, that
when taking for $M.'$ the k-module \hand{$\pi$} placed in degree n and zero in all
other degrees, the corresponding $L_{pt}(X_*') = I_k(M_*')$ embodies the
value of the left derived functor $L\hand{\Gamma}_k$ on $M.'$ , is visibly incorrect if
we do'nt assume moreover that \hand{$\pi$} is projective (flat presumably would be
enough ...). Otherwise, we should first replace \hand{$\pi$} by a projective
(or at any rate flat) resolution\hand{,} which we shift by n to get $M.$ , and \var{\inserted{then}}{\hand{then}}{insertion}{ps_134-7_014-v0005}
take $\hand{\Gamma}_k(M_*)$ to get the correct value of $L\hand{\Gamma}_k(M.')$,

This convinces me that in the question as to whether the maps in (w)
are quasi-isomorphisms (the more crucial one of course being the first
of the two), \uline{we should assume moreover that the components of $X_*$ are}
\uline{described in terms of }\uuline{projective}\uline{ k-modules} $M_n$ , or at any rate k-modules
which are flat. Accordingly, we should make the same restriction on
the semisimplicial schematized model $X_*$ , in order for the description
we gave of ``k-linearization'' as $L_{pt}(X_*)$ (or $L(X_*)$, never mind which)
to be topologically meaningful. Very probably, in the whole schematiza-%
tion set-up, namely in the very definition we gave of $U(k)$ and $U(k)^\bullet$ ,
we should stick to the same restriction. If \var{\xout{atxthexstart}}{}{typewriter deletion}{ps_134-7_014-v0006} I insisted
first (with some inner reluctance, I admit) on taking k-modules M
unrestricted, this was because I was thinking of M, more specifically
of the $M_i$'s occuring in the inductive ``dévissage'' of an object of
$U_m(k)$ (when thinking of the more refined version\var{\deleted{s}}{}{deletion}{ps_134-7_014-v0007} of $U(k)$), as essen-%
tially the homotopy groups of the homotopy type we want to modelize,
or \var{\inserted{rather, as}}{\hand{rather, as}}{insertion}{ps_134-7_014-v0008} the components of the corresponding semisimplicial k-modules (denoted
$M_*'$ some minutes ago). I was still thinking of course,\var{\deleted{\unsure{t}}}{}{deletion}{ps_134-7_014-v0009} be it implicit-%
ly, in terms of Postnikow dévissage, despite yesterday's remark that
to use such dévissage litterally may cause trouble (p.444). Thus the
feeling which gets into the fore now is that \uline{we should kind of forget}
\uline{Postnikov}, \uline{and} \var{\xout{deal}}{}{typewriter deletion}{ps_134-7_014-v0010} \var{\inserted[typed]{work}}{work}{typed insertion}{ps_134-7_014-v0011} \uline{with} \uline{\unsure{semisimplicial|Semisimplicial}} \uline{``schematic''} \uline{models} \uline{buil\var{\over{d}{t}}{\hand{t}}{writing over}{ps_134-7_014-v0012} up}
\uline{with k-modules which are} \uuline{projective}, \uline{or at} any rate flat (namely torsi-%
on free, if $k=\underline{Z}$ ) .

\margin{14}

% Visual collation: both supplied image halves inspected; see CHECK_DECISIONS_134_7_8.tsv.

%% END SOURCE ps_134-7_014
\ifnum\EditionMode=1
\psapparatusstart
\psentry{ps_134-7_014-v0001}{1}{deletion}{\deleted{\illegible{1}}}{}{The author cancels this material by hand; it leaves the corrected running text.}
\psentry{ps_134-7_014-v0002}{2}{insertion}{\inserted{ting}}{\hand{ting}}{The handwritten insertion is incorporated at its indicated position; insertion brackets leave the running text.}
\psentry{ps_134-7_014-v0003}{3}{typewriter deletion}{\xout{MxxxMxxxx}}{}{The typewriter strikes out this material; it leaves the corrected running text.}
\psentry{ps_134-7_014-v0004}{4}{typewriter deletion}{\xout{Mx}}{}{The typewriter strikes out this material; it leaves the corrected running text.}
\psentry{ps_134-7_014-v0005}{5}{insertion}{\inserted{then}}{\hand{then}}{The handwritten insertion is incorporated at its indicated position; insertion brackets leave the running text.}
\psentry{ps_134-7_014-v0006}{6}{typewriter deletion}{\xout{atxthexstart}}{}{The typewriter strikes out this material; it leaves the corrected running text.}
\psentry{ps_134-7_014-v0007}{7}{deletion}{\deleted{s}}{}{The author cancels this material by hand; it leaves the corrected running text.}
\psentry{ps_134-7_014-v0008}{8}{insertion}{\inserted{rather, as}}{\hand{rather, as}}{The handwritten insertion is incorporated at its indicated position; insertion brackets leave the running text.}
\psentry{ps_134-7_014-v0009}{9}{deletion}{\deleted{\unsure{t}}}{}{The author cancels this material by hand; it leaves the corrected running text.}
\psentry{ps_134-7_014-v0010}{10}{typewriter deletion}{\xout{deal}}{}{The typewriter strikes out this material; it leaves the corrected running text.}
\psentry{ps_134-7_014-v0011}{11}{typed insertion}{\inserted[typed]{work}}{work}{The typed insertion is incorporated at its indicated position; insertion brackets leave the running text.}
\psentry{ps_134-7_014-v0012}{12}{writing over}{\over{d}{t}}{\hand{t}}{The superposed handwritten reading replaces the underlying reading.}
\psapparatusend
\fi
\pssource{ps_134-7_015}{134-7/15}{https://grothendieck.umontpellier.fr/web/viewer.html?file=../134-7.pdf\#page=16}
\psnumber{ps_134-7_015-v0001}{1}
\psnumber{ps_134-7_015-v0002}{2}
\psnumber{ps_134-7_015-v0003}{3}
\psnumber{ps_134-7_015-v0004}{4}
\psnumber{ps_134-7_015-v0005}{5}
\psnumber{ps_134-7_015-v0006}{6}
\psnumber{ps_134-7_015-v0007}{7}
%% BEGIN SOURCE ps_134-7_015
% Pursuing Stacks (cotes 134-2 to 134-8): archive cote 134-7, site page 15 (PDF page 16), typescript page 457.
% typescript. Draft compared: ps6.tex:691-749.
% AI-transcribed by Claude Opus 5.5 (claude-opus-5-5, adaptive thinking, effort high), Message Batches API, batch ps3 (msgbatch_01XdFmuVVBAcgc7EFxdtCqsj), 2026-11-12; reviewed against both supplied scan halves; see CHECK_DECISIONS.tsv.
\scanpage{134-7/15}
\mspage{457}

It may be remarked that if M is any k-module,  than the property
that M be projective,  or flat,  can be described in terms of the isomor-%
phism class of the corresponding object X in U(k)\hand{$^\bullet$},  or equivalently, of
the functor W(M) on $\mathrm{Alg}_{/k}$ ,  with values in the category of pointed sets.
Indeed, the isomorphism class of the k-module $\hand{\Gamma}_k(M)$  depends only on the
class of X,  and it is easily seen that M is projective, resp. flat, iff
$\hand{\Gamma}_k(M)$ is.  (The ``only if'' is standard knowledge of commutative algebra,
the ``if'' comes from the fact that M is a direct factor of $\hand{\Gamma}_k(M)$,  hence
projective resp. flat if the latter is.)  Using this,  one even checks
that it is enough to know the isomorphism class of X in U(k)  -  because
two objects of $U(k)^\bullet$  are isomorphic in $U(k)^\bullet$  iff they are in U(k).

* \quad *

*

The second afterthought is  that in the question on page 452,  we
should definitely assume $k=\underline{Z}$ .  Already when asserting hastily (p.453)
that an Eilenberg-Maclane homology group $H_i(\hand{\pi},n;\ \underline{Z})$  should automatically
inherit a structure of a k-module,  whenever \hand{$\pi$} had one,  I was feeling
unconfortable,  because after all the dependance of this group upon
variable \hand{$\pi$} is not at all additive,  thus the operation of k upon this
group stemming from it's operations on \hand{$\pi$} \var{\xout{did}}{}{typewriter deletion}{ps_134-7_015-v0001} \unsure{were|was} not too likely to come
out additive !  To take one example,  take n=1 i.e. we just take ordina-%
ry group homology for \hand{$\pi$} ,  and assume that k is free of finite type over
$\underline{Z}$ ,  and \hand{$\pi$}\var{\inserted{$=M$}}{\hand{$=M$}}{insertion}{ps_134-7_015-v0002} free of finite type over k (for instance \hand{$M$}= k),  hence free
of finite type over $\underline{Z}$ too.  Then it is well known that
\[ H_*(M;\ \underline{Z})\hand{\simeq \bigwedge}_{\underline{Z}}M \qquad , \]
the exterior algebra of M  over $\underline{Z}$ ,  which surely is not endowed with a
structure of a k-module in any natural way !  Thus,  if there was any
such structure (natural or not) on the highest non zero term (correspon-%
ding to the \var{\inserted[typed]{rank d}}{rank d}{typed insertion}{ps_134-7_015-v0003}\var{\xout{rak}}{}{typewriter deletion}{ps_134-7_015-v0004} of M over $\underline{Z}$ )\hand{,} it would follow that we get a ring homomor-%
phism  from  k to  $\underline{Z}\hand{\simeq} \mathrm{End}\hand{_{\underline{Z}}}(H_d\hand{\simeq} \underline{Z})$ ,  and \var{\inserted[typed]{we}}{we}{typed insertion}{ps_134-7_015-v0005}\var{\xout{there}}{}{typewriter deletion}{ps_134-7_015-v0006} may choose k in such a
way that there is no such homomorphism,  for instance  $k = \underline{Z}[T]/(T^2+1)$ .
In this case,  there cannot be \uline{any} isomorphism  between the $H_d$'s of the
two members of \hand{(}u)\var{\over{.}{!}}{\hand{!}}{writing over}{ps_134-7_015-v0007}  Another point,  which I hit upon first,  is that
when  $X_*$  is defined in terms of a semisimplicial object $M_*$ of $\mathrm{Ab}_k$ ,
then the functorial dependance of the first term in\hand{(w}) with respect to

\margin{15}

% Visual collation: both supplied image halves inspected; see CHECK_DECISIONS_134_7_8.tsv.

%% END SOURCE ps_134-7_015
\ifnum\EditionMode=1
\psapparatusstart
\psentry{ps_134-7_015-v0001}{1}{typewriter deletion}{\xout{did}}{}{The typewriter strikes out this material; it leaves the corrected running text.}
\psentry{ps_134-7_015-v0002}{2}{insertion}{\inserted{$=M$}}{\hand{$=M$}}{The handwritten insertion is incorporated at its indicated position; insertion brackets leave the running text.}
\psentry{ps_134-7_015-v0003}{3}{typed insertion}{\inserted[typed]{rank d}}{rank d}{The typed insertion is incorporated at its indicated position; insertion brackets leave the running text.}
\psentry{ps_134-7_015-v0004}{4}{typewriter deletion}{\xout{rak}}{}{The typewriter strikes out this material; it leaves the corrected running text.}
\psentry{ps_134-7_015-v0005}{5}{typed insertion}{\inserted[typed]{we}}{we}{The typed insertion is incorporated at its indicated position; insertion brackets leave the running text.}
\psentry{ps_134-7_015-v0006}{6}{typewriter deletion}{\xout{there}}{}{The typewriter strikes out this material; it leaves the corrected running text.}
\psentry{ps_134-7_015-v0007}{7}{writing over}{\over{.}{!}}{\hand{!}}{The superposed handwritten reading replaces the underlying reading.}
\psapparatusend
\fi
\pssource{ps_134-7_016}{134-7/16}{https://grothendieck.umontpellier.fr/web/viewer.html?file=../134-7.pdf\#page=17}
\psnumber{ps_134-7_016-v0001}{1}
\psnumber{ps_134-7_016-v0002}{2}
\psnumber{ps_134-7_016-v0003}{3}
\psnumber{ps_134-7_016-v0004}{4}
\psnumber{ps_134-7_016-v0005}{5}
\psnumber{ps_134-7_016-v0006}{6}
\psnumber{ps_134-7_016-v0007}{7}
\psnumber{ps_134-7_016-v0008}{8}
\psnumber{ps_134-7_016-v0009}{9}
\psnumber{ps_134-7_016-v0010}{10}
\psnumber{ps_134-7_016-v0011}{11}
\psnumber{ps_134-7_016-v0012}{12}
\psnumber{ps_134-7_016-v0013}{13}
\psnumber{ps_134-7_016-v0014}{14}
%% BEGIN SOURCE ps_134-7_016
% Pursuing Stacks (cotes 134-2 to 134-8): archive cote 134-7, site page 16 (PDF page 17), typescript page 458.
% typescript. Draft compared: ps6.tex:735-796.
% AI-transcribed by Claude Opus 5.5 (claude-opus-5-5, adaptive thinking, effort high), Message Batches API, batch ps3 (msgbatch_01XdFmuVVBAcgc7EFxdtCqsj), 2026-11-12; reviewed against both supplied scan halves; see CHECK_DECISIONS.tsv.
\scanpage{134-7/16}
\mspage{458}

varying $M_{*}$ , is that to a direct sum corresponds the componentwise
tensor product \uline{over} $\underline{Z}$ , whereas for the last term of(w) we have to take
tensor products \uline{over} k (in the middle term, \uline{completed} tensor products
over k) - the two variances are clearly at odds which each other.

Thus, when working with would-be ``k-homotopy types'' as defined
here via semisimp\unsure{l|h}icial objects in $U(k)^{\bullet}$ , \uline{we should altogether drop the}
\uline{idea that the homology groups of the corresponding semisimplicial set}
$X_{*}(k)$ \uline{are k-modules}. I would'nt really look at them as being ``the''
homology groups of the k-homotopy type $X_{*}$ , these should be rather given
via k-linearization \var{\xout{R}}{}{typewriter deletion}{ps_134-7_016-v0001} $L_{pt}$ and they are k-modules, indeed they come
from \var{\over{an}{a}}{\hand{a}}{writing over}{ps_134-7_016-v0002} \var{\inserted{canonical}}{\hand{canonical}}{insertion}{ps_134-7_016-v0003} object of $D.(Ab_{k})$ , namely $L_{pt}(X_{*})$ . In the example just
looked at, presumably we should get the exterior algebra of M \uline{over} k
(not $\underline{Z}$ !) . This makes me suspect even that, except in the case $k=\underline{Z}$ ,
this semisimplicial set $X_{*}(k)$ does'nt make much sense, namely it\unsure{\var{\deleted{'}}{}{deletion}{ps_134-7_016-v0004}}s
homology invariants (and presumably it's homotopy groups too) are not
really relevant for the k-homotopy type $X_{*}$ , \uline{which has invariants of it's}
\uline{own which are completely different}. \uline{Thus I am not} \var{\xout{all}}{}{typewriter deletion}{ps_134-7_016-v0005} \uline{at all convinced}
\margin{X} \uline{any more that the homotopy groups} $\hand{\pi}_{i}(X_{*})$ \uline{carry k-module structures} (as
expected at the beginning, p. 446) - but to clear our mind on that
matter, we should take off from the simplistic example when $X_{*}$ comes
from a semisimplicial k-module $M_{*}$ , in which case of course we have by
Whitehead's isomorphism
\[ \hand{\pi}_{i}(X_{*}(k)) \hand{\simeq} H_{i}(M_{\unsure{*}}.) \quad ( \hand{\simeq} \hand{\pi}_{i}(M_{*})) \quad , \]
where M. is the chain complex associated to $M_{*}$ . Also, \var{\deleted{I am not com-}}{}{deletion}{ps_134-7_016-v0006}\var{\inserted{it is now}}{\hand{it is now}}{insertion}{ps_134-7_016-v0007}
\var{\deleted{pletely clear still}}{}{deletion}{ps_134-7_016-v0008} \var{\inserted{becoming dubious}}{\hand{becoming dubious}}{insertion}{ps_134-7_016-v0009} whether weak equivalences \var{\xout{\unsure{in}}}{}{typewriter deletion}{ps_134-7_016-v0010} for objects $X_{*}$ should
be defined (as we did) via $X_{*}(k)$, as suggested by the \var{\xout{example}}{}{typewriter deletion}{ps_134-7_016-v0011} ``simplis-%
tic'' example above. \var{\inserted{If we}}{\hand{If we}}{insertion}{ps_134-7_016-v0012} \var{\over{D}{d}}{\hand{d}}{writing over}{ps_134-7_016-v0013}efin\unsure{\var{\over{ing}{e}}{\hand{e}}{writing over}{ps_134-7_016-v0014}} homology of $X_{*}$ by the formula
\[ \text{(z)} \qquad LH.(X_{*}) \overset{\text{def}}{=} L_{pt}(X_{*}) \text{ viewed as an object of } D.(Ab_{k}) \ , \]
and, accordingly the homology modules
\[ \text{(z')} \qquad H_{i}(X_{*}) \overset{\text{def}}{=} \hand{\pi}_{i}(L_{pt}(X_{*})) = H_{i}(LH.(X_{*})) \quad , \]
maybe the better idea for defining weak equivalences $X_{*} \hand{\longrightarrow} X'_{*}$ is
by demanding that \uline{they should be transformed into quasi-isomorphisms by}
\uline{the total homology functor} LH. , or equivalently, induce isomorphisms
for the homology modules (z'). If the answer to our crucial question
is affirmative (with the corrections made, including $k = \underline{Z}$ ), then

\margin{16}

% Visual collation: both supplied image halves inspected; see CHECK_DECISIONS_134_7_8.tsv.

%% END SOURCE ps_134-7_016
\ifnum\EditionMode=1
\psapparatusstart
\psentry{ps_134-7_016-v0001}{1}{typewriter deletion}{\xout{R}}{}{The typewriter strikes out this material; it leaves the corrected running text.}
\psentry{ps_134-7_016-v0002}{2}{writing over}{\over{an}{a}}{\hand{a}}{The superposed handwritten reading replaces the underlying reading.}
\psentry{ps_134-7_016-v0003}{3}{insertion}{\inserted{canonical}}{\hand{canonical}}{The handwritten insertion is incorporated at its indicated position; insertion brackets leave the running text.}
\psentry{ps_134-7_016-v0004}{4}{deletion}{\deleted{'}}{}{The author cancels this material by hand; it leaves the corrected running text.}
\psentry{ps_134-7_016-v0005}{5}{typewriter deletion}{\xout{all}}{}{The typewriter strikes out this material; it leaves the corrected running text.}
\psentry{ps_134-7_016-v0006}{6}{deletion}{\deleted{I am not com-}}{}{The author cancels this material by hand; it leaves the corrected running text.}
\psentry{ps_134-7_016-v0007}{7}{insertion}{\inserted{it is now}}{\hand{it is now}}{The handwritten insertion is incorporated at its indicated position; insertion brackets leave the running text.}
\psentry{ps_134-7_016-v0008}{8}{deletion}{\deleted{pletely clear still}}{}{The author cancels this material by hand; it leaves the corrected running text.}
\psentry{ps_134-7_016-v0009}{9}{insertion}{\inserted{becoming dubious}}{\hand{becoming dubious}}{The handwritten insertion is incorporated at its indicated position; insertion brackets leave the running text.}
\psentry{ps_134-7_016-v0010}{10}{typewriter deletion}{\xout{\unsure{in}}}{}{The typewriter strikes out this material; it leaves the corrected running text.}
\psentry{ps_134-7_016-v0011}{11}{typewriter deletion}{\xout{example}}{}{The typewriter strikes out this material; it leaves the corrected running text.}
\psentry{ps_134-7_016-v0012}{12}{insertion}{\inserted{If we}}{\hand{If we}}{The handwritten insertion is incorporated at its indicated position; insertion brackets leave the running text.}
\psentry{ps_134-7_016-v0013}{13}{writing over}{\over{D}{d}}{\hand{d}}{The superposed handwritten reading replaces the underlying reading.}
\psentry{ps_134-7_016-v0014}{14}{writing over}{\over{ing}{e}}{\hand{e}}{The superposed handwritten reading replaces the underlying reading.}
\psapparatusend
\fi
\pssource{ps_134-7_017}{134-7/17}{https://grothendieck.umontpellier.fr/web/viewer.html?file=../134-7.pdf\#page=18}
\psnumber{ps_134-7_017-v0001}{1}
\psnumber{ps_134-7_017-v0002}{2}
\psnumber{ps_134-7_017-v0003}{3}
\psnumber{ps_134-7_017-v0004}{4}
\psnumber{ps_134-7_017-v0005}{5}
\psnumber{ps_134-7_017-v0006}{6}
\psnumber{ps_134-7_017-v0007}{7}
\psnumber{ps_134-7_017-v0008}{8}
\psnumber{ps_134-7_017-v0009}{9}
\psnumber{ps_134-7_017-v0010}{10}
\psnumber{ps_134-7_017-v0011}{11}
\psnumber{ps_134-7_017-v0012}{12}
\psnumber{ps_134-7_017-v0013}{13}
\psnumber{ps_134-7_017-v0014}{14}
\psnumber{ps_134-7_017-v0015}{15}
%% BEGIN SOURCE ps_134-7_017
% Pursuing Stacks (cotes 134-2 to 134-8): archive cote 134-7, site page 17 (PDF page 18), typescript page 459.
% typescript. Draft compared: ps6.tex:783-853.
% AI-transcribed by Claude Opus 5.5 (claude-opus-5-5, adaptive thinking, effort high), Message Batches API, batch ps3 (msgbatch_01XdFmuVVBAcgc7EFxdtCqsj), 2026-11-12; reviewed against both supplied scan halves; see CHECK_DECISIONS.tsv.
\scanpage{134-7/17}
\mspage{459}

in case $k=\underline{Z}$ ,  the new definition just given for weak equivalences is
equivalent to the old one in terms of $X_*(\underline{Z})$  (taking $\underline{Z}$-valued points),
provided at any rate we admit \var{\xout{that}}{}{typewriter deletion}{ps_134-7_017-v0001} or rather assume that $X_*(\underline{Z})$ \var{\xout{is simply}}{}{typewriter deletion}{ps_134-7_017-v0002}
\var{\xout{connected}}{}{typewriter deletion}{ps_134-7_017-v0003} and $X'_*(\underline{Z})$  are simply connected, which will be automatic if
we work in the category of (schematized) models $\hand{\mathcal{M}}_1(k)$,  \var{\deleted{with}}{}{deletion}{ps_134-7_017-v0004} the condi-%
tion \hand{(}t) above  ($X_0 = e$)  being replaced by \hand{(i)}  with n=1  i.e. \var{\inserted{by}}{\hand{by}}{insertion}{ps_134-7_017-v0005}
\margin{($\alpha$)}
\[ X_0 = X_1 = e \qquad , \]
to be on the safe side!  Under this extra assumption at any rate,  I feel
definitely more confident with the \var{\xout{second}}{}{typewriter deletion}{ps_134-7_017-v0006}\var{\inserted[typed]{new}}{new}{typed insertion}{ps_134-7_017-v0007} definition of weak equivalence,
via the homology invariants \hand{(}z'), \var{\inserted{rather than the old one.}}{\hand{rather than the old one.}}{insertion}{ps_134-7_017-v0008} At any rate the question of the two
definitions being equivalent or not should be cleared up,  namely:

\uline{Question} 2.   Let k be any ring,  define the category $U(k)^{\hand{\bullet}}$ of ``pointed
unipotent bundles over k\hand{''}in terms of projective k-modules ,  with maps
defined as in \hand{(}p')  page 453. (This is equally the correct set-up for
question 1 on page 452,  besides the extra condition $k=\underline{Z}$ , as we saw
before.)  Let
\[ u\colon X_* \hand{\longrightarrow} X'_* \]
be a map of semisimplicial objects in $U(k)$,  satisfying \var{\inserted[typed]{both}}{both}{typed insertion}{ps_134-7_017-v0009} the extra assump-%
tions \hand{(}$\alpha$) above.   Then is it true  that the corresponding map \var{\xout{$X_*(k)$}}{}{typewriter deletion}{ps_134-7_017-v0010}
\margin{($\beta$)}
\[ X_*(k) \hand{\longrightarrow} X'_*(k) \]
of semisimplicial sets is a weak equivalence, iff  the map
\[ L_{pt}(X_*) \hand{\longrightarrow} L_{pt}(X'_*) \]
is,  i.e. \var{\inserted{iff}}{\hand{iff}}{insertion}{ps_134-7_017-v0011} u induces an isomorphism for the homology invariants $H_i$
defined in \hand{(}z') above (via the abelianization functor  $L_{pt} = \var{\over{I}{\Gamma}}{\hand{\Gamma}}{writing over}{ps_134-7_017-v0012}_k$ ).

If instead of  \hand{($\alpha$)}  we only assume $X_0=X'_0=e$\var{\deleted{)}}{}{deletion}{ps_134-7_017-v0013} ,  it seems that we
may have to throw in some other extra condition on $X_*$  and $X'_*$  , \var{\xout{\unsure{which}}}{}{typewriter deletion}{ps_134-7_017-v0014}
in order for the definition of weak equivalence in terms of the mere
homology invariants to be reasonable  -  a condition which,  at\var{\xout{\illegible{1}}}{}{typewriter deletion}{ps_134-7_017-v0015} the very
least,  and in case $k=\underline{Z}$ say,  should ensure that the homotopy types
defined by the two terms in  \hand{($\beta$)}  should have \uline{abelian fundamental groups}
(which does'nt look at first sight to be automatic).  At any rate,  the
conditions \hand{(}$\alpha$) above,  which should be viewed as a 1-connectedness assump%
% NOTE: the right edge of the scan cuts this line after "assump"; a hyphen may be hidden
tion,  are natural enough,  and it is natural too to try first to push
through a theory of schematization of homotopy types,  under this assump-%
tion.

\margin{17}

% Visual collation: both supplied image halves inspected; see CHECK_DECISIONS_134_7_8.tsv.

%% END SOURCE ps_134-7_017
\ifnum\EditionMode=1
\psapparatusstart
\psentry{ps_134-7_017-v0001}{1}{typewriter deletion}{\xout{that}}{}{The typewriter strikes out this material; it leaves the corrected running text.}
\psentry{ps_134-7_017-v0002}{2}{typewriter deletion}{\xout{is simply}}{}{The typewriter strikes out this material; it leaves the corrected running text.}
\psentry{ps_134-7_017-v0003}{3}{typewriter deletion}{\xout{connected}}{}{The typewriter strikes out this material; it leaves the corrected running text.}
\psentry{ps_134-7_017-v0004}{4}{deletion}{\deleted{with}}{}{The author cancels this material by hand; it leaves the corrected running text.}
\psentry{ps_134-7_017-v0005}{5}{insertion}{\inserted{by}}{\hand{by}}{The handwritten insertion is incorporated at its indicated position; insertion brackets leave the running text.}
\psentry{ps_134-7_017-v0006}{6}{typewriter deletion}{\xout{second}}{}{The typewriter strikes out this material; it leaves the corrected running text.}
\psentry{ps_134-7_017-v0007}{7}{typed insertion}{\inserted[typed]{new}}{new}{The typed insertion is incorporated at its indicated position; insertion brackets leave the running text.}
\psentry{ps_134-7_017-v0008}{8}{insertion}{\inserted{rather than the old one.}}{\hand{rather than the old one.}}{The handwritten insertion is incorporated at its indicated position; insertion brackets leave the running text.}
\psentry{ps_134-7_017-v0009}{9}{typed insertion}{\inserted[typed]{both}}{both}{The typed insertion is incorporated at its indicated position; insertion brackets leave the running text.}
\psentry{ps_134-7_017-v0010}{10}{typewriter deletion}{\xout{$X_*(k)$}}{}{The typewriter strikes out this material; it leaves the corrected running text.}
\psentry{ps_134-7_017-v0011}{11}{insertion}{\inserted{iff}}{\hand{iff}}{The handwritten insertion is incorporated at its indicated position; insertion brackets leave the running text.}
\psentry{ps_134-7_017-v0012}{12}{writing over}{$\over{I}{\Gamma}$}{$\hand{\Gamma}$}{The superposed handwritten reading replaces the underlying reading.}
\psentry{ps_134-7_017-v0013}{13}{deletion}{\deleted{)}}{}{The author cancels this material by hand; it leaves the corrected running text.}
\psentry{ps_134-7_017-v0014}{14}{typewriter deletion}{\xout{\unsure{which}}}{}{The typewriter strikes out this material; it leaves the corrected running text.}
\psentry{ps_134-7_017-v0015}{15}{typewriter deletion}{\xout{\illegible{1}}}{}{The typewriter strikes out this material; it leaves the corrected running text.}
\psapparatusend
\fi
\pssource{ps_134-7_018}{134-7/18}{https://grothendieck.umontpellier.fr/web/viewer.html?file=../134-7.pdf\#page=19}
\psnumber{ps_134-7_018-v0001}{1}
\psnumber{ps_134-7_018-v0002}{2}
\psnumber{ps_134-7_018-v0003}{3}
\psnumber{ps_134-7_018-v0004}{4}
\psnumber{ps_134-7_018-v0005}{5}
\psnumber{ps_134-7_018-v0006}{6}
%% BEGIN SOURCE ps_134-7_018
% Pursuing Stacks (cotes 134-2 to 134-8): archive cote 134-7, site page 18 (PDF page 19), typescript page 460.
% typescript. Draft compared: ps6.tex:834-896.
% AI-transcribed by Claude Opus 5.5 (claude-opus-5-5, adaptive thinking, effort high), Message Batches API, batch ps3 (msgbatch_01XdFmuVVBAcgc7EFxdtCqsj), 2026-11-12; reviewed against both supplied scan halves; see CHECK_DECISIONS.tsv.
\scanpage{134-7/18}
\mspage{460}

26.8.

112. I am continuing the ``wishful thinking'' abo\var{\over{\unsure{u}}{u}}{\hand{u}}{writing over}{ps_134-7_018-v0001}t schematization of
homotopy types - a wel\var{\deleted{l}}{}{deletion}{ps_134-7_018-v0002}come break in the ``overall review'' on lineariza-%
tion \var{\xout{of}}{}{typewriter deletion}{ps_134-7_018-v0003} (in the context of the modelizer (Cat)), which had been getting
a little fastidious lately !

I'll admit, as one firm hold in all the wishfulness, that in the
``Eilenberg-Maclane case'' of p.453, when moreover $k=\underline{Z}$ and the components
$X_i$ of the semisimplicial unipotent bundle $X_*$ are projective, the two
maps (w) of page 454 are indeed quasiisomorphisms. From this should
follow the similar statement, when $X_*$ comes from a chain complex of
k-modules $M.$ with projective coefficients, by reducing to the case
when only a finite number of components of $M.$ are non zero (by suitable
\margin{!} passage to the limit), and then by \uline{induction} on the number of these
components, using the fact that the three terms in (w) depend on $X_*$ in
a ``multiplicative'' way, namely direct sums being transformed into tensor
products. (NB Under the assumptions of projectivity made, we may as
well express the quasi-isomorphis\unsure{ms}(w) we start with as being semisimpli-%
cial \uline{homotopisms}, and remark that comp\var{\inserted{o}}{\hand{o}}{insertion}{ps_134-7_018-v0004}nentwise tensorproduct of such
homotopisms is again one.) From this, using the relevant spectral se-%
quences in homology, should follow that the maps (w) are still quasi-iso-%
morphisms, whenever $X_*$ can be ``unscrewed'' (``dévissé'') as a finite succes-%
sive fibering with fibers of the type $M_*$ as above. Another passage to
the limit will yield the same result for an infinite dévissage, provided
the fibers $M(i)_*$ ($i=1,2, \ldots$) are ``way-out'' i.e. for given n, only
a finite number of components $M(i)_n$ are non zero (it amounts to the
same to \unsure{demande|dehande} that the sequence $M(i).$ of corresponding chain complexes
be ``way out'' ). This will give already, it seems, a fair number of
cases when (w) are quasi-isomorphisms. (Admittedly, working this out will
involve a fair amount of work, especially for getting the relevant
properties of $L.$ and $L_{pt}$ , which should mimick very closely the known
ones for usual linearization, including spectral sequences or, more
neatly, transitivity isomorphisms in the relevant derived categories...)
The main point here is that those special types of $X_*$ 's (we may call
them \uline{Kan-Postnikov} \uline{complexes} in $U(\var{\over{k}{k}}{\hand{k}}{writing over}{ps_134-7_018-v0005})$ ) are enough in order to modelize,
via the corresponding semisimplicial sets $X_*(\underline{Z})$ , arbitrary \var{\xout{hom}}{}{typewriter deletion}{ps_134-7_018-v0006} pointed
homotopy types with abelian $\hand{\pi}_1$ . This is seen of course using Postnikov

\margin{18}

% Visual collation: both supplied image halves inspected; see CHECK_DECISIONS_134_7_8.tsv.

%% END SOURCE ps_134-7_018
\ifnum\EditionMode=1
\psapparatusstart
\psentry{ps_134-7_018-v0001}{1}{writing over}{\over{\unsure{u}}{u}}{\hand{u}}{The superposed handwritten reading replaces the underlying reading.}
\psentry{ps_134-7_018-v0002}{2}{deletion}{\deleted{l}}{}{The author cancels this material by hand; it leaves the corrected running text.}
\psentry{ps_134-7_018-v0003}{3}{typewriter deletion}{\xout{of}}{}{The typewriter strikes out this material; it leaves the corrected running text.}
\psentry{ps_134-7_018-v0004}{4}{insertion}{\inserted{o}}{\hand{o}}{The handwritten insertion is incorporated at its indicated position; insertion brackets leave the running text.}
\psentry{ps_134-7_018-v0005}{5}{writing over}{$\over{k}{k}$}{$\hand{k}$}{The superposed handwritten reading replaces the underlying reading.}
\psentry{ps_134-7_018-v0006}{6}{typewriter deletion}{\xout{hom}}{}{The typewriter strikes out this material; it leaves the corrected running text.}
\psapparatusend
\fi
\pssource{ps_134-7_019}{134-7/19}{https://grothendieck.umontpellier.fr/web/viewer.html?file=../134-7.pdf\#page=20}
\psnumber{ps_134-7_019-v0001}{1}
\psnumber{ps_134-7_019-v0002}{2}
\psnumber{ps_134-7_019-v0003}{3}
\psnumber{ps_134-7_019-v0004}{4}
\psnumber{ps_134-7_019-v0005}{5}
\psnumber{ps_134-7_019-v0006}{6}
\psnumber{ps_134-7_019-v0007}{7}
\psnumber{ps_134-7_019-v0008}{8}
\psnumber{ps_134-7_019-v0009}{9}
\psnumber{ps_134-7_019-v0010}{10}
\psnumber{ps_134-7_019-v0011}{11}
\psnumber{ps_134-7_019-v0012}{12}
\psnumber{ps_134-7_019-v0013}{13}
\psnumber{ps_134-7_019-v0014}{14}
\psnumber{ps_134-7_019-v0015}{15}
\psnumber{ps_134-7_019-v0016}{16}
\psnumber{ps_134-7_019-v0017}{17}
\psnumber{ps_134-7_019-v0018}{18}
\psnumber{ps_134-7_019-v0019}{19}
\psnumber{ps_134-7_019-v0020}{20}
\psnumber{ps_134-7_019-v0021}{21}
%% BEGIN SOURCE ps_134-7_019
% Pursuing Stacks (cotes 134-2 to 134-8): archive cote 134-7, site page 19 (PDF page 20), typescript page 461.
% typescript. Draft compared: ps6.tex:884-942.
% AI-transcribed by Claude Opus 5.5 (claude-opus-5-5, adaptive thinking, effort high), Message Batches API, batch ps3 (msgbatch_01XdFmuVVBAcgc7EFxdtCqsj), 2026-11-12; reviewed against both supplied scan halves; see CHECK_DECISIONS.tsv.
\scanpage{134-7/19}
\mspage{461}
dévissage of a given homotopy type,  and replacing every homotopy group \hand{$\pi$}$_i$
by a shifted projective resolution (a two-step resolution will do here)
M(i)$_{\hand{\bullet}}$\hand{,}  as indicated on page 456,  and using the corresponding semisimpli-%
cial $\underline{Z}$-modules  M(i)$_*$  as fibers,  in the successive fiberings.  To see
that what is being done on the ``discrete'' level,  working with semisimpli-%
cial \uline{sets},  can be``followed'' in an essentially unique way on the ``schema-%
tic'' level,  we hit now of course upon the key-difficulty,  pointed out
on page 444,  about the Postnikov cohomology group $H^{n+2}(X_{\var{\deleted{n+1}}{}{deletion}{ps_134-7_019-v0001}\hand{(n+1)}},\hand{\pi}_{n+1})$
being isomorphic to the corresponding ``schematic'' group.  Now that this is
so indeed should follow from the homology isomorphisms (w) I hope,
\var{\xout{byx\allowbreak{}passingx\allowbreak{}overx\allowbreak{}tox\allowbreak{}thex\allowbreak{}dual}}{}{typewriter deletion}{ps_134-7_019-v0002} just dualizing the result\var{\inserted{to cohomology}}{\hand{to cohomology}}{insertion}{ps_134-7_019-v0003}.  All this seems to
sound kind of reasonable,  it seems even that for a given homotopy type
in Hot\hand{$_0$}\var{\deleted{\unsure{(W)}}}{}{deletion}{ps_134-7_019-v0004},  we should be able to squeeze out this way a unique \uline{isomorphism
type},  at any rate,  of semisimplicial unipotent bundles  -  but to see
whether it does work,  or if there is some ma\hand{j}or blunder which turns
the whole into nonsense,  will come out only from careful, down-to-earth
work,  which I am not prepared to dive into.

It occured to me that the ``Kan-Postnikov'' complexes in U(k) have
some special features among all possible complexes with $X_0 = e$,  and also
that some extra features \uline{are}  needed,  if we want the maps in (w) to be
quasi-isomorphisms.  I want to dwell upon this a little.  First of all,
the condition  $X_0$=e is indeed essential,  as we see by taking a constant
complex with value $X_0$ ,  then the homology of the three\var{\inserted[typed]{chain}}{chain}{typed insertion}{ps_134-7_019-v0005} complexes
(w)  reduces to degree zero,  and the $H_0$'s are respectively
\[ \underline{Z}^{\hand{(}X_0(\underline{Z})\hand{)}}\var{\inserted{=\underline{Z}^{(M_0)}}}{\hand{=\underline{Z}^{(M_0)}}}{insertion}{ps_134-7_019-v0006} ,\quad L(X_0) = \hand{\Gamma}_{\underline{Z}}{}^{\hand{\wedge}}(M_0) ,\quad L_{pt}(X_0) = \hand{\Gamma}_{\underline{Z}}(M_{\unsure{0}}) , \]
where $M_0$ is the $\underline{Z}$-module giving rise to $X_0$  -  and \var{\deleted{\unsure{all}}}{}{deletion}{ps_134-7_019-v0007}\var{\inserted{none of}}{\hand{none of}}{insertion}{ps_134-7_019-v0008} the two maps \var{\deleted{are}}{}{deletion}{ps_134-7_019-v0009}\hand{is}
\var{\deleted{non}}{}{deletion}{ps_134-7_019-v0010}\var{\inserted{an}}{\hand{an}}{insertion}{ps_134-7_019-v0011} isomorphis\var{\over{ms}{m}}{\hand{m}}{writing over}{ps_134-7_019-v0012}\hand{,} unless $M_0$=0 .

Take now the next simple\var{\inserted{st}}{\hand{st}}{insertion}{ps_134-7_019-v0013} case,  when $X_*$  comes from a monoid object\var{\inserted[typed]{G}}{G}{typed insertion}{ps_134-7_019-v0014}
in U(k) in the usual way;  then what we are after,  in \var{\inserted[typed]{dual}}{dual}{typed insertion}{ps_134-7_019-v0015} terms of cohomology
rather than homology (taking the dual complexes of those in (w)), amounts
essentially to asking whether the usual  discrete cohomology of the
discrete  monoid  G($\underline{Z}$)  can be computed,  using \uline{polynomial} \uline{cochains}
rather than arbitrary ones.  Now this we did admit as ``well known'' in
the most evident case of all,  when G being represented \var{\deleted{by}}{}{deletion}{ps_134-7_019-v0016} as an object
of U($\underline{Z}$)  by a \var{\inserted[typed]{projective (hence free)}}{projective (hence free)}{typed insertion}{ps_134-7_019-v0017} $\underline{Z}$-module M ,  the multiplication law is just \var{\deleted{the}}{}{deletion}{ps_134-7_019-v0018} usual addi-%
tion.  It still looks reasonable enough when the monoid \var{\deleted{law}}{}{deletion}{ps_134-7_019-v0019}\var{\inserted[typed]{G}}{G}{typed insertion}{ps_134-7_019-v0020} is a \uline{group} ,
with M of finite type say.  In this case,  the \var{\inserted{Borel}}{\hand{Borel}}{insertion}{ps_134-7_019-v0021} theory of algebraic
\margin{19}

% Visual collation: both supplied image halves inspected; see CHECK_DECISIONS_134_7_8.tsv.

%% END SOURCE ps_134-7_019
\ifnum\EditionMode=1
\psapparatusstart
\psentry{ps_134-7_019-v0001}{1}{deletion}{$\deleted{n+1}$}{}{The author cancels this material by hand; it leaves the corrected running text.}
\psentry{ps_134-7_019-v0002}{2}{typewriter deletion}{\xout{byx\allowbreak{}passingx\allowbreak{}overx\allowbreak{}tox\allowbreak{}thex\allowbreak{}dual}}{}{The typewriter strikes out this material; it leaves the corrected running text.}
\psentry{ps_134-7_019-v0003}{3}{insertion}{\inserted{to cohomology}}{\hand{to cohomology}}{The handwritten insertion is incorporated at its indicated position; insertion brackets leave the running text.}
\psentry{ps_134-7_019-v0004}{4}{deletion}{\deleted{\unsure{(W)}}}{}{The author cancels this material by hand; it leaves the corrected running text.}
\psentry{ps_134-7_019-v0005}{5}{typed insertion}{\inserted[typed]{chain}}{chain}{The typed insertion is incorporated at its indicated position; insertion brackets leave the running text.}
\psentry{ps_134-7_019-v0006}{6}{insertion}{$\inserted{=\underline{Z}^{(M_0)}}$}{$\hand{=\underline{Z}^{(M_0)}}$}{The handwritten insertion is incorporated at its indicated position; insertion brackets leave the running text.}
\psentry{ps_134-7_019-v0007}{7}{deletion}{\deleted{\unsure{all}}}{}{The author cancels this material by hand; it leaves the corrected running text.}
\psentry{ps_134-7_019-v0008}{8}{insertion}{\inserted{none of}}{\hand{none of}}{The handwritten insertion is incorporated at its indicated position; insertion brackets leave the running text.}
\psentry{ps_134-7_019-v0009}{9}{deletion}{\deleted{are}}{}{The author cancels this material by hand; it leaves the corrected running text.}
\psentry{ps_134-7_019-v0010}{10}{deletion}{\deleted{non}}{}{The author cancels this material by hand; it leaves the corrected running text.}
\psentry{ps_134-7_019-v0011}{11}{insertion}{\inserted{an}}{\hand{an}}{The handwritten insertion is incorporated at its indicated position; insertion brackets leave the running text.}
\psentry{ps_134-7_019-v0012}{12}{writing over}{\over{ms}{m}}{\hand{m}}{The superposed handwritten reading replaces the underlying reading.}
\psentry{ps_134-7_019-v0013}{13}{insertion}{\inserted{st}}{\hand{st}}{The handwritten insertion is incorporated at its indicated position; insertion brackets leave the running text.}
\psentry{ps_134-7_019-v0014}{14}{typed insertion}{\inserted[typed]{G}}{G}{The typed insertion is incorporated at its indicated position; insertion brackets leave the running text.}
\psentry{ps_134-7_019-v0015}{15}{typed insertion}{\inserted[typed]{dual}}{dual}{The typed insertion is incorporated at its indicated position; insertion brackets leave the running text.}
\psentry{ps_134-7_019-v0016}{16}{deletion}{\deleted{by}}{}{The author cancels this material by hand; it leaves the corrected running text.}
\psentry{ps_134-7_019-v0017}{17}{typed insertion}{\inserted[typed]{projective (hence free)}}{projective (hence free)}{The typed insertion is incorporated at its indicated position; insertion brackets leave the running text.}
\psentry{ps_134-7_019-v0018}{18}{deletion}{\deleted{the}}{}{The author cancels this material by hand; it leaves the corrected running text.}
\psentry{ps_134-7_019-v0019}{19}{deletion}{\deleted{law}}{}{The author cancels this material by hand; it leaves the corrected running text.}
\psentry{ps_134-7_019-v0020}{20}{typed insertion}{\inserted[typed]{G}}{G}{The typed insertion is incorporated at its indicated position; insertion brackets leave the running text.}
\psentry{ps_134-7_019-v0021}{21}{insertion}{\inserted{Borel}}{\hand{Borel}}{The handwritten insertion is incorporated at its indicated position; insertion brackets leave the running text.}
\psapparatusend
\fi
\pssource{ps_134-7_020}{134-7/20}{https://grothendieck.umontpellier.fr/web/viewer.html?file=../134-7.pdf\#page=21}
\psnumber{ps_134-7_020-v0001}{1}
\psnumber{ps_134-7_020-v0002}{2}
\psnumber{ps_134-7_020-v0003}{3}
\psnumber{ps_134-7_020-v0004}{4}
\psnumber{ps_134-7_020-v0005}{5}
\psnumber{ps_134-7_020-v0006}{6}
\psnumber{ps_134-7_020-v0007}{7}
\psnumber{ps_134-7_020-v0008}{8}
\psnumber{ps_134-7_020-v0009}{9}
%% BEGIN SOURCE ps_134-7_020
% Pursuing Stacks (cotes 134-2 to 134-8): archive cote 134-7, site page 20 (PDF page 21), typescript page 462.
% typescript. Draft compared: ps6.tex:930-988.
% AI-transcribed by Claude Opus 5.5 (claude-opus-5-5, adaptive thinking, effort high), Message Batches API, batch ps3 (msgbatch_01XdFmuVVBAcgc7EFxdtCqsj), 2026-11-12; reviewed against both supplied scan halves; see CHECK_DECISIONS.tsv.
\scanpage{134-7/20}
\mspage{462}

affine groups over a field (here the field of fractions $\underline{Q}$ of $\underline{Z}$ ) tells
us that $G_{\underline{Q}}$ is a \uline{nilpotent} algebraic group,  and that therefore it admits
a composition series with factors isomorphic to the additive group $G_{a\,\underline{Q}}$ ;
presumably the same dévissage then can be obtained over the base $\underline{Z}$ ,  and
using induction on the length of the composition series,  and the Hoch-%
schild-Serre type of relations (traditionally expressed by a spectral
sequence)  between group cohomology of a group,  quotient group and corres
ponding subgroup,  we should get the wished for quasi-isomorphisms (w).

Take now however the simplest case of a monoid which is'nt a group,
namely the multiplicative law on the affine line,  given by the polynomial
law
\[ W(\underline{Z}) \text{ x } W(\underline{Z}) \hand{\longrightarrow} W(\underline{Z}) \quad : \quad (x,y) \hand{\longmapsto} xy. \]
\margin{\textcircled{$\times$}} the corresponding ``discrete'' monoid is just $\underline{Z}^{\hand{(x)}}$ ,  namely the integers
with multiplication,  it's $H^{1}$ with coefficients in $\underline{Z}$  is just the group
of all homomorphisms
\[ (*) \qquad \underline{Z}^{\hand{(x)}} \hand{\to} \underline{Z}^{\hand{(+)}} = \underline{Z} \quad , \]

and denoting by $\underline{P}$ the set of all primes and using the prime decomposition
of integers,  we find that
\[ \mathrm{Hom}(\underline{Z}^{\hand{(x)}},\underline{Z}^{\hand{(+)}}) \hand{\simeq} \underline{Z}^{P} \quad , \]
to a family $(n_{p})_{p \in P}$ of integers being associated the homomorphism
\[ \hand{\pm\prod_{p\in P}} p^{\hand{\alpha}_{p}} \hand{\mapsto \sum_{p\in P} \alpha}_{p}\, n_{p} \quad . \]
On the other hand,  the schematic $H^{1}$ consists \var{\over{\illegible{1}}{of}}{\hand{of}}{writing over}{ps_134-7_020-v0001} all  homomorphisms (*)
which can be expressed by a polynomial,  hence induce a homomorphism of
algebraic group schemes  $\underline{G}_{m} \hand{\to} \underline{G}_{a}$ ,  and it is well known (and immediately
checked) that there is only the zero homomorphism !

Thus,  it turns out that the assumptions made yesterday on $X_{*}$ ,
in order for the ``linearization theorem'' (!) to hold\hand{,} namely the maps (w)
to be quasi-isomorphisms,  are definitely not strong enough yet\var{\over{.}{\unsure{!}}}{\hand{\unsure{!}}}{writing over}{ps_134-7_020-v0002} One
may think of throwing in the extra condition  $X_{1}=e$ ,  so as to rule out
monoids\var{\inserted{altogether,}}{\hand{altogether,}}{insertion}{ps_134-7_020-v0003} (and even groups, too bad !), \var{\deleted{altogether,}}{}{deletion}{ps_134-7_020-v0004}  but I do'nt think this
helps at all (did'nt try though to make a counterexample\hand{).}\var{\deleted{\xout{still}).}}{}{deletion}{ps_134-7_020-v0005}  On the
other hand,  just restrict\var{\inserted[typed]{ing}}{ing}{typed insertion}{ps_134-7_020-v0006} to Kan-Postnikov complexes seems rather awk-%
ward,   we defini\var{\xout{\illegible{1}}}{}{typewriter deletion}{ps_134-7_020-v0007}tely do'nt want to drag along P\var{\over{\illegible{1}}{o}}{\hand{o}}{writing over}{ps_134-7_020-v0008}stnikiv fibrations as
a compulsory ingredient of the complexes we work with.  The idea which
comes up here is just to \hand{``}drop P\var{\over{\illegible{1}}{o}}{\hand{o}}{writing over}{ps_134-7_020-v0009}stnikow and keep Kan\hand{''} -  namely \uline{introduce}
\uline{a Kan type condition on semisimplicial complexes in U(k)}.  If we mimick

\margin{20}

% Visual collation: both supplied image halves inspected; see CHECK_DECISIONS_134_7_8.tsv.

%% END SOURCE ps_134-7_020
\ifnum\EditionMode=1
\psapparatusstart
\psentry{ps_134-7_020-v0001}{1}{writing over}{\over{\illegible{1}}{of}}{\hand{of}}{The superposed handwritten reading replaces the underlying reading.}
\psentry{ps_134-7_020-v0002}{2}{writing over}{\over{.}{\unsure{!}}}{\hand{\unsure{!}}}{The superposed handwritten reading replaces the underlying reading.}
\psentry{ps_134-7_020-v0003}{3}{insertion}{\inserted{altogether,}}{\hand{altogether,}}{The handwritten insertion is incorporated at its indicated position; insertion brackets leave the running text.}
\psentry{ps_134-7_020-v0004}{4}{deletion}{\deleted{altogether,}}{}{The author cancels this material by hand; it leaves the corrected running text.}
\psentry{ps_134-7_020-v0005}{5}{deletion}{\deleted{\xout{still}).}}{}{The author cancels this material by hand; it leaves the corrected running text.}
\psentry{ps_134-7_020-v0006}{6}{typed insertion}{\inserted[typed]{ing}}{ing}{The typed insertion is incorporated at its indicated position; insertion brackets leave the running text.}
\psentry{ps_134-7_020-v0007}{7}{typewriter deletion}{\xout{\illegible{1}}}{}{The typewriter strikes out this material; it leaves the corrected running text.}
\psentry{ps_134-7_020-v0008}{8}{writing over}{\over{\illegible{1}}{o}}{\hand{o}}{The superposed handwritten reading replaces the underlying reading.}
\psentry{ps_134-7_020-v0009}{9}{writing over}{\over{\illegible{1}}{o}}{\hand{o}}{The superposed handwritten reading replaces the underlying reading.}
\psapparatusend
\fi
\pssource{ps_134-7_021}{134-7/21}{https://grothendieck.umontpellier.fr/web/viewer.html?file=../134-7.pdf\#page=22}
\psnumber{ps_134-7_021-v0001}{1}
\psnumber{ps_134-7_021-v0002}{2}
\psnumber{ps_134-7_021-v0003}{3}
\psnumber{ps_134-7_021-v0004}{4}
\psnumber{ps_134-7_021-v0005}{5}
\psnumber{ps_134-7_021-v0006}{6}
\psnumber{ps_134-7_021-v0007}{7}
\psnumber{ps_134-7_021-v0008}{8}
\psnumber{ps_134-7_021-v0009}{9}
\psnumber{ps_134-7_021-v0010}{10}
\psnumber{ps_134-7_021-v0011}{11}
\psnumber{ps_134-7_021-v0012}{12}
%% BEGIN SOURCE ps_134-7_021
% Pursuing Stacks (cotes 134-2 to 134-8): archive cote 134-7, site page 21 (PDF page 22), typescript page 463.
% typescript. Draft compared: ps6.tex:976-1040.
% AI-transcribed by Claude Opus 5.5 (claude-opus-5-5, adaptive thinking, effort high), Message Batches API, batch ps3 (msgbatch_01XdFmuVVBAcgc7EFxdtCqsj), 2026-11-12; reviewed against both supplied scan halves; see CHECK_DECISIONS.tsv.
\scanpage{134-7/21}
\mspage{463}

formally the usual "discrete" Kan condition, we get that \hand{(}for given pair
of integers $k,n$ with $0\hand{\leq}k\hand{\leq}n$) a certain map from $X_n$ , to a certain
\var{\xout{fiberedxproduct}}{}{typewriter deletion}{ps_134-7_021-v0001} finite projective limit defined in terms of the boundary
maps $X_{n-1}\hand{\to}X_{n-2}$ , should be \var{\xout{surjectivex}}{}{typewriter deletion}{ps_134-7_021-v0002} epimorphic. Now clearly U(k)
is by no means stable under fiber products, except under very special
assumptions (including differ\var{\over{\unsure{an}}{en}}{\hand{en}}{writing over}{ps_134-7_021-v0003}tial transversality conditions, at any
rate), and on the other hand one feels that the notion of "epimorphism"
one will have to work with in U(k) will have to be a lot more exacting
than \var{\over{\illegible{1}}{the}}{\hand{the}}{writing over}{ps_134-7_021-v0004} map $X(k)\hand{\to}Y(k)$ on sections being surjective, or the usual catego-%
rical meaning within U(k), which looks kind of silly here. Even the \var{\inserted[typed]{most}}{most}{typed insertion}{ps_134-7_021-v0005}
exacting surjectivity condition on $X\hand{\to}Y$, namely that it admit a section\hand{\unsure{,}}
does'nt quite satisfy me - what I really want is that X should be a
trivial bundle over Y, more specifically that X is isomorphic to a pro-%
duct $YxZ$, in such a way that the given map $X\hand{\to}Y$ identifies with the
projection $YxZ\hand{\to}Y$. Maybe this is too exacting a condition however, and
hard to check in computational terms sometimes (?), maybe we should be
content with demanding only that $X\hand{\to}Y$ has a section, and moreover is
"smooth" i.e. has everywhere a surjecti\var{\over{\unsure{f}}{v}}{\hand{v}}{writing over}{ps_134-7_021-v0006}e tangent map (which may be
\var{\xout{checked}}{}{typewriter deletion}{ps_134-7_021-v0007} \var{\inserted[typed]{expressed}}{expressed}{typed insertion}{ps_134-7_021-v0008} on the corresponding sheaf on $\mathrm{Aff}_{/k}$ by the familiar condition
of "formal smoothness" namely possibility of lifting sections over
arbitrary infinitesimal neighbourhoods ...). We'll have to choose at any
rate some such \var{\xout{notion}}{}{typewriter deletion}{ps_134-7_021-v0009} strong "surjectivity" notion in U(k), which we'll
call \uline{"submersions"} say. Thus I feel a "Kan complex" in U(k) should have
boundary maps which are submersions. What we should do, is to pin down
some simple \var{\inserted[typed]{"Kan}}{"Kan}{typed insertion}{ps_134-7_021-v0010} condition\hand{"} on a complex $X_*$ , in terms of "submersions",
in such a way as to insure, at any rate

a) that for any pair (\var{\over{k}{n}}{\hand{n}}{writing over}{ps_134-7_021-v0011},k) with $\hand{0\leq}k\hand{\leq}n$, the object $X_{\hand{*}}(n,k)$
"horns of type (n,k) of $X_*$ "\hand{,} expressed by the suitable finite inverse
limits (in terms of boundary maps $X_{n-1}\hand{\to}X_{n-2}$) is representable in U(k),
and

b) the canonical map $X_n\hand{\to}X_*(n,k)$ is a submersion ,

and such of course that \hand{all} Kan-P\hand{\u{o}}stnikov complexes should satisfy this
Kan condition, at the very least.

The first non-trivial case, in view of $X_o=e$, is n=2, in which case
$X_*(2,k)$ is trivially representable by $X_1xX_1$ , and the condition we get
is that the three natural maps coming from boundary maps

\[
\hand{(**)}\qquad
% Hand-drawn arrows are marked blue; typed objects remain plain.
\diagram{\begin{tikzcd} X_2 \ar[color=blue,r,shift left=1.5] \ar[color=blue,r] \ar[color=blue,r,shift right=1.5] \& X_{\var{\over{1}{1}}{\hand{1}}{writing over}{ps_134-7_021-v0012}}\hand{\times X_1} \end{tikzcd}}
\]

\margin{21}

% Visual collation: both supplied image halves inspected; see CHECK_DECISIONS_134_7_8.tsv.

%% END SOURCE ps_134-7_021
\ifnum\EditionMode=1
\psapparatusstart
\psentry{ps_134-7_021-v0001}{1}{typewriter deletion}{\xout{fiberedxproduct}}{}{The typewriter strikes out this material; it leaves the corrected running text.}
\psentry{ps_134-7_021-v0002}{2}{typewriter deletion}{\xout{surjectivex}}{}{The typewriter strikes out this material; it leaves the corrected running text.}
\psentry{ps_134-7_021-v0003}{3}{writing over}{\over{\unsure{an}}{en}}{\hand{en}}{The superposed handwritten reading replaces the underlying reading.}
\psentry{ps_134-7_021-v0004}{4}{writing over}{\over{\illegible{1}}{the}}{\hand{the}}{The superposed handwritten reading replaces the underlying reading.}
\psentry{ps_134-7_021-v0005}{5}{typed insertion}{\inserted[typed]{most}}{most}{The typed insertion is incorporated at its indicated position; insertion brackets leave the running text.}
\psentry{ps_134-7_021-v0006}{6}{writing over}{\over{\unsure{f}}{v}}{\hand{v}}{The superposed handwritten reading replaces the underlying reading.}
\psentry{ps_134-7_021-v0007}{7}{typewriter deletion}{\xout{checked}}{}{The typewriter strikes out this material; it leaves the corrected running text.}
\psentry{ps_134-7_021-v0008}{8}{typed insertion}{\inserted[typed]{expressed}}{expressed}{The typed insertion is incorporated at its indicated position; insertion brackets leave the running text.}
\psentry{ps_134-7_021-v0009}{9}{typewriter deletion}{\xout{notion}}{}{The typewriter strikes out this material; it leaves the corrected running text.}
\psentry{ps_134-7_021-v0010}{10}{typed insertion}{\inserted[typed]{"Kan}}{"Kan}{The typed insertion is incorporated at its indicated position; insertion brackets leave the running text.}
\psentry{ps_134-7_021-v0011}{11}{writing over}{\over{k}{n}}{\hand{n}}{The superposed handwritten reading replaces the underlying reading.}
\psentry{ps_134-7_021-v0012}{12}{writing over}{$\over{1}{1}$}{$\hand{1}$}{The superposed handwritten reading replaces the underlying reading.}
\psapparatusend
\fi
\pssource{ps_134-7_022}{134-7/22}{https://grothendieck.umontpellier.fr/web/viewer.html?file=../134-7.pdf\#page=23}
\psnumber{ps_134-7_022-v0001}{1}
\psnumber{ps_134-7_022-v0002}{2}
\psnumber{ps_134-7_022-v0003}{3}
\psnumber{ps_134-7_022-v0004}{4}
\psnumber{ps_134-7_022-v0005}{5}
\psnumber{ps_134-7_022-v0006}{6}
\psnumber{ps_134-7_022-v0007}{7}
\psnumber{ps_134-7_022-v0008}{8}
\psnumber{ps_134-7_022-v0009}{9}
\psnumber{ps_134-7_022-v0010}{10}
\psnumber{ps_134-7_022-v0011}{11}
\psnumber{ps_134-7_022-v0012}{12}
\psnumber{ps_134-7_022-v0013}{13}
%% BEGIN SOURCE ps_134-7_022
% Pursuing Stacks (cotes 134-2 to 134-8): archive cote 134-7, site page 22 (PDF page 23), typescript page 464.
% typescript. Draft compared: ps6.tex:1023-1079.
% AI-transcribed by Claude Opus 5.5 (claude-opus-5-5, adaptive thinking, effort high), Message Batches API, batch ps3 (msgbatch_01XdFmuVVBAcgc7EFxdtCqsj), 2026-11-12; reviewed against both supplied scan halves; see CHECK_DECISIONS.tsv.
\scanpage{134-7/22}
\mspage{464}

\var{\deleted{and these}}{}{deletion}{ps_134-7_022-v0001} should be submersions for a ``\uline{Kan complex}''. In case $X_*$ is
defined by a monoid object as above, this clearly implies that G is a
group \var{\deleted{(surjectivity \unsure{of} \unsure{(xx)|(**)} on points should be enough, a fortiori it
would suffice that $X_*(k)$ is a Kan complex in the usual sense)}}{}{deletion}{ps_134-7_022-v0002} - which
rules out the counterexample above !

Of course, the very first thing we'll expect from a ``good notion'' of
Kan complexes in U(k), is that for $k=\underline{Z}$ , it should make the linearizati-%
on theorem work, namely the maps in (w) p.454 are quasi-isomorphisms.
The next thing, very close to this one but for arbitrary ground ring k
now, is that a map
\[ X_* \hand{\longrightarrow} X'_* \qquad (\text{ with } X_0=X'_0=X_1=X'_1 = e \quad ) \]
of Kan complexes in U(k) is a homotopism, iff it induces an isomorphism
on the homology modules (z') (p.458) - which sounds reasonable precise
-ly because we are working with Kan complexes. If this is so, the
homotopy category $\mathrm{Hot}_1(k)$ of \var{\xout{simply}}{}{typewriter deletion}{ps_134-7_022-v0003} 1-connected homotopy types over k
may be identified with \var{\over{the}{a}}{\hand{a}}{writing over}{ps_134-7_022-v0004} category of Kan complexes \var{\xout{as usual}}{}{typewriter deletion}{ps_134-7_022-v0005} ``up to
homotopy'' , as usual (but working \var{\inserted{now}}{\hand{now}}{insertion}{ps_134-7_022-v0006} with complexes of unipotent bundles
over k). Third thing, still over arbitrary ground ring, would be
a development of the usual homotopy formalism in the unipotent context,
including (one hopes) homotopy fibers of maps, and P\var{\over{\unsure{a}}{o}}{\hand{o}}{writing over}{ps_134-7_022-v0007}stnikov \var{\xout{fibrations}}{}{typewriter deletion}{ps_134-7_022-v0008}
d\hand{\'e}vissage. Again, it is hard to imagine how to get such d\hand{\'e}vissage,
without getting hold inductively of homotopy invariants $\hand{\pi}_i$ which are
k-modules. This should come out if we are able to define homotopy fibers,
as for an (n-1)-connected k-homotopy type (defined here as one whose
homology invariants $H_i$ are zero for $i\hand{\leqslant}n-1$), $\hand{\pi}_n$ should be no more, no
less than $H_n$ , which is indeed a k-module. Coming back to $k=\underline{Z}$ again,
this should imply that \var{\xout{the}}{}{typewriter deletion}{ps_134-7_022-v0009} for $n\hand{\geqslant}1$ at any rate, the canonical functor
\[ \hand{(***)} \qquad \mathrm{Hot}_n(\underline{Z}) \hand{\longrightarrow} \mathrm{Hot}_n \]
induces a bijection between isomorphism classes of schematic and ordinary
\var{\inserted[typed]{n-connected}}{n-connected}{typed insertion}{ps_134-7_022-v0010} homotopy types - and it will be \var{\over{\unsure{g}}{h}}{\hand{h}}{writing over}{ps_134-7_022-v0011}ard to believe this can be so, without
this functor being actually an equivalence of categories - the expected
apotheosis of the theory ! Maybe to this end, one may even be able to
introduce reasonable internal \uline{Hom}'s within the category of schematic
Kan complexes, \var{\xout{\unsure{and show that they are}}}{}{typewriter deletion}{ps_134-7_022-v0012} \var{\inserted[typed]{in a way}}{in a way}{typed insertion}{ps_134-7_022-v0013} compatible with the familiar notion
in the discrete set-up.

\margin{22}

% Visual collation: both supplied image halves inspected; see CHECK_DECISIONS_134_7_8.tsv.

%% END SOURCE ps_134-7_022
\ifnum\EditionMode=1
\psapparatusstart
\psentry{ps_134-7_022-v0001}{1}{deletion}{\deleted{and these}}{}{The author cancels this material by hand; it leaves the corrected running text.}
\psentry{ps_134-7_022-v0002}{2}{deletion}{\deleted{(surjectivity \unsure{of} \unsure{(xx)|(**)} on points should be enough, a fortiori it
would suffice that $X_*(k)$ is a Kan complex in the usual sense)}}{}{The author cancels this material by hand; it leaves the corrected running text.}
\psentry{ps_134-7_022-v0003}{3}{typewriter deletion}{\xout{simply}}{}{The typewriter strikes out this material; it leaves the corrected running text.}
\psentry{ps_134-7_022-v0004}{4}{writing over}{\over{the}{a}}{\hand{a}}{The superposed handwritten reading replaces the underlying reading.}
\psentry{ps_134-7_022-v0005}{5}{typewriter deletion}{\xout{as usual}}{}{The typewriter strikes out this material; it leaves the corrected running text.}
\psentry{ps_134-7_022-v0006}{6}{insertion}{\inserted{now}}{\hand{now}}{The handwritten insertion is incorporated at its indicated position; insertion brackets leave the running text.}
\psentry{ps_134-7_022-v0007}{7}{writing over}{\over{\unsure{a}}{o}}{\hand{o}}{The superposed handwritten reading replaces the underlying reading.}
\psentry{ps_134-7_022-v0008}{8}{typewriter deletion}{\xout{fibrations}}{}{The typewriter strikes out this material; it leaves the corrected running text.}
\psentry{ps_134-7_022-v0009}{9}{typewriter deletion}{\xout{the}}{}{The typewriter strikes out this material; it leaves the corrected running text.}
\psentry{ps_134-7_022-v0010}{10}{typed insertion}{\inserted[typed]{n-connected}}{n-connected}{The typed insertion is incorporated at its indicated position; insertion brackets leave the running text.}
\psentry{ps_134-7_022-v0011}{11}{writing over}{\over{\unsure{g}}{h}}{\hand{h}}{The superposed handwritten reading replaces the underlying reading.}
\psentry{ps_134-7_022-v0012}{12}{typewriter deletion}{\xout{\unsure{and show that they are}}}{}{The typewriter strikes out this material; it leaves the corrected running text.}
\psentry{ps_134-7_022-v0013}{13}{typed insertion}{\inserted[typed]{in a way}}{in a way}{The typed insertion is incorporated at its indicated position; insertion brackets leave the running text.}
\psapparatusend
\fi
\pssource{ps_134-7_023}{134-7/23}{https://grothendieck.umontpellier.fr/web/viewer.html?file=../134-7.pdf\#page=24}
\psnumber{ps_134-7_023-v0001}{1}
\psnumber{ps_134-7_023-v0002}{2}
\psnumber{ps_134-7_023-v0003}{3}
\psnumber{ps_134-7_023-v0004}{4}
\psnumber{ps_134-7_023-v0005}{5}
\psnumber{ps_134-7_023-v0006}{6}
\psnumber{ps_134-7_023-v0007}{7}
\psnumber{ps_134-7_023-v0008}{8}
%% BEGIN SOURCE ps_134-7_023
% Pursuing Stacks (cotes 134-2 to 134-8): archive cote 134-7, site page 23 (PDF page 24), typescript page 465.
% typescript. Draft compared: ps6.tex:1066-1125.
% AI-transcribed by Claude Opus 5.5 (claude-opus-5-5, adaptive thinking, effort high), Message Batches API, batch ps3 (msgbatch_01XdFmuVVBAcgc7EFxdtCqsj), 2026-11-12; reviewed against both supplied scan halves; see CHECK_DECISIONS.tsv.
\scanpage{134-7/23}
\mspage{465}

If it is possible indeed to construct Postnikov dévissage of a
\var{\xout{Kan}}{}{typewriter deletion}{ps_134-7_023-v0001} schematic Kan complex over any ground ring k, it is clear that this
is compatible with restriction of ground ring, hence it would seem that
formation of the homotopy \var{\xout{module}}{}{typewriter deletion}{ps_134-7_023-v0002} invariants $\hand{\pi}_i$ is compatible with \var{\xout{\illegible{1}}}{}{typewriter deletion}{ps_134-7_023-v0003}
restriction of rings\var{\deleted{\unsure{,}}}{}{deletion}{ps_134-7_023-v0004} (whereas, as we noticed yesterday, the same does
definitely \uline{not} hold for the homology invariants H$_i$). Taking restriction
to the absolute ground ring $\underline{Z}$, this shows that the canonical functor
(***) from schematic to discrete homotopy types is compatible with
taking homotopy groups (but not with homology) - thus the relation
between $X_*$ and the complex of sections $X_*(k)$ seems to be a rather close
one, via the homotopy groups\hand{,} which are the same (and thus, the homotopy
groups of $X_*(k)$ seem to turn out to be k-modules after all !). By the
way, speaking of ``restriction of ground ring'' for Kan complexes was a
little hasty, in view of the projectivity condition on the components,
which a priori \var{\inserted{seems to}}{\hand{seems to}}{insertion}{ps_134-7_023-v0005} oblige\var{\deleted{s}}{}{deletion}{ps_134-7_023-v0006} us to assume k to be a projective k$_o$-module
(for a given ring homomorphism
\[ k_o \hand{\longrightarrow} k \quad ) . \]
Still, the remark about the sections functor $X_* \hand{\longrightarrow} X_*(k)$ makes a
sense, without having to assume k to be a projective $\underline{Z}$-module ! Also,
we feel that, by analogy of what can be done in the linear set up,
when we define a total derived functor
\[ D.(Ab_k) \hand{\longrightarrow} D.(Ab_{k_o}) \]
without any assumption on the ring homomorphism k$_o$ \hand{$\to$} k , a notion of
ring restriction for \var{\inserted[typed]{schematic}}{schematic}{typed insertion}{ps_134-7_023-v0007} homotopy types should make sense without any
restriction, as was surmised yesterday. As in the linear case, we
should allow ourselves to work with schematic complexes which are \uline{not}
projective, but be prepared to take ``resolutions'' \var{\inserted{(in some sense)}}{\hand{(in some sense)}}{insertion}{ps_134-7_023-v0008} of such general comple-%
xes by the more restricted ones. \hand{(with projective components).}

There is no such difficulty in the case of the ring extension func-%
tor, which transforms projective bundles over k$_o$ into projective bundles.
The reflexions above suggest that, whereas ring extension is compatible
with taking total homology invariants LH. , via the corresponding
functor
\[ D.(Ab_{k_o}) \hand{\longrightarrow} D.(Ab_k) \quad , \]
it is compatible too with taking homotopy invariants $\hand{\pi}_i$ separately.

\margin{23}

% Visual collation: both supplied image halves inspected; see CHECK_DECISIONS_134_7_8.tsv.

%% END SOURCE ps_134-7_023
\ifnum\EditionMode=1
\psapparatusstart
\psentry{ps_134-7_023-v0001}{1}{typewriter deletion}{\xout{Kan}}{}{The typewriter strikes out this material; it leaves the corrected running text.}
\psentry{ps_134-7_023-v0002}{2}{typewriter deletion}{\xout{module}}{}{The typewriter strikes out this material; it leaves the corrected running text.}
\psentry{ps_134-7_023-v0003}{3}{typewriter deletion}{\xout{\illegible{1}}}{}{The typewriter strikes out this material; it leaves the corrected running text.}
\psentry{ps_134-7_023-v0004}{4}{deletion}{\deleted{\unsure{,}}}{}{The author cancels this material by hand; it leaves the corrected running text.}
\psentry{ps_134-7_023-v0005}{5}{insertion}{\inserted{seems to}}{\hand{seems to}}{The handwritten insertion is incorporated at its indicated position; insertion brackets leave the running text.}
\psentry{ps_134-7_023-v0006}{6}{deletion}{\deleted{s}}{}{The author cancels this material by hand; it leaves the corrected running text.}
\psentry{ps_134-7_023-v0007}{7}{typed insertion}{\inserted[typed]{schematic}}{schematic}{The typed insertion is incorporated at its indicated position; insertion brackets leave the running text.}
\psentry{ps_134-7_023-v0008}{8}{insertion}{\inserted{(in some sense)}}{\hand{(in some sense)}}{The handwritten insertion is incorporated at its indicated position; insertion brackets leave the running text.}
\psapparatusend
\fi
\pssource{ps_134-7_024}{134-7/24}{https://grothendieck.umontpellier.fr/web/viewer.html?file=../134-7.pdf\#page=25}
\psnumber{ps_134-7_024-v0001}{1}
\psnumber{ps_134-7_024-v0002}{2}
\psnumber{ps_134-7_024-v0003}{3}
\psnumber{ps_134-7_024-v0004}{4}
\psnumber{ps_134-7_024-v0005}{5}
\psnumber{ps_134-7_024-v0006}{6}
\psnumber{ps_134-7_024-v0007}{7}
\psnumber{ps_134-7_024-v0008}{8}
\psnumber{ps_134-7_024-v0009}{9}
\psnumber{ps_134-7_024-v0010}{10}
\psnumber{ps_134-7_024-v0011}{11}
\psnumber{ps_134-7_024-v0012}{12}
\psnumber{ps_134-7_024-v0013}{13}
\psnumber{ps_134-7_024-v0014}{14}
%% BEGIN SOURCE ps_134-7_024
% Pursuing Stacks (cotes 134-2 to 134-8): archive cote 134-7, site page 24 (PDF page 25), typescript page 466.
% typescript. Draft compared: ps6.tex:1106-1176.
% AI-transcribed by Claude Opus 5.5 (claude-opus-5-5, adaptive thinking, effort high), Message Batches API, batch ps3 (msgbatch_01XdFmuVVBAcgc7EFxdtCqsj), 2026-11-12; reviewed against both supplied scan halves; see CHECK_DECISIONS.tsv.
\scanpage{134-7/24}
\mspage{466}

27.8.

\margin{113.} What I was thinking of last night (see last sentence) is that
whereas for total homology (not for the separate H$_i$'s) we have the com-%
prehensive formula
\[ \text{A)} \qquad \mathrm{LH}_{.}(X_{*}\otimes_k k')\hand{\simeq}\mathrm{LH}_{.}(X_{*}) \overset{\hand{L}}{\otimes}_k k' \quad \hand{(}\text{ in } D_{.}(k')) , \]
\hand{(}where $\overset{\hand{L}}{\otimes}_k$ denotes the left derive\hand{d} functor of the ring extension functor
\var{\inserted{for}}{\hand{for}}{insertion}{ps_134-7_024-v0001} $k\hand{\to} k'$, and \var{\inserted[typed]{$X'_{*}=$}}{$X'_{*}=$}{typed insertion}{ps_134-7_024-v0002} $X_{*}\otimes_k k'$ denotes ring extension for the semisimplicial unipo-%
tent bundle $X_{*}$\hand{,)} for homotopy modules we should have the term-by-term
isomorphisms
\[ \text{B)} \qquad \hand{\pi}_i(X_{*}\otimes_k k')\hand{\leftarrow}\hand{\pi}_i(X_{*})\otimes_k k' \quad . \]
\margin{\textcircled{$\leftarrow$}}
This however was pretty rash indeed (it was time to go to sleep I guess !\hand{)}
\hand{\unsure{.}} \var{\over{w}{W}}{\hand{W}}{writing over}{ps_134-7_024-v0003}hereas the map on sections
\[ X_{*}(k)\hand{\longrightarrow} X'_{*}(k') \]
does induce a map
\[ \text{B')} \qquad \hand{\pi}_i(X_{*})\var{\deleted{\illegible{1}}}{}{deletion}{ps_134-7_024-v0004}\hand{\simeq}\hand{\pi}_i(X_{*}(k)) \hand{\longrightarrow} \hand{\pi}_i(X'_{*})\hand{\simeq}\hand{\pi}_i(X'_{*}(k')) \]
which surely is k-linear, and \var{\inserted{hence}}{\hand{hence}}{insertion}{ps_134-7_024-v0005} induces a map B), this map is certainly
\uline{not} an isomorphism without some flatness restriction either on k' over k,
or on the k-modules $\hand{\pi}_j(X_{*})$ for $\hand{j<i}$, as we had noted already three days
ago when looking at the case when $X_{*}$ comes from a chain complex $M_{\hand{\bullet}}$ in
$(\mathrm{Ab}_k)$ (with projective coefficients say), and hence $X'_{*}$ comes from
\[ M_{\var{\over{*}{\bullet}}{\hand{\bullet}}{writing over}{ps_134-7_024-v0006}}' = M_{.}\otimes_k k' \quad . \]
If we look at the description of the k-\uline{module} $\hand{\pi}_i(X_{*}) = \hand{\pi}_i$ in terms of a
Postnikov dévissage of $X_{*}$ , we should recall that the \var{\xout{group object}}{}{typewriter deletion}{ps_134-7_024-v0007}
semisimplicial gro\var{\xout{y}}{}{typewriter deletion}{ps_134-7_024-v0008}up object $M(i)_{*}$ which inters into the picture as the
i.th step fiber is \uline{not} the one defined \uline{directly} (via the Kan-Dold-Puppe
functor) by $\hand{\pi}_i$ placed in degree i , but \var{\inserted{rather}}{\hand{rather}}{insertion}{ps_134-7_024-v0009} by the chain complex \var{\xout{with}}{}{typewriter deletion}{ps_134-7_024-v0010}
$M(i)_{\var{\over{*}{\bullet}}{\hand{\bullet}}{writing over}{ps_134-7_024-v0011}}$ \uline{with projective components}, obtained by taking \var{\inserted[typed]{first}}{first}{typed insertion}{ps_134-7_024-v0012} a shifted projec-%
tive resolution of $\hand{\pi}_i$ . Thus, by ring extension we get from this dévis-%
sage of $X_{*}$ another one of $X'_{*}$ , whose successive fibers are
\[ M(i)'_{*} = M(i)_{*}\otimes_k k' \quad , \]
\var{\xout{with canonic augmentations}}{}{typewriter deletion}{ps_134-7_024-v0013}
\var{\inserted[typed]{corresponding to the chain complexes}}{corresponding to the chain complexes}{typed insertion}{ps_134-7_024-v0014}
\[ M(i)_{.}' = M(i)_{.}\otimes_k k' \quad . \]

\margin{24}

% Visual collation: both supplied image halves inspected; see CHECK_DECISIONS_134_7_8.tsv.

%% END SOURCE ps_134-7_024
\ifnum\EditionMode=1
\psapparatusstart
\psentry{ps_134-7_024-v0001}{1}{insertion}{\inserted{for}}{\hand{for}}{The handwritten insertion is incorporated at its indicated position; insertion brackets leave the running text.}
\psentry{ps_134-7_024-v0002}{2}{typed insertion}{\inserted[typed]{$X'_{*}=$}}{$X'_{*}=$}{The typed insertion is incorporated at its indicated position; insertion brackets leave the running text.}
\psentry{ps_134-7_024-v0003}{3}{writing over}{\over{w}{W}}{\hand{W}}{The superposed handwritten reading replaces the underlying reading.}
\psentry{ps_134-7_024-v0004}{4}{deletion}{$\deleted{\illegible{1}}$}{}{The author cancels this material by hand; it leaves the corrected running text.}
\psentry{ps_134-7_024-v0005}{5}{insertion}{\inserted{hence}}{\hand{hence}}{The handwritten insertion is incorporated at its indicated position; insertion brackets leave the running text.}
\psentry{ps_134-7_024-v0006}{6}{writing over}{$\over{*}{\bullet}$}{$\hand{\bullet}$}{The superposed handwritten reading replaces the underlying reading.}
\psentry{ps_134-7_024-v0007}{7}{typewriter deletion}{\xout{group object}}{}{The typewriter strikes out this material; it leaves the corrected running text.}
\psentry{ps_134-7_024-v0008}{8}{typewriter deletion}{\xout{y}}{}{The typewriter strikes out this material; it leaves the corrected running text.}
\psentry{ps_134-7_024-v0009}{9}{insertion}{\inserted{rather}}{\hand{rather}}{The handwritten insertion is incorporated at its indicated position; insertion brackets leave the running text.}
\psentry{ps_134-7_024-v0010}{10}{typewriter deletion}{\xout{with}}{}{The typewriter strikes out this material; it leaves the corrected running text.}
\psentry{ps_134-7_024-v0011}{11}{writing over}{$\over{*}{\bullet}$}{$\hand{\bullet}$}{The superposed handwritten reading replaces the underlying reading.}
\psentry{ps_134-7_024-v0012}{12}{typed insertion}{\inserted[typed]{first}}{first}{The typed insertion is incorporated at its indicated position; insertion brackets leave the running text.}
\psentry{ps_134-7_024-v0013}{13}{typewriter deletion}{\xout{with canonic augmentations}}{}{The typewriter strikes out this material; it leaves the corrected running text.}
\psentry{ps_134-7_024-v0014}{14}{typed insertion}{\inserted[typed]{corresponding to the chain complexes}}{corresponding to the chain complexes}{The typed insertion is incorporated at its indicated position; insertion brackets leave the running text.}
\psapparatusend
\fi
\pssource{ps_134-7_025}{134-7/25}{https://grothendieck.umontpellier.fr/web/viewer.html?file=../134-7.pdf\#page=26}
\psnumber{ps_134-7_025-v0001}{1}
\psnumber{ps_134-7_025-v0002}{2}
\psnumber{ps_134-7_025-v0003}{3}
\psnumber{ps_134-7_025-v0004}{4}
\psnumber{ps_134-7_025-v0005}{5}
\psnumber{ps_134-7_025-v0006}{6}
\psnumber{ps_134-7_025-v0007}{7}
\psnumber{ps_134-7_025-v0008}{8}
\psnumber{ps_134-7_025-v0009}{9}
\psnumber{ps_134-7_025-v0010}{10}
\psnumber{ps_134-7_025-v0011}{11}
\psnumber{ps_134-7_025-v0012}{12}
\psnumber{ps_134-7_025-v0013}{13}
%% BEGIN SOURCE ps_134-7_025
% Pursuing Stacks (cotes 134-2 to 134-8): archive cote 134-7, site page 25 (PDF page 26), typescript page 467.
% typescript. Draft compared: ps6.tex:1160-1224.
% AI-transcribed by Claude Opus 5.5 (claude-opus-5-5, adaptive thinking, effort high), Message Batches API, batch ps3 (msgbatch_01XdFmuVVBAcgc7EFxdtCqsj), 2026-11-12; reviewed against both supplied scan halves; see CHECK_DECISIONS.tsv.
\scanpage{134-7/25}
\mspage{467}

The latter has \var{\xout{a canonical augmentation}}{}{typewriter deletion}{ps_134-7_025-v0001} $\hand{\pi}_i\otimes_k k'$ as homology module in
degree i (and zero homology in degree $j\hand{<}i$), but the homology modules
\var{\xout{\illegible{1}}}{}{typewriter deletion}{ps_134-7_025-v0002} \var{\xout{\illegible{1}}}{}{typewriter deletion}{ps_134-7_025-v0003}
in degree $j\hand{>}i$ need not be zero. In other words, there is a canonical
\uline{augmentation}
B'') \hfill \var{\xout{\unsure{X'xM(i)'x}}}{}{typewriter deletion}{ps_134-7_025-v0004} $M(i)'_{*}\hand{\longrightarrow} K(\hand{\pi}_i\otimes_k k',i)$ \hfill ,
(where the second member is the semisimplicial \var{\xout{object}}{}{typewriter deletion}{ps_134-7_025-v0005} k'-module defined
by a i-shifted projective resolution of $\hand{\pi}_i(X)\otimes_k k'$), but this augmentation
need not be a quasiisomorphism, unless we make the relevant flatness as-%
sumptions. To sum up, \uline{the dévissage of} \var{\xout{\illegible{1}}}{}{typewriter deletion}{ps_134-7_025-v0006}
\[ X'_{*} = X_{*}\otimes_k k' \]
\uline{deduced from a Postnikov dévissage of $X_{*}$ is} \uuline{not} \uline{a Postnikov dévissage of}
$X'_{*}$ , unless we assume either k' flat over k, or the k-modules $\hand{\pi}_{\var{\over{i}{j}}{\hand{j}}{writing over}{ps_134-7_025-v0007}}(X_{*})$
flat. This at the same time solves the \uline{puzzle} raised on page 444 , and
points towards a \uline{serious shortcoming of the} \var{\inserted[typed]{\hand{(}usual\hand{)}}}{\hand{(}usual\hand{)}}{typed insertion}{ps_134-7_025-v0008} \uline{Postnikov dévissage} - \uline{namely}
\uline{that it is not compatible with ground ring extension}, or, as we would
say in the language of algebraic geometry, that this construction is not
``geometric'' - a harsh thing to say indeed !

At this point the idea comes up that we may define another devissage,
a lot more natural in the spirit of a theory of ``abelianization'' of homo-%
topy types it would seem, and which is ``geometric'' namely compatible
with ring extension. Here we'll have to work though with the ``prohibiti-%
ve'' abelianization functor
\[ U(k)\hand{\longrightarrow} Ab_k \quad , \qquad X\hand{\longmapsto} L(X) \quad ( \hand{=} \hand{\Gamma}^{\hand{\wedge}}_k(M) ) \]
(where M is a k-module ``representing'' the object X in U(k)), as \var{\deleted{we want}}{}{deletion}{ps_134-7_025-v0009}
we'll need the functorial embedding
\margin{X} C) \hfill $X\hand{\longrightarrow} W(L(X))$ \hfill
(where
\[ W:\ Ab_k\hand{\longrightarrow} U(k) \]
denotes the canonical functor from k-modules to unipotent k-bundles);
this is licit anyhow, if we admit that \var{\xout{under}}{}{typewriter deletion}{ps_134-7_025-v0010} the canonical map
\[ L_{pt}(X_{*})\hand{\longrightarrow} L(X_{*}) \]
(cf. page 454 (w)) is a quasiisomorphism i.e. a weak equivalence, for
the complexes $X_{*}$ we are working with. (This of course should hold over
an arbitrary ground ring, not just $\underline{Z}$.) Applying (C) componentwise, we
get for a complex of bundles $X_{*}$ a canonical map into \var{\deleted{the}}{}{deletion}{ps_134-7_025-v0011} \var{\inserted{it's}}{\hand{it's}}{insertion}{ps_134-7_025-v0012} \var{\xout{abelianized}}{}{typewriter deletion}{ps_134-7_025-v0013}

\margin{25}

% Visual collation: both supplied image halves inspected; see CHECK_DECISIONS_134_7_8.tsv.

%% END SOURCE ps_134-7_025
\ifnum\EditionMode=1
\psapparatusstart
\psentry{ps_134-7_025-v0001}{1}{typewriter deletion}{\xout{a canonical augmentation}}{}{The typewriter strikes out this material; it leaves the corrected running text.}
\psentry{ps_134-7_025-v0002}{2}{typewriter deletion}{\xout{\illegible{1}}}{}{The typewriter strikes out this material; it leaves the corrected running text.}
\psentry{ps_134-7_025-v0003}{3}{typewriter deletion}{\xout{\illegible{1}}}{}{The typewriter strikes out this material; it leaves the corrected running text.}
\psentry{ps_134-7_025-v0004}{4}{typewriter deletion}{\xout{\unsure{X'xM(i)'x}}}{}{The typewriter strikes out this material; it leaves the corrected running text.}
\psentry{ps_134-7_025-v0005}{5}{typewriter deletion}{\xout{object}}{}{The typewriter strikes out this material; it leaves the corrected running text.}
\psentry{ps_134-7_025-v0006}{6}{typewriter deletion}{\xout{\illegible{1}}}{}{The typewriter strikes out this material; it leaves the corrected running text.}
\psentry{ps_134-7_025-v0007}{7}{writing over}{$\over{i}{j}$}{$\hand{j}$}{The superposed handwritten reading replaces the underlying reading.}
\psentry{ps_134-7_025-v0008}{8}{typed insertion}{\inserted[typed]{\hand{(}usual\hand{)}}}{\hand{(}usual\hand{)}}{The typed insertion is incorporated at its indicated position; insertion brackets leave the running text.}
\psentry{ps_134-7_025-v0009}{9}{deletion}{\deleted{we want}}{}{The author cancels this material by hand; it leaves the corrected running text.}
\psentry{ps_134-7_025-v0010}{10}{typewriter deletion}{\xout{under}}{}{The typewriter strikes out this material; it leaves the corrected running text.}
\psentry{ps_134-7_025-v0011}{11}{deletion}{\deleted{the}}{}{The author cancels this material by hand; it leaves the corrected running text.}
\psentry{ps_134-7_025-v0012}{12}{insertion}{\inserted{it's}}{\hand{it's}}{The handwritten insertion is incorporated at its indicated position; insertion brackets leave the running text.}
\psentry{ps_134-7_025-v0013}{13}{typewriter deletion}{\xout{abelianized}}{}{The typewriter strikes out this material; it leaves the corrected running text.}
\psapparatusend
\fi
\pssource{ps_134-7_026}{134-7/26}{https://grothendieck.umontpellier.fr/web/viewer.html?file=../134-7.pdf\#page=27}
\psnumber{ps_134-7_026-v0001}{1}
\psnumber{ps_134-7_026-v0002}{2}
\psnumber{ps_134-7_026-v0003}{3}
\psnumber{ps_134-7_026-v0004}{4}
\psnumber{ps_134-7_026-v0005}{5}
\psnumber{ps_134-7_026-v0006}{6}
\psnumber{ps_134-7_026-v0007}{7}
\psnumber{ps_134-7_026-v0008}{8}
\psnumber{ps_134-7_026-v0009}{9}
\psnumber{ps_134-7_026-v0010}{10}
\psnumber{ps_134-7_026-v0011}{11}
\psnumber{ps_134-7_026-v0012}{12}
\psnumber{ps_134-7_026-v0013}{13}
\psnumber{ps_134-7_026-v0014}{14}
%% BEGIN SOURCE ps_134-7_026
% Pursuing Stacks (cotes 134-2 to 134-8): archive cote 134-7, site page 26 (PDF page 27), typescript page 468.
% typescript. Draft compared: ps6.tex:1208-1266.
% AI-transcribed by Claude Opus 5.5 (claude-opus-5-5, adaptive thinking, effort high), Message Batches API, batch ps3 (msgbatch_01XdFmuVVBAcgc7EFxdtCqsj), 2026-11-12; reviewed against both supplied scan halves; see CHECK_DECISIONS.tsv.
\scanpage{134-7/26}
\mspage{468}

abelianization
\[ \text{(C')} \qquad X_{*} \hand{\longrightarrow} W(L(X_{*})) \quad . \]
Postnikovs construction, for \var{\deleted{\unsure{n-}}}{}{deletion}{ps_134-7_026-v0001}\var{\inserted{$(n-1)$-}}{\hand{$(n-1)$-}}{insertion}{ps_134-7_026-v0002}connected $X_{*}$ , consists in composing
this map with the "augmentation"
\[ W(L(X_{*})) \hand{\longrightarrow} K(\hand{\pi}_n,\hand{n}) \var{\deleted{\unsure{.}}}{}{deletion}{ps_134-7_026-v0003} \]
\hand{(}where $\hand{\pi}_n\hand{\simeq \mathrm{H}_n}$ is the first possibly non trivial homology module of the
chain complex corresponding to $L(X_{*})$\hand{)}, and after this \var{\inserted{only}}{\hand{only}}{insertion}{ps_134-7_026-v0004} take the homotopy
fiber, and iterate (the homotopy fiber will be n-connected now). This
process, in a way, breaks the natural abelianizations \var{\xout{\illegible{1}}}{}{typewriter deletion}{ps_134-7_026-v0005} into
pieces, a \var{\xout{little}}{}{typewriter deletion}{ps_134-7_026-v0006} brutal thing to do one will admit, all the more so as
we start with a beautiful complex with projective components, and kind
of destroy it's unmarred harmony by tearing out or breaking off the most
showy part, $\mathrm{H}_n$ to name it, which now looks so lost and awkward we reall\hand{y}
can't just leave it as it is, we first have to take a projective resolu-%
tion of it, choosing it as we may ... \var{\inserted[typed]{But}}{But}{typed insertion}{ps_134-7_026-v0007} We wo'nt do \var{\inserted[typed]{all}}{all}{typed insertion}{ps_134-7_026-v0008} this, will we, \var{\over{\unsure{and}}{and}}{\hand{and}}{writing over}{ps_134-7_026-v0009}
rather keep abelianization and \hand{(}C'\hand{)} as God gave them to us, and take the
homotopy fiber (\unsure{We|we} hope God will give this too ...), and repeat the
process, without even having to care at any stage which $\hand{\pi}_i$\hand{'s} vanish and
which not. Let's call this the "\uline{soft Postnikov dévissage}" , in contrast
to the"brutal"one. In describing the process (which of course makes a
sense in the discrete context as well as in the schematic one), I impli-%
citly admitted that the $X_{*}$ we start with is 1-connected, or for the
very least has abelian $\hand{\pi}_1$ (a notion we'll have to come back to, in the
schematic \var{\xout{\illegible{1}}}{}{typewriter deletion}{ps_134-7_026-v0010} set up). But we may as well apply it to a (discrete)
$K(G,1)$ type, $G$ any discrete group, then it amounts to tak\var{\over{e}{ing}}{\hand{ing}}{writing over}{ps_134-7_026-v0011} the descen-%
ding filtration of $G$ by iterated commutator groups, which is a finite
filtration iff $G$ is solvable. Maybe it would be more natural still to
take the similar descending filtration, suitable for the study of \uline{nilpo-}%
\uline{tent} groups rather than solvable one's, with
\margin{(D)}
\[ G^{(n+1)} = \hand{[}G,\ G^{(n+1)}\hand{]}, \]
where \hand{$[$}A,B\hand{$]$} denotes the subgroup of $G$ generated by commutators \var{\xout{\illegible{1}}}{}{typewriter deletion}{ps_134-7_026-v0012}
\[ (a,b) = aba^{-1}b^{-1}\hand{,} \]
with a in A, b in B. It does'nt seem there is a similar distinction to
make in case we start with a \var{\xout{simply}}{}{typewriter deletion}{ps_134-7_026-v0013} 1-connected $X_{*}$ , more specifically
if $X_0=X_1= e$. There may be some \var{\inserted[typed]{extra}}{extra}{typed insertion}{ps_134-7_026-v0014} caution needed however when we assume
only $X_0= e$ without assuming 1-connectedness, even when $\hand{\pi}_1$ is abelian,

\margin{26}

% Visual collation: both supplied image halves inspected; see CHECK_DECISIONS_134_7_8.tsv.

%% END SOURCE ps_134-7_026
\ifnum\EditionMode=1
\psapparatusstart
\psentry{ps_134-7_026-v0001}{1}{deletion}{\deleted{\unsure{n-}}}{}{The author cancels this material by hand; it leaves the corrected running text.}
\psentry{ps_134-7_026-v0002}{2}{insertion}{\inserted{$(n-1)$-}}{\hand{$(n-1)$-}}{The handwritten insertion is incorporated at its indicated position; insertion brackets leave the running text.}
\psentry{ps_134-7_026-v0003}{3}{deletion}{$\deleted{\unsure{.}}$}{}{The author cancels this material by hand; it leaves the corrected running text.}
\psentry{ps_134-7_026-v0004}{4}{insertion}{\inserted{only}}{\hand{only}}{The handwritten insertion is incorporated at its indicated position; insertion brackets leave the running text.}
\psentry{ps_134-7_026-v0005}{5}{typewriter deletion}{\xout{\illegible{1}}}{}{The typewriter strikes out this material; it leaves the corrected running text.}
\psentry{ps_134-7_026-v0006}{6}{typewriter deletion}{\xout{little}}{}{The typewriter strikes out this material; it leaves the corrected running text.}
\psentry{ps_134-7_026-v0007}{7}{typed insertion}{\inserted[typed]{But}}{But}{The typed insertion is incorporated at its indicated position; insertion brackets leave the running text.}
\psentry{ps_134-7_026-v0008}{8}{typed insertion}{\inserted[typed]{all}}{all}{The typed insertion is incorporated at its indicated position; insertion brackets leave the running text.}
\psentry{ps_134-7_026-v0009}{9}{writing over}{\over{\unsure{and}}{and}}{\hand{and}}{The superposed handwritten reading replaces the underlying reading.}
\psentry{ps_134-7_026-v0010}{10}{typewriter deletion}{\xout{\illegible{1}}}{}{The typewriter strikes out this material; it leaves the corrected running text.}
\psentry{ps_134-7_026-v0011}{11}{writing over}{\over{e}{ing}}{\hand{ing}}{The superposed handwritten reading replaces the underlying reading.}
\psentry{ps_134-7_026-v0012}{12}{typewriter deletion}{\xout{\illegible{1}}}{}{The typewriter strikes out this material; it leaves the corrected running text.}
\psentry{ps_134-7_026-v0013}{13}{typewriter deletion}{\xout{simply}}{}{The typewriter strikes out this material; it leaves the corrected running text.}
\psentry{ps_134-7_026-v0014}{14}{typed insertion}{\inserted[typed]{extra}}{extra}{The typed insertion is incorporated at its indicated position; insertion brackets leave the running text.}
\psapparatusend
\fi
\pssource{ps_134-7_027}{134-7/27}{https://grothendieck.umontpellier.fr/web/viewer.html?file=../134-7.pdf\#page=28}
\psnumber{ps_134-7_027-v0001}{1}
\psnumber{ps_134-7_027-v0002}{2}
\psnumber{ps_134-7_027-v0003}{3}
\psnumber{ps_134-7_027-v0004}{4}
\psnumber{ps_134-7_027-v0005}{5}
\psnumber{ps_134-7_027-v0006}{6}
\psnumber{ps_134-7_027-v0007}{7}
\psnumber{ps_134-7_027-v0008}{8}
\psnumber{ps_134-7_027-v0009}{9}
\psnumber{ps_134-7_027-v0010}{10}
\psnumber{ps_134-7_027-v0011}{11}
\psnumber{ps_134-7_027-v0012}{12}
\psnumber{ps_134-7_027-v0013}{13}
\psnumber{ps_134-7_027-v0014}{14}
\psnumber{ps_134-7_027-v0015}{15}
\psnumber{ps_134-7_027-v0016}{16}
\psnumber{ps_134-7_027-v0017}{17}
\psnumber{ps_134-7_027-v0018}{18}
\psnumber{ps_134-7_027-v0019}{19}
%% BEGIN SOURCE ps_134-7_027
% Pursuing Stacks (cotes 134-2 to 134-8): archive cote 134-7, site page 27 (PDF page 28), typescript page 469.
% typescript. Draft compared: ps6.tex:1253-1313.
% AI-transcribed by Claude Opus 5.5 (claude-opus-5-5, adaptive thinking, effort high), Message Batches API, batch ps3 (msgbatch_01XdFmuVVBAcgc7EFxdtCqsj), 2026-11-12; reviewed against both supplied scan halves; see CHECK_DECISIONS.tsv.
\scanpage{134-7/27}
\mspage{469}

because of the possibility of operation of \hand{$\pi$}$_1$ upon the \hand{$\pi$}$_i$'s. Maybe,
when trying to modelize usual homo\var{\xout{b}}{}{typewriter deletion}{ps_134-7_027-v0001}to\var{\xout{x}}{}{typewriter deletion}{ps_134-7_027-v0002}py types by complexes of unipotent
bundles over $\underline{Z}$, we should restrict to homotopy types which are not only
\var{\xout{zer}}{}{typewriter deletion}{ps_134-7_027-v0003}0-connected and have abelian \hand{$\pi$}$_1$ , but moreover with \hand{$\pi$}$_1$ operating
trivially (or for the very least, in a unipotent way) upon the higher
\hand{$\pi$}$_i$'s. At any rate, as soon as \hand{$\pi$}$_1$ operates non trivially (on itself, or
on the higher \hand{$\pi$}$_i$'s ) there will presumably be two non equivalent ways for
defining soft Postnikov dévissage, corresponding to the two standard
descending commutator group series in a discrete group G. The more relevant
\var{\xout{xxx}}{}{typewriter deletion}{ps_134-7_027-v0004} in view of unipotent schematization would seem to be the ``nilpotent''
one.

Restricting for simplicity to the 1-connected case \unsure{$X_\hand{\bullet}$|$X_0$}$=X_1=e$ , I would
expect soft Postnikov dévissage to be the key \hand{f}or an understanding as well
of the behaviour of the \hand{$\pi$}$_i(X_*)$ modules with respect to ring extension,
as \var{\xout{xxx}}{}{typewriter deletion}{ps_134-7_027-v0005}\var{\inserted[typed]{of}}{of}{typed insertion}{ps_134-7_027-v0006} the full relationship between these invarian ts, and the homology
invarian ts $H_i(X_*)$.

$*$ \quad $*$

$*$

I still should have a look upon complexes $X_*$ in $U(k)^\hand{\bullet}$\hand{/}satisfying
\hand{(}as always in this game\hand{)} $X_0=e$ , but not necessarily $X_1 = e$. Even when
we make the Kan assumption (\var{\deleted{including}}{}{deletion}{ps_134-7_027-v0007}\var{\inserted{plus}}{\hand{plus}}{insertion}{ps_134-7_027-v0008} ``\uline{smoothness}'' of the components,
by which I mean that their k-linearization is projective), I do'nt feel
too sure yet if they fit into a good formalism, for instance \hand{(}when $k=\underline{Z}$\hand{)}
if they satisfy the ``linearization theorem'' (quasiisomorphy for the two
maps in (w) p.454). If we start for instance with a group object\var{\inserted[typed]{G}}{G}{typed insertion}{ps_134-7_027-v0009} of U(k)
and let $X_*$ be the corresponding semisimplicial complex in $U(k)^\hand{\bullet}$ , then
we get an isomorphism
\[ \text{(E)} \qquad \hand{\pi}_1(X_*) \overset{\text{def}}{=} \hand{\pi}_1(X_*(k)) \hand{\simeq} G(k) \quad , \]
which shows us that \hand{$\pi$}$_1(X_*)$ need not be abelian even when $k=\underline{Z}$ , \var{\deleted{all we \xout{sa}}}{}{deletion}{ps_134-7_027-v0010}
\var{\deleted{can safely say,}}{}{deletion}{ps_134-7_027-v0011} \var{\over{i}{I}}{\hand{I}}{writing over}{ps_134-7_027-v0012}f we assume that $X_1 = G$ is ``of finite presentation''
(namely the \var{\inserted[typed]{projective}}{projective}{typed insertion}{ps_134-7_027-v0013} k-module which describes $X_1$ is of finite presentation), \var{\deleted{is}}{}{deletion}{ps_134-7_027-v0014}\var{\inserted{or}}{\hand{or}}{insertion}{ps_134-7_027-v0015}
what amounts to the same, representable by an actual (group) \unsure{scheme|shheme}, \var{\inserted{it}}{\hand{it}}{insertion}{ps_134-7_027-v0016}
is \var{\inserted{true however}}{\hand{true however}}{insertion}{ps_134-7_027-v0017} that G and hence \hand{$\pi$}$_1 = G(k)$ is \uline{nilpotent} (this holds for any k ).
It looks an intriguing question whether \hand{$\pi$}$_1$ is nilpotent under the only
assumption that $X_*$ is a \var{\inserted{smooth}}{\hand{smooth}}{insertion}{ps_134-7_027-v0018} Kan complex with $X_1$ a scheme (wi\var{\over{th}{th}}{\hand{th}}{writing over}{ps_134-7_027-v0019}out assuming
any more X comes from a group object). At any rate, it follows from

\margin{27}

% Visual collation: both supplied image halves inspected; see CHECK_DECISIONS_134_7_8.tsv.

%% END SOURCE ps_134-7_027
\ifnum\EditionMode=1
\psapparatusstart
\psentry{ps_134-7_027-v0001}{1}{typewriter deletion}{\xout{b}}{}{The typewriter strikes out this material; it leaves the corrected running text.}
\psentry{ps_134-7_027-v0002}{2}{typewriter deletion}{\xout{x}}{}{The typewriter strikes out this material; it leaves the corrected running text.}
\psentry{ps_134-7_027-v0003}{3}{typewriter deletion}{\xout{zer}}{}{The typewriter strikes out this material; it leaves the corrected running text.}
\psentry{ps_134-7_027-v0004}{4}{typewriter deletion}{\xout{xxx}}{}{The typewriter strikes out this material; it leaves the corrected running text.}
\psentry{ps_134-7_027-v0005}{5}{typewriter deletion}{\xout{xxx}}{}{The typewriter strikes out this material; it leaves the corrected running text.}
\psentry{ps_134-7_027-v0006}{6}{typed insertion}{\inserted[typed]{of}}{of}{The typed insertion is incorporated at its indicated position; insertion brackets leave the running text.}
\psentry{ps_134-7_027-v0007}{7}{deletion}{\deleted{including}}{}{The author cancels this material by hand; it leaves the corrected running text.}
\psentry{ps_134-7_027-v0008}{8}{insertion}{\inserted{plus}}{\hand{plus}}{The handwritten insertion is incorporated at its indicated position; insertion brackets leave the running text.}
\psentry{ps_134-7_027-v0009}{9}{typed insertion}{\inserted[typed]{G}}{G}{The typed insertion is incorporated at its indicated position; insertion brackets leave the running text.}
\psentry{ps_134-7_027-v0010}{10}{deletion}{\deleted{all we \xout{sa}}}{}{The author cancels this material by hand; it leaves the corrected running text.}
\psentry{ps_134-7_027-v0011}{11}{deletion}{\deleted{can safely say,}}{}{The author cancels this material by hand; it leaves the corrected running text.}
\psentry{ps_134-7_027-v0012}{12}{writing over}{\over{i}{I}}{\hand{I}}{The superposed handwritten reading replaces the underlying reading.}
\psentry{ps_134-7_027-v0013}{13}{typed insertion}{\inserted[typed]{projective}}{projective}{The typed insertion is incorporated at its indicated position; insertion brackets leave the running text.}
\psentry{ps_134-7_027-v0014}{14}{deletion}{\deleted{is}}{}{The author cancels this material by hand; it leaves the corrected running text.}
\psentry{ps_134-7_027-v0015}{15}{insertion}{\inserted{or}}{\hand{or}}{The handwritten insertion is incorporated at its indicated position; insertion brackets leave the running text.}
\psentry{ps_134-7_027-v0016}{16}{insertion}{\inserted{it}}{\hand{it}}{The handwritten insertion is incorporated at its indicated position; insertion brackets leave the running text.}
\psentry{ps_134-7_027-v0017}{17}{insertion}{\inserted{true however}}{\hand{true however}}{The handwritten insertion is incorporated at its indicated position; insertion brackets leave the running text.}
\psentry{ps_134-7_027-v0018}{18}{insertion}{\inserted{smooth}}{\hand{smooth}}{The handwritten insertion is incorporated at its indicated position; insertion brackets leave the running text.}
\psentry{ps_134-7_027-v0019}{19}{writing over}{\over{th}{th}}{\hand{th}}{The superposed handwritten reading replaces the underlying reading.}
\psapparatusend
\fi
\pssource{ps_134-7_028}{134-7/28}{https://grothendieck.umontpellier.fr/web/viewer.html?file=../134-7.pdf\#page=29}
\psnumber{ps_134-7_028-v0001}{1}
\psnumber{ps_134-7_028-v0002}{2}
\psnumber{ps_134-7_028-v0003}{3}
\psnumber{ps_134-7_028-v0004}{4}
\psnumber{ps_134-7_028-v0005}{5}
\psnumber{ps_134-7_028-v0006}{6}
\psnumber{ps_134-7_028-v0007}{7}
\psnumber{ps_134-7_028-v0008}{8}
\psnumber{ps_134-7_028-v0009}{9}
\psnumber{ps_134-7_028-v0010}{10}
\psnumber{ps_134-7_028-v0011}{11}
\psnumber{ps_134-7_028-v0012}{12}
\psnumber{ps_134-7_028-v0013}{13}
\psnumber{ps_134-7_028-v0014}{14}
\psnumber{ps_134-7_028-v0015}{15}
\psnumber{ps_134-7_028-v0016}{16}
\psnumber{ps_134-7_028-v0017}{17}
\psnumber{ps_134-7_028-v0018}{18}
\psnumber{ps_134-7_028-v0019}{19}
\psnumber{ps_134-7_028-v0020}{20}
\psnumber{ps_134-7_028-v0021}{21}
\psnumber{ps_134-7_028-v0022}{22}
\psnumber{ps_134-7_028-v0023}{23}
\psnumber{ps_134-7_028-v0024}{24}
%% BEGIN SOURCE ps_134-7_028
% Pursuing Stacks (cotes 134-2 to 134-8): archive cote 134-7, site page 28 (PDF page 29), typescript page 470.
% typescript. Draft compared: ps6.tex:1298-1359.
% AI-transcribed by Claude Opus 5.5 (claude-opus-5-5, adaptive thinking, effort high), Message Batches API, batch ps3 (msgbatch_01XdFmuVVBAcgc7EFxdtCqsj), 2026-11-12; reviewed against both supplied scan halves; see CHECK_DECISIONS.tsv.
\scanpage{134-7/28}
\mspage{470}

\var{\deleted{from}}{}{deletion}{ps_134-7_028-v0001} the Kan condition that $\hand{\pi}_1$ may be interpreted as a quotient set of \var{\xout{X}}{}{typewriter deletion}{ps_134-7_028-v0002}
\var{\inserted[typed]{$E \overset{\mathrm{def}}{=}$}}{$E \overset{\mathrm{def}}{=}$}{typed insertion}{ps_134-7_028-v0003}$X_1(k)$ (without having to pass to the full free group generated by this
set), with a set of relations
\[ \text{(F)} \qquad z_i = x_iy_i \quad , \quad i \text{ in } X_2(k)=I, \]
indexed by $X_2(k)$, where
\[ i\hand{\mapsto} x_i \ , \ i\hand{\mapsto} y_i \ , \ i\hand{\mapsto} z_i \]
are the three boundary maps, remembering moreover \var{\inserted[typed]{the Kan condition}}{the Kan condition}{typed insertion}{ps_134-7_028-v0004} that the three maps
% Hand-drawn arrows are marked blue; typed objects remain plain.
\[ \diagram{\begin{tikzcd}[cramped,sep=small] I \ar[color=blue,r]\ar[color=blue,r,shift left=1.5]\ar[color=blue,r,shift right=1.5] \& E\mathrm{x}E \end{tikzcd}} \ , \quad i\hand{\mapsto}(x_i,y_i) \ , \quad i\hand{\mapsto}(x_i,z_i) \ , \quad i\hand{\mapsto}(y_i,z_i) \]
are \uline{surjective} , which implies indeed that \var{\xout{xxx}}{}{typewriter deletion}{ps_134-7_028-v0005} any element of the group \var{\inserted{$\pi$}}{\hand{$\pi$}}{insertion}{ps_134-7_028-v0006}
described by the set of generators E and relations (F) \var{\xout{xx}}{}{typewriter deletion}{ps_134-7_028-v0007} comes from an
element in E. Replacing k by any k-algebra k', we see that we have a
presheaf\var{\xout{\illegible{1}}}{}{typewriter deletion}{ps_134-7_028-v0008}
\[ k'\hand{\mapsto} \hand{\pi}_1(X_*(k')) = \hand{\pi}_1(X_*\otimes_k k') \]
on the category $\mathrm{Aff}_{/k}$ of affine schemes over k, with values in the
category of groups, which may be viewed (as a presheaf of sets) as a
quotient presheaf of the presheaf on $\mathrm{Aff}_{/k}$ defined by $X_1$ . \var{\xout{In order}}{}{typewriter deletion}{ps_134-7_028-v0009} We
feel that this \var{\xout{\illegible{1}}}{}{typewriter deletion}{ps_134-7_028-v0010}\var{\inserted[typed]{presheaf}}{presheaf}{typed insertion}{ps_134-7_028-v0011}will fit into a reasonable ``schematic'' set up,
only if it turns out to be a \uline{sheaf}, and \var{\xout{even}}{}{typewriter deletion}{ps_134-7_028-v0012} more exactingly still, if
this sheaf is \var{\xout{representable}}{}{typewriter deletion}{ps_134-7_028-v0013}\var{\inserted[typed]{isomorphic}}{isomorphic}{typed insertion}{ps_134-7_028-v0014} (as a sheaf of sets) to one stemming from an
object of U(k) i.e. if it is isomorphic to a sheaf W(M), for suitable
k-module M (not necessarily a projective one). If we denote by G this
object of U(k), it will be endowed with a group structure, and \var{\inserted{it is}}{\hand{it is}}{insertion}{ps_134-7_028-v0015} \uline{this}
\uline{group object of} U(k), \uline{rather than just the} \var{\inserted[typed]{set-theoretic}}{set-theoretic}{typed insertion}{ps_134-7_028-v0016} \uline{group of it's sections} i.e.
\uline{of k-valued ``points''}, \var{\inserted{which}}{\hand{which}}{insertion}{ps_134-7_028-v0017} \uline{merits to be viewed as the ``true''} $\hand{\pi}_1(X_*)$. To say
it differently, whereas the higher $\hand{\pi}_i(X_*) \overset{\mathrm{def}}{=} \hand{\pi}_i(X_*(k))$ \var{\xout{xxx}}{}{typewriter deletion}{ps_134-7_028-v0018} (for $i\hand{\geq} 2$),
in the cases considered so far, should be viewed as being not mere abeli\var{\xout{a}}{}{typewriter deletion}{ps_134-7_028-v0019}
\hand{-}an groups, but moreover endowed with a natural \var{\over{k}{k}}{\hand{k}}{writing over}{ps_134-7_028-v0020}-module structure,
in the case when i=1 i.e. for the fundamental group $\hand{\pi}_1(X_*)$, the natu-%
ral structure to expect on this (possibly non commutative) group is a
\var{\inserted{``unipotent}}{\hand{``unipotent}}{insertion}{ps_134-7_028-v0021} ``schematic'' structure, namely essentially a \var{\inserted[typed]{pointed}}{pointed}{typed insertion}{ps_134-7_028-v0022} ``parametrisation'' of this
group by elements of a suitable k-module $M_1$ , in such a way that the
composition law is expressed in terms of a polynomial law, making sense
therefore not only for k-valued points i.e. \var{\inserted{for}}{\hand{for}}{insertion}{ps_134-7_028-v0023} elements of M, but \var{\inserted{for}}{\hand{for}}{insertion}{ps_134-7_028-v0024} k'-valued
points as well \hand{(}for any k-algebra k'\hand{),} namely defining a group law on

\margin{28}

% Visual collation: both supplied image halves inspected; see CHECK_DECISIONS_134_7_8.tsv.

%% END SOURCE ps_134-7_028
\ifnum\EditionMode=1
\psapparatusstart
\psentry{ps_134-7_028-v0001}{1}{deletion}{\deleted{from}}{}{The author cancels this material by hand; it leaves the corrected running text.}
\psentry{ps_134-7_028-v0002}{2}{typewriter deletion}{\xout{X}}{}{The typewriter strikes out this material; it leaves the corrected running text.}
\psentry{ps_134-7_028-v0003}{3}{typed insertion}{\inserted[typed]{$E \overset{\mathrm{def}}{=}$}}{$E \overset{\mathrm{def}}{=}$}{The typed insertion is incorporated at its indicated position; insertion brackets leave the running text.}
\psentry{ps_134-7_028-v0004}{4}{typed insertion}{\inserted[typed]{the Kan condition}}{the Kan condition}{The typed insertion is incorporated at its indicated position; insertion brackets leave the running text.}
\psentry{ps_134-7_028-v0005}{5}{typewriter deletion}{\xout{xxx}}{}{The typewriter strikes out this material; it leaves the corrected running text.}
\psentry{ps_134-7_028-v0006}{6}{insertion}{\inserted{$\pi$}}{\hand{$\pi$}}{The handwritten insertion is incorporated at its indicated position; insertion brackets leave the running text.}
\psentry{ps_134-7_028-v0007}{7}{typewriter deletion}{\xout{xx}}{}{The typewriter strikes out this material; it leaves the corrected running text.}
\psentry{ps_134-7_028-v0008}{8}{typewriter deletion}{\xout{\illegible{1}}}{}{The typewriter strikes out this material; it leaves the corrected running text.}
\psentry{ps_134-7_028-v0009}{9}{typewriter deletion}{\xout{In order}}{}{The typewriter strikes out this material; it leaves the corrected running text.}
\psentry{ps_134-7_028-v0010}{10}{typewriter deletion}{\xout{\illegible{1}}}{}{The typewriter strikes out this material; it leaves the corrected running text.}
\psentry{ps_134-7_028-v0011}{11}{typed insertion}{\inserted[typed]{presheaf}}{presheaf}{The typed insertion is incorporated at its indicated position; insertion brackets leave the running text.}
\psentry{ps_134-7_028-v0012}{12}{typewriter deletion}{\xout{even}}{}{The typewriter strikes out this material; it leaves the corrected running text.}
\psentry{ps_134-7_028-v0013}{13}{typewriter deletion}{\xout{representable}}{}{The typewriter strikes out this material; it leaves the corrected running text.}
\psentry{ps_134-7_028-v0014}{14}{typed insertion}{\inserted[typed]{isomorphic}}{isomorphic}{The typed insertion is incorporated at its indicated position; insertion brackets leave the running text.}
\psentry{ps_134-7_028-v0015}{15}{insertion}{\inserted{it is}}{\hand{it is}}{The handwritten insertion is incorporated at its indicated position; insertion brackets leave the running text.}
\psentry{ps_134-7_028-v0016}{16}{typed insertion}{\inserted[typed]{set-theoretic}}{set-theoretic}{The typed insertion is incorporated at its indicated position; insertion brackets leave the running text.}
\psentry{ps_134-7_028-v0017}{17}{insertion}{\inserted{which}}{\hand{which}}{The handwritten insertion is incorporated at its indicated position; insertion brackets leave the running text.}
\psentry{ps_134-7_028-v0018}{18}{typewriter deletion}{\xout{xxx}}{}{The typewriter strikes out this material; it leaves the corrected running text.}
\psentry{ps_134-7_028-v0019}{19}{typewriter deletion}{\xout{a}}{}{The typewriter strikes out this material; it leaves the corrected running text.}
\psentry{ps_134-7_028-v0020}{20}{writing over}{\over{k}{k}}{\hand{k}}{The superposed handwritten reading replaces the underlying reading.}
\psentry{ps_134-7_028-v0021}{21}{insertion}{\inserted{``unipotent}}{\hand{``unipotent}}{The handwritten insertion is incorporated at its indicated position; insertion brackets leave the running text.}
\psentry{ps_134-7_028-v0022}{22}{typed insertion}{\inserted[typed]{pointed}}{pointed}{The typed insertion is incorporated at its indicated position; insertion brackets leave the running text.}
\psentry{ps_134-7_028-v0023}{23}{insertion}{\inserted{for}}{\hand{for}}{The handwritten insertion is incorporated at its indicated position; insertion brackets leave the running text.}
\psentry{ps_134-7_028-v0024}{24}{insertion}{\inserted{for}}{\hand{for}}{The handwritten insertion is incorporated at its indicated position; insertion brackets leave the running text.}
\psapparatusend
\fi
\pssource{ps_134-7_029}{134-7/29}{https://grothendieck.umontpellier.fr/web/viewer.html?file=../134-7.pdf\#page=30}
\psnumber{ps_134-7_029-v0001}{1}
\psnumber{ps_134-7_029-v0002}{2}
\psnumber{ps_134-7_029-v0003}{3}
\psnumber{ps_134-7_029-v0004}{4}
\psnumber{ps_134-7_029-v0005}{5}
\psnumber{ps_134-7_029-v0006}{6}
\psnumber{ps_134-7_029-v0007}{7}
\psnumber{ps_134-7_029-v0008}{8}
%% BEGIN SOURCE ps_134-7_029
% Pursuing Stacks (cotes 134-2 to 134-8): archive cote 134-7, site page 29 (PDF page 30), typescript page 471.
% typescript. Draft compared: ps6.tex:1342-1405.
% AI-transcribed by Claude Opus 5.5 (claude-opus-5-5, adaptive thinking, effort high), Message Batches API, batch ps3 (msgbatch_01XdFmuVVBAcgc7EFxdtCqsj), 2026-11-12; reviewed against both supplied scan halves; see CHECK_DECISIONS.tsv.
\scanpage{134-7/29}
\mspage{471}

$W(M)(k') = M\otimes_k k'$ . If henceforth we denote by $\hand{\pi}_1(X_*)$ this group
object of $U(k)$, the relevant formula now is
\[ \text{(G)} \qquad \hand{\pi}_1(X_*(k'))\hand{\simeq} \hand{\pi}_1(X_*)(k') \qquad , \]
a group isomorphism functorial with respect to variable k-algebra k',
which will imply the corresponding isomorphism
\[ \text{(G')} \qquad \hand{\pi}_1(X_*\otimes_k k')\hand{\simeq} \hand{\pi}_1(X_*)\otimes_k k' \]
of group objects in $U(k')$, i.e. formation of the ``schematic'' $\hand{\pi}_1$ is compa-%
tible with ground ring extension $k \quad k'$ (provided \var{\over{it}{$\pi_1$}}{\hand{$\pi_1$}}{writing over}{ps_134-7_029-v0001} exists\var{\deleted{\unsure{x}}}{}{deletion}{ps_134-7_029-v0002}, for a given
$X_*$ ).

We will expect \var{\xout{that}}{}{typewriter deletion}{ps_134-7_029-v0003} the map of passage to quotient
\[ \text{(H)} \qquad X_1\hand{\to}\hand{\pi}_1(X_*) = G \]
to be ``epimorphic'' in a very strong sense, stronger even than just in
the sense of presheaves, the first thought that comes to mind here is
that it should be a``\uline{submersion}'' , in the sense suggested in yesterday's
reflections in connection with the description of the Kan condition. If
however we want to be able to get for $\hand{\pi}_1(X_*(k)) = \hand{\pi}_1(X_*(k))$ \uline{any} abelian
group given beforehand, in the case $k=\underline{Z}$ say, without demanding that it
be a projective k-module, and still get it via an $X_*$ with \uline{smooth}
components, this shows that when defining a notion of ``submersion'' for
objects of $U(k)$ which may not be smooth, we should not be quite as deman-%
ding \var{\deleted{in the definition}}{}{deletion}{ps_134-7_029-v0004} as suggested yesterday (cf. page 463), but find
\var{\deleted{one}}{}{deletion}{ps_134-7_029-v0005} \var{\inserted{a definition}}{\hand{a definition}}{insertion}{ps_134-7_029-v0006} which will include also any map $X \hand{\to} Y$ coming from an epimorphism
$M\hand{\to} N$ of k-modules (which will allow us to take $X_1$ as associated to
a projective k-module admitting the given $\hand{\pi}_1$ as it's quotient). One
idea that comes to mind here, is to take this property as the \uline{definition}
of a submersion, as an arrow in $U(k)$ which \var{\xout{comes from}}{}{typewriter deletion}{ps_134-7_029-v0007} is isomorphic to
one obtained from an epimorphism in $Ab_k$ . This\hand{,} \var{\deleted{at present,}}{}{deletion}{ps_134-7_029-v0008} of the
three definitions that have come to my mind so far for this notion, is
the one which looks the most convincing to me. I would'nt expect too
much from a complex $X_*$ , even a smooth one and satisfying the Kan condi-%
tion, unless (in terms of the three boundary maps from $X_2$ to $X_1$) it
does'nt give rise, as just explained, to a group object $G = \hand{\pi}_1(X_*)$
in $U(k)$, together with a \uline{submersion} (H). Thus definitely, when de-%
fining a shhematic model category $\hand{\mathcal{M}}_n(k)$ of n-connected ss complexes
of unipotent bundles over k, I feel like insisting in case n=0 at least
upon this extra condition (plus of course $X_o = e$). The still more strin-%
gent condition one may think of, in order to have schematic models as

\margin{29}

% Visual collation: both supplied image halves inspected; see CHECK_DECISIONS_134_7_8.tsv.

%% END SOURCE ps_134-7_029
\ifnum\EditionMode=1
\psapparatusstart
\psentry{ps_134-7_029-v0001}{1}{writing over}{\over{it}{$\pi_1$}}{\hand{$\pi_1$}}{The superposed handwritten reading replaces the underlying reading.}
\psentry{ps_134-7_029-v0002}{2}{deletion}{\deleted{\unsure{x}}}{}{The author cancels this material by hand; it leaves the corrected running text.}
\psentry{ps_134-7_029-v0003}{3}{typewriter deletion}{\xout{that}}{}{The typewriter strikes out this material; it leaves the corrected running text.}
\psentry{ps_134-7_029-v0004}{4}{deletion}{\deleted{in the definition}}{}{The author cancels this material by hand; it leaves the corrected running text.}
\psentry{ps_134-7_029-v0005}{5}{deletion}{\deleted{one}}{}{The author cancels this material by hand; it leaves the corrected running text.}
\psentry{ps_134-7_029-v0006}{6}{insertion}{\inserted{a definition}}{\hand{a definition}}{The handwritten insertion is incorporated at its indicated position; insertion brackets leave the running text.}
\psentry{ps_134-7_029-v0007}{7}{typewriter deletion}{\xout{comes from}}{}{The typewriter strikes out this material; it leaves the corrected running text.}
\psentry{ps_134-7_029-v0008}{8}{deletion}{\deleted{at present,}}{}{The author cancels this material by hand; it leaves the corrected running text.}
\psapparatusend
\fi
\pssource{ps_134-7_030}{134-7/30}{https://grothendieck.umontpellier.fr/web/viewer.html?file=../134-7.pdf\#page=31}
\psnumber{ps_134-7_030-v0001}{1}
\psnumber{ps_134-7_030-v0002}{2}
\psnumber{ps_134-7_030-v0003}{3}
\psnumber{ps_134-7_030-v0004}{4}
\psnumber{ps_134-7_030-v0005}{5}
\psnumber{ps_134-7_030-v0006}{6}
\psnumber{ps_134-7_030-v0007}{7}
\psnumber{ps_134-7_030-v0008}{8}
\psnumber{ps_134-7_030-v0009}{9}
\psnumber{ps_134-7_030-v0010}{10}
\psnumber{ps_134-7_030-v0011}{11}
\psnumber{ps_134-7_030-v0012}{12}
\psnumber{ps_134-7_030-v0013}{13}
\psnumber{ps_134-7_030-v0014}{14}
\psnumber{ps_134-7_030-v0015}{15}
\psnumber{ps_134-7_030-v0016}{16}
%% BEGIN SOURCE ps_134-7_030
% Pursuing Stacks (cotes 134-2 to 134-8): archive cote 134-7, site page 30 (PDF page 31), typescript page 472.
% typescript. Draft compared: ps6.tex:1393-1452.
% AI-transcribed by Claude Opus 5.5 (claude-opus-5-5, adaptive thinking, effort high), Message Batches API, batch ps3 (msgbatch_01XdFmuVVBAcgc7EFxdtCqsj), 2026-11-12; reviewed against both supplied scan halves; see CHECK_DECISIONS.tsv.
\scanpage{134-7/30}
\mspage{472}

close as one may w\var{\xout{h}}{}{typewriter deletion}{ps_134-7_030-v0001}ish to k-modules,  is to demand that moreover G is
isomorphic to the object of U(k) defined by a k-module $M_1$ ,  the group
law moreover coming from the addition law in M.  This condition is stron-%
ger still than merely demanding that G be commutative,  even when $M_1$ is
free of rank one,  \var{\deleted{as}}{}{deletion}{ps_134-7_030-v0002} \var{\inserted{because}}{\hand{because}}{insertion}{ps_134-7_030-v0003} one knows that over a non-perfect field k  there
may be "forms" of the additive group $\underline{G}_a$ which are not isomorphic to $\underline{G}_a$ ;
presumably there should be similar examples over $\underline{Z}$ to,  with rank larger
than one however.

We feel however that the case when  $G = \hand{\pi}_1(X_*)$  is a non-linear
or even a non abelian group object of U(k),  is still worth of interest.
The \var{\xout{xxxx}}{}{typewriter deletion}{ps_134-7_030-v0004}\var{\inserted[typed]{first}}{first}{typed insertion}{ps_134-7_030-v0005} test it would seem,  to check if we do have a good notion indeed,
is to see if it does satisfy to the ``linearization theorem'' in case $k=\underline{Z}$ ,
i.e. the maps (w) on  page 454 are quasi-isomorphisms.  Another key test,
which now makes sense for arbitrary k,  is whether for a smooth Kan com\var{\xout{p}}{}{typewriter deletion}{ps_134-7_030-v0006}-%
plex in \var{\xout{M}}{}{typewriter deletion}{ps_134-7_030-v0007}$\hand{\mathcal{M}_0}(k)$  (i.e. satisfying the extra assumption involving G),
$X_*$  is homotopic to a bundle over K(G,1),  with a 1-connected fiber or
more specifically,  a fiber $Y_*$  satisfying  $Y_0=Y_1=e$ .  Among other featu-%
res to expect is a natural operation of the group object G on the k-modu-%
les $\hand{\pi}_i(X_*)$,  as well as $H_i(X_*)$ .  If however we wish,  for $k=\underline{Z}$ ,  to
use models in $\hand{\mathcal{M}_0}(\underline{Z})$  for describing possibly homotopy types with nilpo-%
tent $\hand{\pi}_1$ say, \var{\inserted{and moreover unipotent action on the $\pi_i$}}{\hand{and moreover unipotent action on the $\pi_i$}}{insertion}{ps_134-7_030-v0008} and devise a corresponding equivalence between suitable
homotopy categories,  we should first investigate the question of the
relationship between nilpotent discrete groups,  and group objects of
U(k)  -  a question already touched upon earlier in our reflection on
linearization (see end of section 94),, and of separate interest.

\var{\deleted{\illegible{1}}}{}{deletion}{ps_134-7_030-v0009}

\hand{114.}  \var{\xout{Afterxthxx}}{}{typewriter deletion}{ps_134-7_030-v0010}\var{\inserted[typed]{During these}}{During these}{typed insertion}{ps_134-7_030-v0011} four\var{\xout{th}}{}{typewriter deletion}{ps_134-7_030-v0012} days of reflection on schematization of homotopy
types,  \var{\xout{finally}}{}{typewriter deletion}{ps_134-7_030-v0013} a relatively coherent picture \var{\over{w}{h}}{\hand{h}}{writing over}{ps_134-7_030-v0014}as gradually \var{\inserted{been}}{\hand{been}}{insertion}{ps_134-7_030-v0015} emerging from
darkness.  How far this image reflects substantial reality,  not just
daydreaming,  I would be at a loss to tell now.  Maybe some substantial
corrections will have to be made still,  besides getting in other ideas
for a more complete picture  -  I would be amazed at any rate if every-%
thing should turn out as just nonsense !  If it does'nt,  there is surely
a lot of work ahead to get everything straightened out and ready-to-use.
I will leave it at that I suppose,  for the time being  -  maybe just
finish this digression by a quick review of the set up\hand{,} and \var{\inserted{of}}{\hand{of}}{insertion}{ps_134-7_030-v0016} some main
questions which have come out.

\margin{30}

% Visual collation: both supplied image halves inspected; see CHECK_DECISIONS_134_7_8.tsv.

%% END SOURCE ps_134-7_030
\ifnum\EditionMode=1
\psapparatusstart
\psentry{ps_134-7_030-v0001}{1}{typewriter deletion}{\xout{h}}{}{The typewriter strikes out this material; it leaves the corrected running text.}
\psentry{ps_134-7_030-v0002}{2}{deletion}{\deleted{as}}{}{The author cancels this material by hand; it leaves the corrected running text.}
\psentry{ps_134-7_030-v0003}{3}{insertion}{\inserted{because}}{\hand{because}}{The handwritten insertion is incorporated at its indicated position; insertion brackets leave the running text.}
\psentry{ps_134-7_030-v0004}{4}{typewriter deletion}{\xout{xxxx}}{}{The typewriter strikes out this material; it leaves the corrected running text.}
\psentry{ps_134-7_030-v0005}{5}{typed insertion}{\inserted[typed]{first}}{first}{The typed insertion is incorporated at its indicated position; insertion brackets leave the running text.}
\psentry{ps_134-7_030-v0006}{6}{typewriter deletion}{\xout{p}}{}{The typewriter strikes out this material; it leaves the corrected running text.}
\psentry{ps_134-7_030-v0007}{7}{typewriter deletion}{\xout{M}}{}{The typewriter strikes out this material; it leaves the corrected running text.}
\psentry{ps_134-7_030-v0008}{8}{insertion}{\inserted{and moreover unipotent action on the $\pi_i$}}{\hand{and moreover unipotent action on the $\pi_i$}}{The handwritten insertion is incorporated at its indicated position; insertion brackets leave the running text.}
\psentry{ps_134-7_030-v0009}{9}{deletion}{\deleted{\illegible{1}}}{}{The author cancels this material by hand; it leaves the corrected running text.}
\psentry{ps_134-7_030-v0010}{10}{typewriter deletion}{\xout{Afterxthxx}}{}{The typewriter strikes out this material; it leaves the corrected running text.}
\psentry{ps_134-7_030-v0011}{11}{typed insertion}{\inserted[typed]{During these}}{During these}{The typed insertion is incorporated at its indicated position; insertion brackets leave the running text.}
\psentry{ps_134-7_030-v0012}{12}{typewriter deletion}{\xout{th}}{}{The typewriter strikes out this material; it leaves the corrected running text.}
\psentry{ps_134-7_030-v0013}{13}{typewriter deletion}{\xout{finally}}{}{The typewriter strikes out this material; it leaves the corrected running text.}
\psentry{ps_134-7_030-v0014}{14}{writing over}{\over{w}{h}}{\hand{h}}{The superposed handwritten reading replaces the underlying reading.}
\psentry{ps_134-7_030-v0015}{15}{insertion}{\inserted{been}}{\hand{been}}{The handwritten insertion is incorporated at its indicated position; insertion brackets leave the running text.}
\psentry{ps_134-7_030-v0016}{16}{insertion}{\inserted{of}}{\hand{of}}{The handwritten insertion is incorporated at its indicated position; insertion brackets leave the running text.}
\psapparatusend
\fi
\pssource{ps_134-7_031}{134-7/31}{https://grothendieck.umontpellier.fr/web/viewer.html?file=../134-7.pdf\#page=32}
\psnumber{ps_134-7_031-v0001}{1}
\psnumber{ps_134-7_031-v0002}{2}
\psnumber{ps_134-7_031-v0003}{3}
\psnumber{ps_134-7_031-v0004}{4}
\psnumber{ps_134-7_031-v0005}{5}
%% BEGIN SOURCE ps_134-7_031
% Pursuing Stacks (cotes 134-2 to 134-8): archive cote 134-7, site page 31 (PDF page 32), typescript page 473.
% typescript. Draft compared: ps6.tex:1438-1502.
% AI-transcribed by Claude Opus 5.5 (claude-opus-5-5, adaptive thinking, effort high), Message Batches API, batch ps3 (msgbatch_01XdFmuVVBAcgc7EFxdtCqsj), 2026-11-12; reviewed against both supplied scan halves; see CHECK_DECISIONS.tsv.
\scanpage{134-7/31}
\mspage{473}

For a given ground ring k,  the basic category\var{\inserted{we'll use}}{\hand{we'll use}}{insertion}{ps_134-7_031-v0001} of ``schematic'' objects
over k is the category of \uline{unipotent} \uline{bundles} over k,  which may be defined
as the category of \var{\xout{xxx}}{}{typewriter deletion}{ps_134-7_031-v0002} functors  from Alg$_{/k}$ to \var{\xout{sets of the type W(}}{}{typewriter deletion}{ps_134-7_031-v0003} (Sets)
isomorphic to functors of the type
\[ W(M) = ( k' \hand{\longmapsto} M\otimes_k k' ) \qquad , \]
where M is any k-module.  We do not restrict here M to be projective or
flat,  as we definitely want to have,  for a ring homomorphism $k \hand{\to} k'$ ,
a problemless functor  ``\uline{restriction of scalars}''
\[ U(k') \hand{\to} U(k) \qquad , \]
inserting in the commutative diagram
% Hand-drawn arrows are marked blue; typed objects remain plain.
\diagram{
(I) \qquad
\begin{tikzcd}
Ab_{k'} \arrow[color=blue,r, "\text{restr.}"] \arrow[color=blue,d, "W_{k'}"'] \& Ab_k \arrow[color=blue,d, "W_k"] \\
U(k') \arrow[color=blue,r, "\text{restr.}"] \& U(k)
\end{tikzcd}
\qquad .
}
Another reason is that we want that the k-modules of the type $\hand{\pi}_i(X_{\hand{*}})$
which will come out should be eligible for defining objects in U(k).
We are more specifically interested though in ``\uline{smooth}'' objects of U(k),
namely those which correspond to \uline{projective} k-modules.  \hand{(}We prefer to
call them ``smooth'' rather than ``projective'', in order to avoid confusion
with the notion of a projective object in the usual categorical sense
for U(k).)\hand{\unsure{)}}  Another relevant notion is the notion of a \uline{submersion} ,
namely a map in U(k) \var{\inserted{isomorphic to one}}{\hand{isomorphic to one}}{insertion}{ps_134-7_031-v0004} coming from an epimorphism \var{\xout{in}}{}{typewriter deletion}{ps_134-7_031-v0005} $M \hand{\to} N$  in Ab$_k$ .

(If the latter can be chosen to have a projective kernel,  we may speak
of a smooth submersion.)  The ring restriction functor transforms sub-%
mersions into submersions,  and also smooth objects into smooth ones
provided k' is projective as a module over k.    We also have a \uline{ring}
\uline{extension functor} from U(k) to U(k'),  giving rise to a diagram (I')
similar to (I) above,  it transforms submersions into submersions,
smooth objects into smooth ones.

The smoothness condition is likely to come in in two ways,  one
is via \uline{flatness}  (we may call an object of U(k) ``flat'' when it is iso-%
morphic to some W(M) with M a flat k-module),  whereas projectivity is
needed in order to insure that in certain cases,  weak equivalences are
homotopisms.  Flatness is the kind of condition which insures the vali-%
dity of ``naive'' universal coefficients formulae for homotopy or homology
objects ,  whereas projectivity may be needed in case of such formulae
for cohomology rather than homology.

\margin{31}

% Visual collation: both supplied image halves inspected; see CHECK_DECISIONS_134_7_8.tsv.

%% END SOURCE ps_134-7_031
\ifnum\EditionMode=1
\psapparatusstart
\psentry{ps_134-7_031-v0001}{1}{insertion}{\inserted{we'll use}}{\hand{we'll use}}{The handwritten insertion is incorporated at its indicated position; insertion brackets leave the running text.}
\psentry{ps_134-7_031-v0002}{2}{typewriter deletion}{\xout{xxx}}{}{The typewriter strikes out this material; it leaves the corrected running text.}
\psentry{ps_134-7_031-v0003}{3}{typewriter deletion}{\xout{sets of the type W(}}{}{The typewriter strikes out this material; it leaves the corrected running text.}
\psentry{ps_134-7_031-v0004}{4}{insertion}{\inserted{isomorphic to one}}{\hand{isomorphic to one}}{The handwritten insertion is incorporated at its indicated position; insertion brackets leave the running text.}
\psentry{ps_134-7_031-v0005}{5}{typewriter deletion}{\xout{in}}{}{The typewriter strikes out this material; it leaves the corrected running text.}
\psapparatusend
\fi
\pssource{ps_134-7_032}{134-7/32}{https://grothendieck.umontpellier.fr/web/viewer.html?file=../134-7.pdf\#page=33}
\psnumber{ps_134-7_032-v0001}{1}
\psnumber{ps_134-7_032-v0002}{2}
\psnumber{ps_134-7_032-v0003}{3}
\psnumber{ps_134-7_032-v0004}{4}
\psnumber{ps_134-7_032-v0005}{5}
\psnumber{ps_134-7_032-v0006}{6}
\psnumber{ps_134-7_032-v0007}{7}
%% BEGIN SOURCE ps_134-7_032
% Pursuing Stacks (cotes 134-2 to 134-8): archive cote 134-7, site page 32 (PDF page 33), typescript page 474.
% typescript. Draft compared: ps6.tex:1488-1560.
% AI-transcribed by Claude Opus 5.5 (claude-opus-5-5, adaptive thinking, effort high), Message Batches API, batch ps3 (msgbatch_01XdFmuVVBAcgc7EFxdtCqsj), 2026-11-12; reviewed against both supplied scan halves; see CHECK_DECISIONS.tsv.
\scanpage{134-7/32}
\mspage{474}

The description I just recalled of U(k) is the one most intuitive to
my mind,  other people may prefer the more computational one given on
page 451 in terms of $\hand{\Gamma}_k^{\hand{\wedge}}\,(M)$  (endowed with it's augmentation to k and it's
diagonal map),  which is of importance in it's own right.  It shows
the existence of a canonical k-\uline{linearization} \uline{functor}
\[ \text{(J)} \qquad L : U(k) \hand{\longrightarrow} Ab_k \quad , \]
giving rise to the commutative diagram (up to can. isom.)
\var{\deleted{\xout{$Ab_k$x\allowbreak{}x\allowbreak{}x\allowbreak{}x\allowbreak{}x\allowbreak{}x\allowbreak{}$W$x\allowbreak{}x\allowbreak{}x\allowbreak{}x\allowbreak{}x\allowbreak{}x\allowbreak{}x\allowbreak{}$Ab_k$x\allowbreak{}x\allowbreak{}W(k)}}}{}{deletion}{ps_134-7_032-v0001}
% Hand-drawn arrows are marked blue; typed objects remain plain.
\diagram{
\text{(J')}\qquad
\begin{tikzcd}
Ab_k \arrow[color=blue,rr, "W"] \arrow[color=blue,dr, "\hand{\Gamma}_k^{\hand{\wedge}}"'] \& \& U(k) \arrow[color=blue,dl, "L"] \\
\& Ab_k \&
\end{tikzcd}
\quad \var{\xout{,}}{}{typewriter deletion}{ps_134-7_032-v0002}
}
where
\[ \hand{\Gamma}_k^{\hand{\wedge}}\,(M) = \hand{\prod}_{i\hand{\geq}0} \hand{\Gamma}_k^{\,i}(M) \qquad . \]
This linearization is not quite compatible with ring extension,  it beco-%
mes so only when we view it as a functor with values,  not just in the
category $Ab_k$ of k-modules,  but of separated and complete \var{\xout{k-modules for}}{}{typewriter deletion}{ps_134-7_032-v0003}
linearly topologized k-modules,  the ring extension functor for these
being \uline{completed} tensor product.  This is a little (or big ?) technical
\margin{x} drawback for this notion of linearization.  We have a canonical embedding
\[ \text{(K)} \qquad \hand{x \mapsto \exp(x) :} \qquad X \hand{\longrightarrow} \var{\xout{\unsure{W(L(X))}xxx}}{}{typewriter deletion}{ps_134-7_032-v0004}\ W^{\hand{\wedge}}\,(L(X)) \qquad \hand{\text{(cf. p.450, (n))}} \]
functorial in X,  \var{\xout{which satisfies a}}{}{typewriter deletion}{ps_134-7_032-v0005} where for a topological k-module M
as above,  described as a filtering inverse limit of discrete ones $M_i$ ,
we define
\[ \text{(K')} \qquad \var{\xout{W}}{}{typewriter deletion}{ps_134-7_032-v0006}\ W^{\hand{\wedge}}\,(M) = (k' \hand{\longmapsto} \underset{\hand{\leftarrow\, i}}{\underline{\lim}}\ M_{\unsure{i}} \otimes_k k' \ ) \quad . \]
The map (K) has a universal property with respect to all possible maps
$X \hand{\to} W^{\hand{\wedge}}\,(M)$  with M a linearly topologized separated and complete k-module,
which acounts for it's role as ``linearization''.  It should be noted here
that the map (C) on page (467) does'nt quite exist,  we have corrected
this point here -  definitely we cannot in (K) replace $W^{\hand{\wedge}}$ by W .  Of cour-%
se linearization L (or it's variant $L_{pt}$)  does'nt commute in any sense
whatever to restriction of ground ring.

The image of  X  in $W^{\hand{\wedge}}\,(L(X))$  is caracterized by the simple
formulae (q) p. 451 .  Maps from X to Y  may be described as just continu-%
ous \var{\inserted[typed]{k-linear}}{k-linear}{typed insertion}{ps_134-7_032-v0007} maps from  L(X) to L(Y),  compatible with augmentations and diag\unsure{o}nal
maps.

\margin{32}

% Visual collation: both supplied image halves inspected; see CHECK_DECISIONS_134_7_8.tsv.

%% END SOURCE ps_134-7_032
\ifnum\EditionMode=1
\psapparatusstart
\psentry{ps_134-7_032-v0001}{1}{deletion}{\deleted{\xout{$Ab_k$x\allowbreak{}x\allowbreak{}x\allowbreak{}x\allowbreak{}x\allowbreak{}x\allowbreak{}$W$x\allowbreak{}x\allowbreak{}x\allowbreak{}x\allowbreak{}x\allowbreak{}x\allowbreak{}x\allowbreak{}$Ab_k$x\allowbreak{}x\allowbreak{}W(k)}}}{}{The author cancels this material by hand; it leaves the corrected running text.}
\psentry{ps_134-7_032-v0002}{2}{typewriter deletion}{\xout{,}}{}{The typewriter strikes out this material; it leaves the corrected running text.}
\psentry{ps_134-7_032-v0003}{3}{typewriter deletion}{\xout{k-modules for}}{}{The typewriter strikes out this material; it leaves the corrected running text.}
\psentry{ps_134-7_032-v0004}{4}{typewriter deletion}{$\xout{\unsure{W(L(X))}xxx}$}{}{The typewriter strikes out this material; it leaves the corrected running text.}
\psentry{ps_134-7_032-v0005}{5}{typewriter deletion}{\xout{which satisfies a}}{}{The typewriter strikes out this material; it leaves the corrected running text.}
\psentry{ps_134-7_032-v0006}{6}{typewriter deletion}{$\xout{W}$}{}{The typewriter strikes out this material; it leaves the corrected running text.}
\psentry{ps_134-7_032-v0007}{7}{typed insertion}{\inserted[typed]{k-linear}}{k-linear}{The typed insertion is incorporated at its indicated position; insertion brackets leave the running text.}
\psapparatusend
\fi
\pssource{ps_134-7_033}{134-7/33}{https://grothendieck.umontpellier.fr/web/viewer.html?file=../134-7.pdf\#page=34}
\psnumber{ps_134-7_033-v0001}{1}
\psnumber{ps_134-7_033-v0002}{2}
\psnumber{ps_134-7_033-v0003}{3}
\psnumber{ps_134-7_033-v0004}{4}
\psnumber{ps_134-7_033-v0005}{5}
\psnumber{ps_134-7_033-v0006}{6}
\psnumber{ps_134-7_033-v0007}{7}
\psnumber{ps_134-7_033-v0008}{8}
\psnumber{ps_134-7_033-v0009}{9}
\psnumber{ps_134-7_033-v0010}{10}
%% BEGIN SOURCE ps_134-7_033
% Pursuing Stacks (cotes 134-2 to 134-8): archive cote 134-7, site page 33 (PDF page 34), typescript page 475.
% typescript. Draft compared: ps6.tex:1543-1620.
% AI-transcribed by Claude Opus 5.5 (claude-opus-5-5, adaptive thinking, effort high), Message Batches API, batch ps3 (msgbatch_01XdFmuVVBAcgc7EFxdtCqsj), 2026-11-12; reviewed against both supplied scan halves; see CHECK_DECISIONS.tsv.
\scanpage{134-7/33}
\mspage{475}

We'll more specifically work in the category $U(k)^\hand{\bullet}$ of \uline{pointed} objec\unsure{ts}
of $U(k)$, namely objects endowed with a section over the final object $e$,
the so-called \uline{pointed} \uline{unipotent bundles}. We now have a functor
\[ W^\hand{\bullet} : \quad Ab_k \hand{\longrightarrow} U(k)^\hand{\bullet} \tag{L} \]
deduced from $W$ using the fact that $W(0)=e$ , and a ``\uline{pointed linearizati-}%
\uline{on functor}''
\[ L^\hand{\bullet} \text{or } L_{pt} \;:\; U(k)^\hand{\bullet} \hand{\longrightarrow} Ab_k \quad , \tag{M} \]
giving rise to a commutative diagram similar to (J')
% Hand-drawn arrows are marked blue; typed objects remain plain.
\diagram{
\begin{tikzcd}
Ab_k \arrow[color=blue,rr, "W^\hand{\bullet}"] \arrow[color=blue,dr, "\hand{\Gamma}_k"'] \& \& U(k)^\hand{\bullet} \arrow[color=blue,dl, "L^\hand{\bullet}"] \\
 \& Ab_k \& ;
\end{tikzcd}
\quad\textnormal{(M')}
}
the notation $L^\hand{\bullet}$ seems here the most coherent one, but may bring about
confusion with the similar notation for some cochain complex say, there-%
fore we had first used the alternative notation $L_{pt}$ , to which one may
still come back if needed. This time the functor $L^\hand{\bullet}$ commutes to ring
extension without any grain of salt. We have of course a canonical
embedding
\[ L^\hand{\bullet}(X) \hand{\hookrightarrow} L(X) \tag{N} \]
\var{\deleted{corresponding to}}{}{deletion}{ps_134-7_033-v0001} \var{\inserted{defined via}}{\hand{defined via}}{insertion}{ps_134-7_033-v0002} the corresponding embedding for an object $M$ of $Ab_k$
\[ \hand{\Gamma}_k(M) \hand{\hookrightarrow} \hand{\Gamma}_k{}^{\hand{\wedge}}(M) \quad , \tag{N'} \]
by which $L(X)$ may be viewed as the completion of $L^\hand{\bullet}(X)$ with respect to
the topology it induces on it, which is a \uline{canonical} \uline{topology} on \var{\deleted{\unsure{X|it}}}{}{deletion}{ps_134-7_033-v0003}\hand{$L^\bullet(X)$}.
% NOTE: a typed ``X'' and ``Y'' stand between this line and the next, under ``the'' and ``topology''
$X$ \quad $Y$
Maps in $U(k)^\hand{\bullet}$ correspond to \uline{k-linear} \uline{maps}
\[ L^\hand{\bullet}(X) \hand{\longrightarrow} L^\hand{\bullet}(Y) \]
which are moreover \uline{continous}, \var{\inserted[typed]{and}}{and}{typed insertion}{ps_134-7_033-v0004} \uline{compatible with} \var{\inserted[typed]{co}}{co}{typed insertion}{ps_134-7_033-v0005}\uline{augmentation}\var{\deleted{\unsure{s}}}{}{deletion}{ps_134-7_033-v0006}\var{\xout{as well as}}{}{typewriter deletion}{ps_134-7_033-v0007}
(i.e. transforms 1 into 1) \uline{as well as with augmentations and diagonal}
\uline{maps}. I wonder if there is any simple \hand{c}haracterization of submersions
\var{\inserted[typed]{in $U(k)^\hand{\bullet}$}}{in $U(k)^\hand{\bullet}$}{typed insertion}{ps_134-7_033-v0008} in terms of the corresponding map between the linearizations. At any
rate, an object \var{\inserted[typed]{$X$}}{$X$}{typed insertion}{ps_134-7_033-v0009} of $U(k)^\hand{\bullet}$ is smooth resp. flat iff $L^\hand{\bullet}(X)$ is a projecti-%
ve resp. a flat $k$-module.

For any natural integer, we want now to define a model category
\[ \hand{\mathcal{M}}_n(k) \hand{\subset} \uline{Hom}(\hand{\Delta^{\circ}}, U(k)^\hand{\bullet}) \quad , \tag{O} \]
which should be a full subcategory of the category of semisimplicial
\var{\deleted{\xout{objects in U(k)}}}{}{deletion}{ps_134-7_033-v0010}

\margin{33}

% Visual collation: both supplied image halves inspected; see CHECK_DECISIONS_134_7_8.tsv.

%% END SOURCE ps_134-7_033
\ifnum\EditionMode=1
\psapparatusstart
\psentry{ps_134-7_033-v0001}{1}{deletion}{\deleted{corresponding to}}{}{The author cancels this material by hand; it leaves the corrected running text.}
\psentry{ps_134-7_033-v0002}{2}{insertion}{\inserted{defined via}}{\hand{defined via}}{The handwritten insertion is incorporated at its indicated position; insertion brackets leave the running text.}
\psentry{ps_134-7_033-v0003}{3}{deletion}{\deleted{\unsure{X|it}}}{}{The author cancels this material by hand; it leaves the corrected running text.}
\psentry{ps_134-7_033-v0004}{4}{typed insertion}{\inserted[typed]{and}}{and}{The typed insertion is incorporated at its indicated position; insertion brackets leave the running text.}
\psentry{ps_134-7_033-v0005}{5}{typed insertion}{\inserted[typed]{co}}{co}{The typed insertion is incorporated at its indicated position; insertion brackets leave the running text.}
\psentry{ps_134-7_033-v0006}{6}{deletion}{\deleted{\unsure{s}}}{}{The author cancels this material by hand; it leaves the corrected running text.}
\psentry{ps_134-7_033-v0007}{7}{typewriter deletion}{\xout{as well as}}{}{The typewriter strikes out this material; it leaves the corrected running text.}
\psentry{ps_134-7_033-v0008}{8}{typed insertion}{\inserted[typed]{in $U(k)^\hand{\bullet}$}}{in $U(k)^\hand{\bullet}$}{The typed insertion is incorporated at its indicated position; insertion brackets leave the running text.}
\psentry{ps_134-7_033-v0009}{9}{typed insertion}{\inserted[typed]{$X$}}{$X$}{The typed insertion is incorporated at its indicated position; insertion brackets leave the running text.}
\psentry{ps_134-7_033-v0010}{10}{deletion}{\deleted{\xout{objects in U(k)}}}{}{The author cancels this material by hand; it leaves the corrected running text.}
\psapparatusend
\fi
\pssource{ps_134-7_034}{134-7/34}{https://grothendieck.umontpellier.fr/web/viewer.html?file=../134-7.pdf\#page=35}
\psnumber{ps_134-7_034-v0001}{1}
\psnumber{ps_134-7_034-v0002}{2}
\psnumber{ps_134-7_034-v0003}{3}
\psnumber{ps_134-7_034-v0004}{4}
\psnumber{ps_134-7_034-v0005}{5}
%% BEGIN SOURCE ps_134-7_034
% Pursuing Stacks (cotes 134-2 to 134-8): archive cote 134-7, site page 34 (PDF page 35), typescript page 476.
% typescript. Draft compared: ps6.tex:1601-1686.
% AI-transcribed by Claude Opus 5.5 (claude-opus-5-5, adaptive thinking, effort high), Message Batches API, batch ps3 (msgbatch_01XdFmuVVBAcgc7EFxdtCqsj), 2026-11-12; reviewed against both supplied scan halves; see CHECK_DECISIONS.tsv.
\scanpage{134-7/34}
\mspage{476}

objects in $U(k)^\hand{\bullet}$ . We'll get a functor
\[ \text{(P)}\qquad X_{*}\hand{\longmapsto}X_{*}(k) \ : \ \hand{\mathcal{M}}_{n}(k)\hand{\longrightarrow} \ \uline{Hom}(\hand{\Delta}^{o},\mathrm{Sets}^{\hand{\bullet}}) \]
from this category to the category of semisimplicial pointed sets. \var{\xout{Ifxn}}{}{typewriter deletion}{ps_134-7_034-v0001}
For $n\hand{\geq}1$ , the only condition it seems to impose upon $X_{*}$ in the
second member of (O), i.e. \var{\inserted{upon}}{\hand{upon}}{insertion}{ps_134-7_034-v0002} a semisimplicial pointed unipotent bundle over
k, in order to belong to $\hand{\mathcal{M}}_{n}(k)$ , is
\[ (Q_{n})\qquad X_{i} = e \quad\text{for}\quad i \hand{\leq} n \ . \]
This for $n=0$ reduces to the common condition
\[ (Q_{o})\qquad X_{o} = e \ , \]
which definitely is not enough though to get a category of ``models''
$\hand{\mathcal{M}_{\bullet}}(k)$ whose objects should have the kind of properties we are after.
There are various kind of \var{\inserted[typed]{extra}}{extra}{typed insertion}{ps_134-7_034-v0003} restrictions one may want to impose, accor-%
ding to the type of situations one wants to describe, some hints along
these lines are given on pa\var{\over{\unsure{g}}{g}}{\hand{g}}{writing over}{ps_134-7_034-v0004}es 469-472. For a preliminary study, the
case $n\hand{\geq} 1$, and more specifically, the case $n=1$, is quite enough, the
latter corresponding to the restrictions
\[ (Q_{1})\qquad X_{o} = X_{1} = e \ . \]

From P we get a functor
\[ (P')\qquad \hand{\mathcal{M}}_{n}(k) \hand{\longrightarrow} \mathrm{Hot}^{\hand{\bullet}}_{n} \quad = \quad \begin{array}{l}\text{category of pointed n-connected}\\ \text{homotopy types ,}\end{array} \]
we define \var{\xout{thexnotion}}{}{typewriter deletion}{ps_134-7_034-v0005} a map in $\hand{\mathcal{M}}_{n}(k)$ to be a \uline{weak} \uline{equivalence} if it's
image by (P) is i.e. it's image by(P')is an isomorphism, and locali-%
sing by weak equivalences we get the category
\[ (R)\qquad \mathrm{Hot}_{n}(k) \]
of ``n-connected schematic homotopy types over k'' , together with a
``sections functor'' induced by (P)
\[ (R')\qquad \mathrm{Hot}_{n}(k) \hand{\longrightarrow} \mathrm{Hot}_{n} \ . \]
One main point in our definitions is that \uline{we hope this functor to be an}
\uline{equivalence of categories}, \uline{in the case when $k=\underline{Z}$} , and of course $n\hand{\geq} 1$.

The description just given of categories $\mathrm{Hot}_{n}(k)$ is suitable for
defining functors of restriction of ground ring for $k \hand{\rightarrow} k'$
\[ (S)\qquad \mathrm{Hot}_{n}(k')\hand{\longrightarrow} \mathrm{Hot}_{n}(k) \ , \]
compatible with the sections functors (R') for k and k'. It is'nt
directly suited though for describing ring extension - as a matter of
\margin{34}

% Visual collation: both supplied image halves inspected; see CHECK_DECISIONS_134_7_8.tsv.

%% END SOURCE ps_134-7_034
\ifnum\EditionMode=1
\psapparatusstart
\psentry{ps_134-7_034-v0001}{1}{typewriter deletion}{\xout{Ifxn}}{}{The typewriter strikes out this material; it leaves the corrected running text.}
\psentry{ps_134-7_034-v0002}{2}{insertion}{\inserted{upon}}{\hand{upon}}{The handwritten insertion is incorporated at its indicated position; insertion brackets leave the running text.}
\psentry{ps_134-7_034-v0003}{3}{typed insertion}{\inserted[typed]{extra}}{extra}{The typed insertion is incorporated at its indicated position; insertion brackets leave the running text.}
\psentry{ps_134-7_034-v0004}{4}{writing over}{\over{\unsure{g}}{g}}{\hand{g}}{The superposed handwritten reading replaces the underlying reading.}
\psentry{ps_134-7_034-v0005}{5}{typewriter deletion}{\xout{thexnotion}}{}{The typewriter strikes out this material; it leaves the corrected running text.}
\psapparatusend
\fi
\pssource{ps_134-7_035}{134-7/35}{https://grothendieck.umontpellier.fr/web/viewer.html?file=../134-7.pdf\#page=36}
\psnumber{ps_134-7_035-v0001}{1}
\psnumber{ps_134-7_035-v0002}{2}
\psnumber{ps_134-7_035-v0003}{3}
\psnumber{ps_134-7_035-v0004}{4}
\psnumber{ps_134-7_035-v0005}{5}
\psnumber{ps_134-7_035-v0006}{6}
\psnumber{ps_134-7_035-v0007}{7}
\psnumber{ps_134-7_035-v0008}{8}
\psnumber{ps_134-7_035-v0009}{9}
\psnumber{ps_134-7_035-v0010}{10}
%% BEGIN SOURCE ps_134-7_035
% Pursuing Stacks (cotes 134-2 to 134-8): archive cote 134-7, site page 35 (PDF page 36), typescript page 477.
% typescript. Draft compared: ps6.tex:1668-1740.
% AI-transcribed by Claude Opus 5.5 (claude-opus-5-5, adaptive thinking, effort high), Message Batches API, batch ps3 (msgbatch_01XdFmuVVBAcgc7EFxdtCqsj), 2026-11-12; reviewed against both supplied scan halves; see CHECK_DECISIONS.tsv.
\scanpage{134-7/35}
\mspage{477}
fact,  ring extension for homotopy types (an operation of greater interest
than ring restriction surely) is \uline{not} expressed,  in general,  by just
performing the trivial ring extension operation
\[ X_{*} \hand{\mapsto} X_{*} \otimes_{k} k' \]
on models in $\hand{\pi}_{n}(k)$,  unless we assume k' to be flat over k say  -  but
even in this case it is by no means clear a priori that the operation
above transforms weak equivalences into weak equivalences.  This is very
clearly \var{\deleted{suggested}}{}{deletion}{ps_134-7_035-v0001}\var{\inserted{shown}}{\hand{shown}}{insertion}{ps_134-7_035-v0002} by the linear analogon,  the categories $\hand{\mathcal{M}}_{n}(k)$  being
replaced by \var{\inserted[typed]{the}}{the}{typed insertion}{ps_134-7_035-v0003} categories of chain complexes in  $Ab_{k}$ say,  or by $Comp^{-}(Ab_{k})$
or the like.  \var{\xout{In}}{}{typewriter deletion}{ps_134-7_035-v0004} In order to correctly describe ground ring extension on
homotopy types,  we'll have first to take a suitable``resolution'' of $X_{*}$ ,
namely replace $X_{*}$  by some  $K_{*}$  say,  endowed with a weak equivalence
\[ K_{*} \hand{\to} X_{*} \quad , \]
and $K_{*}$ satisfying some extra assumptions.  Maybe flatness of the compo-%
nent would be enough here.  For other purposes,  we may have to use
resolutions which are even smooth (componentwise),  or which satisfy a
suitable \uline{Kan condition}  (of a type outlined on page 463),  or both.
Our expectation is that,  when we restrict to the subcategory
\[ \var{\deleted{\unsure{SK_{n}(k)}}}{}{deletion}{ps_134-7_035-v0005} \ \hand{sK\mathcal{M}_{n}(k)} \]
of the model category $\hand{\mathcal{M}}_{n}(k)$ made up with smooth Kan complexes,  \uline{that}
\uline{the category} $Hot_{n}(k)$  \uline{may be described simply in  terms of such}  sK-%
\uline{complexes} \uline{``up to h\hand{$\smile$}omotopy''},  as usual.  If this is so,  the ground ring
extension functor \var{\deleted{(S)}}{}{deletion}{ps_134-7_035-v0006}  follows trivially from a corresponding functor on
the \var{\xout{\unsure{Kss}}}{}{typewriter deletion}{ps_134-7_035-v0007} sK-model categories
\[ \text{(T)} \qquad sK\hand{\mathcal{M}}_{n}(k) \hand{\longrightarrow} sK\hand{\mathcal{M}}_{n}(k') \quad , \quad \hand{X_{*} \longmapsto X_{*} \otimes_{k} k'} \hand{,} \]
hence
\[ \text{(T')} \qquad Hot_{n}(k) \hand{\longrightarrow} Hot_{n}(k') \quad . \]

From the \var{\inserted[typed]{sections}}{sections}{typed insertion}{ps_134-7_035-v0008} functor  (R')  we get homotopy invariants  $\hand{\pi}_{i}$ for an object
in  $Hot_{n}(k)$,  but the relevant k-module structure on these is not appea-%
rant on this definition.  We have a better hold,  via linear algebra over
k,  upon  homology invariants of $X_{*}$  ,  which are k-linear objects,  and
are definitely distinct (unless $k=\underline{Z}$ )  from the homology invariants of
$X_{*}(k)$ ,  which for gene ral k are definitely of little interest it seems.
The definition of homology goes via the pointed linearization functor \hand{(M)}\hand{:}
\[ \text{(U)} \qquad \hand{\Big\{} \begin{array}{l} LH.(X_{*}) \overset{def}{=} \var{\xout{\unsure{KK}}}{}{typewriter deletion}{ps_134-7_035-v0009} \ L_{\hand{pt}}\hand{(X_{*})} \var{\inserted{, viewed as an object in D.(Ab_{k}),}}{\hand{, viewed as an object in D.(Ab_{k}),}}{insertion}{ps_134-7_035-v0010} \\ \hand{H_{i}(X_{*}) \overset{def}{=} H_{i}(LH.(X_{*})) = \pi_{i}(L_{pt}(X_{*}))} \ \hand{in Ab_{k},} \end{array} \]
\margin{35}

% Visual collation: both supplied image halves inspected; see CHECK_DECISIONS_134_7_8.tsv.

%% END SOURCE ps_134-7_035
\ifnum\EditionMode=1
\psapparatusstart
\psentry{ps_134-7_035-v0001}{1}{deletion}{\deleted{suggested}}{}{The author cancels this material by hand; it leaves the corrected running text.}
\psentry{ps_134-7_035-v0002}{2}{insertion}{\inserted{shown}}{\hand{shown}}{The handwritten insertion is incorporated at its indicated position; insertion brackets leave the running text.}
\psentry{ps_134-7_035-v0003}{3}{typed insertion}{\inserted[typed]{the}}{the}{The typed insertion is incorporated at its indicated position; insertion brackets leave the running text.}
\psentry{ps_134-7_035-v0004}{4}{typewriter deletion}{\xout{In}}{}{The typewriter strikes out this material; it leaves the corrected running text.}
\psentry{ps_134-7_035-v0005}{5}{deletion}{$\deleted{\unsure{SK_{n}(k)}}$}{}{The author cancels this material by hand; it leaves the corrected running text.}
\psentry{ps_134-7_035-v0006}{6}{deletion}{\deleted{(S)}}{}{The author cancels this material by hand; it leaves the corrected running text.}
\psentry{ps_134-7_035-v0007}{7}{typewriter deletion}{\xout{\unsure{Kss}}}{}{The typewriter strikes out this material; it leaves the corrected running text.}
\psentry{ps_134-7_035-v0008}{8}{typed insertion}{\inserted[typed]{sections}}{sections}{The typed insertion is incorporated at its indicated position; insertion brackets leave the running text.}
\psentry{ps_134-7_035-v0009}{9}{typewriter deletion}{$\xout{\unsure{KK}}$}{}{The typewriter strikes out this material; it leaves the corrected running text.}
\psentry{ps_134-7_035-v0010}{10}{insertion}{$\inserted{, viewed as an object in D.(Ab_{k}),}$}{$\hand{, viewed as an object in D.(Ab_{k}),}$}{The handwritten insertion is incorporated at its indicated position; insertion brackets leave the running text.}
\psapparatusend
\fi
\pssource{ps_134-7_036}{134-7/36}{https://grothendieck.umontpellier.fr/web/viewer.html?file=../134-7.pdf\#page=37}
\psnumber{ps_134-7_036-v0001}{1}
\psnumber{ps_134-7_036-v0002}{2}
\psnumber{ps_134-7_036-v0003}{3}
\psnumber{ps_134-7_036-v0004}{4}
\psnumber{ps_134-7_036-v0005}{5}
\psnumber{ps_134-7_036-v0006}{6}
\psnumber{ps_134-7_036-v0007}{7}
\psnumber{ps_134-7_036-v0008}{8}
\psnumber{ps_134-7_036-v0009}{9}
\psnumber{ps_134-7_036-v0010}{10}
\psnumber{ps_134-7_036-v0011}{11}
\psnumber{ps_134-7_036-v0012}{12}
\psnumber{ps_134-7_036-v0013}{13}
\psnumber{ps_134-7_036-v0014}{14}
\psnumber{ps_134-7_036-v0015}{15}
\psnumber{ps_134-7_036-v0016}{16}
\psnumber{ps_134-7_036-v0017}{17}
\psnumber{ps_134-7_036-v0018}{18}
%% BEGIN SOURCE ps_134-7_036
% Pursuing Stacks (cotes 134-2 to 134-8): archive cote 134-7, site page 36 (PDF page 37), typescript page 478.
% typescript. Draft compared: ps6.tex:1723-1793.
% AI-transcribed by Claude Opus 5.5 (claude-opus-5-5, adaptive thinking, effort high), Message Batches API, batch ps3 (msgbatch_01XdFmuVVBAcgc7EFxdtCqsj), 2026-11-12; reviewed against both supplied scan halves; see CHECK_DECISIONS.tsv.
\scanpage{134-7/36}
\mspage{478}

where the L in LH. suggests that we are taking something similar to a
total left derived functor, and where definitely in the right hand mem-%
ber we had to write L$_{\mathrm{pt}}$ and not L$^{\hand{\bullet}}$ , in order not to \var{\xout{sink into}}{}{typewriter deletion}{ps_134-7_036-v0001} get sunk
into a morass of confusion ! In the formulae (U) we should assume \var{\over{\unsure{n}}{h}}{\hand{h}}{writing over}{ps_134-7_036-v0002}ow-
\unsure{ever} that X$_{*}$ is a smooth Kan complex, which will imply (if indeed \var{\xout{H}}{}{typewriter deletion}{ps_134-7_036-v0003}
Hot$_n$(k) may be described in terms of sK\hand{$\mathcal{M}$}$_n$(k) as said above) that LH.
may be viewed as a functor
\[ \text{(U')} \qquad \mathrm{Hot}_n(k) \hand{\longrightarrow} \mathrm{D.}(\mathrm{Ab}_k) \quad , \]
and likewise the H$_i$'s are functors from Hot$_n$(k) to Ab$_k$ . In order to
compute these homology invariants for an arbitrary complex in \hand{$\mathcal{M}$}$_n$(k),
we'll first have to resolve it by a sK complex, and then apply (U).

We expect \uline{that a map} $X_{*} \hand{\longrightarrow} Y_{*}$ \uline{in} \hand{$\mathcal{M}$}$_n$(k) \uline{is a} \var{\xout{quasiis}}{}{typewriter deletion}{ps_134-7_036-v0004} \uline{weak equiva-}%
\uline{lence iff the corresponding map for} LH. \uline{is a quasiisomorphism}, in other
words we expect the functor (U') to be ``conservative'': a map in the
first category is an isomorphism iff it's image in the second one is.
A second main feature we expect from linearization, is that in the case
k=\underline{Z} it corresponds to the usual abelianization of homotopy types, This
statement, when made more specific as in (w) page 454, decomposes into
two distinct ones. O\var{\over{\unsure{n}}{n}}{\hand{n}}{writing over}{ps_134-7_036-v0005}e is of significance over an arbitrary ring k,
and states that \uline{for a} sK \uline{complex} X$_{*}$ , \uline{the inclusion} (coming from (N))
\[ \text{(V)} \qquad \mathrm{L}_{\mathrm{pt}}(X_{*}) \hand{\longrightarrow} \mathrm{L}(X_{*}) \] \margin{x}
\uline{from} L$_{\mathrm{pt}}$ \uline{into it's completion}, \uline{when viewed as a map of chain complexes}
\uline{in Ab$_k$} (using the simplicial differenti\var{\xout{\unsure{t}}}{}{typewriter deletion}{ps_134-7_036-v0006}al operator, or passing to the
corresponding ``normalized'' chain complexes first) \uline{is a quasi-isomorphism}.
\var{\over{w}{W}}{\hand{W}}{writing over}{ps_134-7_036-v0007}hether this is \var{\inserted{always}}{\hand{always}}{insertion}{ps_134-7_036-v0008} so or not, \var{\over{not}{\unsure{or}}}{\hand{\unsure{or}}}{writing over}{ps_134-7_036-v0009} \var{\xout{\unsure{xx}}}{}{typewriter deletion}{ps_134-7_036-v0010} \var{\inserted[typed]{whether}}{whether}{typed insertion}{ps_134-7_036-v0011} noetherian conditions on k or some \var{\xout{fin}}{}{typewriter deletion}{ps_134-7_036-v0012}
finiteness conditions for the components of X$_{*}$ are needed, looks like a
rather standard question of linear homological algebra ! O\var{\over{\unsure{nx}}{n}}{\hand{n}}{writing over}{ps_134-7_036-v0013} the other
hand, using the exponential embedding (K) for sections, we get another
map \var{\xout{of}}{}{typewriter deletion}{ps_134-7_036-v0014} of semisimplicial k-modules
\[ \text{(V')} \qquad k^{(X\hand{_{*}(k)})} \hand{\longrightarrow} \mathrm{L}(X_{*}) \quad , \] \margin{x}
and here \uline{the question again} (the expectation I might say?\hand{)} \uline{is} \var{\deleted{\xout{that}}}{}{deletion}{ps_134-7_036-v0015} \var{\inserted[typed]{whether}}{whether}{typed insertion}{ps_134-7_036-v0016} this
\uline{is \var{\deleted{again}}{}{deletion}{ps_134-7_036-v0017} \hand{a} quasiisomorphism}. This would just mean \var{\xout{that}}{}{typewriter deletion}{ps_134-7_036-v0018} (if coupled with
quasi-isomorphy of (V)) that the homology invariants (U) are just the
usual homology invariants of the discrete homotopy type modelled by
X$_{*}$(k), but with coefficients not in \underline{Z} , but in k. We certainly do ex-%
pect this to be true for k = \underline{Z} - which was the content of ``question 1''

\margin{36}

% Visual collation: both supplied image halves inspected; see CHECK_DECISIONS_134_7_8.tsv.

%% END SOURCE ps_134-7_036
\ifnum\EditionMode=1
\psapparatusstart
\psentry{ps_134-7_036-v0001}{1}{typewriter deletion}{\xout{sink into}}{}{The typewriter strikes out this material; it leaves the corrected running text.}
\psentry{ps_134-7_036-v0002}{2}{writing over}{\over{\unsure{n}}{h}}{\hand{h}}{The superposed handwritten reading replaces the underlying reading.}
\psentry{ps_134-7_036-v0003}{3}{typewriter deletion}{\xout{H}}{}{The typewriter strikes out this material; it leaves the corrected running text.}
\psentry{ps_134-7_036-v0004}{4}{typewriter deletion}{\xout{quasiis}}{}{The typewriter strikes out this material; it leaves the corrected running text.}
\psentry{ps_134-7_036-v0005}{5}{writing over}{\over{\unsure{n}}{n}}{\hand{n}}{The superposed handwritten reading replaces the underlying reading.}
\psentry{ps_134-7_036-v0006}{6}{typewriter deletion}{\xout{\unsure{t}}}{}{The typewriter strikes out this material; it leaves the corrected running text.}
\psentry{ps_134-7_036-v0007}{7}{writing over}{\over{w}{W}}{\hand{W}}{The superposed handwritten reading replaces the underlying reading.}
\psentry{ps_134-7_036-v0008}{8}{insertion}{\inserted{always}}{\hand{always}}{The handwritten insertion is incorporated at its indicated position; insertion brackets leave the running text.}
\psentry{ps_134-7_036-v0009}{9}{writing over}{\over{not}{\unsure{or}}}{\hand{\unsure{or}}}{The superposed handwritten reading replaces the underlying reading.}
\psentry{ps_134-7_036-v0010}{10}{typewriter deletion}{\xout{\unsure{xx}}}{}{The typewriter strikes out this material; it leaves the corrected running text.}
\psentry{ps_134-7_036-v0011}{11}{typed insertion}{\inserted[typed]{whether}}{whether}{The typed insertion is incorporated at its indicated position; insertion brackets leave the running text.}
\psentry{ps_134-7_036-v0012}{12}{typewriter deletion}{\xout{fin}}{}{The typewriter strikes out this material; it leaves the corrected running text.}
\psentry{ps_134-7_036-v0013}{13}{writing over}{\over{\unsure{nx}}{n}}{\hand{n}}{The superposed handwritten reading replaces the underlying reading.}
\psentry{ps_134-7_036-v0014}{14}{typewriter deletion}{\xout{of}}{}{The typewriter strikes out this material; it leaves the corrected running text.}
\psentry{ps_134-7_036-v0015}{15}{deletion}{\deleted{\xout{that}}}{}{The author cancels this material by hand; it leaves the corrected running text.}
\psentry{ps_134-7_036-v0016}{16}{typed insertion}{\inserted[typed]{whether}}{whether}{The typed insertion is incorporated at its indicated position; insertion brackets leave the running text.}
\psentry{ps_134-7_036-v0017}{17}{deletion}{\deleted{again}}{}{The author cancels this material by hand; it leaves the corrected running text.}
\psentry{ps_134-7_036-v0018}{18}{typewriter deletion}{\xout{that}}{}{The typewriter strikes out this material; it leaves the corrected running text.}
\psapparatusend
\fi
\pssource{ps_134-7_037}{134-7/37}{https://grothendieck.umontpellier.fr/web/viewer.html?file=../134-7.pdf\#page=38}
\psnumber{ps_134-7_037-v0001}{1}
\psnumber{ps_134-7_037-v0002}{2}
\psnumber{ps_134-7_037-v0003}{3}
\psnumber{ps_134-7_037-v0004}{4}
\psnumber{ps_134-7_037-v0005}{5}
\psnumber{ps_134-7_037-v0006}{6}
\psnumber{ps_134-7_037-v0007}{7}
\psnumber{ps_134-7_037-v0008}{8}
\psnumber{ps_134-7_037-v0009}{9}
%% BEGIN SOURCE ps_134-7_037
% Pursuing Stacks (cotes 134-2 to 134-8): archive cote 134-7, site page 37 (PDF page 38), typescript page 479.
% typescript. Draft compared: ps6.tex:1781-1849.
% AI-transcribed by Claude Opus 5.5 (claude-opus-5-5, adaptive thinking, effort high), Message Batches API, batch ps3 (msgbatch_01XdFmuVVBAcgc7EFxdtCqsj), 2026-11-12; reviewed against both supplied scan halves; see CHECK_DECISIONS.tsv.
\scanpage{134-7/37}
\mspage{479}

on page 452 (taken up again on page 454 and following).   Of course in
case we do'nt assume $k=\underline{Z}$ ,  writing $\underline{Z}^{(X_*(k))}$ instead of $k^{(X_*(k))}$ as
we did in (w) p. 454 now looks kind of silly,  and the idea in this \var{\xout{k-li-}}{}{typewriter deletion}{ps_134-7_037-v0001}
k-linear context to take k-valued homology of $X_*(k)$ rather than $\underline{Z}$-valued
one is evident enough !  However I was confused by the misconception that
the \uline{integral} homology of $X_*(k)$ should carry k-linear structure, as
this was what I expected too for the invariants $\hand{\pi}_i(X_*(k))$ (which seems
to turn out to be correct).  This misconception was corrected a few
pages later (p.457) but still I kept dragging along the silly first
member of (w) .  Anyhow, it \var{\inserted{now just}}{\hand{now just}}{insertion}{ps_134-7_037-v0002} occurs to me that except in case $k=\underline{Z}$ , it
is definitely \uline{false} that (V') is a quasi-isomorphism, except in some
wholly trivial cases.  Indeed, let $H_i$ be the first non vanishing homolo-%
gy invariant (U) of $X_*$ (or more safely still, take i=n+1), then
we definitely expect to have a canonical isomorphism \hand{of $k$-modules}
\[ \text{(V)} \qquad \hand{\pi}_i(X_*(k)) \hand{\simeq} H_i(X_*) = H_i \]
but we equally have by Hopf's theorem, as the lower $\hand{\pi}_j$'s of $X_*(k)$ are
zero
\[ \hand{\pi}_i(X_*(k)) = H_i(X_*(k), \underline{Z}) \quad , \]
hence \var{\xout{(posing S = $X_*(k)$)}}{}{typewriter deletion}{ps_134-7_037-v0003}
\[ \text{(V')} \qquad H_i \hand{\simeq} H_i(S,\underline{Z}) \quad , \quad \text{where} \quad S = X_*(k) \quad , \]
which is not compatible with the guess that $H_i \hand{\simeq} H_i(S,k)$ ($\hand{\simeq} H_i \hand{\otimes}_{\underline{Z}} k$ ) .
Presumably the isomorphism (V') above is induced by the first
map in (w) above, but (except for $k=\underline{Z}$ ) we should expect in (w) to
have an isomorphism only for the lowest dimensional homology groups
which are occuring in the two first members.  Anyhow, it appears after
all that this map in (w) is the more reasonable one \var{\inserted{compared to}}{\hand{compared to}}{insertion}{ps_134-7_037-v0004} \var{\deleted{than}}{}{deletion}{ps_134-7_037-v0005} (V') above, as
it yi\var{\over{\unsure{le}}{el}}{\hand{el}}{writing over}{ps_134-7_037-v0006}ds an isomorphism on homology in the key dimension n+1, whereas
(V') apparently will practically \uline{never} give an isomorphism.

What is mainly lacking still in this review of the \var{\deleted{main}}{}{deletion}{ps_134-7_037-v0007} expected \var{\inserted{main}}{\hand{main}}{insertion}{ps_134-7_037-v0008}
features of schematization of homotopy types, is description of the
k-module structure on the homotopy groups
\[ \text{(W)} \qquad \hand{\pi}_i(X_*) \overset{\mathrm{def}}{=} \hand{\pi}_i(X_*(k)) \quad , \]
or preferably still, a direct description of those invariants as
k-modules, working within the model category $\hand{\mathcal{M}}_n(k)$.  This, as sugges-%
ted in yesterday's notes\var{\inserted{(p. 464)}}{\hand{(p. 464)}}{insertion}{ps_134-7_037-v0009}, may be achieved \uline{by developping a theory of}
\uline{Postnikov dévissage within} $\hand{\mathcal{M}}_n(k)$ \uline{and using} (V) in order to pull oursel-%
ves by the bootstraps, defining homotopy finally in terms of homology.

\margin{37}

% Visual collation: both supplied image halves inspected; see CHECK_DECISIONS_134_7_8.tsv.

%% END SOURCE ps_134-7_037
\ifnum\EditionMode=1
\psapparatusstart
\psentry{ps_134-7_037-v0001}{1}{typewriter deletion}{\xout{k-li-}}{}{The typewriter strikes out this material; it leaves the corrected running text.}
\psentry{ps_134-7_037-v0002}{2}{insertion}{\inserted{now just}}{\hand{now just}}{The handwritten insertion is incorporated at its indicated position; insertion brackets leave the running text.}
\psentry{ps_134-7_037-v0003}{3}{typewriter deletion}{\xout{(posing S = $X_*(k)$)}}{}{The typewriter strikes out this material; it leaves the corrected running text.}
\psentry{ps_134-7_037-v0004}{4}{insertion}{\inserted{compared to}}{\hand{compared to}}{The handwritten insertion is incorporated at its indicated position; insertion brackets leave the running text.}
\psentry{ps_134-7_037-v0005}{5}{deletion}{\deleted{than}}{}{The author cancels this material by hand; it leaves the corrected running text.}
\psentry{ps_134-7_037-v0006}{6}{writing over}{\over{\unsure{le}}{el}}{\hand{el}}{The superposed handwritten reading replaces the underlying reading.}
\psentry{ps_134-7_037-v0007}{7}{deletion}{\deleted{main}}{}{The author cancels this material by hand; it leaves the corrected running text.}
\psentry{ps_134-7_037-v0008}{8}{insertion}{\inserted{main}}{\hand{main}}{The handwritten insertion is incorporated at its indicated position; insertion brackets leave the running text.}
\psentry{ps_134-7_037-v0009}{9}{insertion}{\inserted{(p. 464)}}{\hand{(p. 464)}}{The handwritten insertion is incorporated at its indicated position; insertion brackets leave the running text.}
\psapparatusend
\fi
\pssource{ps_134-7_038}{134-7/38}{https://grothendieck.umontpellier.fr/web/viewer.html?file=../134-7.pdf\#page=39}
\psnumber{ps_134-7_038-v0001}{1}
\psnumber{ps_134-7_038-v0002}{2}
\psnumber{ps_134-7_038-v0003}{3}
\psnumber{ps_134-7_038-v0004}{4}
\psnumber{ps_134-7_038-v0005}{5}
\psnumber{ps_134-7_038-v0006}{6}
\psnumber{ps_134-7_038-v0007}{7}
\psnumber{ps_134-7_038-v0008}{8}
\psnumber{ps_134-7_038-v0009}{9}
\psnumber{ps_134-7_038-v0010}{10}
\psnumber{ps_134-7_038-v0011}{11}
\psnumber{ps_134-7_038-v0012}{12}
\psnumber{ps_134-7_038-v0013}{13}
\psnumber{ps_134-7_038-v0014}{14}
\psnumber{ps_134-7_038-v0015}{15}
\psnumber{ps_134-7_038-v0016}{16}
\psnumber{ps_134-7_038-v0017}{17}
\psnumber{ps_134-7_038-v0018}{18}
\psnumber{ps_134-7_038-v0019}{19}
%% BEGIN SOURCE ps_134-7_038
% Pursuing Stacks (cotes 134-2 to 134-8): archive cote 134-7, site page 38 (PDF page 39), typescript page 480.
% typescript. Draft compared: ps6.tex:1836-1904.
% AI-transcribed by Claude Opus 5.5 (claude-opus-5-5, adaptive thinking, effort high), Message Batches API, batch ps3 (msgbatch_01XdFmuVVBAcgc7EFxdtCqsj), 2026-11-12; reviewed against both supplied scan halves; see CHECK_DECISIONS.tsv.
\scanpage{134-7/38}
\mspage{480}

At this point it should be noted that the devissage we'll have to use
here is the "brutal" one,  which we frowned upon earlier today !  To
develop such a formalism,  it seems essential to work with smooth Kan
complexes and projective resolutions of the\var{\inserted[typed]{k-modules}}{k-modules}{typed insertion}{ps_134-7_038-v0001} $\hand{\pi}_{i}$ as they appear one by one,
\var{\xout{endowed with}}{}{typewriter deletion}{ps_134-7_038-v0002}  Whether we want to describe\var{\xout{soft}}{}{typewriter deletion}{ps_134-7_038-v0003} "hard" or "soft" Postnikov
dévissage (see p.468 for the latter),  one common key step is the
\uline{linearization map}  coming from the exponential map (K)  applied component-%
wise
\[
(W) \qquad X_{\ast} \hand{\longrightarrow} \var{\xout{L(X_{\ast})}}{}{typewriter deletion}{ps_134-7_038-v0004}\ W^{\hand{\wedge}}(L(X_{\ast})) \ ,
\]
which we would like to look upon as \var{\xout{being}}{}{typewriter deletion}{ps_134-7_038-v0005} defining a homotopy class of
maps in $\hand{\mathcal{M}}_n(k)$
\[
\margin{X}\ (W\ ?) \qquad X_{\ast} \hand{\longrightarrow} W(L_{pt}(X_{\ast})) \quad ,
\]
where the second member moreover is endowed with it's natural abelian
group structure (it's components are abelian group objects of U(k) and the
simplicial maps are additive).  To pass from (W) to (W ?) ,  it is felt
that the essential step is that (V) above be a quasi-isomorphism,\var{\xout{\unsure{x}}}{}{typewriter deletion}{ps_134-7_038-v0006} hence\hand{,}
applying the functor W\hand{,} we should get a \var{\xout{quasi-isomorphism}}{}{typewriter deletion}{ps_134-7_038-v0007} weak equivalence\hand{,}
hence an isomorphism in the \var{\xout{suitable}}{}{typewriter deletion}{ps_134-7_038-v0008} derived category Hot$_n$(k),  hence
(W)  implies (W ?).  The \var{\xout{main}}{}{typewriter deletion}{ps_134-7_038-v0009} main \var{\deleted{trouble here}}{}{deletion}{ps_134-7_038-v0010}\var{\inserted{flaw in this ``argument''}}{\hand{flaw in this ``argument''}}{insertion}{ps_134-7_038-v0011} comes from the W$^{\hand{\wedge}}$ in the
second member of (W),  which is'nt quite the same as W definitely.  Thus
some\var{\inserted[typed]{further}}{further}{typed insertion}{ps_134-7_038-v0012} amount of work will be needed presumably to get (W ?) from (W).
Of course,  we cant possibly just keep (W) as it is,  as for getting
devissage we need a map in $\hand{\mathcal{M}}_n$(k),  whereas the map (W) is\var{\inserted{just}}{\hand{just}}{insertion}{ps_134-7_038-v0013} a map of
semisimplicial sheaves on Aff$_{/k}$  , \var{\deleted{but}}{}{deletion}{ps_134-7_038-v0014}\var{\inserted{where}}{\hand{where}}{insertion}{ps_134-7_038-v0015} the second member is \uline{not} in
$\hand{\mathcal{M}}_n$(k)  i.e. it's components are not in U(k).  Once we got (W ?)
factoring (W) up to homotopy (NB of course we assume $X_{\ast}$ to be an sK
complex in all this),  we still need a reasonable notion of homotopy
fibers of maps in $\hand{\mathcal{M}}_n$(k),  in order to push through the inductive step.

\qquad Thus,  a large part of the weight of the work ahead may well lie
upon \uline{developping the standard homotopy constructions} \uline{within the model}
\uline{category} $\hand{\mathcal{M}}_n$(k),  as contemplated on page 464.  This should be fun,  if
it can be done indeed\var{\over{.}{!}}{\hand{!}}{writing over}{ps_134-7_038-v0016}  One difficulty here seems to be that \var{\xout{there}}{}{typewriter deletion}{ps_134-7_038-v0017} Quil-%
len's  standard machines  wo'nt work,\var{\inserted[typed]{not}}{not}{typed insertion}{ps_134-7_038-v0018} "telles quelles" at any rate,  becau-%
se of the category U(k) failing to be stable under finite limits - it
does'nt even have fi\var{\over{t}{b}}{\hand{b}}{writing over}{ps_134-7_038-v0019}er products.  But I think I'll stop my ponderings
on schematization here ...

\margin{38}

% Visual collation: both supplied image halves inspected; see CHECK_DECISIONS_134_7_8.tsv.

%% END SOURCE ps_134-7_038
\ifnum\EditionMode=1
\psapparatusstart
\psentry{ps_134-7_038-v0001}{1}{typed insertion}{\inserted[typed]{k-modules}}{k-modules}{The typed insertion is incorporated at its indicated position; insertion brackets leave the running text.}
\psentry{ps_134-7_038-v0002}{2}{typewriter deletion}{\xout{endowed with}}{}{The typewriter strikes out this material; it leaves the corrected running text.}
\psentry{ps_134-7_038-v0003}{3}{typewriter deletion}{\xout{soft}}{}{The typewriter strikes out this material; it leaves the corrected running text.}
\psentry{ps_134-7_038-v0004}{4}{typewriter deletion}{$\xout{L(X_{\ast})}$}{}{The typewriter strikes out this material; it leaves the corrected running text.}
\psentry{ps_134-7_038-v0005}{5}{typewriter deletion}{\xout{being}}{}{The typewriter strikes out this material; it leaves the corrected running text.}
\psentry{ps_134-7_038-v0006}{6}{typewriter deletion}{\xout{\unsure{x}}}{}{The typewriter strikes out this material; it leaves the corrected running text.}
\psentry{ps_134-7_038-v0007}{7}{typewriter deletion}{\xout{quasi-isomorphism}}{}{The typewriter strikes out this material; it leaves the corrected running text.}
\psentry{ps_134-7_038-v0008}{8}{typewriter deletion}{\xout{suitable}}{}{The typewriter strikes out this material; it leaves the corrected running text.}
\psentry{ps_134-7_038-v0009}{9}{typewriter deletion}{\xout{main}}{}{The typewriter strikes out this material; it leaves the corrected running text.}
\psentry{ps_134-7_038-v0010}{10}{deletion}{\deleted{trouble here}}{}{The author cancels this material by hand; it leaves the corrected running text.}
\psentry{ps_134-7_038-v0011}{11}{insertion}{\inserted{flaw in this ``argument''}}{\hand{flaw in this ``argument''}}{The handwritten insertion is incorporated at its indicated position; insertion brackets leave the running text.}
\psentry{ps_134-7_038-v0012}{12}{typed insertion}{\inserted[typed]{further}}{further}{The typed insertion is incorporated at its indicated position; insertion brackets leave the running text.}
\psentry{ps_134-7_038-v0013}{13}{insertion}{\inserted{just}}{\hand{just}}{The handwritten insertion is incorporated at its indicated position; insertion brackets leave the running text.}
\psentry{ps_134-7_038-v0014}{14}{deletion}{\deleted{but}}{}{The author cancels this material by hand; it leaves the corrected running text.}
\psentry{ps_134-7_038-v0015}{15}{insertion}{\inserted{where}}{\hand{where}}{The handwritten insertion is incorporated at its indicated position; insertion brackets leave the running text.}
\psentry{ps_134-7_038-v0016}{16}{writing over}{\over{.}{!}}{\hand{!}}{The superposed handwritten reading replaces the underlying reading.}
\psentry{ps_134-7_038-v0017}{17}{typewriter deletion}{\xout{there}}{}{The typewriter strikes out this material; it leaves the corrected running text.}
\psentry{ps_134-7_038-v0018}{18}{typed insertion}{\inserted[typed]{not}}{not}{The typed insertion is incorporated at its indicated position; insertion brackets leave the running text.}
\psentry{ps_134-7_038-v0019}{19}{writing over}{\over{t}{b}}{\hand{b}}{The superposed handwritten reading replaces the underlying reading.}
\psapparatusend
\fi
\pssource{ps_134-7_039}{134-7/39}{https://grothendieck.umontpellier.fr/web/viewer.html?file=../134-7.pdf\#page=40}
\psnumber{ps_134-7_039-v0001}{1}
\psnumber{ps_134-7_039-v0002}{2}
\psnumber{ps_134-7_039-v0003}{3}
\psnumber{ps_134-7_039-v0004}{4}
\psnumber{ps_134-7_039-v0005}{5}
\psnumber{ps_134-7_039-v0006}{6}
\psnumber{ps_134-7_039-v0007}{7}
\psnumber{ps_134-7_039-v0008}{8}
\psnumber{ps_134-7_039-v0009}{9}
\psnumber{ps_134-7_039-v0010}{10}
%% BEGIN SOURCE ps_134-7_039
% Pursuing Stacks (cotes 134-2 to 134-8): archive cote 134-7, site page 39 (PDF page 40), typescript page 481.
% typescript. Draft compared: ps6.tex:1886-1953.
% AI-transcribed by Claude Opus 5.5 (claude-opus-5-5, adaptive thinking, effort high), Message Batches API, batch ps3 (msgbatch_01XdFmuVVBAcgc7EFxdtCqsj), 2026-11-12; reviewed against both supplied scan halves; see CHECK_DECISIONS.tsv.
\scanpage{134-7/39}
\mspage{481}

28.8.

115.   For the last four days,  while reflecting on ``schematization'',
each time I think I am going to be through with that unforeseen green
apple within an hour or two\hand{,} and get back to ``l'ordre du jour'' -  and
overnight something else still appears I feel I should still look into
just a little;  and there I am again,  sure enough,  with some extra refle\unsure{c}%
% NOTE: the right edge of the scan cuts the end of this line after "refle"; a hyphen may be lost there
tion on ``schematic linearization'' which I had'nt quite understood yet,
it appears to me now.   These last days I had given up numbering formulas
as usual by arab ciphers (1), (2)  etc,  as I did'nt want to ``cut'' the
numbering of that unending ``review'' of sections 104 to 109 which was'nt
quite finished yet got it only till formula (136).  But now I will stop
this nonsense with numberings (a) (b) and (A) (B),  after all even if
there are in-between ``arab'' formulas now,  this does'nt prevent me,when
it comes to it,to start a \var{\inserted[typed]{``review''}}{``review''}{typed insertion}{ps_134-7_039-v0001} section with formula (137) and go on till
(1000) if I like...   And now to the schematic linearization functor
again,  for unipotent bundles !

When writing up that schematization program yesterday,  some techni-%
cal difficulties appeared at the end (see page before) \var{\deleted{\unsure{in}}}{}{deletion}{ps_134-7_039-v0002}\var{\inserted[typed]{for}}{for}{typed insertion}{ps_134-7_039-v0003} a proper under-%
standing of the relationship between the two linearizations $L_{pt}$ and $L$ ,
\var{\deleted{needed}}{}{deletion}{ps_134-7_039-v0004} in order to define,  in a suitable derived category,  a map
\[ X_{*}\hand{\to} W(L_{pt}(X_{*})) \]
using the canonical term-by-term \var{\xout{map}}{}{typewriter deletion}{ps_134-7_039-v0005} exponential map
\[ X_{*}\hand{\to} W^{\hand{\wedge}}L(X_{*}) \quad . \]
It seems to me \var{\inserted{that}}{\hand{that}}{insertion}{ps_134-7_039-v0006} the exact significance of the object $L(X_{*})$ or \var{\xout{$W^{\hand{\wedge}}L(X$}}{}{typewriter deletion}{ps_134-7_039-v0007}
$W^{\hand{\wedge}}L(X_{*})$ \var{\xout{wa}}{}{typewriter deletion}{ps_134-7_039-v0008}is'nt quite understood yet,  and \var{\inserted{that}}{\hand{that}}{insertion}{ps_134-7_039-v0009} the confusion which occured
between which usual kind of linearization we should compare this with,
whether  $X\hand{\mapsto} \underline{Z}^{(X)}$  as I did first ,  or  $X\hand{\mapsto} k^{(X)}$  as it occured to me
yesterday (pages 478-479)\hand{,} is quite typical of this lack of understan-%
ding.  It now occured to me that neither term,  for a general ground
ring $k$ (namely not assuming $k=\underline{Z}$ ),  is reasonable,  whereas the reaso-%
nable ``usual'' kind of linearization comparing with $L(X)$ \var{\xout{or $L(X_{*})$}}{}{typewriter deletion}{ps_134-7_039-v0010}
(when $X$ is a unipotent bundle or a ss complex of such)  is
\[ \hand{(1)} \qquad X\hand{\mapsto} \underline{O}_{k}^{(X)} \quad , \]
where $\underline{O}_{k}$  is the basic \uline{quasicoherent}  sheaf of rings over $k$  i.e. over
$\mathrm{Aff}_{/k}$ ,  given by the tautological functor
\[ \hand{(2)} \qquad \hand{\underline{O}_{k}:}\ (\mathrm{Aff}_{\hand{/k}})^{\hand{o}} \hand{\xrightarrow{\approx}} \mathrm{Alg}_{\hand{/k}} \hand{\longrightarrow} (\mathrm{Rings}) \ , \ \hand{k'\mapsto k',} \]
\margin{39}

% Visual collation: both supplied image halves inspected; see CHECK_DECISIONS_134_7_8.tsv.

%% END SOURCE ps_134-7_039
\ifnum\EditionMode=1
\psapparatusstart
\psentry{ps_134-7_039-v0001}{1}{typed insertion}{\inserted[typed]{``review''}}{``review''}{The typed insertion is incorporated at its indicated position; insertion brackets leave the running text.}
\psentry{ps_134-7_039-v0002}{2}{deletion}{\deleted{\unsure{in}}}{}{The author cancels this material by hand; it leaves the corrected running text.}
\psentry{ps_134-7_039-v0003}{3}{typed insertion}{\inserted[typed]{for}}{for}{The typed insertion is incorporated at its indicated position; insertion brackets leave the running text.}
\psentry{ps_134-7_039-v0004}{4}{deletion}{\deleted{needed}}{}{The author cancels this material by hand; it leaves the corrected running text.}
\psentry{ps_134-7_039-v0005}{5}{typewriter deletion}{\xout{map}}{}{The typewriter strikes out this material; it leaves the corrected running text.}
\psentry{ps_134-7_039-v0006}{6}{insertion}{\inserted{that}}{\hand{that}}{The handwritten insertion is incorporated at its indicated position; insertion brackets leave the running text.}
\psentry{ps_134-7_039-v0007}{7}{typewriter deletion}{\xout{$W^{\hand{\wedge}}L(X$}}{}{The typewriter strikes out this material; it leaves the corrected running text.}
\psentry{ps_134-7_039-v0008}{8}{typewriter deletion}{\xout{wa}}{}{The typewriter strikes out this material; it leaves the corrected running text.}
\psentry{ps_134-7_039-v0009}{9}{insertion}{\inserted{that}}{\hand{that}}{The handwritten insertion is incorporated at its indicated position; insertion brackets leave the running text.}
\psentry{ps_134-7_039-v0010}{10}{typewriter deletion}{\xout{or $L(X_{*})$}}{}{The typewriter strikes out this material; it leaves the corrected running text.}
\psapparatusend
\fi
\pssource{ps_134-7_040}{134-7/40}{https://grothendieck.umontpellier.fr/web/viewer.html?file=../134-7.pdf\#page=41}
\psnumber{ps_134-7_040-v0001}{1}
\psnumber{ps_134-7_040-v0002}{2}
\psnumber{ps_134-7_040-v0003}{3}
\psnumber{ps_134-7_040-v0004}{4}
\psnumber{ps_134-7_040-v0005}{5}
\psnumber{ps_134-7_040-v0006}{6}
\psnumber{ps_134-7_040-v0007}{7}
\psnumber{ps_134-7_040-v0008}{8}
\psnumber{ps_134-7_040-v0009}{9}
%% BEGIN SOURCE ps_134-7_040
% Pursuing Stacks (cotes 134-2 to 134-8): archive cote 134-7, site page 40 (PDF page 41), typescript page 482.
% typescript. Draft compared: ps6.tex:1935-2003.
% AI-transcribed by Claude Opus 5.5 (claude-opus-5-5, adaptive thinking, effort high), Message Batches API, batch ps3 (msgbatch_01XdFmuVVBAcgc7EFxdtCqsj), 2026-11-12; reviewed against both supplied scan halves; see CHECK_DECISIONS.tsv.
\scanpage{134-7/40}
\mspage{482}

associating to any affine scheme \var{\inserted[typed]{S=}}{S=}{typed insertion}{ps_134-7_040-v0001} Spec(k') over k the ring of sections
of  it\var{\deleted{'}}{}{deletion}{ps_134-7_040-v0002}s usual Zariski structure sheaf,  which ring is canonically isomor-%
phic to k' itself !  The operation (1) is the usual linearization opera-%
tion with respect to this sheaf of rings,  working in the topos of
fpqc sheaves of sets over k which we described at some length in section
111  (p.447).  As I was fearing that working in such a thing would cause
anguish to a number of prospective readers,  I took pains to translate
unipotent bundles from the geometric language which is the suggestive
one,  to the language of commutative algebra which is more liable to hide
than to disclose geometrical meaning;  so much so that in the process I
m\var{\xout{x}}{}{typewriter deletion}{ps_134-7_040-v0003}yself lost contact somewhat with the geometric flavor,  and more speci-%
fically still with this basic fact,  that in our context of unipotent
bundles and complexes of such,  the ``natural'' coefficients for cohomology
\margin{$X(n)_*,\hand{\pi}_{n+1}$} (such as the $H^{n+2}(\unsure{X(n)_*},\hand{\pi}_{n+1})$ groups occuring in Posthikov dévissage) are
by no means``\uline{discrete}''\hand{/}one's such as $\underline{Z}$ or k ,  but \uline{quasi}-\uline{coherent} one's
namely provided by quasi-coherent sheaves of \var{\xout{modules}}{}{typewriter deletion}{ps_134-7_040-v0004} $\underline{O}_k$-modules or com-%
plexes of such.  Thus in the above Postnikov obstruction group, $\hand{\pi}_{n+1}$
does \uline{not} stand as a constant group of coefficients (if it was,  this
would drag us into the niceties and difficulties of etale cohomology
for the components $X(n)_i$ of the semisimplicial unipotent bundle $X(n)_*$);
but using the k-module structure of $\hand{\pi}_{n+1}$ for defining a \uline{quasi}-\uline{coherent}
\hand{s}heaf of modules  $W(\hand{\pi}_{n+1})$  ``over k'' i.e.  over $Aff_{/k}$ ,  it is this
``continuous'' sheaf (or ``vector bundle'') over k, lifted of course to the
various components \var{\xout{\unsure{at}}}{}{typewriter deletion}{ps_134-7_040-v0005} $X(n)_i$ ,  which yields the correct answer.  This
was kind of clear in my mind the very first day when I started reflecting
on schematization, \var{\inserted{even}}{\hand{even}}{insertion}{ps_134-7_040-v0006} before \var{\deleted{even}}{}{deletion}{ps_134-7_040-v0007} introducing formally unipotent bundles
(pages 443-444),  but this instinctive understanding later became dulled
somewhat,  largely due it seems to me to the concession I had made to
\var{\xout{commutative}}{}{typewriter deletion}{ps_134-7_040-v0008} algebra,  giving up to some extent the language of geometry.

Let's recall that the operation (1) may be defined as the solution
of a universal problem,  namely sending the non-linear object X into
a ``linear'' one namely into a sheaf of $\underline{O}_k$-modules (or a $\underline{O}_k$-Module,  as
we'll simply say).  This is expressed by a canonical \var{\xout{map}}{}{typewriter deletion}{ps_134-7_040-v0009} map of sheaves
of sets
\[ (3) \qquad X \hand{\longrightarrow} \underline{O}_k^{(X)} \]
(which I am tempted to call the ``\uline{exponential}'' \uline{map} for X, and denote by
a corresponding symbol such as $exp_X$ ) ,  giving rise,  for every Module
(over $\underline{O}_k$)  to a corresponding map which is \uline{bijective}

\margin{40}

% Visual collation: both supplied image halves inspected; see CHECK_DECISIONS_134_7_8.tsv.

%% END SOURCE ps_134-7_040
\ifnum\EditionMode=1
\psapparatusstart
\psentry{ps_134-7_040-v0001}{1}{typed insertion}{\inserted[typed]{S=}}{S=}{The typed insertion is incorporated at its indicated position; insertion brackets leave the running text.}
\psentry{ps_134-7_040-v0002}{2}{deletion}{\deleted{'}}{}{The author cancels this material by hand; it leaves the corrected running text.}
\psentry{ps_134-7_040-v0003}{3}{typewriter deletion}{\xout{x}}{}{The typewriter strikes out this material; it leaves the corrected running text.}
\psentry{ps_134-7_040-v0004}{4}{typewriter deletion}{\xout{modules}}{}{The typewriter strikes out this material; it leaves the corrected running text.}
\psentry{ps_134-7_040-v0005}{5}{typewriter deletion}{\xout{\unsure{at}}}{}{The typewriter strikes out this material; it leaves the corrected running text.}
\psentry{ps_134-7_040-v0006}{6}{insertion}{\inserted{even}}{\hand{even}}{The handwritten insertion is incorporated at its indicated position; insertion brackets leave the running text.}
\psentry{ps_134-7_040-v0007}{7}{deletion}{\deleted{even}}{}{The author cancels this material by hand; it leaves the corrected running text.}
\psentry{ps_134-7_040-v0008}{8}{typewriter deletion}{\xout{commutative}}{}{The typewriter strikes out this material; it leaves the corrected running text.}
\psentry{ps_134-7_040-v0009}{9}{typewriter deletion}{\xout{map}}{}{The typewriter strikes out this material; it leaves the corrected running text.}
\psapparatusend
\fi
\pssource{ps_134-7_041}{134-7/41}{https://grothendieck.umontpellier.fr/web/viewer.html?file=../134-7.pdf\#page=42}
\psnumber{ps_134-7_041-v0001}{1}
\psnumber{ps_134-7_041-v0002}{2}
\psnumber{ps_134-7_041-v0003}{3}
\psnumber{ps_134-7_041-v0004}{4}
\psnumber{ps_134-7_041-v0005}{5}
%% BEGIN SOURCE ps_134-7_041
% Pursuing Stacks (cotes 134-2 to 134-8): archive cote 134-7, site page 41 (PDF page 42), typescript page 483.
% typescript. Draft compared: ps6.tex:1985-2059.
% AI-transcribed by Claude Opus 5.5 (claude-opus-5-5, adaptive thinking, effort high), Message Batches API, batch ps3 (msgbatch_01XdFmuVVBAcgc7EFxdtCqsj), 2026-11-12; reviewed against both supplied scan halves; see CHECK_DECISIONS.tsv.
\scanpage{134-7/41}
\mspage{483}

(4) \[ \mathrm{Hom}_{\underline{O}_k}( \underline{O}_k^{(X)}, F ) \hand{\simeq} \mathrm{Hom}(\hand{X},F\hand{)} \hand{(}\var{\deleted{\unsure{def}}}{}{deletion}{ps_134-7_041-v0001}\hand{\simeq} \hand{\Gamma}(X,F_X)\hand{)}, \]
where in the last member (included as a more geometric interpretation of
the second) $F_X$ denotes the restriction of $F$ to the object $X$, more accu-%
rately to the topos (or site) induced \var{\xout{up}}{}{typewriter deletion}{ps_134-7_041-v0002}on $X$ b\hand{y} the ambient topos (or
site) we are working in. Thus we may indeed view (1) as the most per-%
fect notion of linearization, as far as generality goes - it makes
sense of course in any ringed topos (without even a commutativity assump-%
tion !). The only trouble is that, even for such a down to earth $X$ as
a unipotent bundle, \var{\xout{say}}{}{typewriter deletion}{ps_134-7_041-v0003} the standard affine line $E_k^1$ say, the sheaf
$\underline{O}_k^{(X)}$ in (1) is not quasi-coherent and therefore not too amenable it
seems to computations - thus we get easily from (4) a canonical map
(for general $X$)
\[ k^{(X(k))} \hand{\to} \hand{\Gamma}(k, \underline{O}_k^{(X)}) \ (= \underline{O}_k^{(X)}\ (k)) \]
(where the $\hand{\Gamma}$ in the second member denotes sections over $k$ i.e. value of
a functor on $\mathrm{Alg}_{/k}$ on the initial object $k$ , and remembering in the
first member that the ring of sections of $\underline{O}_k$ is $k$ ), but I would be at
a loss to make a guess \var{\xout{when}}{}{typewriter deletion}{ps_134-7_041-v0004}\var{\inserted[typed]{as}}{as}{typed insertion}{ps_134-7_041-v0005} for reasonable conditions for this map to
be an isomorphism ! This may seem a prohibitive ``contra'' against using
at all such huge sheaves as $\underline{O}_k^{(X)}$ , the point though is that in most
questions where such linearizations are introduced (mainly questions
where interest lies in computing cohomology invariants), one is practi-%
cally never interested in taking the groups of sections of these, but
rather in looking at their maps into sheaves of modules $F$ precisely,
which is achieved by (4), or taking more generally their $\mathrm{Ext}^i$ with
such an $F$, which is achieved by the similar familiar formula

(5) \[ \mathrm{Ext}^i_{\underline{O}_k}(\underline{O}_k^{(X)}, F ) \hand{\simeq} H^i(X,F_X) \quad , \]
more neatly

(5') \[ \mathrm{RHom}_{\underline{O}_k}( \underline{O}_k^{(X)}, F ) \hand{\simeq} \mathrm{R}\hand{\Gamma}_X(F_X) \quad , \]
valid of course again for any ringed topos. In the present context
however, the ``coefficients'' $F$ we are interested in, as was just pointed
out, are not arbitrary $\underline{O}_k$-Modules, but rather \uline{quasi}-\uline{coherent} ones.
Thus, if we get a variant of (1)

(6) \[ X \hand{\longrightarrow} L(X) \]
with $L(X)$ some \uline{quasi}-\uline{coherent} Module, giving rise to (4), or even to
(5) and (5'), this would be for us a perfectly good substitue for (3),
which would deserve the name of a ``\uline{quasicoherent} \uline{envelope}'' of $X$. Of

\margin{41}

% Visual collation: both supplied image halves inspected; see CHECK_DECISIONS_134_7_8.tsv.

%% END SOURCE ps_134-7_041
\ifnum\EditionMode=1
\psapparatusstart
\psentry{ps_134-7_041-v0001}{1}{deletion}{$\deleted{\unsure{def}}$}{}{The author cancels this material by hand; it leaves the corrected running text.}
\psentry{ps_134-7_041-v0002}{2}{typewriter deletion}{\xout{up}}{}{The typewriter strikes out this material; it leaves the corrected running text.}
\psentry{ps_134-7_041-v0003}{3}{typewriter deletion}{\xout{say}}{}{The typewriter strikes out this material; it leaves the corrected running text.}
\psentry{ps_134-7_041-v0004}{4}{typewriter deletion}{\xout{when}}{}{The typewriter strikes out this material; it leaves the corrected running text.}
\psentry{ps_134-7_041-v0005}{5}{typed insertion}{\inserted[typed]{as}}{as}{The typed insertion is incorporated at its indicated position; insertion brackets leave the running text.}
\psapparatusend
\fi
\pssource{ps_134-7_042}{134-7/42}{https://grothendieck.umontpellier.fr/web/viewer.html?file=../134-7.pdf\#page=43}
\psnumber{ps_134-7_042-v0001}{1}
\psnumber{ps_134-7_042-v0002}{2}
\psnumber{ps_134-7_042-v0003}{3}
\psnumber{ps_134-7_042-v0004}{4}
\psnumber{ps_134-7_042-v0005}{5}
\psnumber{ps_134-7_042-v0006}{6}
\psnumber{ps_134-7_042-v0007}{7}
%% BEGIN SOURCE ps_134-7_042
% Pursuing Stacks (cotes 134-2 to 134-8): archive cote 134-7, site page 42 (PDF page 43), typescript page 484.
% typescript. Draft compared: ps6.tex:2038-2124.
% AI-transcribed by Claude Opus 5.5 (claude-opus-5-5, adaptive thinking, effort high), Message Batches API, batch ps3 (msgbatch_01XdFmuVVBAcgc7EFxdtCqsj), 2026-11-12; reviewed against both supplied scan halves; see CHECK_DECISIONS.tsv.
\scanpage{134-7/42}
\mspage{484}
course,  this Module L(X)  would be unique up to unique isomorphism,  as
the solution of a universal problem embodied by (4),  namely as the
quasi-coherent Module representing the functor
\[ (7) \qquad F \hand{\longmapsto} \mathrm{Hom}(X,F) \hand{\simeq} \hand{\Gamma}(X,F_X) \]
on the category this time of all \uline{quasicoherent} $\underline{O}_k$-Modules.

For the unipotent schematization story,  we are more specifically
interested in the case when X comes from a quasicoherent Module itself,
by forgetting it's module structure.  Now \var{\xout{the}}{}{typewriter deletion}{ps_134-7_042-v0001} as well known,  the
functor
\[ (8) \qquad M \hand{\mapsto} W(M) \quad : \quad \mathrm{Ab}_k = (k\text{-Mod}) \qquad \begin{array}[t]{l} \text{category of quasicoherent} \\ \underline{O}_k\text{-Modules, } \hand{\mathrm{Qucoh}(k) \text{ say}} \end{array} \]
is an equivalence of categories.  Thus for X defined by such an M,  the
question of representability of (7)  within the category of quasicoherent
Modules,  amounts to the similar question in $\mathrm{Ab}_k$  for the functor
\[ (7') \qquad N \hand{\longmapsto} \mathrm{Hom}(W(M),W(N)) \qquad , \]
where the Hom denotes homomorphisms of sheaves \uline{of sets} of course.  Now,
as suggested first\hand{,} somewhat vaguely still\hand{,} in section 111 (page 450),
we have an alternative expression of this functor\hand{,} \var{\xout{as}}{}{typewriter deletion}{ps_134-7_042-v0002} \hand{via}
\[ (9) \qquad \mathrm{Hom}(W(M),W(N)) \hand{\simeq} \mathrm{Homcont}_k(\hand{\Gamma}_k^{\hand{\wedge}}(M),N) \hand{\simeq} \var{\xout{\underset{\hand{\longrightarrow}}{\mathrm{lim}} \hand{\Gamma}_k(M)(i}}{}{typewriter deletion}{ps_134-7_042-v0003} \]
\margin{\textcircled{$\mathrm{Hom}_k($}} \hand{\unsure{$\nearrow$|caret}}
\[ \underset{i}{\underset{\hand{\longrightarrow}}{\mathrm{lim}}} \mathrm{Hom}\var{\xout{\mathrm{cont}}}{}{typewriter deletion}{ps_134-7_042-v0004}_k(\hand{\Gamma}_k(M)(i),N) \quad , \]
where in the second member $\mathrm{Homcont}_k$ denotes the set of k-homomorphisms
\var{\inserted[typed]{\xout{product}}}{}{typed insertion}{ps_134-7_042-v0005}
which are continuous \var{\xout{for the topology induces by the product of}}{}{typewriter deletion}{ps_134-7_042-v0006} on
\[ (10) \qquad \hand{\Gamma}_k^{\hand{\wedge}}(M) = \hand{\prod_{i \geqslant 0}} \hand{\Gamma}_k^{\,i}(M) \]
\unsure{(}endowed with the product of discrete topologies\unsure{)},  and in the third
we have written
\[ (10') \qquad \hand{\Gamma}_k(M)(i) = \hand{\prod}_{j \hand{\leqslant} i} \hand{\Gamma}_k^{\,j}(M) \]
for the product of the i first factors occuring in (10).   The map (9)
is \var{\xout{induced by}}{}{typewriter deletion}{ps_134-7_042-v0007} deduced in evident way from the exponential map
\[ (11) \qquad M \hand{\longrightarrow} W^{\hand{\wedge}}\hand{\Gamma}_k^{\hand{\wedge}}(M) \overset{\mathrm{def}}{=} \underset{i}{\underset{\hand{\longleftarrow}}{\mathrm{lim}}} W(\hand{\Gamma}_k(M)(i)) \qquad . \]
(NB.  The relation between the description (9) of maps $W(M) \hand{\rightarrow} W(N)$ with
the description given p.451  in terms of maps (p) from $\hand{\Gamma}_k^{\hand{\wedge}}(M)$ to
$\hand{\Gamma}_k^{\hand{\wedge}}(N)$ ,  is by associating to such a map f it's composition with the
projection of the target upon it's factor N ...).  An uncorrect way of
expressing (9),  which I slipped into in section 111 and kind of remained
in till now,  is by pretending that the k-module  $\hand{\Gamma}_k^{\hand{\wedge}}(M)$ represents the

\margin{42}

% Visual collation: both supplied image halves inspected; see CHECK_DECISIONS_134_7_8.tsv.

%% END SOURCE ps_134-7_042
\ifnum\EditionMode=1
\psapparatusstart
\psentry{ps_134-7_042-v0001}{1}{typewriter deletion}{\xout{the}}{}{The typewriter strikes out this material; it leaves the corrected running text.}
\psentry{ps_134-7_042-v0002}{2}{typewriter deletion}{\xout{as}}{}{The typewriter strikes out this material; it leaves the corrected running text.}
\psentry{ps_134-7_042-v0003}{3}{typewriter deletion}{$\xout{\underset{\hand{\longrightarrow}}{\mathrm{lim}} \hand{\Gamma}_k(M)(i}$}{}{The typewriter strikes out this material; it leaves the corrected running text.}
\psentry{ps_134-7_042-v0004}{4}{typewriter deletion}{$\xout{\mathrm{cont}}$}{}{The typewriter strikes out this material; it leaves the corrected running text.}
\psentry{ps_134-7_042-v0005}{5}{typed insertion}{\inserted[typed]{\xout{product}}}{}{The typed insertion is incorporated at its indicated position; insertion brackets leave the running text.}
\psentry{ps_134-7_042-v0006}{6}{typewriter deletion}{\xout{for the topology induces by the product of}}{}{The typewriter strikes out this material; it leaves the corrected running text.}
\psentry{ps_134-7_042-v0007}{7}{typewriter deletion}{\xout{induced by}}{}{The typewriter strikes out this material; it leaves the corrected running text.}
\psapparatusend
\fi
\pssource{ps_134-7_043}{134-7/43}{https://grothendieck.umontpellier.fr/web/viewer.html?file=../134-7.pdf\#page=44}
\psnumber{ps_134-7_043-v0001}{1}
\psnumber{ps_134-7_043-v0002}{2}
\psnumber{ps_134-7_043-v0003}{3}
\psnumber{ps_134-7_043-v0004}{4}
\psnumber{ps_134-7_043-v0005}{5}
\psnumber{ps_134-7_043-v0006}{6}
\psnumber{ps_134-7_043-v0007}{7}
\psnumber{ps_134-7_043-v0008}{8}
%% BEGIN SOURCE ps_134-7_043
% Pursuing Stacks (cotes 134-2 to 134-8): archive cote 134-7, site page 43 (PDF page 44), typescript page 485.
% typescript. Draft compared: ps6.tex:2108-2184.
% AI-transcribed by Claude Opus 5.5 (claude-opus-5-5, adaptive thinking, effort high), Message Batches API, batch ps3 (msgbatch_01XdFmuVVBAcgc7EFxdtCqsj), 2026-11-12; reviewed against both supplied scan halves; see CHECK_DECISIONS.tsv.
\scanpage{134-7/43}
\mspage{485}
functor (7'), this is clearly false, as we do not have any canonical
map from W(M) into $W(\hand{\Gamma}_k^{\hand{\wedge}}(M))$ , only into $W^{\hand{\wedge}}\hand{\Gamma}_k^{\hand{\wedge}}(M)$ \var{\xout{wh}}{}{typewriter deletion}{ps_134-7_043-v0001} - we have an
embedding
\[ (12) \qquad W(\hand{\Gamma}_k^{\hand{\wedge}}(M)) \hand{\hookrightarrow} W^{\hand{\wedge}}\hand{\Gamma}_k^{\hand{\wedge}}(M) \quad , \]
but it is clear that in general, the exponential map (11) does not
factor through the first term in (12). (It does of course when we look
at sections over k only, but when we go over to a general k', we hit
into the trouble that formation of inverse limits des not commute with
ring extension $\hand{\otimes}_k k'$ !). We may however express (9) by stating that the
functor (7') is ``\uline{prorepresentable}'' by the \uline{pro}-\uline{object}
\[ (13) \qquad \mathrm{Pr}\bullet\, \hand{\Gamma}_k(M) \overset{\mathrm{def}}{=} (\hand{\Gamma}_k(M)(i))_{i \hand{\geqslant} 0} \quad \text{in} \quad \mathrm{Pro}(\mathrm{Ab}_k) \quad , \]
this is even a \uline{strict} proobject (the transition morphisms are epimor-%
phisms), which implies that the functor it prorepresents is represen-%
table iff this projective system is ``essentially constant'' in the most
trivial sense, which means here
\[ \hand{\Gamma}_k^{\,i}(M) \;=\; 0 \quad \text{for large } i \quad , \]
a condition which presumably is satisfied only for $M=0$ ! Thus, the
``correct'' interpretation of non-pointed \var{\inserted[typed]{\hand{(}quasi-coherent\hand{)}}}{\hand{(}quasi-coherent\hand{)}}{typed insertion}{ps_134-7_043-v0002} linearization seems to me to
be the corresponding functor, which I would like now to call $L_k$ or
simply L as before \hand{but with slightly different meaning:}
\[ (14) \qquad L \text{ or } L_k \;:\; U(k) \hand{\longrightarrow} \mathrm{Pr}\bullet(\mathrm{Ab}_k) \hand{\underline{\approx}} \mathrm{Pr}\bullet(\mathrm{Qucoh}(k)) \quad , \]
where \var{\xout{Qu}}{}{typewriter deletion}{ps_134-7_043-v0003} Qucoh(k) is defined in (8). In computational terms, I
would like to view \var{\xout{Qucohx}}{}{typewriter deletion}{ps_134-7_043-v0004}\var{\inserted{L(X)}}{\hand{L(X)}}{insertion}{ps_134-7_043-v0005} (for a unipotent bundle X) \var{\deleted{\hand{\unsure{L(X)}}}}{}{deletion}{ps_134-7_043-v0006} to be a
pro-k-module, but in terms of geometric intuition, I would see it
rather as a pro-$\underline{O}_k$-Module, i.e. essentially as an inverse system of
quasicoherent modules. It is in these latter terms that the constructi-%
on we just gave generalises to unipotent bundles over arbitrary ground
schemes, not necessarily affine ones. As for the ``pointed'' quasicohe-%
rent linearization functor
\[ (15) \qquad L_{pt} \text{ or } L_{k\ pt} \;:\; U(k)^{\hand{\bullet}} \hand{\longrightarrow} (\mathrm{Ab}_k) \hand{\underline{\approx}} \mathrm{Qucoh}(k) \quad , \]
which I like best to view as taking values $L_{pt}(X)$ which are quasicoherent
sheaves, it maps into L by
\[ (1\var{\over{\unsure{5}}{6}}{\hand{6}}{writing over}{ps_134-7_043-v0007}) \qquad L_{pt}(X) \hand{\hookrightarrow} L(X) \quad , \]
interpreting objects of \var{\xout{categoriexx}}{}{typewriter deletion}{ps_134-7_043-v0008} a category as special cases of pro-%
objects. We'll denote by
\[ \hand{(17)} \qquad WL(X) \hand{\in} \mathrm{Pro}(U(k)) \]
\margin{43}

% Visual collation: both supplied image halves inspected; see CHECK_DECISIONS_134_7_8.tsv.

%% END SOURCE ps_134-7_043
\ifnum\EditionMode=1
\psapparatusstart
\psentry{ps_134-7_043-v0001}{1}{typewriter deletion}{\xout{wh}}{}{The typewriter strikes out this material; it leaves the corrected running text.}
\psentry{ps_134-7_043-v0002}{2}{typed insertion}{\inserted[typed]{\hand{(}quasi-coherent\hand{)}}}{\hand{(}quasi-coherent\hand{)}}{The typed insertion is incorporated at its indicated position; insertion brackets leave the running text.}
\psentry{ps_134-7_043-v0003}{3}{typewriter deletion}{\xout{Qu}}{}{The typewriter strikes out this material; it leaves the corrected running text.}
\psentry{ps_134-7_043-v0004}{4}{typewriter deletion}{\xout{Qucohx}}{}{The typewriter strikes out this material; it leaves the corrected running text.}
\psentry{ps_134-7_043-v0005}{5}{insertion}{\inserted{L(X)}}{\hand{L(X)}}{The handwritten insertion is incorporated at its indicated position; insertion brackets leave the running text.}
\psentry{ps_134-7_043-v0006}{6}{deletion}{\deleted{\hand{\unsure{L(X)}}}}{}{The author cancels this material by hand; it leaves the corrected running text.}
\psentry{ps_134-7_043-v0007}{7}{writing over}{$\over{\unsure{5}}{6}$}{$\hand{6}$}{The superposed handwritten reading replaces the underlying reading.}
\psentry{ps_134-7_043-v0008}{8}{typewriter deletion}{\xout{categoriexx}}{}{The typewriter strikes out this material; it leaves the corrected running text.}
\psapparatusend
\fi
\pssource{ps_134-7_044}{134-7/44}{https://grothendieck.umontpellier.fr/web/viewer.html?file=../134-7.pdf\#page=45}
\psnumber{ps_134-7_044-v0001}{1}
\psnumber{ps_134-7_044-v0002}{2}
\psnumber{ps_134-7_044-v0003}{3}
\psnumber{ps_134-7_044-v0004}{4}
\psnumber{ps_134-7_044-v0005}{5}
\psnumber{ps_134-7_044-v0006}{6}
\psnumber{ps_134-7_044-v0007}{7}
\psnumber{ps_134-7_044-v0008}{8}
\psnumber{ps_134-7_044-v0009}{9}
\psnumber{ps_134-7_044-v0010}{10}
%% BEGIN SOURCE ps_134-7_044
% Pursuing Stacks (cotes 134-2 to 134-8): archive cote 134-7, site page 44 (PDF page 45), typescript page 486.
% typescript. Draft compared: ps6.tex:2160-2237.
% AI-transcribed by Claude Opus 5.5 (claude-opus-5-5, adaptive thinking, effort high), Message Batches API, batch ps3 (msgbatch_01XdFmuVVBAcgc7EFxdtCqsj), 2026-11-12; reviewed against both supplied scan halves; see CHECK_DECISIONS.tsv.
\scanpage{134-7/44}
\mspage{486}

the pro-unipotent bundle defined in terms of L(X) via the canonical
extension
\[ \mathrm{Pro}(W) \text{ or simply } W : \quad \mathrm{Pro}(\mathrm{Ab}_k) \hand{\longrightarrow} \mathrm{Pro}(U(k)) \]
of W \var{\xout{as}}{}{typewriter deletion}{ps_134-7_044-v0001} (cf. (8)) to proobjects. Thus, instead of the map (6) which
does'nt quite exist, we get a canonical ``exponential'' map
\[ (18) \qquad X \hand{\longrightarrow} WL(X) \quad \var{\xout{\unsure{*}}}{}{typewriter deletion}{ps_134-7_044-v0002} \]
in Pro U(k) , which has of course little chance to factor through
\[ (16') \qquad WL_{pt}(X) \hand{\longrightarrow} WL(X) \]
deduced from (16) by applying W. It is via this map (18) that we may
declare that L(X) prorepresents the functor (7') - it may be viewed
as the universal map of the type
\[ X \hand{\longrightarrow} W(N) \quad , \]
where now N is (not just a k-module\var{\xout{s}}{}{typewriter deletion}{ps_134-7_044-v0003}, but) a variable object in Pr$\bullet$ Ab$_k$.
Whereas the proobject L(X) is of a ``k-linear'' nature and may be viewed
as the (\uline{quasi-coherent}) \uline{k-linearization} \uline{of} \uline{the unipotent bundle} X,
the proobject WL(X) \var{\xout{has lost}}{}{typewriter deletion}{ps_134-7_044-v0004} of U(k) has lost \var{\over{the}{\unsure{it's|its}}}{\hand{\unsure{it's|its}}}{writing over}{ps_134-7_044-v0005} k-linear nature,
we would rather view it as the canonical ``abelianization'' of X , retai-%
ning mainly it's additive structure (plus maybe operation of k on it,
which is a lot weaker though than a structure of an $\underline{O}_k$-module ...).

I would like now to examine if the \var{\inserted[typed]{\hand{(}quasi-coherent\hand{)}}}{\hand{(}quasi-coherent\hand{)}}{typed insertion}{ps_134-7_044-v0006} proobject L(X), which has
been obtained as the suitable quasi-coherent substitute \var{\xout{a}}{}{typewriter deletion}{ps_134-7_044-v0007}for $\underline{O}_k^{(X)}$
in order to get the basic isomorphism (4) for quasicoherent F, may
serve the same puppose for RHom$_{\underline{O}_k}$ , in \var{\deleted{relation}}{}{deletion}{ps_134-7_044-v0008}\var{\inserted{analogy}}{\hand{analogy}}{insertion}{ps_134-7_044-v0009} to (5),(5'). Quite
generally, if
\[ \hand{\Gamma}=(\hand{\Gamma}_{\hand{\alpha}}) \]
is any pro-$\underline{O}_k$-module, let's define for any Module F
\[ (19) \qquad \mathrm{RHom}_{\underline{O}_k}(\,\hand{\Gamma}, F) \overset{\mathrm{def}}{=} \mathrm{Hom}_{\underline{O}_k}(\,\hand{\Gamma} , C^{\hand{\bullet}}(F)) \quad , \]
where C$^{\hand{\bullet}}$(F) is an injective resolution of F - thus the definition
extends the usual one \var{\xout{for}}{}{typewriter deletion}{ps_134-7_044-v0010} when $\hand{\Gamma}$ is just an $\underline{O}_k$-Module . Our expectati-%
on now would be
\[ (20) \qquad \mathrm{RHom}_{\underline{O}_k}(L(X), F\,) \hand{\xrightarrow{\sim}} R\hand{\Gamma}_X(F_X) \quad , \]
giving rise to
\[ (20') \qquad \mathrm{Ext}^{\hand{n}}_{\underline{O}_k}(L(X),F) \hand{\xrightarrow{\sim}} H^{\hand{n}}(X,F_X) \quad , \]

\margin{44}

% Visual collation: both supplied image halves inspected; see CHECK_DECISIONS_134_7_8.tsv.

%% END SOURCE ps_134-7_044
\ifnum\EditionMode=1
\psapparatusstart
\psentry{ps_134-7_044-v0001}{1}{typewriter deletion}{\xout{as}}{}{The typewriter strikes out this material; it leaves the corrected running text.}
\psentry{ps_134-7_044-v0002}{2}{typewriter deletion}{$\xout{\unsure{*}}$}{}{The typewriter strikes out this material; it leaves the corrected running text.}
\psentry{ps_134-7_044-v0003}{3}{typewriter deletion}{\xout{s}}{}{The typewriter strikes out this material; it leaves the corrected running text.}
\psentry{ps_134-7_044-v0004}{4}{typewriter deletion}{\xout{has lost}}{}{The typewriter strikes out this material; it leaves the corrected running text.}
\psentry{ps_134-7_044-v0005}{5}{writing over}{\over{the}{\unsure{it's|its}}}{\hand{\unsure{it's|its}}}{The superposed handwritten reading replaces the underlying reading.}
\psentry{ps_134-7_044-v0006}{6}{typed insertion}{\inserted[typed]{\hand{(}quasi-coherent\hand{)}}}{\hand{(}quasi-coherent\hand{)}}{The typed insertion is incorporated at its indicated position; insertion brackets leave the running text.}
\psentry{ps_134-7_044-v0007}{7}{typewriter deletion}{\xout{a}}{}{The typewriter strikes out this material; it leaves the corrected running text.}
\psentry{ps_134-7_044-v0008}{8}{deletion}{\deleted{relation}}{}{The author cancels this material by hand; it leaves the corrected running text.}
\psentry{ps_134-7_044-v0009}{9}{insertion}{\inserted{analogy}}{\hand{analogy}}{The handwritten insertion is incorporated at its indicated position; insertion brackets leave the running text.}
\psentry{ps_134-7_044-v0010}{10}{typewriter deletion}{\xout{for}}{}{The typewriter strikes out this material; it leaves the corrected running text.}
\psapparatusend
\fi
\pssource{ps_134-7_045}{134-7/45}{https://grothendieck.umontpellier.fr/web/viewer.html?file=../134-7.pdf\#page=46}
\psnumber{ps_134-7_045-v0001}{1}
\psnumber{ps_134-7_045-v0002}{2}
\psnumber{ps_134-7_045-v0003}{3}
\psnumber{ps_134-7_045-v0004}{4}
\psnumber{ps_134-7_045-v0005}{5}
\psnumber{ps_134-7_045-v0006}{6}
\psnumber{ps_134-7_045-v0007}{7}
\psnumber{ps_134-7_045-v0008}{8}
\psnumber{ps_134-7_045-v0009}{9}
\psnumber{ps_134-7_045-v0010}{10}
\psnumber{ps_134-7_045-v0011}{11}
\psnumber{ps_134-7_045-v0012}{12}
\psnumber{ps_134-7_045-v0013}{13}
\psnumber{ps_134-7_045-v0014}{14}
\psnumber{ps_134-7_045-v0015}{15}
\psnumber{ps_134-7_045-v0016}{16}
\psnumber{ps_134-7_045-v0017}{17}
\psnumber{ps_134-7_045-v0018}{18}
%% BEGIN SOURCE ps_134-7_045
% Pursuing Stacks (cotes 134-2 to 134-8): archive cote 134-7, site page 45 (PDF page 46), typescript page 487.
% typescript. Draft compared: ps6.tex:2215-2280.
% AI-transcribed by Claude Opus 5.5 (claude-opus-5-5, adaptive thinking, effort high), Message Batches API, batch ps3 (msgbatch_01XdFmuVVBAcgc7EFxdtCqsj), 2026-11-12; reviewed against both supplied scan halves; see CHECK_DECISIONS.tsv.
\scanpage{134-7/45}
\mspage{487}

for any unipotent bundle X over k,  and any quasicoherent sheaf
\[ F = W(N) \quad , \]
where N is any k-module.  Of course, (19)  yields (for general $\hand{\Gamma}$ )
\[ (19') \qquad \mathrm{Ext}^{i}_{\underline{O}_{\hand{k}}}( \hand{\Gamma} , F) \hand{\simeq} \underset{\hand{\alpha}}{\underset{\hand{\longrightarrow}}{\mathrm{lim}}} \mathrm{Ext}^{i}_{\underline{O}_k}( \hand{\Gamma}_{\hand{\alpha}} , F) \quad , \]
so that (20')  may be rewritten more explicitly,  if  $X \hand{\simeq} W(M)$, \hand{as}
\[ (21) \qquad H^{n}(X,W(N)_X) \hand{\simeq} \var{\xout{\mathrm{lim}}}{}{typewriter deletion}{ps_134-7_045-v0001}_{\var{\deleted{\hand{i}}}{}{deletion}{ps_134-7_045-v0002}} \hand{\bigoplus_{i\geq 0}} \mathrm{Ext}^{n}_{k}(\hand{\Gamma}_k^{\ i}(M),N) \quad . \]
At any rate,  we have a canonical map (20) in $D^{+}(Ab_k)$,  hence maps
(20'),  in view of the isomorphism (5)  and the canonical map
\[ (\hand{*}) \qquad \underline{O}_k^{(X)} \hand{\longrightarrow} L(X) \]
deduced from (18),  and the question now is whether th\var{\over{is}{ese}}{\hand{ese}}{writing over}{ps_134-7_045-v0003} \var{\over{is}{are}}{\hand{are}}{writing over}{ps_134-7_045-v0004} isomor-%
phism\hand{s}. \var{\deleted{The question \inserted[typed]{(20)} makes a sense for any X,  we are more specifically
interested though in the case when X is a unipotent bundle,  to which
\xout{We may of course assume X = W(M)}
case I will restrict.}}{}{deletion}{ps_134-7_045-v0005} We may of course assume  in (21)
\[ X = W(M) \quad , \quad n \hand{\geq} 1 \ . \]
If M is projective,  so are the modules $\hand{\Gamma}_k^{\ i}(M)$,  and hence \var{\xout{\unsure{for n>0}}}{}{typewriter deletion}{ps_134-7_045-v0006}
the second  member in (21) is zero,  so we should check the first mem-%
ber is too.  This \var{\over[typed]{\unsure{a}}{i}}{i}{typewriter overwriting}{ps_134-7_045-v0007}s clear when M is of finite type,  hence X represen-%
table by an affine scheme,  who\var{\over[typed]{\unsure{e}}{s}}{s}{typewriter overwriting}{ps_134-7_045-v0008}e quasi\var{\over[typed]{\unsure{o}}{c}}{c}{typewriter overwriting}{ps_134-7_045-v0009}oherent cohomology is well known
therefore to vanish in dim. $n \hand{>} 0$.  The  general case should be a con-%
sequence of this,  representing X as the filtering direct limit of
\var{\over[typed]{p}{i}}{i}{typewriter overwriting}{ps_134-7_045-v0010}t's submodules which are projective of finite type  -  this should work
at any rate when M is free with a basis which is at most countable,
using the standard so-called ``Mittag\var{\deleted{\unsure{e}}}{}{deletion}{ps_134-7_045-v0011}-Leffler'' argument for passage to
limit.   Thus in case M projective,  (21)  and hence (20) seems OK
indeed.  When M is \uline{not} projective however,  \var{\over[typed]{w}{t}}{t}{typewriter overwriting}{ps_134-7_045-v0012}here must be some k-module
N such that  $\mathrm{Ext}^{1}_{k}(M,N) \hand{\neq} 0$ ,  and hence the second member of (19) is
non zero for n=1,  which should imply rather unexpectedly
\[ H^{1}(X,F) = H^{1}(W(M), W(N)_{W(M)}) \hand{\neq} 0 \quad , \]
whereas till this very moment I \var{\over{has}{had}}{\hand{had}}{writing over}{ps_134-7_045-v0013} \var{\inserted{been}}{\hand{been}}{insertion}{ps_134-7_045-v0014} under the impression that quasi-cohe%
rent cohomology of quasiunipotent bundles should be zero,  just as
for affine schemes !  Maybe it has been familiar to Larry Breen for a
long time that this is \uline{not} so\var{\over{!}{?}}{\hand{?}}{writing over}{ps_134-7_045-v0015}  Maybe \var{\inserted{also}}{\hand{also}}{insertion}{ps_134-7_045-v0016} for what we want to do it is'nt
really basic to find out whether (21) is true in full generality,  as
for the purpose of studying P\var{\over{\unsure{a}}{o}}{\hand{o}}{writing over}{ps_134-7_045-v0017}stnikov type dévissage,  the \var{\deleted{coefficients}}{}{deletion}{ps_134-7_045-v0018}

\margin{45}

% Visual collation: both supplied image halves inspected; see CHECK_DECISIONS_134_7_8.tsv.

%% END SOURCE ps_134-7_045
\ifnum\EditionMode=1
\psapparatusstart
\psentry{ps_134-7_045-v0001}{1}{typewriter deletion}{$\xout{\mathrm{lim}}$}{}{The typewriter strikes out this material; it leaves the corrected running text.}
\psentry{ps_134-7_045-v0002}{2}{deletion}{$\deleted{\hand{i}}$}{}{The author cancels this material by hand; it leaves the corrected running text.}
\psentry{ps_134-7_045-v0003}{3}{writing over}{\over{is}{ese}}{\hand{ese}}{The superposed handwritten reading replaces the underlying reading.}
\psentry{ps_134-7_045-v0004}{4}{writing over}{\over{is}{are}}{\hand{are}}{The superposed handwritten reading replaces the underlying reading.}
\psentry{ps_134-7_045-v0005}{5}{deletion}{\deleted{The question \inserted[typed]{(20)} makes a sense for any X,  we are more specifically
interested though in the case when X is a unipotent bundle,  to which
\xout{We may of course assume X = W(M)}
case I will restrict.}}{}{The author cancels this material by hand; it leaves the corrected running text.}
\psentry{ps_134-7_045-v0006}{6}{typewriter deletion}{\xout{\unsure{for n>0}}}{}{The typewriter strikes out this material; it leaves the corrected running text.}
\psentry{ps_134-7_045-v0007}{7}{typewriter overwriting}{\over[typed]{\unsure{a}}{i}}{i}{The superposed typed reading replaces the underlying reading.}
\psentry{ps_134-7_045-v0008}{8}{typewriter overwriting}{\over[typed]{\unsure{e}}{s}}{s}{The superposed typed reading replaces the underlying reading.}
\psentry{ps_134-7_045-v0009}{9}{typewriter overwriting}{\over[typed]{\unsure{o}}{c}}{c}{The superposed typed reading replaces the underlying reading.}
\psentry{ps_134-7_045-v0010}{10}{typewriter overwriting}{\over[typed]{p}{i}}{i}{The superposed typed reading replaces the underlying reading.}
\psentry{ps_134-7_045-v0011}{11}{deletion}{\deleted{\unsure{e}}}{}{The author cancels this material by hand; it leaves the corrected running text.}
\psentry{ps_134-7_045-v0012}{12}{typewriter overwriting}{\over[typed]{w}{t}}{t}{The superposed typed reading replaces the underlying reading.}
\psentry{ps_134-7_045-v0013}{13}{writing over}{\over{has}{had}}{\hand{had}}{The superposed handwritten reading replaces the underlying reading.}
\psentry{ps_134-7_045-v0014}{14}{insertion}{\inserted{been}}{\hand{been}}{The handwritten insertion is incorporated at its indicated position; insertion brackets leave the running text.}
\psentry{ps_134-7_045-v0015}{15}{writing over}{\over{!}{?}}{\hand{?}}{The superposed handwritten reading replaces the underlying reading.}
\psentry{ps_134-7_045-v0016}{16}{insertion}{\inserted{also}}{\hand{also}}{The handwritten insertion is incorporated at its indicated position; insertion brackets leave the running text.}
\psentry{ps_134-7_045-v0017}{17}{writing over}{\over{\unsure{a}}{o}}{\hand{o}}{The superposed handwritten reading replaces the underlying reading.}
\psentry{ps_134-7_045-v0018}{18}{deletion}{\deleted{coefficients}}{}{The author cancels this material by hand; it leaves the corrected running text.}
\psapparatusend
\fi
\pssource{ps_134-7_046}{134-7/46}{https://grothendieck.umontpellier.fr/web/viewer.html?file=../134-7.pdf\#page=47}
\psnumber{ps_134-7_046-v0001}{1}
\psnumber{ps_134-7_046-v0002}{2}
\psnumber{ps_134-7_046-v0003}{3}
\psnumber{ps_134-7_046-v0004}{4}
\psnumber{ps_134-7_046-v0005}{5}
\psnumber{ps_134-7_046-v0006}{6}
\psnumber{ps_134-7_046-v0007}{7}
\psnumber{ps_134-7_046-v0008}{8}
\psnumber{ps_134-7_046-v0009}{9}
\psnumber{ps_134-7_046-v0010}{10}
\psnumber{ps_134-7_046-v0011}{11}
\psnumber{ps_134-7_046-v0012}{12}
\psnumber{ps_134-7_046-v0013}{13}
\psnumber{ps_134-7_046-v0014}{14}
\psnumber{ps_134-7_046-v0015}{15}
\psnumber{ps_134-7_046-v0016}{16}
\psnumber{ps_134-7_046-v0017}{17}
%% BEGIN SOURCE ps_134-7_046
% Pursuing Stacks (cotes 134-2 to 134-8): archive cote 134-7, site page 46 (PDF page 47), typescript page 488.
% typescript. Draft compared: ps6.tex:2268-2332.
% AI-transcribed by Claude Opus 5.5 (claude-opus-5-5, adaptive thinking, effort high), Message Batches API, batch ps3 (msgbatch_01XdFmuVVBAcgc7EFxdtCqsj), 2026-11-12; reviewed against both supplied scan halves; see CHECK_DECISIONS.tsv.
\scanpage{134-7/46}
\mspage{488}

unipotent bundles $X$ we are going to work with will be smooth i.e. $M$ pro-%
jective (and we may even get away with free $M$'s, if we need so). The
natural idea here for getting (20) via (21) in full generality, is to
use a projective resolution of $M$ (even a free one), but I'll not try to
work this out now. The main impression which remains is that for the
more relevant cases (involving cohomology \var{\xout{\unsure{inva}}}{}{typewriter deletion}{ps_134-7_046-v0001} groups of a \uline{smooth}
unipotent bundle\var{\inserted[typed]{$X$}}{$X$}{typed insertion}{ps_134-7_046-v0002} at any rate, with quasicoherent coefficients), the
quasicoherent "pro"-linearization $L(X)$ is just as good for computing
cohomology invariants, as the forbidding $\underline{O}_k^{(X)}$ modules we were shrinking
from.

If now we take an $X_*$ instead of just $X$, namely a ss complex in
$U(k)$, and assuming the components $X_n$ to be smooth (to be safe), the
isomorphisms (20') should give rise to isomorphisms (for $F=W(N)$)
\[ (22) \qquad H^n(X_*, F ) \hand{\simeq} \operatorname{Ext}^n_{\underline{O}_k}(L(X_*), F) \hand{\simeq} \operatorname{Ext}^n_k(L(X_*),N) , \]
where the $\operatorname{Ext}^n$ should be viewed as Hyperext functors (not term by
term), and where \var{\inserted{in the last member}}{\hand{in the last member}}{insertion}{ps_134-7_046-v0003} $L(X_*)$ may be interpreted as a chain complex in
Pro $(\mathrm{Ab}_k)$. The first member of (22) is the kind of group occuring as
obstruction group in the Postnikov-type devissage of $X$\hand{$_*$}into linear
structures $W(M(i)_*)$. The chain-pro complex $L(X_*)$ may still look a
\margin{x} little forbidding, ou\var{\over{t}{r}}{\hand{r}}{writing over}{ps_134-7_046-v0004} hope though now is that in the "pointed" case
we are really interested in, we may replace $L(X_*)$ by $L$\var{\deleted{\illegible{1}}}{}{deletion}{ps_134-7_046-v0005}$_{pt}(X_*)$, which
is just a true honest chain complex in $\mathrm{Ab}_k$ . Now from (16) we get indeed
a canonical map
\[ (23) \qquad \operatorname{Ext}^n_k(L(X_*), N ) \hand{\longrightarrow} \operatorname{Ext}^n_k(L_{pt}(X_*), N ) \]
and we hope that this is an isomorphism, under suitable assumptions on
$X_*$ , the most basic one I can think of now being
\[ X_o = e . \]
The map (23) w\var{\over{\unsure{i}}{a}}{\hand{a}}{writing over}{ps_134-7_046-v0006}s \var{\inserted{defined as the}}{\hand{defined as the}}{insertion}{ps_134-7_046-v0007} transposed of a map of chain complexes in \var{\deleted{$\mathrm{Ab}_k$}}{}{deletion}{ps_134-7_046-v0008} \hand{Pro $\mathrm{Ab}_k$}
\[ (24) \qquad L_{pt}(X_*) \hand{\longrightarrow} L(X_*) \quad , \]
\var{\xout{by applying}}{}{typewriter deletion}{ps_134-7_046-v0009} deduced from (16) by applying it componentwise. We recog-%
nize here,\var{\deleted{\illegible{1}}}{}{deletion}{ps_134-7_046-v0010} but with a different interpretation (which seems to me
"the correct" one), the second map in the often referred\hand{-}to diagram (w)
of page 454 , or (V\hand{)} in yesterday's reflections (p.478). To say that
it gives rise to isomorphisms (23), for any $n$ and any \var{\over[typed]{m}{M}}{M}{typewriter overwriting}{ps_134-7_046-v0011}odule $N$, should
be equivalent to saying that (24) is a quasiisomorphism\var{\deleted{s}}{}{deletion}{ps_134-7_046-v0012} - but to
make sure \var{\inserted{I \deleted{expected}}}{\hand{I }}{insertion}{ps_134-7_046-v0013} \var{\deleted{entails}}{}{deletion}{ps_134-7_046-v0014}, \var{\inserted[typed]{should demand \deleted{\unsure{a}}}}{should demand }{typed insertion}{ps_134-7_046-v0015} \var{\xout{and prove will}}{}{typewriter deletion}{ps_134-7_046-v0016} \var{\deleted{amount}}{}{deletion}{ps_134-7_046-v0017} a little work on foundational
matters on pro-complexes I guess\hand{;} \hand{Also, to see}

\margin{46}

% Visual collation: both supplied image halves inspected; see CHECK_DECISIONS_134_7_8.tsv.

%% END SOURCE ps_134-7_046
\ifnum\EditionMode=1
\psapparatusstart
\psentry{ps_134-7_046-v0001}{1}{typewriter deletion}{\xout{\unsure{inva}}}{}{The typewriter strikes out this material; it leaves the corrected running text.}
\psentry{ps_134-7_046-v0002}{2}{typed insertion}{\inserted[typed]{$X$}}{$X$}{The typed insertion is incorporated at its indicated position; insertion brackets leave the running text.}
\psentry{ps_134-7_046-v0003}{3}{insertion}{\inserted{in the last member}}{\hand{in the last member}}{The handwritten insertion is incorporated at its indicated position; insertion brackets leave the running text.}
\psentry{ps_134-7_046-v0004}{4}{writing over}{\over{t}{r}}{\hand{r}}{The superposed handwritten reading replaces the underlying reading.}
\psentry{ps_134-7_046-v0005}{5}{deletion}{\deleted{\illegible{1}}}{}{The author cancels this material by hand; it leaves the corrected running text.}
\psentry{ps_134-7_046-v0006}{6}{writing over}{\over{\unsure{i}}{a}}{\hand{a}}{The superposed handwritten reading replaces the underlying reading.}
\psentry{ps_134-7_046-v0007}{7}{insertion}{\inserted{defined as the}}{\hand{defined as the}}{The handwritten insertion is incorporated at its indicated position; insertion brackets leave the running text.}
\psentry{ps_134-7_046-v0008}{8}{deletion}{\deleted{$\mathrm{Ab}_k$}}{}{The author cancels this material by hand; it leaves the corrected running text.}
\psentry{ps_134-7_046-v0009}{9}{typewriter deletion}{\xout{by applying}}{}{The typewriter strikes out this material; it leaves the corrected running text.}
\psentry{ps_134-7_046-v0010}{10}{deletion}{\deleted{\illegible{1}}}{}{The author cancels this material by hand; it leaves the corrected running text.}
\psentry{ps_134-7_046-v0011}{11}{typewriter overwriting}{\over[typed]{m}{M}}{M}{The superposed typed reading replaces the underlying reading.}
\psentry{ps_134-7_046-v0012}{12}{deletion}{\deleted{s}}{}{The author cancels this material by hand; it leaves the corrected running text.}
\psentry{ps_134-7_046-v0013}{13}{insertion}{\inserted{I \deleted{expected}}}{\hand{I }}{The handwritten insertion is incorporated at its indicated position; insertion brackets leave the running text.}
\psentry{ps_134-7_046-v0014}{14}{deletion}{\deleted{entails}}{}{The author cancels this material by hand; it leaves the corrected running text.}
\psentry{ps_134-7_046-v0015}{15}{typed insertion}{\inserted[typed]{should demand \deleted{\unsure{a}}}}{should demand }{The typed insertion is incorporated at its indicated position; insertion brackets leave the running text.}
\psentry{ps_134-7_046-v0016}{16}{typewriter deletion}{\xout{and prove will}}{}{The typewriter strikes out this material; it leaves the corrected running text.}
\psentry{ps_134-7_046-v0017}{17}{deletion}{\deleted{amount}}{}{The author cancels this material by hand; it leaves the corrected running text.}
\psapparatusend
\fi
\pssource{ps_134-7_047}{134-7/47}{https://grothendieck.umontpellier.fr/web/viewer.html?file=../134-7.pdf\#page=48}
\psnumber{ps_134-7_047-v0001}{1}
\psnumber{ps_134-7_047-v0002}{2}
\psnumber{ps_134-7_047-v0003}{3}
\psnumber{ps_134-7_047-v0004}{4}
\psnumber{ps_134-7_047-v0005}{5}
\psnumber{ps_134-7_047-v0006}{6}
\psnumber{ps_134-7_047-v0007}{7}
\psnumber{ps_134-7_047-v0008}{8}
\psnumber{ps_134-7_047-v0009}{9}
\psnumber{ps_134-7_047-v0010}{10}
\psnumber{ps_134-7_047-v0011}{11}
\psnumber{ps_134-7_047-v0012}{12}
\psnumber{ps_134-7_047-v0013}{13}
\psnumber{ps_134-7_047-v0014}{14}
\psnumber{ps_134-7_047-v0015}{15}
\psnumber{ps_134-7_047-v0016}{16}
\psnumber{ps_134-7_047-v0017}{17}
\psnumber{ps_134-7_047-v0018}{18}
\psnumber{ps_134-7_047-v0019}{19}
\psnumber{ps_134-7_047-v0020}{20}
\psnumber{ps_134-7_047-v0021}{21}
%% BEGIN SOURCE ps_134-7_047
% Pursuing Stacks (cotes 134-2 to 134-8): archive cote 134-7, site page 47 (PDF page 48), typescript page 489.
% typescript. Draft compared: ps6.tex:2320-2387.
% AI-transcribed by Claude Opus 5.5 (claude-opus-5-5, adaptive thinking, effort high), Message Batches API, batch ps3 (msgbatch_01XdFmuVVBAcgc7EFxdtCqsj), 2026-11-12; reviewed against both supplied scan halves; see CHECK_DECISIONS.tsv.
\scanpage{134-7/47}
\mspage{\var{\xout{(25)}}{}{typewriter deletion}{ps_134-7_047-v0001} 489}

if the \var{\xout{fact}}{}{typewriter deletion}{ps_134-7_047-v0002}\var{\inserted[typed]{assumption}}{assumption}{typed insertion}{ps_134-7_047-v0003} that (24) is a quasi-isomorphism should imply the same state-%
ment with $L(X_*)$ replaced by it's componentwise projective limit - name-%
ly \hand{that} (V) on p. 478 is, which we'll need of course in case $k=\underline{Z}$ for the
so-called ``linearization theorem''. \var{\xout{These}}{}{typewriter deletion}{ps_134-7_047-v0004} Thus we get three \var{\xout{related}}{}{typewriter deletion}{ps_134-7_047-v0005}
isomorphism or quasi-isomorphism statement\hand{s,} concerning (23),(24) and
(V) in yesterday's notes, which are at any rate closely related, and
which \var{\xout{I feel}}{}{typewriter deletion}{ps_134-7_047-v0006} one hopes to be true, be\hand{c}ause this\var{\inserted{seems}}{\hand{seems}}{insertion}{ps_134-7_047-v0007} \var{\deleted{is}}{}{deletion}{ps_134-7_047-v0008} needed for a schema-%
tization theory of homotopy types to work. But I should confess I have
not tried even to get any clue as to why this should be true, under the
only assumptions say that the components of $X_*$ should be smooth (and
possibly the Kan assumption ?), plus $X_0=e$ say.

Now to the second ingredient of the looked-for ``linearization theo-%
rem'', which previously was the first map in (w) p.454, or (V') on page
478, involving maps of \var{\xout{either Z(Xx(k))}}{}{typewriter deletion}{ps_134-7_047-v0009}
\[ \text{either } \underline{Z}^{(X_{\hand{*}}(k))} \quad \text{or} \quad k^{(X_{\hand{*}}(k))} \]
into what was previously called $L(X_*)$, and which we \var{\xout{now}}{}{typewriter deletion}{ps_134-7_047-v0010} would now
rather denote\var{\deleted{s}}{}{deletion}{ps_134-7_047-v0011} by
\[ \underset{\hand{\longleftarrow}}{\lim}\, L(X_*) \quad ( \hand{\simeq} \hand{\Gamma}_k^{\hand{\wedge}}\,(M_{\hand{*}}) \ \text{if} \ X_{\hand{*}} = W(M_{\hand{*}}) ) \ . \]
We made sure that, unless $k=\underline{Z}$ , none of the two \var{\over{\unsure{b}}{h}}{\hand{h}}{writing over}{ps_134-7_047-v0012}ad any chance to be
a quasi-isomorphism. The only positive thing that came out in this
direction was that the first one of these maps would induce an isomor-%
phism on $H_i$ in the critical dimension (namely $i=n+1$ in the n-connected
case). We now understand why, \var{\xout{ex}}{}{typewriter deletion}{ps_134-7_047-v0013} for $k \hand{\neq} \underline{Z}$ , we would not get any
actual quasiisomorphism - namely the ``correct'' naive linearization
which compares reasonably with $L(X_*)$ should not be relative to a
\uline{constant} ring such as $\underline{Z}$ or $k$ , but relative to $\underline{O}_k$ , via the map
($*$) (p. 487) giving rise now to
\margin{\textcircled{$(X_*)$}}
\[ (25) \qquad \underline{O}_k^{(X_{\hand{*}})} \hand{\longrightarrow} L(X_*) \quad . \]
The more reasonable question now, making a good sense really for any
ground ring $k$, is whether this map (under the usual assumptions\var{\inserted{say}}{\hand{say}}{insertion}{ps_134-7_047-v0014} on $X_*$)
is a quasi-isomorphism. I would'nt really bet it is, as I have some
doubts as to whether the \var{\xout{co}}{}{typewriter deletion}{ps_134-7_047-v0015}homology sheaves of the first member (both
\var{\inserted{members}}{\hand{members}}{insertion}{ps_134-7_047-v0016} of course being viewed as chain complexes of $\underline{O}_k$-Modules\var{\inserted[typed]{or ``pro'' such}}{or ``pro'' such}{typed insertion}{ps_134-7_047-v0017}) are quasi-%
coherent - but for the time being I am not sure either if those of
the second member are essentially constant proobjects ! But even if
(25) is'nt a quasiisomorphism, it does behave like one for all practi-%
cal purposes of computing\var{\inserted[typed]{quasi-coherent}}{quasi-coherent}{typed insertion}{ps_134-7_047-v0018} cohomology\hand{,} it would seem, as \var{\xout{leads to}}{}{typewriter deletion}{ps_134-7_047-v0019}\var{\inserted[typed]{this boils}}{this boils}{typed insertion}{ps_134-7_047-v0020}
\var{\inserted{down indeed to the isomorphisms (20) or (20').}}{\hand{down indeed to the isomorphisms (20) or (20').}}{insertion}{ps_134-7_047-v0021}

\margin{47}

% Visual collation: both supplied image halves inspected; see CHECK_DECISIONS_134_7_8.tsv.

%% END SOURCE ps_134-7_047
\ifnum\EditionMode=1
\psapparatusstart
\psentry{ps_134-7_047-v0001}{1}{typewriter deletion}{\xout{(25)}}{}{The typewriter strikes out this material; it leaves the corrected running text.}
\psentry{ps_134-7_047-v0002}{2}{typewriter deletion}{\xout{fact}}{}{The typewriter strikes out this material; it leaves the corrected running text.}
\psentry{ps_134-7_047-v0003}{3}{typed insertion}{\inserted[typed]{assumption}}{assumption}{The typed insertion is incorporated at its indicated position; insertion brackets leave the running text.}
\psentry{ps_134-7_047-v0004}{4}{typewriter deletion}{\xout{These}}{}{The typewriter strikes out this material; it leaves the corrected running text.}
\psentry{ps_134-7_047-v0005}{5}{typewriter deletion}{\xout{related}}{}{The typewriter strikes out this material; it leaves the corrected running text.}
\psentry{ps_134-7_047-v0006}{6}{typewriter deletion}{\xout{I feel}}{}{The typewriter strikes out this material; it leaves the corrected running text.}
\psentry{ps_134-7_047-v0007}{7}{insertion}{\inserted{seems}}{\hand{seems}}{The handwritten insertion is incorporated at its indicated position; insertion brackets leave the running text.}
\psentry{ps_134-7_047-v0008}{8}{deletion}{\deleted{is}}{}{The author cancels this material by hand; it leaves the corrected running text.}
\psentry{ps_134-7_047-v0009}{9}{typewriter deletion}{\xout{either Z(Xx(k))}}{}{The typewriter strikes out this material; it leaves the corrected running text.}
\psentry{ps_134-7_047-v0010}{10}{typewriter deletion}{\xout{now}}{}{The typewriter strikes out this material; it leaves the corrected running text.}
\psentry{ps_134-7_047-v0011}{11}{deletion}{\deleted{s}}{}{The author cancels this material by hand; it leaves the corrected running text.}
\psentry{ps_134-7_047-v0012}{12}{writing over}{\over{\unsure{b}}{h}}{\hand{h}}{The superposed handwritten reading replaces the underlying reading.}
\psentry{ps_134-7_047-v0013}{13}{typewriter deletion}{\xout{ex}}{}{The typewriter strikes out this material; it leaves the corrected running text.}
\psentry{ps_134-7_047-v0014}{14}{insertion}{\inserted{say}}{\hand{say}}{The handwritten insertion is incorporated at its indicated position; insertion brackets leave the running text.}
\psentry{ps_134-7_047-v0015}{15}{typewriter deletion}{\xout{co}}{}{The typewriter strikes out this material; it leaves the corrected running text.}
\psentry{ps_134-7_047-v0016}{16}{insertion}{\inserted{members}}{\hand{members}}{The handwritten insertion is incorporated at its indicated position; insertion brackets leave the running text.}
\psentry{ps_134-7_047-v0017}{17}{typed insertion}{\inserted[typed]{or ``pro'' such}}{or ``pro'' such}{The typed insertion is incorporated at its indicated position; insertion brackets leave the running text.}
\psentry{ps_134-7_047-v0018}{18}{typed insertion}{\inserted[typed]{quasi-coherent}}{quasi-coherent}{The typed insertion is incorporated at its indicated position; insertion brackets leave the running text.}
\psentry{ps_134-7_047-v0019}{19}{typewriter deletion}{\xout{leads to}}{}{The typewriter strikes out this material; it leaves the corrected running text.}
\psentry{ps_134-7_047-v0020}{20}{typed insertion}{\inserted[typed]{this boils}}{this boils}{The typed insertion is incorporated at its indicated position; insertion brackets leave the running text.}
\psentry{ps_134-7_047-v0021}{21}{insertion}{\inserted{down indeed to the isomorphisms (20) or (20').}}{\hand{down indeed to the isomorphisms (20) or (20').}}{The handwritten insertion is incorporated at its indicated position; insertion brackets leave the running text.}
\psapparatusend
\fi
\pssource{ps_134-7_048}{134-7/48}{https://grothendieck.umontpellier.fr/web/viewer.html?file=../134-7.pdf\#page=49}
\psnumber{ps_134-7_048-v0001}{1}
\psnumber{ps_134-7_048-v0002}{2}
\psnumber{ps_134-7_048-v0003}{3}
\psnumber{ps_134-7_048-v0004}{4}
\psnumber{ps_134-7_048-v0005}{5}
\psnumber{ps_134-7_048-v0006}{6}
\psnumber{ps_134-7_048-v0007}{7}
\psnumber{ps_134-7_048-v0008}{8}
\psnumber{ps_134-7_048-v0009}{9}
\psnumber{ps_134-7_048-v0010}{10}
\psnumber{ps_134-7_048-v0011}{11}
%% BEGIN SOURCE ps_134-7_048
% Pursuing Stacks (cotes 134-2 to 134-8): archive cote 134-7, site page 48 (PDF page 49), typescript page 490.
% typescript. Draft compared: ps6.tex:2369-2431.
% AI-transcribed by Claude Opus 5.5 (claude-opus-5-5, adaptive thinking, effort high), Message Batches API, batch ps3 (msgbatch_01XdFmuVVBAcgc7EFxdtCqsj), 2026-11-12; reviewed against both supplied scan halves; see CHECK_DECISIONS.tsv.
\scanpage{134-7/48}
\mspage{490}

7.9.

116.   For ten days I hav'nt written any notes,  and the time when I
stopped looks a lot more remote still.  For two days still after I last
wrote on the notes,  I kept pondering about schematization of homotopy
types  -  it were rather lively days\var{\over{.}{,}}{\hand{,}}{writing over}{ps_134-7_048-v0001}  \var{\over{F}{f}}{\hand{f}}{writing over}{ps_134-7_048-v0002}irst day I found the amazingly
simple description of the homotopy groups in the schematic set-up\hand{,} which
got me quite excited;  next day\hand{,} from a phone call to Illusie\hand{,} it turned
out that the key assumption in all my ponderings on schematization,
namely the ``abelianization theorem'' (sic) asserting isomorphism between
discrete and schematic (namely quasi-coherent) \var{\xout{co}}{}{typewriter deletion}{ps_134-7_048-v0003}homology (or equivalent-%
ly,  cohomology) invariants,  was definitely \uline{false};  consequently,  the
canonical functor from schematic to discrete homotopy types turns out
definitely not to be an equivalence.  This completely overthrows the
idyllic picture in my mind about the relationship between schematic  and
discrete homotopy types  -  but the reflection on schematic homotopy types
``in their own right'' had by then proceeded far enough, the very day befo-%
re,\hand{so} that my faith in the relevance of schematic homotopy types was'nt
seriously shaken  -  rather,  I got excited at drawing a systematic
``bilan'' from the evidence \var{\inserted[typed]{now}}{now}{typed insertion}{ps_134-7_048-v0004} \var{\deleted{at}}{}{deletion}{ps_134-7_048-v0005} hand,about the prospects of developping a
theory of ``schematic homotopy types'' satisfying some basic formal proper-%
ties,  whether or not such theory \var{\over{\unsure{is}}{be}}{\hand{be}}{writing over}{ps_134-7_048-v0006} based on semisimplicial unipotent
bundles as \var{\deleted{objects}}{}{deletion}{ps_134-7_048-v0007} models,  or \var{\inserted{on}}{\hand{on}}{insertion}{ps_134-7_048-v0008} any other kind of models making a sense
over arbitrary ground rings.  Before reverting to a review of the more
formal properties of abeli\var{\over{\unsure{a}}{a}}{\hand{a}}{writing over}{ps_134-7_048-v0009}nization in the context of the basic modelizer
(Cat),  I would like still to write down with some care what had thus
come up monday and tuesday last week...

The next days I felt a great fatigue in all my body,  and I then
stopped (till yesterday) any involvment in mathematical reflection.  I
am glad I followed this time the hint that had come to me through my
body,  rather than brush it aside and go on rushing ahead with the work
I was so intensely involved in,  as had been a rule in my life for \var{\xout{\illegible{1}}}{}{typewriter deletion}{ps_134-7_048-v0010}
many years.  This time,  I understood that the reluctance of the body
to follow that forward rush,  even though I was taking good care of
myself with sleep and food,  had strong reasons,  which had nothing to do
with neither sleep or food or with my general way of life.  Rather,
during the weeks before and also during those very days,  a number of
things had occured in my life,  not all visibly related and of differing
weight and magnitude,  to none of which I had really devoted serious
reflection,  nor even a minimum of time and attention \var{\inserted{needed}}{\hand{needed}}{insertion}{ps_134-7_048-v0011} for giving me a

\margin{48}

% Visual collation: both supplied image halves inspected; see CHECK_DECISIONS_134_7_8.tsv.

%% END SOURCE ps_134-7_048
\ifnum\EditionMode=1
\psapparatusstart
\psentry{ps_134-7_048-v0001}{1}{writing over}{\over{.}{,}}{\hand{,}}{The superposed handwritten reading replaces the underlying reading.}
\psentry{ps_134-7_048-v0002}{2}{writing over}{\over{F}{f}}{\hand{f}}{The superposed handwritten reading replaces the underlying reading.}
\psentry{ps_134-7_048-v0003}{3}{typewriter deletion}{\xout{co}}{}{The typewriter strikes out this material; it leaves the corrected running text.}
\psentry{ps_134-7_048-v0004}{4}{typed insertion}{\inserted[typed]{now}}{now}{The typed insertion is incorporated at its indicated position; insertion brackets leave the running text.}
\psentry{ps_134-7_048-v0005}{5}{deletion}{\deleted{at}}{}{The author cancels this material by hand; it leaves the corrected running text.}
\psentry{ps_134-7_048-v0006}{6}{writing over}{\over{\unsure{is}}{be}}{\hand{be}}{The superposed handwritten reading replaces the underlying reading.}
\psentry{ps_134-7_048-v0007}{7}{deletion}{\deleted{objects}}{}{The author cancels this material by hand; it leaves the corrected running text.}
\psentry{ps_134-7_048-v0008}{8}{insertion}{\inserted{on}}{\hand{on}}{The handwritten insertion is incorporated at its indicated position; insertion brackets leave the running text.}
\psentry{ps_134-7_048-v0009}{9}{writing over}{\over{\unsure{a}}{a}}{\hand{a}}{The superposed handwritten reading replaces the underlying reading.}
\psentry{ps_134-7_048-v0010}{10}{typewriter deletion}{\xout{\illegible{1}}}{}{The typewriter strikes out this material; it leaves the corrected running text.}
\psentry{ps_134-7_048-v0011}{11}{insertion}{\inserted{needed}}{\hand{needed}}{The handwritten insertion is incorporated at its indicated position; insertion brackets leave the running text.}
\psapparatusend
\fi
\pssource{ps_134-7_049}{134-7/49}{https://grothendieck.umontpellier.fr/web/viewer.html?file=../134-7.pdf\#page=50}
\psnumber{ps_134-7_049-v0001}{1}
\psnumber{ps_134-7_049-v0002}{2}
\psnumber{ps_134-7_049-v0003}{3}
\psnumber{ps_134-7_049-v0004}{4}
\psnumber{ps_134-7_049-v0005}{5}
\psnumber{ps_134-7_049-v0006}{6}
\psnumber{ps_134-7_049-v0007}{7}
\psnumber{ps_134-7_049-v0008}{8}
\psnumber{ps_134-7_049-v0009}{9}
\psnumber{ps_134-7_049-v0010}{10}
\psnumber{ps_134-7_049-v0011}{11}
\psnumber{ps_134-7_049-v0012}{12}
\psnumber{ps_134-7_049-v0013}{13}
\psnumber{ps_134-7_049-v0014}{14}
\psnumber{ps_134-7_049-v0015}{15}
\psnumber{ps_134-7_049-v0016}{16}
\psnumber{ps_134-7_049-v0017}{17}
\psnumber{ps_134-7_049-v0018}{18}
\psnumber{ps_134-7_049-v0019}{19}
\psnumber{ps_134-7_049-v0020}{20}
\psnumber{ps_134-7_049-v0021}{21}
\psnumber{ps_134-7_049-v0022}{22}
\psnumber{ps_134-7_049-v0023}{23}
\psnumber{ps_134-7_049-v0024}{24}
\psnumber{ps_134-7_049-v0025}{25}
\psnumber{ps_134-7_049-v0026}{26}
%% BEGIN SOURCE ps_134-7_049
% Pursuing Stacks (cotes 134-2 to 134-8): archive cote 134-7, site page 49 (PDF page 50), typescript page 491.
% typescript. Draft compared: ps6.tex:2419-2478.
% AI-transcribed by Claude Opus 5.5 (claude-opus-5-5, adaptive thinking, effort high), Message Batches API, batch ps3 (msgbatch_01XdFmuVVBAcgc7EFxdtCqsj), 2026-11-12; reviewed against both supplied scan halves; see CHECK_DECISIONS.tsv.
\scanpage{134-7/49}
\mspage{491}

chance to let these things and their meaning ``enter''.  In lack of this,
there was little chance my response to current events \var{\inserted[typed]{would}}{would}{typed insertion}{ps_134-7_049-v0001} \var{\deleted{to}}{}{deletion}{ps_134-7_049-v0002} be any better
than purely mechanical,  and my interaction with some of the people I
love \var{\inserted{would}}{\hand{would}}{insertion}{ps_134-7_049-v0003} \var{\deleted{to}}{}{deletion}{ps_134-7_049-v0004} be in any way creative.  There was this need for being attentive,
an urgent need springing from life itself and which I was about to ignore
-  and their was this drive,  this impatience driving me recklessly
ahead,  with no look left nor right.  Of course,  I did know about the
need, ``somewhere''  -  and in my head too I kind of knew,  but the head
was prejudiced as usual and would take no notice,  not of the need and
not of the conflict between an urgent need and a powerful\hand{,} \var{\xout{craving}}{}{typewriter deletion}{ps_134-7_049-v0005}
ego-inve\var{\over{\unsure{s}}{s}}{\hand{s}}{writing over}{ps_134-7_049-v0006}ted passion.  The head was prejudiced and foolish  -  so \var{\over[typed]{\unsure{i}}{i}}{i}{typewriter overwriting}{ps_134-7_049-v0007}t
was my body \var{\deleted{who}}{}{deletion}{ps_134-7_049-v0008} finally \var{\inserted{who}}{\hand{who}}{insertion}{ps_134-7_049-v0009} told me:  now you stop this non-sense and \var{\inserted{you}}{\hand{you}}{insertion}{ps_134-7_049-v0010} take
care of what you well know you better take care \uline{first},  and \uline{now}!  And
his language was strong and simple enough and caus\var{\over{e}{ed}}{\hand{ed}}{writing over}{ps_134-7_049-v0011} me to listen.

Thus the main work I was involved in for these last seven days was
to let a number of things ``enter''  -  mainly things that were being revea-%
led through the death of my granddaughter Ella.  It surely \uline{was}\hand{|}``work''  ,
taking up the\hand{|}largest part of my nights and my days,  -  so much so that
I can't really say there is any less fatigue now than seven days ago. It
seems to me though it is'nt quite the same fatigue  -  this time it is
the fatigue coming from work done,  not from work shunned.  The\hand{``}work done\hand{''}
was'nt really done by me,  I feel,  rather work taking place \var{\inserted[typed]{with}}{with}{typed insertion}{ps_134-7_049-v0012}in me,  and
``my'' main contribution has been to allow it to take place,by providing
the necessary time and qui\var{\over{\unsure{t}}{et}}{\hand{et}}{writing over}{ps_134-7_049-v0013}ness;  and of course,  also,  to allow the
outcome of this work to become conscious knowledge,  rather than burying
it away in some dark corner of the mind.  The \var{\deleted{mat}}{}{deletion}{ps_134-7_049-v0014} rough material\hand{,}as well
as the outcome of this work\hand{,}have not been this time new facts \var{\over{\unsure{a}}{o}}{\hand{o}}{writing over}{ps_134-7_049-v0015}r new in-%
sights;  rather\hand{,} things which I had come to\var{\deleted{\unsure{o}}}{}{deletion}{ps_134-7_049-v0016} perceive and notice for some
time already,  over the last two years,  without \var{\deleted{allowing}}{}{deletion}{ps_134-7_049-v0017} \var{\deleted{giving}}{}{deletion}{ps_134-7_049-v0018} \var{\inserted{granting}}{\hand{granting}}{insertion}{ps_134-7_049-v0019} them the
proper we\var{\over{\unsure{ie}}{ei}}{\hand{ei}}{writing over}{ps_134-7_049-v0020}ght and perspective\var{\over{;}{\unsure{,}}}{\hand{\unsure{,}}}{writing over}{ps_134-7_049-v0021}- somehow as if I did'nt quite believe what
I was unmistakeably perceiving,  or \var{\deleted{not}}{}{deletion}{ps_134-7_049-v0022} \var{\inserted{did'nt}}{\hand{did'nt}}{insertion}{ps_134-7_049-v0023} take it quite seriously.  Such a
thing\hand{,}I noticed\hand{,}happens quite often,  not only with me,  and takes care
of making even the most lively perseptions inocuous,  by disconnecting
them at all price from the image of reality and of ourselves we are car-%
rying around with us,  \hand{\unsure{**}} \var{\deleted{which}}{}{deletion}{ps_134-7_049-v0024} \var{\inserted{that the image}}{\hand{that the image}}{insertion}{ps_134-7_049-v0025} remain\hand{s} static\hand{,} \var{\deleted{and almost entirely}}{}{deletion}{ps_134-7_049-v0026} unaffected
by any kind of ``information'' flowing around or through us.

\margin{49}

% Visual collation: both supplied image halves inspected; see CHECK_DECISIONS_134_7_8.tsv.

%% END SOURCE ps_134-7_049
\ifnum\EditionMode=1
\psapparatusstart
\psentry{ps_134-7_049-v0001}{1}{typed insertion}{\inserted[typed]{would}}{would}{The typed insertion is incorporated at its indicated position; insertion brackets leave the running text.}
\psentry{ps_134-7_049-v0002}{2}{deletion}{\deleted{to}}{}{The author cancels this material by hand; it leaves the corrected running text.}
\psentry{ps_134-7_049-v0003}{3}{insertion}{\inserted{would}}{\hand{would}}{The handwritten insertion is incorporated at its indicated position; insertion brackets leave the running text.}
\psentry{ps_134-7_049-v0004}{4}{deletion}{\deleted{to}}{}{The author cancels this material by hand; it leaves the corrected running text.}
\psentry{ps_134-7_049-v0005}{5}{typewriter deletion}{\xout{craving}}{}{The typewriter strikes out this material; it leaves the corrected running text.}
\psentry{ps_134-7_049-v0006}{6}{writing over}{\over{\unsure{s}}{s}}{\hand{s}}{The superposed handwritten reading replaces the underlying reading.}
\psentry{ps_134-7_049-v0007}{7}{typewriter overwriting}{\over[typed]{\unsure{i}}{i}}{i}{The superposed typed reading replaces the underlying reading.}
\psentry{ps_134-7_049-v0008}{8}{deletion}{\deleted{who}}{}{The author cancels this material by hand; it leaves the corrected running text.}
\psentry{ps_134-7_049-v0009}{9}{insertion}{\inserted{who}}{\hand{who}}{The handwritten insertion is incorporated at its indicated position; insertion brackets leave the running text.}
\psentry{ps_134-7_049-v0010}{10}{insertion}{\inserted{you}}{\hand{you}}{The handwritten insertion is incorporated at its indicated position; insertion brackets leave the running text.}
\psentry{ps_134-7_049-v0011}{11}{writing over}{\over{e}{ed}}{\hand{ed}}{The superposed handwritten reading replaces the underlying reading.}
\psentry{ps_134-7_049-v0012}{12}{typed insertion}{\inserted[typed]{with}}{with}{The typed insertion is incorporated at its indicated position; insertion brackets leave the running text.}
\psentry{ps_134-7_049-v0013}{13}{writing over}{\over{\unsure{t}}{et}}{\hand{et}}{The superposed handwritten reading replaces the underlying reading.}
\psentry{ps_134-7_049-v0014}{14}{deletion}{\deleted{mat}}{}{The author cancels this material by hand; it leaves the corrected running text.}
\psentry{ps_134-7_049-v0015}{15}{writing over}{\over{\unsure{a}}{o}}{\hand{o}}{The superposed handwritten reading replaces the underlying reading.}
\psentry{ps_134-7_049-v0016}{16}{deletion}{\deleted{\unsure{o}}}{}{The author cancels this material by hand; it leaves the corrected running text.}
\psentry{ps_134-7_049-v0017}{17}{deletion}{\deleted{allowing}}{}{The author cancels this material by hand; it leaves the corrected running text.}
\psentry{ps_134-7_049-v0018}{18}{deletion}{\deleted{giving}}{}{The author cancels this material by hand; it leaves the corrected running text.}
\psentry{ps_134-7_049-v0019}{19}{insertion}{\inserted{granting}}{\hand{granting}}{The handwritten insertion is incorporated at its indicated position; insertion brackets leave the running text.}
\psentry{ps_134-7_049-v0020}{20}{writing over}{\over{\unsure{ie}}{ei}}{\hand{ei}}{The superposed handwritten reading replaces the underlying reading.}
\psentry{ps_134-7_049-v0021}{21}{writing over}{\over{;}{\unsure{,}}}{\hand{\unsure{,}}}{The superposed handwritten reading replaces the underlying reading.}
\psentry{ps_134-7_049-v0022}{22}{deletion}{\deleted{not}}{}{The author cancels this material by hand; it leaves the corrected running text.}
\psentry{ps_134-7_049-v0023}{23}{insertion}{\inserted{did'nt}}{\hand{did'nt}}{The handwritten insertion is incorporated at its indicated position; insertion brackets leave the running text.}
\psentry{ps_134-7_049-v0024}{24}{deletion}{\deleted{which}}{}{The author cancels this material by hand; it leaves the corrected running text.}
\psentry{ps_134-7_049-v0025}{25}{insertion}{\inserted{that the image}}{\hand{that the image}}{The handwritten insertion is incorporated at its indicated position; insertion brackets leave the running text.}
\psentry{ps_134-7_049-v0026}{26}{deletion}{\deleted{and almost entirely}}{}{The author cancels this material by hand; it leaves the corrected running text.}
\psapparatusend
\fi
\pssource{ps_134-7_050}{134-7/50}{https://grothendieck.umontpellier.fr/web/viewer.html?file=../134-7.pdf\#page=51}
\psnumber{ps_134-7_050-v0001}{1}
\psnumber{ps_134-7_050-v0002}{2}
\psnumber{ps_134-7_050-v0003}{3}
\psnumber{ps_134-7_050-v0004}{4}
\psnumber{ps_134-7_050-v0005}{5}
\psnumber{ps_134-7_050-v0006}{6}
\psnumber{ps_134-7_050-v0007}{7}
%% BEGIN SOURCE ps_134-7_050
% Pursuing Stacks (cotes 134-2 to 134-8): archive cote 134-7, site page 50 (PDF page 51), typescript page 492.
% typescript. Draft compared: ps6.tex:2458-2524.
% AI-transcribed by Claude Opus 5.5 (claude-opus-5-5, adaptive thinking, effort high), Message Batches API, batch ps3 (msgbatch_01XdFmuVVBAcgc7EFxdtCqsj), 2026-11-12; reviewed against both supplied scan halves; see CHECK_DECISIONS.tsv.
\scanpage{134-7/50}
\mspage{492}

10.9.

117. Maybe the best will be to write up (and possibly develop some)
my reflections \var{\xout{after the notes}}{}{typewriter deletion}{ps_134-7_050-v0001}\hand{(}of the two days after I stopped with the
notes\hand{)} roughly in the order as they occured.

There were some somewhat technical afterthoughts. One was about the
logical difficulty coming from the site Aff$_{/k}$ on page 446 not being an
$\underline{U}$-site (where $\underline{U}$ is the universe we are working in), hence strictly spea-%
king the category of all sheaves on this site is possibly not even a
$\underline{U}$-category (i.e. the Hom objects need not be small i.e. with cardinal in
$\underline{U}$ ), still less a topos, and hence the standard panoply of notions and
constructions in a topos does not apply. This does'nt look really serious
though, one way out is to limit beforehand the ``size'' of the unipotent
bundles we are allowing i.e. of the k-modules describing them, in terms
say of cardinality of a family of generators of the latter - and then
restrict accordingly the size of the k-algebras k' taken as ``arguments''
for our sheaves i.e. as objects of the basic site we\var{\inserted{are}}{\hand{are}}{insertion}{ps_134-7_050-v0002} work\var{\over{in.}{ing}}{\hand{ing}}{writing over}{ps_134-7_050-v0003}\var{\inserted{on.}}{\hand{on.}}{insertion}{ps_134-7_050-v0004} For instanc\hand{e,}
when working with unipotent bundles of finite type only (i.e. correspon-%
ding to k-modules of finite type - a rather interesting and natural
finiteness \var{\xout{requi}}{}{typewriter deletion}{ps_134-7_050-v0005} condition\var{\inserted[typed]{anyhow}}{anyhow}{typed insertion}{ps_134-7_050-v0006} on the components of a schematic model $X_{\ast}$ ),
it is appropriate to work on the ``fppf site'',where the arguments k' are
k-algebras of \uline{finite presentation}. If we should be unwilling to be limi\hand{ted}
by a fixed size restriction on the unipotent bundles we are working with
(and hence also on the corresponding homotopy types), we may have to work
with a hierarchy of size restrictions and passage from one to any other
less stringent one - a technical nuisance to be sure, if we do'nt find
a more elegant way out, but surely not a substantial difficulty. At the
present heuristic stage of reflections, it does'nt seem worth while real-%
ly to dwell on such questions any longer.

Another afterthought is about the functor
\[ M \hand{\longmapsto} W(M) \quad : \quad \mathrm{Ab}_k =(k\text{-Mod}) \hand{\longrightarrow} (\underline{O}_k\text{-Mod}) \]
from k-modules to \var{\xout{quasi}}{}{typewriter deletion}{ps_134-7_050-v0007} $\underline{O}_k$-Modules - a fully faithful functor we know,
whose essential image by definition consists of the so-called \uline{quasi}-\uline{cohe-}%
\uline{rent} sheaves. A little caution is needed, as this functor is right
exact, but \uline{not} \uline{exact} i.e. it does not commute with formation of kernels,
because for a k-algebra k' which is'nt flat, the functor
\[ M \hand{\longmapsto} M\hand{\otimes}_k k' \]

\margin{50}

% Visual collation: both supplied image halves inspected; see CHECK_DECISIONS_134_7_8.tsv.

%% END SOURCE ps_134-7_050
\ifnum\EditionMode=1
\psapparatusstart
\psentry{ps_134-7_050-v0001}{1}{typewriter deletion}{\xout{after the notes}}{}{The typewriter strikes out this material; it leaves the corrected running text.}
\psentry{ps_134-7_050-v0002}{2}{insertion}{\inserted{are}}{\hand{are}}{The handwritten insertion is incorporated at its indicated position; insertion brackets leave the running text.}
\psentry{ps_134-7_050-v0003}{3}{writing over}{\over{in.}{ing}}{\hand{ing}}{The superposed handwritten reading replaces the underlying reading.}
\psentry{ps_134-7_050-v0004}{4}{insertion}{\inserted{on.}}{\hand{on.}}{The handwritten insertion is incorporated at its indicated position; insertion brackets leave the running text.}
\psentry{ps_134-7_050-v0005}{5}{typewriter deletion}{\xout{requi}}{}{The typewriter strikes out this material; it leaves the corrected running text.}
\psentry{ps_134-7_050-v0006}{6}{typed insertion}{\inserted[typed]{anyhow}}{anyhow}{The typed insertion is incorporated at its indicated position; insertion brackets leave the running text.}
\psentry{ps_134-7_050-v0007}{7}{typewriter deletion}{\xout{quasi}}{}{The typewriter strikes out this material; it leaves the corrected running text.}
\psapparatusend
\fi
\pssource{ps_134-7_051}{134-7/51}{https://grothendieck.umontpellier.fr/web/viewer.html?file=../134-7.pdf\#page=52}
\psnumber{ps_134-7_051-v0001}{1}
\psnumber{ps_134-7_051-v0002}{2}
\psnumber{ps_134-7_051-v0003}{3}
\psnumber{ps_134-7_051-v0004}{4}
\psnumber{ps_134-7_051-v0005}{5}
\psnumber{ps_134-7_051-v0006}{6}
\psnumber{ps_134-7_051-v0007}{7}
\psnumber{ps_134-7_051-v0008}{8}
\psnumber{ps_134-7_051-v0009}{9}
\psnumber{ps_134-7_051-v0010}{10}
\psnumber{ps_134-7_051-v0011}{11}
%% BEGIN SOURCE ps_134-7_051
% Pursuing Stacks (cotes 134-2 to 134-8): archive cote 134-7, site page 51 (PDF page 52), typescript page 493.
% typescript. Draft compared: ps6.tex:2507-2578.
% AI-transcribed by Claude Opus 5.5 (claude-opus-5-5, adaptive thinking, effort high), Message Batches API, batch ps4 (msgbatch_012Aoc1NuskATPsyBeKSYK7J), 2026-11-12; reviewed against both supplied scan halves; see CHECK_DECISIONS.tsv.
\scanpage{134-7/51}
\mspage{493}

does'nt.  When we wrote down a formula such as (21) on page 487,  we were
implicitly making use of the assumption that for k-modules M,N we have
a canonical isomorphism
\[ (1) \qquad \mathrm{RHom}_{k}(M,N) \hand{\xleftarrow{\sim}} \mathrm{RHom}_{\underline{O}_k}(W(M),W(N)) \ , \]
at any rate such a formula is needed if we want to view (21) as a more
explicit way of writing (20) or (20').  Using a free resolution L. of M,
and the fact that
\[ \mathrm{Ext}^{i}_{\underline{O}_k}(\underline{O}_k, \var{\deleted{M}}{}{deletion}{ps_134-7_051-v0001} W(N)) \hand{\simeq} H^{i}(\mathrm{Spec}(k), W(N)) \ = 0 \ \text{if} \ \ i \hand{>} 0, \]
stemming from the cohomology properties of flat descent,  we easily get
a map (1),  but I did not check that this map is an isomorphism,  the dif-%
ficulty coming  from the fact that  W(L.)  need not be a resolution of
W(M),  unless M is flat.  This perplexity already arises in the fppf
context  -  surely Larry Breen should know the answer.  For what we are
\var{\deleted{reall}}{}{deletion}{ps_134-7_051-v0002} after here \var{\deleted{though}}{}{deletion}{ps_134-7_051-v0003} it does'nt seem to  matter too much,  as the com-%
putations of loc. cit.  were of interest mainly (maybe exclusively) in the
case when M is projective or at any rate flat,  i.e. when working with
\uline{flat} unipotent bundles   -  in which case (1) is indeed an isomorphism.

The interpretation of polynomial maps between k-modules M,N  in terms
of the topological augmented coalgebras  $\hand{\Gamma}_k{}^{\hand{\wedge}}$ associated to these  may seem
a little forbidding to some readers.  In the all-important and typical
case when M \var{\xout{is}}{}{typewriter deletion}{ps_134-7_051-v0004} \var{\inserted[typed]{and N are}}{and N are}{typed insertion}{ps_134-7_051-v0005} projective and of finite type,  things become quite evident
though,  by just dualizing the more familiar concepts around polynomial
functions and homomorphisms between rings of such.  Thus W(M) is just a
usual vector-bundle,\var{\xout{\unsure{t}}}{}{typewriter deletion}{ps_134-7_051-v0006}hence also a\var{\xout{n}}{}{typewriter deletion}{ps_134-7_051-v0007} true honest affine scheme over k,
whose affine ring is the ring of polynomial functions on M,  which can
be identified with  $\var{\over{\mathrm{Sym}}{\mathrm{Sym}}}{\hand{\mathrm{Sym}}}{writing over}{ps_134-7_051-v0008}^{\hand{*}}_{\hand{k}}(M^{\hand{\vee}})$,  the symmetric algebra on the dual module :
\[ (2) \qquad W(M) \hand{\simeq} \mathrm{Spec}(\mathrm{Sym}^{\hand{*}}_{\hand{k}}(M\ )) \ , \]
and similarly for W(N).  Polynomial maps from M to N i.e. maps from the
k-scheme W(M) to the k-scheme W(N),   just correspond to k-algebra homomor-%
phisms
\[ \hand{(3)} \qquad \mathrm{Sym}^{*}_{\hand{k}}(N) \hand{\longrightarrow} \mathrm{Sym}^{*}_{\hand{k}}(M) \]
(irrespective of the graded structures).  As each of these algebras\var{\inserted{as a k-module,}}{\hand{as a k-module,}}{insertion}{ps_134-7_051-v0009} is
a filtering direct limit of \var{\inserted{it's}}{\hand{it's}}{insertion}{ps_134-7_051-v0010} projective \var{\inserted{sub}}{\hand{sub}}{insertion}{ps_134-7_051-v0011}modules of finite type such as
\[ \hand{(4)} \qquad \mathrm{Sym}_{k}(M)(i) \ = \ \hand{\bigoplus}_{j \hand{\leq} i} \mathrm{Sym}_{k}{}^{j}(M) \qquad , \]

\margin{51}

% Visual collation: both supplied image halves inspected; see CHECK_DECISIONS_134_7_8.tsv.

%% END SOURCE ps_134-7_051
\ifnum\EditionMode=1
\psapparatusstart
\psentry{ps_134-7_051-v0001}{1}{deletion}{$\deleted{M}$}{}{The author cancels this material by hand; it leaves the corrected running text.}
\psentry{ps_134-7_051-v0002}{2}{deletion}{\deleted{reall}}{}{The author cancels this material by hand; it leaves the corrected running text.}
\psentry{ps_134-7_051-v0003}{3}{deletion}{\deleted{though}}{}{The author cancels this material by hand; it leaves the corrected running text.}
\psentry{ps_134-7_051-v0004}{4}{typewriter deletion}{\xout{is}}{}{The typewriter strikes out this material; it leaves the corrected running text.}
\psentry{ps_134-7_051-v0005}{5}{typed insertion}{\inserted[typed]{and N are}}{and N are}{The typed insertion is incorporated at its indicated position; insertion brackets leave the running text.}
\psentry{ps_134-7_051-v0006}{6}{typewriter deletion}{\xout{\unsure{t}}}{}{The typewriter strikes out this material; it leaves the corrected running text.}
\psentry{ps_134-7_051-v0007}{7}{typewriter deletion}{\xout{n}}{}{The typewriter strikes out this material; it leaves the corrected running text.}
\psentry{ps_134-7_051-v0008}{8}{writing over}{$\over{\mathrm{Sym}}{\mathrm{Sym}}$}{$\hand{\mathrm{Sym}}$}{The superposed handwritten reading replaces the underlying reading.}
\psentry{ps_134-7_051-v0009}{9}{insertion}{\inserted{as a k-module,}}{\hand{as a k-module,}}{The handwritten insertion is incorporated at its indicated position; insertion brackets leave the running text.}
\psentry{ps_134-7_051-v0010}{10}{insertion}{\inserted{it's}}{\hand{it's}}{The handwritten insertion is incorporated at its indicated position; insertion brackets leave the running text.}
\psentry{ps_134-7_051-v0011}{11}{insertion}{\inserted{sub}}{\hand{sub}}{The handwritten insertion is incorporated at its indicated position; insertion brackets leave the running text.}
\psapparatusend
\fi
\pssource{ps_134-7_052}{134-7/52}{https://grothendieck.umontpellier.fr/web/viewer.html?file=../134-7.pdf\#page=53}
\psnumber{ps_134-7_052-v0001}{1}
\psnumber{ps_134-7_052-v0002}{2}
\psnumber{ps_134-7_052-v0003}{3}
\psnumber{ps_134-7_052-v0004}{4}
\psnumber{ps_134-7_052-v0005}{5}
\psnumber{ps_134-7_052-v0006}{6}
\psnumber{ps_134-7_052-v0007}{7}
%% BEGIN SOURCE ps_134-7_052
% Pursuing Stacks (cotes 134-2 to 134-8): archive cote 134-7, site page 52 (PDF page 53), typescript page 494.
% typescript. Draft compared: ps6.tex:2559-2630.
% AI-transcribed by Claude Opus 5.5 (claude-opus-5-5, adaptive thinking, effort high), Message Batches API, batch ps4 (msgbatch_012Aoc1NuskATPsyBeKSYK7J), 2026-11-12; reviewed against both supplied scan halves; see CHECK_DECISIONS.tsv.
\scanpage{134-7/52}
\mspage{494}

the dual module
\[\hand{\Gamma}_k^{\hand{\wedge}}(M)\hand{\simeq}(\mathrm{Sym}_k^{*}(M))\hand{^{\vee}}\hand{\simeq}\ \hand{\prod}_j\ (\mathrm{Sym}_k^{j}(M)\hand{^{\vee}}\ \hand{\simeq}\ \hand{\Gamma}_k^{\,j}(M))\]
\unsure{may|mgy} be viewed as the inverse limit of the duals of those submodules
(compare p.484 (11) ), and topologized accordingly; linear maps (3)
may be interpreted in terms of continuous maps between the dual structures
(or, equivalently, between the corresponding proobjects)
\[(3')\qquad \hand{\Gamma}_k^{\hand{\uparrow}}(M)\hand{\longrightarrow}\hand{\Gamma}_k^{\hand{\wedge}}(N)\quad ,\]
and compatibility of (3) with multiplication and units is expressed by
compatibility of (3') with``comultiplication'' i.e. diagonal maps, and
with augmentations. On the other hand, maps $W(M)\hand{\rightarrow} W(N)$ respecting the
``pointed structures'' coming from zero sections, correspond to maps (3')
transforming 1 into 1 , besides the other requirements. Such maps,
it turns out, automatically \var{\xout{transform}}{}{typewriter deletion}{ps_134-7_052-v0001} induce a map between the submodules
\[(4)\qquad \hand{\Gamma}_k(M)\hand{\longrightarrow}\hand{\Gamma}_k(N)\quad ,\]
which is rather evident indeed, if we \var{\xout{recall that}}{}{typewriter deletion}{ps_134-7_052-v0002} remind ourselves of \var{\xout{th}}{}{typewriter deletion}{ps_134-7_052-v0003}
the fact that the submodule $\hand{\Gamma}_k(M)$ may be viewed as the\var{\inserted{(topological)}}{\hand{(topological)}}{insertion}{ps_134-7_052-v0004}dual of $\mathrm{Sym}_k^{*}(M)$\hand{,}
topologized by the powers of the augmentation ideal \var{\deleted{\xout{Sym}\hand{$^{+}$}}}{}{deletion}{ps_134-7_052-v0005}
\[\mathrm{Sym}_k^{+}(M)\ =\ \hand{\bigoplus}_{i>0}\ \mathrm{Sym}_k^{i}(M)\quad ,\]
or equivalently of the corresponding adic completion
\[(5)\qquad \mathrm{Sym}_k^{\overset{\hand{\illegible{1}}}{*}\hand{\wedge}}(M)\ =\ \underset{i}{\underset{\hand{\longleftarrow}}{\mathrm{lim}}}\ \mathrm{Sym}_k^{*}(M)(i)\hand{\simeq}\ \hand{\prod}_{\hand{j\geq 0}}\mathrm{Sym}_k^{j}(M)\quad ,\]
and correspondingly for N. The``pointed'' assumption on \var{\deleted{(3)}}{}{deletion}{ps_134-7_052-v0006} a map $W(M)\hand{\rightarrow}W(N)$\hand{,}
in terms of the corresponding homomorphism of k-algebras (3), just
translates into compatibility with the augmentations or, equivalently,
with the corresponding ideals $\mathrm{Sym}_k^{+}$ , which implies that it induces a
homomorphism on the corresponding adic rings, and hence by duality a
homomorphism (4) on their duals.

These reminders bring near that working with the\var{\inserted[typed]{(discrete)}}{(discrete)}{typed insertion}{ps_134-7_052-v0007}k-algebras
$\mathrm{Sym}_k^{*}(M)$, or equivalently with the \uline{topological} coalgebras $\hand{\Gamma}_k^{\hand{\wedge}}(M)$, amounts
to working with ``unipotent bundles'' (projective and of finite type),
which are just usual schemes (of a rather particular structure of course);
whereas working with the \uline{topological} k-algebras $\mathrm{Sym}_k^{*\hand{\wedge}}(M)$, or equivalent
\hand{-}ly with the discrete coalgebras $\hand{\Gamma}_k(M)$ , amounts to working with \uline{formal}
\uline{schemes} , namely essentially, with the formal completions of the former
along the zero sections.( Of course, the topological objects just consi-

\margin{52}

% Visual collation: both supplied image halves inspected; see CHECK_DECISIONS_134_7_8.tsv.

%% END SOURCE ps_134-7_052
\ifnum\EditionMode=1
\psapparatusstart
\psentry{ps_134-7_052-v0001}{1}{typewriter deletion}{\xout{transform}}{}{The typewriter strikes out this material; it leaves the corrected running text.}
\psentry{ps_134-7_052-v0002}{2}{typewriter deletion}{\xout{recall that}}{}{The typewriter strikes out this material; it leaves the corrected running text.}
\psentry{ps_134-7_052-v0003}{3}{typewriter deletion}{\xout{th}}{}{The typewriter strikes out this material; it leaves the corrected running text.}
\psentry{ps_134-7_052-v0004}{4}{insertion}{\inserted{(topological)}}{\hand{(topological)}}{The handwritten insertion is incorporated at its indicated position; insertion brackets leave the running text.}
\psentry{ps_134-7_052-v0005}{5}{deletion}{\deleted{\xout{Sym}\hand{$^{+}$}}}{}{The author cancels this material by hand; it leaves the corrected running text.}
\psentry{ps_134-7_052-v0006}{6}{deletion}{\deleted{(3)}}{}{The author cancels this material by hand; it leaves the corrected running text.}
\psentry{ps_134-7_052-v0007}{7}{typed insertion}{\inserted[typed]{(discrete)}}{(discrete)}{The typed insertion is incorporated at its indicated position; insertion brackets leave the running text.}
\psapparatusend
\fi
\pssource{ps_134-7_053}{134-7/53}{https://grothendieck.umontpellier.fr/web/viewer.html?file=../134-7.pdf\#page=54}
\psnumber{ps_134-7_053-v0001}{1}
\psnumber{ps_134-7_053-v0002}{2}
\psnumber{ps_134-7_053-v0003}{3}
\psnumber{ps_134-7_053-v0004}{4}
\psnumber{ps_134-7_053-v0005}{5}
\psnumber{ps_134-7_053-v0006}{6}
\psnumber{ps_134-7_053-v0007}{7}
\psnumber{ps_134-7_053-v0008}{8}
\psnumber{ps_134-7_053-v0009}{9}
%% BEGIN SOURCE ps_134-7_053
% Pursuing Stacks (cotes 134-2 to 134-8): archive cote 134-7, site page 53 (PDF page 54), typescript page 495.
% typescript. Draft compared: ps6.tex:2617-2677.
% AI-transcribed by Claude Opus 5.5 (claude-opus-5-5, adaptive thinking, effort high), Message Batches API, batch ps4 (msgbatch_012Aoc1NuskATPsyBeKSYK7J), 2026-11-12; reviewed against both supplied scan halves; see CHECK_DECISIONS.tsv.
\scanpage{134-7/53}
\mspage{495}

may be equally viewed as being \uline{pro}-\uline{modules} endowed with suitable extra
structure.) Cor\-res\-pond\-ing\-ly, we will expect the cohomology invariants
constructed in terms of (the apparently more forbidding) $\hand{\Gamma}_k{}^{\hand{\wedge}}$ to express
quasi-coherent cohomology of the corresponding \uline{schemes}, or semisimplicial
systems of such; whereas, working with the (apparently more anodine !)
$\hand{\Gamma}_k$ will lead to cohomology invariants of \uline{formal schemes} and semisimplici-%
al systems of such. Both types of invariants are of interest it would
seem, the one however which looks the more relevant in connection with
studying ordinary homotopy types in terms of schematic one's, is surely
the first. On the other hand, there is'nt any reason whatever to
believe that under fairly general conditions, these two types of invari-%
ants are going to be isomorphic, by the evident map from ``schematic'' to
``formal'' quasicoherent cohomology. To say it differently, I do no longer
expect that underreasonably wide assumptions, the map (24) of p.488
\[
(6)\qquad L_{pt}(X_{*})\hand{\longrightarrow} L(X_{*})
\]
is a quasi-isomorphism, nor behaves like one with respect to taking
$\mathrm{Ext}^{i}$ 's with values in a quasicoherent module, as I was \var{\inserted[typed]{hastily}}{hastily}{typed insertion}{ps_134-7_053-v0001} surmising for
about one week, while loosing track of the geometric meaning of the
algebraic objects I was playing around with. To give just one example,
take $X_{*}$ to be the standard semisimplicial \var{\xout{object}}{}{typewriter deletion}{ps_134-7_053-v0002} unipotent bundle associa-%
ted to the group-object $\underline{G}_a$ in $\hand{\mathcal{U}}(k)$, namely just the usual affine line
with addition law. The Ext$^{i}$'s of the two members of (6) with values in
the k-module k may be interpreted as either schematic or formal Hochschild
cohomology of the additive group, with coefficients in \var{\xout{\illegible{1}}}{}{typewriter deletion}{ps_134-7_053-v0003} $\underline{O}_k = \underline{G}_a$ .
The map of the former into the latter is not \var{\inserted{always}}{\hand{always}}{insertion}{ps_134-7_053-v0004} an isomorphism, already
in dimension 1 , where the two groups to be compared are just the groups
of endomorphisms of $\underline{G}_a$ , and of the corresponding formal group. If k is
of char. p $\hand{>}$ 0, a prime, then the latter group can be described as the
group of all formal power series of the type
\[
F(t) = \hand{\sum}_{i\hand{\geq} 0} c_i t^{(p^{i})} ,
\]
whereas for the first group we must restrict to those F which are polyno-%
mials i.e. \var{\xout{all}}{}{typewriter deletion}{ps_134-7_053-v0005} only a finite number of the coefficients $c_i$ are non-zero.

One may of course object to this example, because the $X_{*}$ we are
working with is not simply connected, and \var{\deleted{that}}{}{deletion}{ps_134-7_053-v0006} \var{\inserted{because}}{\hand{because}}{insertion}{ps_134-7_053-v0007} the example does not apply
\var{\deleted{either}}{}{deletion}{ps_134-7_053-v0008} over a ring such as $\underline{Z}$ , which is the one we are interested in most
of all. I am convinced \var{\inserted{now}}{\hand{now}}{insertion}{ps_134-7_053-v0009} however that even when assuming $X_1 = X_0 = e$
and k=$\underline{Z}$ , (6) is very far from being a quasiisomorphism, even for such

\margin{53}

% Visual collation: both supplied image halves inspected; see CHECK_DECISIONS_134_7_8.tsv.

%% END SOURCE ps_134-7_053
\ifnum\EditionMode=1
\psapparatusstart
\psentry{ps_134-7_053-v0001}{1}{typed insertion}{\inserted[typed]{hastily}}{hastily}{The typed insertion is incorporated at its indicated position; insertion brackets leave the running text.}
\psentry{ps_134-7_053-v0002}{2}{typewriter deletion}{\xout{object}}{}{The typewriter strikes out this material; it leaves the corrected running text.}
\psentry{ps_134-7_053-v0003}{3}{typewriter deletion}{\xout{\illegible{1}}}{}{The typewriter strikes out this material; it leaves the corrected running text.}
\psentry{ps_134-7_053-v0004}{4}{insertion}{\inserted{always}}{\hand{always}}{The handwritten insertion is incorporated at its indicated position; insertion brackets leave the running text.}
\psentry{ps_134-7_053-v0005}{5}{typewriter deletion}{\xout{all}}{}{The typewriter strikes out this material; it leaves the corrected running text.}
\psentry{ps_134-7_053-v0006}{6}{deletion}{\deleted{that}}{}{The author cancels this material by hand; it leaves the corrected running text.}
\psentry{ps_134-7_053-v0007}{7}{insertion}{\inserted{because}}{\hand{because}}{The handwritten insertion is incorporated at its indicated position; insertion brackets leave the running text.}
\psentry{ps_134-7_053-v0008}{8}{deletion}{\deleted{either}}{}{The author cancels this material by hand; it leaves the corrected running text.}
\psentry{ps_134-7_053-v0009}{9}{insertion}{\inserted{now}}{\hand{now}}{The handwritten insertion is incorporated at its indicated position; insertion brackets leave the running text.}
\psapparatusend
\fi
\pssource{ps_134-7_054}{134-7/54}{https://grothendieck.umontpellier.fr/web/viewer.html?file=../134-7.pdf\#page=55}
\psnumber{ps_134-7_054-v0001}{1}
\psnumber{ps_134-7_054-v0002}{2}
\psnumber{ps_134-7_054-v0003}{3}
\psnumber{ps_134-7_054-v0004}{4}
\psnumber{ps_134-7_054-v0005}{5}
\psnumber{ps_134-7_054-v0006}{6}
\psnumber{ps_134-7_054-v0007}{7}
\psnumber{ps_134-7_054-v0008}{8}
\psnumber{ps_134-7_054-v0009}{9}
\psnumber{ps_134-7_054-v0010}{10}
\psnumber{ps_134-7_054-v0011}{11}
\psnumber{ps_134-7_054-v0012}{12}
\psnumber{ps_134-7_054-v0013}{13}
%% BEGIN SOURCE ps_134-7_054
% Pursuing Stacks (cotes 134-2 to 134-8): archive cote 134-7, site page 54 (PDF page 55), typescript page 496.
% typescript. Draft compared: ps6.tex:2663-2741.
% AI-transcribed by Claude Opus 5.5 (claude-opus-5-5, adaptive thinking, effort high), Message Batches API, batch ps4 (msgbatch_012Aoc1NuskATPsyBeKSYK7J), 2026-11-12; reviewed against both supplied scan halves; see CHECK_DECISIONS.tsv.
\scanpage{134-7/54}
\mspage{496}

basic structures as $K(\underline{Z},n)$ with $n\hand{\geqslant} 2$. At any rate, the \hand{``}way-out\hand{''} argument
I have finally been thinking of, in order to try and check (6) is a quasi%
isomorphism, rests on vanishing assumptions which (as I was informed by
Illusie the next day) are wholly unrealistic. This finally clears up,
it would seem, a tenacious misconception which has been sticking to my
first heuristic ponderings about the homology and cohomology formalism
for schematic homotopy types: \uline{one should be very careful not to substitu\var{\xout{t}}{}{typewriter deletion}{ps_134-7_054-v0001}}%
\uline{te the ``pointed'' linearization} $L_{pt}(W(M))\hand{\simeq}\hand{\Gamma}_k(M)$ \uline{for the nonpointed}
\uline{one} $L(W(M))\hand{\simeq}\hand{\Gamma}_k{}^{\hand{\wedge}}(M)$, \uline{in computing homology and cohomology invariants}
\uline{of schematic homotopy types}. To say it differently, in order to be
able to \var{\xout{define}}{}{typewriter deletion}{ps_134-7_054-v0002} compute (or just define) the ``schematic'' homology and
cohomology invariants, we do need as a model a fullfledged semisimplicial
unipotent bundle, not just the corresponding formal one, giving rise
to invariants of \var{\deleted{their}}{}{deletion}{ps_134-7_054-v0003}\var{\inserted{it's}}{\hand{it's}}{insertion}{ps_134-7_054-v0004} own, namely ``formal'' homology and cohomology,
which are definitely distinct from the former.

\var{\deleted{\xout{It is all the more remarkable that}}}{}{deletion}{ps_134-7_054-v0005}
118. It is all the more remarkable, in view of the preceding findings,
that the homotopy invariants
\[ \hand{\pi}_i(X_*) \ , \quad i \hand{\geqslant} 0 , \]
of a \var{\inserted[typed]{pointed}}{pointed}{typed insertion}{ps_134-7_054-v0006} semisimplicial \var{\inserted[typed]{unipotent}}{unipotent}{typed insertion}{ps_134-7_054-v0007} bundle $X_*$ (still assuming the components $X_n$ to be
smooth, i.e. \var{\inserted{to}}{\hand{to}}{insertion}{ps_134-7_054-v0008} correspond\var{\deleted{ing}}{}{deletion}{ps_134-7_054-v0009} to projective modules) turn out to be invari-%
ants of the corresponding ``formal'' object, and, more startling still,
of the corresponding ``infinitesimal'' object of order 1. M\var{\over[typed]{g}{o}}{o}{typewriter overwriting}{ps_134-7_054-v0010}re specifical-%
ly, consider the ``Lie functor'' or ``tangent space at the origin''
\[ (7) \qquad \mathrm{Lie}: U(k)^{\hand{\bullet}} \hand{\longrightarrow} Ab_k \ , \qquad X = W(M) \hand{\longmapsto} M \ , \]
which we'll need only for the time being for smooth $X$, when the geometri
meaning of it is clear. This functor transforms semisimplicial \var{\inserted[typed]{pointed}}{pointed}{typed insertion}{ps_134-7_054-v0011} bundles
into semisimplicial $k$-modules $\mathrm{Lie}(X_*)$ , thus we should get, besides
abelianization, another remarkable functor\hand{,} from $\hand{\mathcal{M}_1}(k)$ \var{\inserted{say}}{\hand{say}}{insertion}{ps_134-7_054-v0012} to $D.(Ab_k)$ :
\[ (8) \qquad \mathrm{Lie}: \hand{\mathcal{M}_1}(k) \hand{\longrightarrow} D.(Ab_k) \ \var{\inserted[typed]{via}}{via}{typed insertion}{ps_134-7_054-v0013}, \ X_* \hand{\longmapsto} \mathrm{Lie}(X_*) \ , \]
granting that $X_* \hand{\mapsto} \mathrm{Lie}(X_*)$ transforms quasiisomorphisms into quasiiso-%
morphisms. Now that this must be so follows from the really startling
formula
\[ (9) \qquad \hand{\pi}_i(X_*) \hand{\simeq} \hand{\pi}_i(\mathrm{Lie}(X_*)) \quad (\hand{\simeq} H_i \text{ of the associated chain} \]
\[ \text{complex in } Ab_k \ ) \ , \]
where the left hand side, I recall, is defined as
\[ \hand{(10)} \qquad \hand{\pi}_i(X_*) = \hand{\pi}_i(X_*(k)) \ . \]

\margin{54}

% Visual collation: both supplied image halves inspected; see CHECK_DECISIONS_134_7_8.tsv.

%% END SOURCE ps_134-7_054
\ifnum\EditionMode=1
\psapparatusstart
\psentry{ps_134-7_054-v0001}{1}{typewriter deletion}{\xout{t}}{}{The typewriter strikes out this material; it leaves the corrected running text.}
\psentry{ps_134-7_054-v0002}{2}{typewriter deletion}{\xout{define}}{}{The typewriter strikes out this material; it leaves the corrected running text.}
\psentry{ps_134-7_054-v0003}{3}{deletion}{\deleted{their}}{}{The author cancels this material by hand; it leaves the corrected running text.}
\psentry{ps_134-7_054-v0004}{4}{insertion}{\inserted{it's}}{\hand{it's}}{The handwritten insertion is incorporated at its indicated position; insertion brackets leave the running text.}
\psentry{ps_134-7_054-v0005}{5}{deletion}{\deleted{\xout{It is all the more remarkable that}}}{}{The author cancels this material by hand; it leaves the corrected running text.}
\psentry{ps_134-7_054-v0006}{6}{typed insertion}{\inserted[typed]{pointed}}{pointed}{The typed insertion is incorporated at its indicated position; insertion brackets leave the running text.}
\psentry{ps_134-7_054-v0007}{7}{typed insertion}{\inserted[typed]{unipotent}}{unipotent}{The typed insertion is incorporated at its indicated position; insertion brackets leave the running text.}
\psentry{ps_134-7_054-v0008}{8}{insertion}{\inserted{to}}{\hand{to}}{The handwritten insertion is incorporated at its indicated position; insertion brackets leave the running text.}
\psentry{ps_134-7_054-v0009}{9}{deletion}{\deleted{ing}}{}{The author cancels this material by hand; it leaves the corrected running text.}
\psentry{ps_134-7_054-v0010}{10}{typewriter overwriting}{\over[typed]{g}{o}}{o}{The superposed typed reading replaces the underlying reading.}
\psentry{ps_134-7_054-v0011}{11}{typed insertion}{\inserted[typed]{pointed}}{pointed}{The typed insertion is incorporated at its indicated position; insertion brackets leave the running text.}
\psentry{ps_134-7_054-v0012}{12}{insertion}{\inserted{say}}{\hand{say}}{The handwritten insertion is incorporated at its indicated position; insertion brackets leave the running text.}
\psentry{ps_134-7_054-v0013}{13}{typed insertion}{$\inserted[typed]{via}$}{$via$}{The typed insertion is incorporated at its indicated position; insertion brackets leave the running text.}
\psapparatusend
\fi
\pssource{ps_134-7_055}{134-7/55}{https://grothendieck.umontpellier.fr/web/viewer.html?file=../134-7.pdf\#page=56}
\psnumber{ps_134-7_055-v0001}{1}
\psnumber{ps_134-7_055-v0002}{2}
\psnumber{ps_134-7_055-v0003}{3}
\psnumber{ps_134-7_055-v0004}{4}
\psnumber{ps_134-7_055-v0005}{5}
\psnumber{ps_134-7_055-v0006}{6}
\psnumber{ps_134-7_055-v0007}{7}
\psnumber{ps_134-7_055-v0008}{8}
\psnumber{ps_134-7_055-v0009}{9}
\psnumber{ps_134-7_055-v0010}{10}
\psnumber{ps_134-7_055-v0011}{11}
\psnumber{ps_134-7_055-v0012}{12}
\psnumber{ps_134-7_055-v0013}{13}
\psnumber{ps_134-7_055-v0014}{14}
\psnumber{ps_134-7_055-v0015}{15}
%% BEGIN SOURCE ps_134-7_055
% Pursuing Stacks (cotes 134-2 to 134-8): archive cote 134-7, site page 55 (PDF page 56), typescript page 497.
% typescript. Draft compared: ps6.tex:2719-2783.
% AI-transcribed by Claude Opus 5.5 (claude-opus-5-5, adaptive thinking, effort high), Message Batches API, batch ps4 (msgbatch_012Aoc1NuskATPsyBeKSYK7J), 2026-11-12; reviewed against both supplied scan halves; see CHECK_DECISIONS.tsv.
\scanpage{134-7/55}
\mspage{497}

I do'nt have,  I must confess,  any direct description of \var{\inserted{such}}{\hand{such}}{insertion}{ps_134-7_055-v0001} an isomorphism
(9),  valid for any semisimplicial bundle $X_{*}$  say satisfying the assumti-%
ons
\[ X_o=X_l= e \ , \quad X_n \text{ smooth for any n } \hand{\unsure{,}}\text{(maybe flat is enough) ,} \]
plus possibly (if needed) a Kan type condition.  However,  we have
such isomorphisms (9) in a tautological way,  when $X_{*}$  comes \var{\xout{\unsure{fxxix}}}{}{typewriter deletion}{ps_134-7_055-v0002} in the
usual way from a chain complex in Ab$_k$  with projective components,  hence
also when $X_{*}$  admits a Postnikov-type dévissage into ``abelian'' pieces as
above.  If we admit that any $X_{*}$  satisfying the assumptions is homotopic
to one admitting such a dévissage,  the isomorphisms (9) should follow,
except of course that extra work would be still needed to get naturality
of (9).

I have the feeling however that,  besides the specific abelianiza-%
tion functor in the schematic context,  \uline{formula} (9) \uline{should be} \var{\xout{viewed as}}{}{typewriter deletion}{ps_134-7_055-v0003} \var{\inserted[typed]{made a cor-}}{made a cor-}{typed insertion}{ps_134-7_055-v0004}
\var{\xout{being ``the'' natural \uline{definition}}}{}{typewriter deletion}{ps_134-7_055-v0005} \var{\inserted[typed]{\uline{nerstone of a theory of schematic}}}{\uline{nerstone of a theory of schematic}}{typed insertion}{ps_134-7_055-v0006}  \uline{homotopy types},  \uline{and serve as}``\uline{the}''
\uline{natural definition of the homotopy invariants} of a model $X_{*}$ ,  within the
context of schematic models and without any need a priori to \var{\xout{subordinate}}{}{typewriter deletion}{ps_134-7_055-v0007}
\var{\xout{these}}{}{typewriter deletion}{ps_134-7_055-v0008} tie them up with,  let alone subordinate their study\hand{\unsure{,}} to,  \var{\xout{corres-}}{}{typewriter deletion}{ps_134-7_055-v0009}
\var{\xout{ponding discr homotopy}}{}{typewriter deletion}{ps_134-7_055-v0010}invariants of  the corresponding discrete homotopy
type $X_{*}(k)$.  Accordingly,  \var{\xout{quasi}}{}{typewriter deletion}{ps_134-7_055-v0011} weak equivalences should be defined \var{\xout{as}}{}{typewriter deletion}{ps_134-7_055-v0012}
\var{\xout{maps}}{}{typewriter deletion}{ps_134-7_055-v0013} (for semisimplicial bundles satisfying the suitable assumptions at
any rate)  as maps inducing isomorphisms for the $\hand{\pi}_i$ invariants,  namely
inducing quasi-isomorphisms for the corresponding Lie chain complexes.
It is immediately checked that this implies that the corresponding map
for ``formal homology'' namely
\[ L_{pt}(X_{*}) \hand{\to} L_{pt}(Y_{*}) \]
is then a quasi-isomorphism too,  when viewed as a map of chain complexes,
and hopefully the same should hold for ``schematic homology''
\[ L(X_{*}) \hand{\to} L(Y_{*}) \ \ , \]
and of course one would expect the converse statements to hold too.

I would like to comment a little on the significance of formula (9).
\var{\xout{To my k}}{}{typewriter deletion}{ps_134-7_055-v0014} As far as I know,  this is the only fairly general formula,  not
reducing to an ``abelian'' case,  where the homotopy groups $\hand{\pi}_i$  appear
as just the $H_i$  invariants of a suitable chain complex,  defined up to
unique isomorphism in the relevant derived category.  This chain complex
comes here,  moreover,  with an amazingly simple description,  of immedia-%
te geometrical significance,  and suggestive of \var{\xout{very p}}{}{typewriter deletion}{ps_134-7_055-v0015} relationships with

\margin{55}

% Visual collation: both supplied image halves inspected; see CHECK_DECISIONS_134_7_8.tsv.

%% END SOURCE ps_134-7_055
\ifnum\EditionMode=1
\psapparatusstart
\psentry{ps_134-7_055-v0001}{1}{insertion}{\inserted{such}}{\hand{such}}{The handwritten insertion is incorporated at its indicated position; insertion brackets leave the running text.}
\psentry{ps_134-7_055-v0002}{2}{typewriter deletion}{\xout{\unsure{fxxix}}}{}{The typewriter strikes out this material; it leaves the corrected running text.}
\psentry{ps_134-7_055-v0003}{3}{typewriter deletion}{\xout{viewed as}}{}{The typewriter strikes out this material; it leaves the corrected running text.}
\psentry{ps_134-7_055-v0004}{4}{typed insertion}{\inserted[typed]{made a cor-}}{made a cor-}{The typed insertion is incorporated at its indicated position; insertion brackets leave the running text.}
\psentry{ps_134-7_055-v0005}{5}{typewriter deletion}{\xout{being ``the'' natural \uline{definition}}}{}{The typewriter strikes out this material; it leaves the corrected running text.}
\psentry{ps_134-7_055-v0006}{6}{typed insertion}{\inserted[typed]{\uline{nerstone of a theory of schematic}}}{\uline{nerstone of a theory of schematic}}{The typed insertion is incorporated at its indicated position; insertion brackets leave the running text.}
\psentry{ps_134-7_055-v0007}{7}{typewriter deletion}{\xout{subordinate}}{}{The typewriter strikes out this material; it leaves the corrected running text.}
\psentry{ps_134-7_055-v0008}{8}{typewriter deletion}{\xout{these}}{}{The typewriter strikes out this material; it leaves the corrected running text.}
\psentry{ps_134-7_055-v0009}{9}{typewriter deletion}{\xout{corres-}}{}{The typewriter strikes out this material; it leaves the corrected running text.}
\psentry{ps_134-7_055-v0010}{10}{typewriter deletion}{\xout{ponding discr homotopy}}{}{The typewriter strikes out this material; it leaves the corrected running text.}
\psentry{ps_134-7_055-v0011}{11}{typewriter deletion}{\xout{quasi}}{}{The typewriter strikes out this material; it leaves the corrected running text.}
\psentry{ps_134-7_055-v0012}{12}{typewriter deletion}{\xout{as}}{}{The typewriter strikes out this material; it leaves the corrected running text.}
\psentry{ps_134-7_055-v0013}{13}{typewriter deletion}{\xout{maps}}{}{The typewriter strikes out this material; it leaves the corrected running text.}
\psentry{ps_134-7_055-v0014}{14}{typewriter deletion}{\xout{To my k}}{}{The typewriter strikes out this material; it leaves the corrected running text.}
\psentry{ps_134-7_055-v0015}{15}{typewriter deletion}{\xout{very p}}{}{The typewriter strikes out this material; it leaves the corrected running text.}
\psapparatusend
\fi
\pssource{ps_134-7_056}{134-7/56}{https://grothendieck.umontpellier.fr/web/viewer.html?file=../134-7.pdf\#page=57}
\psnumber{ps_134-7_056-v0001}{1}
\psnumber{ps_134-7_056-v0002}{2}
\psnumber{ps_134-7_056-v0003}{3}
\psnumber{ps_134-7_056-v0004}{4}
\psnumber{ps_134-7_056-v0005}{5}
\psnumber{ps_134-7_056-v0006}{6}
\psnumber{ps_134-7_056-v0007}{7}
\psnumber{ps_134-7_056-v0008}{8}
\psnumber{ps_134-7_056-v0009}{9}
\psnumber{ps_134-7_056-v0010}{10}
\psnumber{ps_134-7_056-v0011}{11}
\psnumber{ps_134-7_056-v0012}{12}
\psnumber{ps_134-7_056-v0013}{13}
\psnumber{ps_134-7_056-v0014}{14}
\psnumber{ps_134-7_056-v0015}{15}
\psnumber{ps_134-7_056-v0016}{16}
\psnumber{ps_134-7_056-v0017}{17}
\psnumber{ps_134-7_056-v0018}{18}
\psnumber{ps_134-7_056-v0019}{19}
\psnumber{ps_134-7_056-v0020}{20}
\psnumber{ps_134-7_056-v0021}{21}
%% BEGIN SOURCE ps_134-7_056
% Pursuing Stacks (cotes 134-2 to 134-8): archive cote 134-7, site page 56 (PDF page 57), typescript page 498.
% typescript. Draft compared: ps6.tex:2771-2829.
% AI-transcribed by Claude Opus 5.5 (claude-opus-5-5, adaptive thinking, effort high), Message Batches API, batch ps4 (msgbatch_012Aoc1NuskATPsyBeKSYK7J), 2026-11-12; reviewed against both supplied scan halves; see CHECK_DECISIONS.tsv.
\scanpage{134-7/56}
\mspage{498}

the homology invariants a lot more precise\var{\inserted[typed]{presumably,}}{presumably,}{typed insertion}{ps_134-7_056-v0001} than those currently used so
far.  Of course,  the significance of (9) for the study of the usual,
``discrete'' homotopy types,  will be subordinated to haw difficult it will
turn out for such a homotopy type to be  a) realizable by a schematic
one over $\underline{Z}$,  and  b) to get hold of a\var{\xout{n explicit}}{}{typewriter deletion}{ps_134-7_056-v0002} more or less explicit
description of \var{\xout{a}}{}{typewriter deletion}{ps_134-7_056-v0003} such a schematic homotopy type,  via say a semisimplicial
bundle as a model.  One is of course thinking more specifically of the
case of the spheres $S^n$ ,  the first case (besides the trivial $S^1$ case)
being $S^2$ ,  the sequence of homotopy groups of which (as far as I know)
is not understood as yet.  Viewing the spheres $S^n$ as successive suspensi-%
ons of $S^1$ ,  which $S^1$ is fitting nicely into the formalism of schematic
homotopy types as a $K(\underline{Z},1)$  (except that the \var{\inserted[typed]{1-connectedness}}{1-connectedness}{typed insertion}{ps_134-7_056-v0004} condition $X_1=e$ \var{\xout{will}}{}{typewriter deletion}{ps_134-7_056-v0005} \var{\inserted[typed]{is}}{is}{typed insertion}{ps_134-7_056-v0006} not \var{\xout{be}}{}{typewriter deletion}{ps_134-7_056-v0007}
satisfied),  this brings near the question of defining \var{\inserted{the suspension}}{\hand{the suspension}}{insertion}{ps_134-7_056-v0008} \var{\deleted{such}}{}{deletion}{ps_134-7_056-v0009} operation\var{\deleted{s}}{}{deletion}{ps_134-7_056-v0010}
in the \var{\inserted[typed]{relevant}}{relevant}{typed insertion}{ps_134-7_056-v0011} derived category $\hand{\mathcal{M}}_o(k)$ or $\hand{\mathcal{M}}_1(k)$\unsure{,} \hand{(}whereas before we had met with
the ``dual'' question of constructing homotopy fibers of maps,  such as
loop spaces\hand{)}.  Thus it would seem that \uline{a breakthrough in getting hold of}
\uline{the standard homotopy constructions within the schematic context},  assu-%
ming that these constructions do still make a sense,  \uline{may well mean a}
\uline{significant advance} \var{\xout{in}}{}{typewriter deletion}{ps_134-7_056-v0012} \uline{for the understanding of the homotopy groups of}
\uline{spheres}.  This \var{\deleted{at the same time}}{}{deletion}{ps_134-7_056-v0013} looks like a very strong motivation for
trying to \var{\xout{develop the}}{}{typewriter deletion}{ps_134-7_056-v0014} carry through those constructions (possibly even
construct a corresponding ``derivator'' embodying any kinds of finite
``integration'' and ``cointegration'' operations on schematic homotopy types)\hand{,}
and \var{\inserted{(at the same time)}}{\hand{(at the same time)}}{insertion}{ps_134-7_056-v0015} warns us that such work,  if at all feasible,  will most probably be
a highly non trivial one.

The question of``schematizing'' the homotopy types of $S^2$ and $S^3$ reminds
me of a fact which struck me a long time ago (maybe J.P. Serre or someone
else pointed it out to me first).  Namely,  in some respects there does'nt
seem to be \var{\xout{a}}{}{typewriter deletion}{ps_134-7_056-v0016} really satisfactory algebraic models for these homotopy types\hand{,}
taking into account the basic relationship between the two,  namely:
\var{\over{t}{\unsure{T}}}{\hand{\unsure{T}}}{writing over}{ps_134-7_056-v0017}he 2-sphere (or,  equivalently,  the projective complex line)  is a
\uline{homogeneous space} under the quaternionic \uline{group} $S^3$ (or,  equivalently,
under the complex linear group $S\unsure{l}(2,\underline{\unsure{C}})$).  This \var{\xout{sit}}{}{typewriter deletion}{ps_134-7_056-v0018} relationship,  and
it's manifold ``avatars'' in the realms of discrete groups, Lie groups,
algebraic  groups or group schemes etc,  is one of the \var{\inserted[typed]{few}}{few}{typed insertion}{ps_134-7_056-v0019} key situations
met with\hand{,} and of equal basic significance\hand{,} in the most diverse quarters\var{\inserted{in mathematics,}}{\hand{in mathematics,}}{insertion}{ps_134-7_056-v0020}
from topology to arithmetics.   Thus $S\unsure{l}(2)$ \var{\xout{plays the r\^ole of \uline{the}}}{}{typewriter deletion}{ps_134-7_056-v0021} as a
simple Lie (or algebraic) group of minimum rank 1,  plays the  role of

\margin{56}

% Visual collation: both supplied image halves inspected; see CHECK_DECISIONS_134_7_8.tsv.

%% END SOURCE ps_134-7_056
\ifnum\EditionMode=1
\psapparatusstart
\psentry{ps_134-7_056-v0001}{1}{typed insertion}{\inserted[typed]{presumably,}}{presumably,}{The typed insertion is incorporated at its indicated position; insertion brackets leave the running text.}
\psentry{ps_134-7_056-v0002}{2}{typewriter deletion}{\xout{n explicit}}{}{The typewriter strikes out this material; it leaves the corrected running text.}
\psentry{ps_134-7_056-v0003}{3}{typewriter deletion}{\xout{a}}{}{The typewriter strikes out this material; it leaves the corrected running text.}
\psentry{ps_134-7_056-v0004}{4}{typed insertion}{\inserted[typed]{1-connectedness}}{1-connectedness}{The typed insertion is incorporated at its indicated position; insertion brackets leave the running text.}
\psentry{ps_134-7_056-v0005}{5}{typewriter deletion}{\xout{will}}{}{The typewriter strikes out this material; it leaves the corrected running text.}
\psentry{ps_134-7_056-v0006}{6}{typed insertion}{\inserted[typed]{is}}{is}{The typed insertion is incorporated at its indicated position; insertion brackets leave the running text.}
\psentry{ps_134-7_056-v0007}{7}{typewriter deletion}{\xout{be}}{}{The typewriter strikes out this material; it leaves the corrected running text.}
\psentry{ps_134-7_056-v0008}{8}{insertion}{\inserted{the suspension}}{\hand{the suspension}}{The handwritten insertion is incorporated at its indicated position; insertion brackets leave the running text.}
\psentry{ps_134-7_056-v0009}{9}{deletion}{\deleted{such}}{}{The author cancels this material by hand; it leaves the corrected running text.}
\psentry{ps_134-7_056-v0010}{10}{deletion}{\deleted{s}}{}{The author cancels this material by hand; it leaves the corrected running text.}
\psentry{ps_134-7_056-v0011}{11}{typed insertion}{\inserted[typed]{relevant}}{relevant}{The typed insertion is incorporated at its indicated position; insertion brackets leave the running text.}
\psentry{ps_134-7_056-v0012}{12}{typewriter deletion}{\xout{in}}{}{The typewriter strikes out this material; it leaves the corrected running text.}
\psentry{ps_134-7_056-v0013}{13}{deletion}{\deleted{at the same time}}{}{The author cancels this material by hand; it leaves the corrected running text.}
\psentry{ps_134-7_056-v0014}{14}{typewriter deletion}{\xout{develop the}}{}{The typewriter strikes out this material; it leaves the corrected running text.}
\psentry{ps_134-7_056-v0015}{15}{insertion}{\inserted{(at the same time)}}{\hand{(at the same time)}}{The handwritten insertion is incorporated at its indicated position; insertion brackets leave the running text.}
\psentry{ps_134-7_056-v0016}{16}{typewriter deletion}{\xout{a}}{}{The typewriter strikes out this material; it leaves the corrected running text.}
\psentry{ps_134-7_056-v0017}{17}{writing over}{\over{t}{\unsure{T}}}{\hand{\unsure{T}}}{The superposed handwritten reading replaces the underlying reading.}
\psentry{ps_134-7_056-v0018}{18}{typewriter deletion}{\xout{sit}}{}{The typewriter strikes out this material; it leaves the corrected running text.}
\psentry{ps_134-7_056-v0019}{19}{typed insertion}{\inserted[typed]{few}}{few}{The typed insertion is incorporated at its indicated position; insertion brackets leave the running text.}
\psentry{ps_134-7_056-v0020}{20}{insertion}{\inserted{in mathematics,}}{\hand{in mathematics,}}{The handwritten insertion is incorporated at its indicated position; insertion brackets leave the running text.}
\psentry{ps_134-7_056-v0021}{21}{typewriter deletion}{\xout{plays the r\^ole of \uline{the}}}{}{The typewriter strikes out this material; it leaves the corrected running text.}
\psapparatusend
\fi
\pssource{ps_134-7_057}{134-7/57}{https://grothendieck.umontpellier.fr/web/viewer.html?file=../134-7.pdf\#page=58}
\psnumber{ps_134-7_057-v0001}{1}
\psnumber{ps_134-7_057-v0002}{2}
\psnumber{ps_134-7_057-v0003}{3}
\psnumber{ps_134-7_057-v0004}{4}
\psnumber{ps_134-7_057-v0005}{5}
\psnumber{ps_134-7_057-v0006}{6}
\psnumber{ps_134-7_057-v0007}{7}
\psnumber{ps_134-7_057-v0008}{8}
\psnumber{ps_134-7_057-v0009}{9}
\psnumber{ps_134-7_057-v0010}{10}
\psnumber{ps_134-7_057-v0011}{11}
%% BEGIN SOURCE ps_134-7_057
% Pursuing Stacks (cotes 134-2 to 134-8): archive cote 134-7, site page 57 (PDF page 58), typescript page 499.
% typescript. Draft compared: ps6.tex:2817-2869.
% AI-transcribed by Claude Opus 5.5 (claude-opus-5-5, adaptive thinking, effort high), Message Batches API, batch ps4 (msgbatch_012Aoc1NuskATPsyBeKSYK7J), 2026-11-12; reviewed against both supplied scan halves; see CHECK_DECISIONS.tsv.
\scanpage{134-7/57}
\mspage{499}

\uline{the} basic building block for building up the most general semisimple
groups,  whereas $P^{1}_{\hand{\underline{\unsure{C}}}}$ may be viewed as being the most significant homoge-%
neous space under this group,  namely the first and most elementary case
of flag manifolds.  In view of this significance of $S^2$ and $S^3$,  it is all
the more remarkable that no simple, non plethoric semisimplicial model
say is available,  in terms say of a semisimplicial group having the
homotopy type of \var{\xout{\unsure{z}}}{}{typewriter deletion}{ps_134-7_057-v0001}
\[ S^3 \hand{\sim} S\unsure{l}(2,\underline{C} )= G \hand{,} \]
and a subgroup having the homotopy type of
\[ S^1 \hand{\sim} \underline{C}^{\hand{*}} = G\unsure{l}(1,\underline{C}) \hand{\sim} K(\underline{Z},1) \]
and playing the part of a Borel subgroup or a maximal torus,  in such a
way that the quotient will have the homotopy type of
\[ S^2 = P^1_{\underline{C}} \hand{\simeq} S^3/S^1 \hand{\simeq} G/B \quad . \]
It would be tempting now to try and construct such a model of the situa-%
tion in terms of semisimplicial unipotent bund\var{\over{\illegible{1}}{les}}{\hand{les}}{writing over}{ps_134-7_057-v0002} over $\underline{Z}$  - which would
at the same time display the homotopy groups of $S^2$ and $S^3$ (not much of a
difference\hand{!}) via formula (9).  All the more tempting of course,  as it is
felt that the geometric objects and their relationship, \var{\xout{which}}{}{typewriter deletion}{ps_134-7_057-v0003} the homotopy
shadow of which we want to modelize \var{\deleted{\illegible{1}}}{}{deletion}{ps_134-7_057-v0004} schematically,  are themselves\var{\inserted{already,}}{\hand{already,}}{insertion}{ps_134-7_057-v0005}
basically, \var{\inserted{most beautiful}}{\hand{most beautiful}}{insertion}{ps_134-7_057-v0006} schemes over Spec($\underline{Z}$)\hand{!}

\margin{!\ \rule{0.4pt}{3\baselineskip}}Another way of getting a display of the homotopy groups of $S^2$ and
$S^3$ would be in terms of a (discrete) model of the situation above,  in
terms of ``hemispherical complexes'' rather than semisimplicial ones.  On
the other hand,  there is no reason why a theory of schematic homotopy
types could not be carried through as well,  using hemispherical comple-%
xes rather than semisimplicial ones.  The latter \var{\inserted{kind of complexes}}{\hand{kind of complexes}}{insertion}{ps_134-7_057-v0007} have the advantage \var{\deleted{of}}{}{deletion}{ps_134-7_057-v0008}
that they have become \var{\xout{familiar}}{}{typewriter deletion}{ps_134-7_057-v0009}  thoroughly familiar \var{\xout{by}}{}{typewriter deletion}{ps_134-7_057-v0010} \var{\inserted[typed]{through}}{through}{typed insertion}{ps_134-7_057-v0011} constant use
by topologists and homotopy people for thirty years or so  -  the former
however and newcomers,  have the advantage of still greater formal simpli-%
city (just two boundary operations,  and just one degeneracy),  and more
important still,  of allowing for a direct computational description of
the homotopy invariants,  in the discrete set-up.  When working with
hemispherical complexes of unipotent bundles as models for schematic
homotopy types,  we'll get then \uline{two} highly differen\hand{t} descriptions of the
homotopy invariants $\hand{\pi}_i$ ,  one by the ``infinitesimal'' formula(9)  interpre-%
ting them as derived functors of the Lie functor,  the other one via the

\margin{57}

% Visual collation: both supplied image halves inspected; see CHECK_DECISIONS_134_7_8.tsv.

%% END SOURCE ps_134-7_057
\ifnum\EditionMode=1
\psapparatusstart
\psentry{ps_134-7_057-v0001}{1}{typewriter deletion}{\xout{\unsure{z}}}{}{The typewriter strikes out this material; it leaves the corrected running text.}
\psentry{ps_134-7_057-v0002}{2}{writing over}{\over{\illegible{1}}{les}}{\hand{les}}{The superposed handwritten reading replaces the underlying reading.}
\psentry{ps_134-7_057-v0003}{3}{typewriter deletion}{\xout{which}}{}{The typewriter strikes out this material; it leaves the corrected running text.}
\psentry{ps_134-7_057-v0004}{4}{deletion}{\deleted{\illegible{1}}}{}{The author cancels this material by hand; it leaves the corrected running text.}
\psentry{ps_134-7_057-v0005}{5}{insertion}{\inserted{already,}}{\hand{already,}}{The handwritten insertion is incorporated at its indicated position; insertion brackets leave the running text.}
\psentry{ps_134-7_057-v0006}{6}{insertion}{\inserted{most beautiful}}{\hand{most beautiful}}{The handwritten insertion is incorporated at its indicated position; insertion brackets leave the running text.}
\psentry{ps_134-7_057-v0007}{7}{insertion}{\inserted{kind of complexes}}{\hand{kind of complexes}}{The handwritten insertion is incorporated at its indicated position; insertion brackets leave the running text.}
\psentry{ps_134-7_057-v0008}{8}{deletion}{\deleted{of}}{}{The author cancels this material by hand; it leaves the corrected running text.}
\psentry{ps_134-7_057-v0009}{9}{typewriter deletion}{\xout{familiar}}{}{The typewriter strikes out this material; it leaves the corrected running text.}
\psentry{ps_134-7_057-v0010}{10}{typewriter deletion}{\xout{by}}{}{The typewriter strikes out this material; it leaves the corrected running text.}
\psentry{ps_134-7_057-v0011}{11}{typed insertion}{\inserted[typed]{through}}{through}{The typed insertion is incorporated at its indicated position; insertion brackets leave the running text.}
\psapparatusend
\fi
\pssource{ps_134-7_058}{134-7/58}{https://grothendieck.umontpellier.fr/web/viewer.html?file=../134-7.pdf\#page=59}
\psnumber{ps_134-7_058-v0001}{1}
\psnumber{ps_134-7_058-v0002}{2}
\psnumber{ps_134-7_058-v0003}{3}
\psnumber{ps_134-7_058-v0004}{4}
\psnumber{ps_134-7_058-v0005}{5}
\psnumber{ps_134-7_058-v0006}{6}
\psnumber{ps_134-7_058-v0007}{7}
\psnumber{ps_134-7_058-v0008}{8}
\psnumber{ps_134-7_058-v0009}{9}
\psnumber{ps_134-7_058-v0010}{10}
\psnumber{ps_134-7_058-v0011}{11}
\psnumber{ps_134-7_058-v0012}{12}
\psnumber{ps_134-7_058-v0013}{13}
\psnumber{ps_134-7_058-v0014}{14}
\psnumber{ps_134-7_058-v0015}{15}
\psnumber{ps_134-7_058-v0016}{16}
%% BEGIN SOURCE ps_134-7_058
% Pursuing Stacks (cotes 134-2 to 134-8): archive cote 134-7, site page 58 (PDF page 59), typescript page 500.
% typescript. Draft compared: ps6.tex:2858-2921.
% AI-transcribed by Claude Opus 5.5 (claude-opus-5-5, adaptive thinking, effort high), Message Batches API, batch ps4 (msgbatch_012Aoc1NuskATPsyBeKSYK7J), 2026-11-12; reviewed against both supplied scan halves; see CHECK_DECISIONS.tsv.
\scanpage{134-7/58}
\mspage{500}

hemispherical set $X_{\underline{*}}(k)$, handled as if it came from an actual $\hand{\infty}$-groupoid
(\var{\inserted[typed]{by}}{by}{typed insertion}{ps_134-7_058-v0001} taking it's \var{\xout{\unsure{time}}}{}{typewriter deletion}{ps_134-7_058-v0002}\var{\inserted[typed]{source}}{source}{typed insertion}{ps_134-7_058-v0003} and target operations etc) even if it does not. (I con-%
fess I did not check that this process does correctly describe the homoto\unsure{-}%
py groups of a hemispherical set\hand{,} even \var{\over{\unsure{with}}{with}}{\hand{with}}{writing over}{ps_134-7_058-v0004}out assuming it comes from
a\hand{n} $\hand{\infty}$-groupoid or an $\hand{\infty}$-stack, but never mind for the time being ...)
Among the interesting things still ahead (once we get a little \var{\xout{more}}{}{typewriter deletion}{ps_134-7_058-v0005} accus-%
tomed working with hemispherical complexes) is to try and understand how
these two descriptions relate to each other, which may be \var{\deleted{a way}}{}{deletion}{ps_134-7_058-v0006}\var{\inserted{one means}}{\hand{one means}}{insertion}{ps_134-7_058-v0007} for \var{\deleted{coming}}{}{deletion}{ps_134-7_058-v0008}
\var{\deleted{to}}{}{deletion}{ps_134-7_058-v0009} a better understanding of the basic formula (9), in the context of
hemispheric schematic mod\var{\over{\unsure{les}}{els}}{\hand{els}}{writing over}{ps_134-7_058-v0010}.


26.9.

119. After the last notes (of September 10) I was a little sick for a
few days, then I was taken by current tasks from professional and family
life, which left little leisure for mathematical reflection, except once
or twice for a couple of hours, by way of recreation. It would seem now
that in the days and weeks ahead, there will be more time to go on with
the notes, and I feel eager indeed to push ahead. Also I more or less
promised the publisher, Pierre Bérès, that a first volume would be ready
for the printer by the end of this calendar year, and I would like to
keep promise.

I still have to tie in with the reflections and happenings of the end
of last month, as I started upon with the last notes (of Sept. 10). Next
thing then to report upon is the ``coup de th\'e\hand{\^{a}}tre'' occuring through \var{\xout{my}}{}{typewriter deletion}{ps_134-7_058-v0011}
the phone call to Luc Illusie. \var{\over{\unsure{Then}}{When}}{\hand{When}}{writing over}{ps_134-7_058-v0012} I told him about what by then still
looked to me as the key assumption for a theory of schematization of homo-%
topy types, namely that the homology of $K(\hand{\pi},n)$ should be computable \hand{in}
terms of derived functors of the ``divided power algebra''-functor $\hand{\Gamma}_{\underline{Z}}$ , he
at once felt rather sceptical, and later he called me back to tell me it
was definitely false. He could not give me an explicit counterexample for
$H_{i}(n,\underline{Z};\underline{Z})$ say with given n and i, rather he said that when suitably ``sta-%
bilizing'' the assumption I had in mind, it went against result\hand{s} of Larry
Breen on $\mathrm{Ext}^{i}$ functors of $\underline{G}_{a}$ with itself over prime fields $\underline{F}_{p}$ . I do'nt
know if I am going some day to dive into Illusie's argument and into Larry
Breen's results - \var{\deleted{at any rate}}{}{deletion}{ps_134-7_058-v0013}\var{\inserted{however}}{\hand{however}}{insertion}{ps_134-7_058-v0014}, even before I got Illusie's confirmation
that definitely my assumption was wrong, I convinced myself that at any
rate it was false for $n=1$ \var{\over{,}{.}}{\hand{.}}{writing over}{ps_134-7_058-v0015} \var{\inserted{This is}}{\hand{This is}}{insertion}{ps_134-7_058-v0016} a non simply connected case and hence not

\margin{58}

% Visual collation: both supplied image halves inspected; see CHECK_DECISIONS_134_7_8.tsv.

%% END SOURCE ps_134-7_058
\ifnum\EditionMode=1
\psapparatusstart
\psentry{ps_134-7_058-v0001}{1}{typed insertion}{\inserted[typed]{by}}{by}{The typed insertion is incorporated at its indicated position; insertion brackets leave the running text.}
\psentry{ps_134-7_058-v0002}{2}{typewriter deletion}{\xout{\unsure{time}}}{}{The typewriter strikes out this material; it leaves the corrected running text.}
\psentry{ps_134-7_058-v0003}{3}{typed insertion}{\inserted[typed]{source}}{source}{The typed insertion is incorporated at its indicated position; insertion brackets leave the running text.}
\psentry{ps_134-7_058-v0004}{4}{writing over}{\over{\unsure{with}}{with}}{\hand{with}}{The superposed handwritten reading replaces the underlying reading.}
\psentry{ps_134-7_058-v0005}{5}{typewriter deletion}{\xout{more}}{}{The typewriter strikes out this material; it leaves the corrected running text.}
\psentry{ps_134-7_058-v0006}{6}{deletion}{\deleted{a way}}{}{The author cancels this material by hand; it leaves the corrected running text.}
\psentry{ps_134-7_058-v0007}{7}{insertion}{\inserted{one means}}{\hand{one means}}{The handwritten insertion is incorporated at its indicated position; insertion brackets leave the running text.}
\psentry{ps_134-7_058-v0008}{8}{deletion}{\deleted{coming}}{}{The author cancels this material by hand; it leaves the corrected running text.}
\psentry{ps_134-7_058-v0009}{9}{deletion}{\deleted{to}}{}{The author cancels this material by hand; it leaves the corrected running text.}
\psentry{ps_134-7_058-v0010}{10}{writing over}{\over{\unsure{les}}{els}}{\hand{els}}{The superposed handwritten reading replaces the underlying reading.}
\psentry{ps_134-7_058-v0011}{11}{typewriter deletion}{\xout{my}}{}{The typewriter strikes out this material; it leaves the corrected running text.}
\psentry{ps_134-7_058-v0012}{12}{writing over}{\over{\unsure{Then}}{When}}{\hand{When}}{The superposed handwritten reading replaces the underlying reading.}
\psentry{ps_134-7_058-v0013}{13}{deletion}{\deleted{at any rate}}{}{The author cancels this material by hand; it leaves the corrected running text.}
\psentry{ps_134-7_058-v0014}{14}{insertion}{\inserted{however}}{\hand{however}}{The handwritten insertion is incorporated at its indicated position; insertion brackets leave the running text.}
\psentry{ps_134-7_058-v0015}{15}{writing over}{\over{,}{.}}{\hand{.}}{The superposed handwritten reading replaces the underlying reading.}
\psentry{ps_134-7_058-v0016}{16}{insertion}{\inserted{This is}}{\hand{This is}}{The handwritten insertion is incorporated at its indicated position; insertion brackets leave the running text.}
\psapparatusend
\fi
\pssource{ps_134-7_059}{134-7/59}{https://grothendieck.umontpellier.fr/web/viewer.html?file=../134-7.pdf\#page=60}
\psnumber{ps_134-7_059-v0001}{1}
\psnumber{ps_134-7_059-v0002}{2}
\psnumber{ps_134-7_059-v0003}{3}
\psnumber{ps_134-7_059-v0004}{4}
\psnumber{ps_134-7_059-v0005}{5}
\psnumber{ps_134-7_059-v0006}{6}
\psnumber{ps_134-7_059-v0007}{7}
\psnumber{ps_134-7_059-v0008}{8}
\psnumber{ps_134-7_059-v0009}{9}
\psnumber{ps_134-7_059-v0010}{10}
\psnumber{ps_134-7_059-v0011}{11}
\psnumber{ps_134-7_059-v0012}{12}
\psnumber{ps_134-7_059-v0013}{13}
%% BEGIN SOURCE ps_134-7_059
% Pursuing Stacks (cotes 134-2 to 134-8): archive cote 134-7, site page 59 (PDF page 60), typescript page 501.
% typescript. Draft compared: ps6.tex:2907-2961.
% AI-transcribed by Claude Opus 5.5 (claude-opus-5-5, adaptive thinking, effort high), Message Batches API, batch ps4 (msgbatch_012Aoc1NuskATPsyBeKSYK7J), 2026-11-12; reviewed against both supplied scan halves; see CHECK_DECISIONS.tsv.
\scanpage{134-7/59}
\mspage{501}
entirely conclusive maybe,  but still it was enough to shake my confidence
that the assumption was OK.  T\var{\over[typed]{h}{h}}{h}{typewriter overwriting}{ps_134-7_059-v0001}e counterexample is in terms of cohomology
\[ H^2(1,\underline{Z};\ \underline{Z}) = H^2(\underline{Z},\underline{Z}) = H^2(B_{\underline{Z}},\ \underline{Z}) \]
rather than homology, as usual.   As the classifying space $B_{\underline{Z}} = S^1$ of $\underline{Z}$
is one dimension\hand{al,} it's cohomology is zero in dimensions $i \hand{\geq} 2$.  On the \var{\xout{oth}}{}{typewriter deletion}{ps_134-7_059-v0002}
other hand,  $H^2$ classifies central extensions of $\underline{Z}$ by $\underline{Z}$ ,  and an\var{\deleted{direct}}{}{deletion}{ps_134-7_059-v0003}
immediate direct argument shows \var{\deleted{too}}{}{deletion}{ps_134-7_059-v0004}\var{\inserted{indeed}}{\hand{indeed}}{insertion}{ps_134-7_059-v0005} that such extensions (indeed, any exten%
\hand{-}sion of $\underline{Z}$ by any group) split.  If we take the schematic $H^2$,  defined by
Hochschild cochains which are polynomial functions,  we get the classifica-%
tion of central extensions of $\underline{G}_a$ \var{\deleted{\xout{\illegible{1}}}}{}{deletion}{ps_134-7_059-v0006}by $\underline{G}_a$ \var{\xout{\illegible{1}}}{}{typewriter deletion}{ps_134-7_059-v0007} ,  as group schemes over
the ring of integers.  Now it is e\hand{a}sy to find such an extension which does
\uline{not} split,  the first one one may think of being the group scheme represen-%
ting the functor
\[ k \hand{\mapsto} W_2(k) = \begin{array}[t]{l} \var{\xout{\text{ring}}}{}{typewriter deletion}{ps_134-7_059-v0008}\var{\inserted[typed]{\text{group}}}{\text{group}}{typed insertion}{ps_134-7_059-v0009}\text{ of truncated power series }\ 1+at+bt^2 \\ \text{in }\ k\hand{[}t\hand{]}/(\hand{t^3}), \end{array} \]
where a,b are parameters in the commutative ring k,  the group structure
being multiplication.  These parameters define in an evident way \var{\over[typed]{\unsure{a}}{a}}{a}{typewriter overwriting}{ps_134-7_059-v0010} struc-%
ture of an extension of $\underline{G}_a$ by $\underline{G}_a$ upon $W_2$ ,  a splitting of which would
correspond to a group homomorphism
\[ \underline{G}_a \hand{\longrightarrow} W_2 \ , \quad a \hand{\longmapsto} 1 + at + P(a)t^2 \quad , \]
with
\[ P \hand{\in}\ Z\hand{[}t\hand{]} \ . \]
Expressing compatibility with the group laws gives the condition
\[ P(a+a') = P(a) + P(a') + aa' \quad , \]
which has\hand{,} as unique solutions  in  $\underline{Q}$\hand{$[$}t\hand{$]$}\hand{,} expressions
\[ P(t) = t(t-1)/2 \ + \ ct \]
with c in $\underline{Q}$ ,  none of which has coefficients in $\underline{Z}$ . \var{\deleted{\unsure{(}}}{}{deletion}{ps_134-7_059-v0011}This argument shows
in fact that for given ring k, $(W_2)_k$ is a split extension iff 2 is inver-%
tible in k,  in which case a splitting is given by
\[ a \hand{\longmapsto} 1 + at +(a(a-1)/2)t \quad . \]

This example brings near one plausible ``reason'' why the expected
comparison statement about discrete and schematic linearization could not
reasonably hold true,  and in particular\var{\inserted{why}}{\hand{why}}{insertion}{ps_134-7_059-v0012} we should'nt expect discrete and
schematic Hochs\var{\over[typed]{\unsure{c}}{ch}}{ch}{typewriter overwriting}{ps_134-7_059-v0013}ild cohomology (for group schemes over $\underline{Z}$  such as $\underline{G}_a$

\margin{59}

% Visual collation: both supplied image halves inspected; see CHECK_DECISIONS_134_7_8.tsv.

%% END SOURCE ps_134-7_059
\ifnum\EditionMode=1
\psapparatusstart
\psentry{ps_134-7_059-v0001}{1}{typewriter overwriting}{\over[typed]{h}{h}}{h}{The superposed typed reading replaces the underlying reading.}
\psentry{ps_134-7_059-v0002}{2}{typewriter deletion}{\xout{oth}}{}{The typewriter strikes out this material; it leaves the corrected running text.}
\psentry{ps_134-7_059-v0003}{3}{deletion}{\deleted{direct}}{}{The author cancels this material by hand; it leaves the corrected running text.}
\psentry{ps_134-7_059-v0004}{4}{deletion}{\deleted{too}}{}{The author cancels this material by hand; it leaves the corrected running text.}
\psentry{ps_134-7_059-v0005}{5}{insertion}{\inserted{indeed}}{\hand{indeed}}{The handwritten insertion is incorporated at its indicated position; insertion brackets leave the running text.}
\psentry{ps_134-7_059-v0006}{6}{deletion}{\deleted{\xout{\illegible{1}}}}{}{The author cancels this material by hand; it leaves the corrected running text.}
\psentry{ps_134-7_059-v0007}{7}{typewriter deletion}{\xout{\illegible{1}}}{}{The typewriter strikes out this material; it leaves the corrected running text.}
\psentry{ps_134-7_059-v0008}{8}{typewriter deletion}{$\xout{\text{ring}}$}{}{The typewriter strikes out this material; it leaves the corrected running text.}
\psentry{ps_134-7_059-v0009}{9}{typed insertion}{$\inserted[typed]{\text{group}}$}{$\text{group}$}{The typed insertion is incorporated at its indicated position; insertion brackets leave the running text.}
\psentry{ps_134-7_059-v0010}{10}{typewriter overwriting}{\over[typed]{\unsure{a}}{a}}{a}{The superposed typed reading replaces the underlying reading.}
\psentry{ps_134-7_059-v0011}{11}{deletion}{\deleted{\unsure{(}}}{}{The author cancels this material by hand; it leaves the corrected running text.}
\psentry{ps_134-7_059-v0012}{12}{insertion}{\inserted{why}}{\hand{why}}{The handwritten insertion is incorporated at its indicated position; insertion brackets leave the running text.}
\psentry{ps_134-7_059-v0013}{13}{typewriter overwriting}{\over[typed]{\unsure{c}}{ch}}{ch}{The superposed typed reading replaces the underlying reading.}
\psapparatusend
\fi
\pssource{ps_134-7_060}{134-7/60}{https://grothendieck.umontpellier.fr/web/viewer.html?file=../134-7.pdf\#page=61}
\psnumber{ps_134-7_060-v0001}{1}
\psnumber{ps_134-7_060-v0002}{2}
\psnumber{ps_134-7_060-v0003}{3}
\psnumber{ps_134-7_060-v0004}{4}
\psnumber{ps_134-7_060-v0005}{5}
\psnumber{ps_134-7_060-v0006}{6}
\psnumber{ps_134-7_060-v0007}{7}
\psnumber{ps_134-7_060-v0008}{8}
\psnumber{ps_134-7_060-v0009}{9}
\psnumber{ps_134-7_060-v0010}{10}
\psnumber{ps_134-7_060-v0011}{11}
\psnumber{ps_134-7_060-v0012}{12}
%% BEGIN SOURCE ps_134-7_060
% Pursuing Stacks (cotes 134-2 to 134-8): archive cote 134-7, site page 60 (PDF page 61), typescript page 502.
% typescript. Draft compared: ps6.tex:2948-3007.
% AI-transcribed by Claude Opus 5.5 (claude-opus-5-5, adaptive thinking, effort high), Message Batches API, batch ps4 (msgbatch_012Aoc1NuskATPsyBeKSYK7J), 2026-11-12; reviewed against both supplied scan halves; see CHECK_DECISIONS.tsv.
\scanpage{134-7/60}
\mspage{502}

or successive extensions of such) to give the same result.  Namely \var{\deleted{the}}{}{deletion}{ps_134-7_060-v0001}
the latter is computed in terms of cochains which are polynomial functi\var{\over{on}{on}}{\hand{on}}{writing over}{ps_134-7_060-v0002}s
\uline{with coefficients in $\underline{Z}$} ,  whereas there exist polynomial functtions \uline{with}
\uline{coefficients in $\underline{Q}$} \var{\inserted{(not in $\underline{Z}$)}}{\hand{(not in $\underline{Z}$)}}{insertion}{ps_134-7_060-v0003} which\hand{,} however\hand{,} give rise to integer-valued functions on
the group of integer-valued points\var{\over{,}{.}}{\hand{.}}{writing over}{ps_134-7_060-v0004} \var{\deleted{and notably}}{}{deletion}{ps_134-7_060-v0005} \var{\inserted{Such are}}{\hand{Such are}}{insertion}{ps_134-7_060-v0006} the binomial expressions
\[ P_n(t) = t(t-1)\dots(t-n+1)/n! \qquad \hand{(for n \in \underline{N}).} \]

These (in the case of just one variable t) are known to form a basis of
the $\underline{Z}$-module of all integer-valued functions on $\underline{Z}$ ,  and there is a corres-%
ponding basis for integer-valued functions on $\underline{Z}^r$\hand{,} for any natural integer
r.  Thus the hope still remains that a sweeping comparison theorem for
discrete versus ``schematic'' \var{\xout{(\illegible{1})}}{}{typewriter deletion}{ps_134-7_060-v0007} linearization might hold true,
provided it is expressed in such a way that the ``schematic models'' we are
working with should be built up with ``schemes'' (of sorts) described in
terms of spectra not of polynomial algebra\unsure{s}  $\underline{Z}\hand{[}t\hand{]}$ and tensor powers of
these,  but rather of  ``binomial algebra\unsure{s}''  $\underline{Z}\hand{\langle t\rangle}$ built up with the binomial
expressions above,  and tensor powers  of such.  If we want to develop a
corresponding notion of homotopy types over a general gound ring k,  we
should then require upon k an extra structure of a ``binomial ring'' (as
introduced in the Riemann-Roch Seminar SGA 6 in some talks of Berthelot),
nam\hand{e}ly a ring endowed with operations
\[ x \hand{\mapsto} \binom{x}{n} \ : \ k \hand{\to} k \qquad (n \hand{\in} \hand{\underline{N}}) \ , \]
satisfying the formal properties of the binomial functions  $x \hand{\mapsto} P_n(x)$
in the case $k=\underline{Z}$ or $\underline{Q}$ .  Whereas linearization of homotopy types via De
Rham complexes with devided powers relies on a ``commutative algebra with
d\var{\over{\unsure{e}}{e}}{\hand{e}}{writing over}{ps_134-7_060-v0008}vided powers'' \var{\xout{and corresponding}}{}{typewriter deletion}{ps_134-7_060-v0009} \hand{(which} was developped extensively by Berthe-%
lot and others for the needs of crystalline cohomology\hand{)}\hand{,}  linearization via
unipotent bundles (assuming it can be done in such a way as to insure that
any discrete homotopy type can be``schematized''in an essentially unique way\hand{)}
might well rely on the development of a ``binomial commutative algebra''
and a corresponding notion of ``binomial schemes''.  There should be a lot
of fun ahead developping the necessary algebraic machinery,  which may
prove of interest in it's own right.  It should be realized however that \margin{See comments next section (p.506-507)}
for a ring k to admit a binomial structure is a rather strong restriction
\var{\xout{\unsure{such}}}{}{typewriter deletion}{ps_134-7_060-v0010}  -  thus for a given prime p,  no field of char.p (except possibly
the prime field ?) admits such a structure.  This remark may temper some-%
what the enthousiasm for pushing in this direction, \var{\inserted{even}}{\hand{even}}{insertion}{ps_134-7_060-v0011} granting that a
``binomial comparison theorem'' for di\var{\deleted{\unsure{t|e}}}{}{deletion}{ps_134-7_060-v0012}screte versus ``binomial'' linearizati-%
on holds true.

\margin{60}

% Visual collation: both supplied image halves inspected; see CHECK_DECISIONS_134_7_8.tsv.

%% END SOURCE ps_134-7_060
\ifnum\EditionMode=1
\psapparatusstart
\psentry{ps_134-7_060-v0001}{1}{deletion}{\deleted{the}}{}{The author cancels this material by hand; it leaves the corrected running text.}
\psentry{ps_134-7_060-v0002}{2}{writing over}{\over{on}{on}}{\hand{on}}{The superposed handwritten reading replaces the underlying reading.}
\psentry{ps_134-7_060-v0003}{3}{insertion}{\inserted{(not in $\underline{Z}$)}}{\hand{(not in $\underline{Z}$)}}{The handwritten insertion is incorporated at its indicated position; insertion brackets leave the running text.}
\psentry{ps_134-7_060-v0004}{4}{writing over}{\over{,}{.}}{\hand{.}}{The superposed handwritten reading replaces the underlying reading.}
\psentry{ps_134-7_060-v0005}{5}{deletion}{\deleted{and notably}}{}{The author cancels this material by hand; it leaves the corrected running text.}
\psentry{ps_134-7_060-v0006}{6}{insertion}{\inserted{Such are}}{\hand{Such are}}{The handwritten insertion is incorporated at its indicated position; insertion brackets leave the running text.}
\psentry{ps_134-7_060-v0007}{7}{typewriter deletion}{\xout{(\illegible{1})}}{}{The typewriter strikes out this material; it leaves the corrected running text.}
\psentry{ps_134-7_060-v0008}{8}{writing over}{\over{\unsure{e}}{e}}{\hand{e}}{The superposed handwritten reading replaces the underlying reading.}
\psentry{ps_134-7_060-v0009}{9}{typewriter deletion}{\xout{and corresponding}}{}{The typewriter strikes out this material; it leaves the corrected running text.}
\psentry{ps_134-7_060-v0010}{10}{typewriter deletion}{\xout{\unsure{such}}}{}{The typewriter strikes out this material; it leaves the corrected running text.}
\psentry{ps_134-7_060-v0011}{11}{insertion}{\inserted{even}}{\hand{even}}{The handwritten insertion is incorporated at its indicated position; insertion brackets leave the running text.}
\psentry{ps_134-7_060-v0012}{12}{deletion}{\deleted{\unsure{t|e}}}{}{The author cancels this material by hand; it leaves the corrected running text.}
\psapparatusend
\fi
\pssource{ps_134-7_061}{134-7/61}{https://grothendieck.umontpellier.fr/web/viewer.html?file=../134-7.pdf\#page=62}
\psnumber{ps_134-7_061-v0001}{1}
\psnumber{ps_134-7_061-v0002}{2}
\psnumber{ps_134-7_061-v0003}{3}
\psnumber{ps_134-7_061-v0004}{4}
\psnumber{ps_134-7_061-v0005}{5}
\psnumber{ps_134-7_061-v0006}{6}
\psnumber{ps_134-7_061-v0007}{7}
\psnumber{ps_134-7_061-v0008}{8}
\psnumber{ps_134-7_061-v0009}{9}
\psnumber{ps_134-7_061-v0010}{10}
\psnumber{ps_134-7_061-v0011}{11}
\psnumber{ps_134-7_061-v0012}{12}
\psnumber{ps_134-7_061-v0013}{13}
\psnumber{ps_134-7_061-v0014}{14}
\psnumber{ps_134-7_061-v0015}{15}
%% BEGIN SOURCE ps_134-7_061
% Pursuing Stacks (cotes 134-2 to 134-8): archive cote 134-7, site page 61 (PDF page 62), typescript page 503.
% typescript. Draft compared: ps6.tex:2992-3051.
% AI-transcribed by Claude Opus 5.5 (claude-opus-5-5, adaptive thinking, effort high), Message Batches API, batch ps4 (msgbatch_012Aoc1NuskATPsyBeKSYK7J), 2026-11-12; reviewed against both supplied scan halves; see CHECK_DECISIONS.tsv.
\scanpage{134-7/61}
\mspage{503}

Maybe it is worthwhile still to give a down-to-earth formulation of
such a comparison statement. For any free $\underline{Z}$-module $M$\var{\inserted[typed]{of finite type}}{of finite type}{typed insertion}{ps_134-7_061-v0001}, let
\[ \mathrm{Symbin}^{\hand{\unsure{*}}}_{\underline{Z}}(M) \hand{\subset} \mathrm{Sym}^{*}_{\underline{Q}}(M_{\underline{Q}}) \qquad \text{(where } M_{\underline{Q}} = M\otimes_{\underline{Z}}\underline{Q} \text{ )} \]
be the subalgebra of the algebra of polynomial functions on $M_{\underline{Q}}^{\hand{\vee}}\hand{\simeq}(M^{\hand{\vee}})_{\underline{Q}}$
which are integgal - valued on $M^{\hand{\vee}}$ = dual module of M. Now let $L_{*}$
be any semisimplicial $\underline{Z}$-module whose components are free of finite type,
and consider the canonical map of cosemisimplicial $\underline{Z}$-modules
\[ \mathrm{Symbin}_{\underline{Z}}(L_{*}^{\hand{\vee}}) \hand{\to} \mathrm{Maps}(L_{*}, \underline{Z}) \quad , \]
described componentwise in an obvious way. The question is whether this
is a weak equivalence i.e. induces a \var{\xout{chain equiva}}{}{typewriter deletion}{ps_134-7_061-v0002} quasi-isomorphism for
the associated cochain complexes, under the extra assumption that $L_{*}$ is
0-connected i.e. the associated chain complex has zero $H_{0}$ (and possibly,
if needed, assuming even 1-connectedness i.e. $H_{0}$ \uline{and} $H_{1}$ of the associated
chain complex are \var{\inserted{both}}{\hand{both}}{insertion}{ps_134-7_061-v0003} zero). Presumably, by easy devissage arguments one
should be able to reduce to the case when $L_{*}$ is a $K(\underline{Z},n)$ type, \var{\xout{and}}{}{typewriter deletion}{ps_134-7_061-v0004} and
more specifically still, that it is the semisimplicial abelian complex
associated to the chain complex reduced to $\underline{Z}$ placed in degree n (where
$n\hand{\geq}1$). Thus, the question is whether Eilenberg-Maclane cohomology (with
coefficients in $\underline{Z}$) for $K(\hand{\pi},n)$ types (or more specifically, $K(\underline{Z},n)$ types)
can be \var{\xout{computed}}{}{typewriter deletion}{ps_134-7_061-v0005}\var{\inserted[typed]{expressed}}{expressed}{typed insertion}{ps_134-7_061-v0006} in terms of derived functors of the $\mathrm{Symbin}(M^{\hand{\vee}})$ functor .
At any rate, whether Symbin is just the right functor to fit in or not,
it looks like an interesting question whether Eilenberg-Maclane cohomology
(or, more relevant still, homology) can be expressed in ter\var{\over{\unsure{n}}{m}}{\hand{m}}{writing over}{ps_134-7_061-v0007}s of the deri-%
ved functor of a suitable non-additive contravariant (resp. covariant)
functor\var{\inserted[typed]{B}}{B}{typed insertion}{ps_134-7_061-v0008} from (Ab) (or from the subcategory of free $\underline{Z}$-modules) to itself.
If so, there should be a \var{\xout{(possibly somewhat sophisticated)}}{}{typewriter deletion}{ps_134-7_061-v0009} way of defi-%
ning a (possibly somewhat sophisticated) notion of ``B-schematic homotopy
types'' (over a ground ring k endowed with suitable extra structure, such
as a binomial structure ), \var{\inserted{in terms of ``unipotent B-bundles'',}}{\hand{in terms of ``unipotent B-bundles'',}}{insertion}{ps_134-7_061-v0010} in such a way that any \var{\xout{\illegible{1}}}{}{typewriter deletion}{ps_134-7_061-v0011} ``discrete'' namely
usual homotopy type, satisfying a suitable 1-connectedness restriction,
admits an essentially unique ``B-schematization''.

I do'nt feel like pursuing these questions here, which would take me
too far off the main line of investigation I've been out for. At any rate,
whether or not Eilenberg-Maclane \var{\xout{co}}{}{typewriter deletion}{ps_134-7_061-v0012}homology may be expressed in terms of
the \var{\inserted{total}}{\hand{total}}{insertion}{ps_134-7_061-v0013} \var{\xout{total}}{}{typewriter deletion}{ps_134-7_061-v0014}\var{\inserted[typed]{left}}{left}{typed insertion}{ps_134-7_061-v0015} derived functor of a suitable functor from (Ab) to itself, it
would seem that the somewhat naive approach towards schematic homotopy
types we have been following, valid over an arbitrary (commutative)
ground ring k without any extra structure needed on it, is worthwhile

\margin{61}

% Visual collation: both supplied image halves inspected; see CHECK_DECISIONS_134_7_8.tsv.

%% END SOURCE ps_134-7_061
\ifnum\EditionMode=1
\psapparatusstart
\psentry{ps_134-7_061-v0001}{1}{typed insertion}{\inserted[typed]{of finite type}}{of finite type}{The typed insertion is incorporated at its indicated position; insertion brackets leave the running text.}
\psentry{ps_134-7_061-v0002}{2}{typewriter deletion}{\xout{chain equiva}}{}{The typewriter strikes out this material; it leaves the corrected running text.}
\psentry{ps_134-7_061-v0003}{3}{insertion}{\inserted{both}}{\hand{both}}{The handwritten insertion is incorporated at its indicated position; insertion brackets leave the running text.}
\psentry{ps_134-7_061-v0004}{4}{typewriter deletion}{\xout{and}}{}{The typewriter strikes out this material; it leaves the corrected running text.}
\psentry{ps_134-7_061-v0005}{5}{typewriter deletion}{\xout{computed}}{}{The typewriter strikes out this material; it leaves the corrected running text.}
\psentry{ps_134-7_061-v0006}{6}{typed insertion}{\inserted[typed]{expressed}}{expressed}{The typed insertion is incorporated at its indicated position; insertion brackets leave the running text.}
\psentry{ps_134-7_061-v0007}{7}{writing over}{\over{\unsure{n}}{m}}{\hand{m}}{The superposed handwritten reading replaces the underlying reading.}
\psentry{ps_134-7_061-v0008}{8}{typed insertion}{\inserted[typed]{B}}{B}{The typed insertion is incorporated at its indicated position; insertion brackets leave the running text.}
\psentry{ps_134-7_061-v0009}{9}{typewriter deletion}{\xout{(possibly somewhat sophisticated)}}{}{The typewriter strikes out this material; it leaves the corrected running text.}
\psentry{ps_134-7_061-v0010}{10}{insertion}{\inserted{in terms of ``unipotent B-bundles'',}}{\hand{in terms of ``unipotent B-bundles'',}}{The handwritten insertion is incorporated at its indicated position; insertion brackets leave the running text.}
\psentry{ps_134-7_061-v0011}{11}{typewriter deletion}{\xout{\illegible{1}}}{}{The typewriter strikes out this material; it leaves the corrected running text.}
\psentry{ps_134-7_061-v0012}{12}{typewriter deletion}{\xout{co}}{}{The typewriter strikes out this material; it leaves the corrected running text.}
\psentry{ps_134-7_061-v0013}{13}{insertion}{\inserted{total}}{\hand{total}}{The handwritten insertion is incorporated at its indicated position; insertion brackets leave the running text.}
\psentry{ps_134-7_061-v0014}{14}{typewriter deletion}{\xout{total}}{}{The typewriter strikes out this material; it leaves the corrected running text.}
\psentry{ps_134-7_061-v0015}{15}{typed insertion}{\inserted[typed]{left}}{left}{The typed insertion is incorporated at its indicated position; insertion brackets leave the running text.}
\psapparatusend
\fi
\pssource{ps_134-7_062}{134-7/62}{https://grothendieck.umontpellier.fr/web/viewer.html?file=../134-7.pdf\#page=63}
\psnumber{ps_134-7_062-v0001}{1}
\psnumber{ps_134-7_062-v0002}{2}
\psnumber{ps_134-7_062-v0003}{3}
\psnumber{ps_134-7_062-v0004}{4}
\psnumber{ps_134-7_062-v0005}{5}
\psnumber{ps_134-7_062-v0006}{6}
\psnumber{ps_134-7_062-v0007}{7}
\psnumber{ps_134-7_062-v0008}{8}
\psnumber{ps_134-7_062-v0009}{9}
\psnumber{ps_134-7_062-v0010}{10}
\psnumber{ps_134-7_062-v0011}{11}
\psnumber{ps_134-7_062-v0012}{12}
\psnumber{ps_134-7_062-v0013}{13}
\psnumber{ps_134-7_062-v0014}{14}
\psnumber{ps_134-7_062-v0015}{15}
\psnumber{ps_134-7_062-v0016}{16}
\psnumber{ps_134-7_062-v0017}{17}
\psnumber{ps_134-7_062-v0018}{18}
\psnumber{ps_134-7_062-v0019}{19}
\psnumber{ps_134-7_062-v0020}{20}
\psnumber{ps_134-7_062-v0021}{21}
%% BEGIN SOURCE ps_134-7_062
% Pursuing Stacks (cotes 134-2 to 134-8): archive cote 134-7, site page 62 (PDF page 63), typescript page 504.
% typescript. Draft compared: ps6.tex:3039-3098.
% AI-transcribed by Claude Opus 5.5 (claude-opus-5-5, adaptive thinking, effort high), Message Batches API, batch ps4 (msgbatch_012Aoc1NuskATPsyBeKSYK7J), 2026-11-12; reviewed against both supplied scan halves; see CHECK_DECISIONS.tsv.
\scanpage{134-7/62}
\mspage{504}

pursuing even for the \var{\inserted{mere}}{\hand{mere}}{insertion}{ps_134-7_062-v0001} sake of studying ordinary homotopy types.  The main
reason for feeling this way is the amazingly simple description of the
homotopy modules $\hand{\pi}_{i}$ of a homotopy type defined in terms of a semisimplicial
(or hemispherical) unipotent bundle,  \var{\xout{which}}{}{typewriter deletion}{ps_134-7_062-v0002} as derived functors \var{\xout{of the}}{}{typewriter deletion}{ps_134-7_062-v0003}
(so to say) of the Lie functor (cf. previous section 118).  T\var{\over[typed]{h}{h}}{h}{typewriter overwriting}{ps_134-7_062-v0004}e main test
for deciding whether there is indeed a rewarding new tool to be dug out,
is to see whether or not in the model categories \var{\xout{\illegible{1}}}{}{typewriter deletion}{ps_134-7_062-v0005} $\hand{\mathcal{M}}_{n}(k)$ we have been
working with so far,  \var{\xout{\illegible{1}}}{}{typewriter deletion}{ps_134-7_062-v0006} the standart homotopy constructions (around loop-%
spaces and suspensions) make a sense,  and in such a way of course that
the canonical functor from schematic to discrete homotopy types should com-%
mute to these operations.  It may well turn out that to get a handy forma-%
lism,  one will have still to modify more or less the conceptual set-up
of unipotent bundles I've been tentatively working with so far.  I already
\var{\inserted{lately}}{\hand{lately}}{insertion}{ps_134-7_062-v0007} hit upon suggestions of such modifications \var{\deleted{lately}}{}{deletion}{ps_134-7_062-v0008}, and presumahly  I'm
going to discuss this still,  before leaving the topic of schematizations.

Another reason which makes me feel that there should exist a notion of
homotopy types over more general groundrings \var{\inserted[typed]{$k$}}{$k$}{typed insertion}{ps_134-7_062-v0009} than $\underline{Z}$ ,  is that for a number
of rings,  such a notion has been known for quite a while.  If I got it \var{\xout{rig}}{}{typewriter deletion}{ps_134-7_062-v0010}
right,  already in the late sixties (even before I withdrew from the mathe-%
matical milieu) I heard about such things as\var{\xout{s}}{}{typewriter deletion}{ps_134-7_062-v0011} homotopy types over residue\var{\xout{s}}{}{typewriter deletion}{ps_134-7_062-v0012}
\var{\inserted[typed]{class}}{class}{typed insertion}{ps_134-7_062-v0013} rings $\underline{Z}/n\underline{Z}$ ,  or over rings(such as $\underline{Q}$) which are localizations of $\underline{Z}$ ,  or
over rings such as $\underline{Z}^{\hand{\wedge}}$ \var{\xout{\illegible{1}}}{}{typewriter deletion}{ps_134-7_062-v0014} or $\underline{Z}_{p}$ (p-adic integers)  which are completions
of $\underline{Z}$ with respect to a suitable  \var{\xout{\unsure{padic topolgy}}}{}{typewriter deletion}{ps_134-7_062-v0015} linear topology.  Last
week,  which was the first time I was at the university \var{\xout{since}}{}{typewriter deletion}{ps_134-7_062-v0016} after the
summer vacations,  I took from the library the Bousfield-Kan \var{\inserted[typed]{Lecture Notes}}{Lecture Notes}{typed insertion}{ps_134-7_062-v0017} book on
homotopy limits (which had been pointed out to me by Tim Porter in June,
when he had taken the trouble to tell me about``shape theory'' and it's rela-%
tions to \var{\inserted[typed]{(filtering)}}{(filtering)}{typed insertion}{ps_134-7_062-v0018} \var{\deleted{homotopy}}{}{deletion}{ps_134-7_062-v0019} limits), and while glancing through it,  I noticed there
\var{\xout{was}}{}{typewriter deletion}{ps_134-7_062-v0020} is a systematic treatment of such homotopy \uline{types.}\hand{$\rfloor\!\ulcorner$}

At this very minute
I had a closer look upon the introduction of part I,  it turns out that
\var{\xout{they}}{}{typewriter deletion}{ps_134-7_062-v0021} Bousfield and Kan are working with an arbitrary (commutative) ground-%
ring k,  and they are defining corresponding k-completion  $k_{\infty}X$  of a homo-%
topy type k,  rather than a notion of ``homotopy type over k''.  But the two
kind of notions are surely closely related,  the k-completion of BK presu-%
mably should have more or less the meaning of ground \uline{ring extension} $\underline{Z}\hand{\to}$ k. \margin{definitely not, in general!}
At any rate,  for 1-connected spaces and k=$\underline{Z}$  the completion operation
seems to be no more no less than just the identity,  thus it would seem
that the implicit notion of ``homotopy type over $\underline{Z}$ ''  should be just the
ordinary ``discrete'' notion of homotopy types  -  unlike the notion of

\margin{62}

% Visual collation: both supplied image halves inspected; see CHECK_DECISIONS_134_7_8.tsv.

%% END SOURCE ps_134-7_062
\ifnum\EditionMode=1
\psapparatusstart
\psentry{ps_134-7_062-v0001}{1}{insertion}{\inserted{mere}}{\hand{mere}}{The handwritten insertion is incorporated at its indicated position; insertion brackets leave the running text.}
\psentry{ps_134-7_062-v0002}{2}{typewriter deletion}{\xout{which}}{}{The typewriter strikes out this material; it leaves the corrected running text.}
\psentry{ps_134-7_062-v0003}{3}{typewriter deletion}{\xout{of the}}{}{The typewriter strikes out this material; it leaves the corrected running text.}
\psentry{ps_134-7_062-v0004}{4}{typewriter overwriting}{\over[typed]{h}{h}}{h}{The superposed typed reading replaces the underlying reading.}
\psentry{ps_134-7_062-v0005}{5}{typewriter deletion}{\xout{\illegible{1}}}{}{The typewriter strikes out this material; it leaves the corrected running text.}
\psentry{ps_134-7_062-v0006}{6}{typewriter deletion}{\xout{\illegible{1}}}{}{The typewriter strikes out this material; it leaves the corrected running text.}
\psentry{ps_134-7_062-v0007}{7}{insertion}{\inserted{lately}}{\hand{lately}}{The handwritten insertion is incorporated at its indicated position; insertion brackets leave the running text.}
\psentry{ps_134-7_062-v0008}{8}{deletion}{\deleted{lately}}{}{The author cancels this material by hand; it leaves the corrected running text.}
\psentry{ps_134-7_062-v0009}{9}{typed insertion}{\inserted[typed]{$k$}}{$k$}{The typed insertion is incorporated at its indicated position; insertion brackets leave the running text.}
\psentry{ps_134-7_062-v0010}{10}{typewriter deletion}{\xout{rig}}{}{The typewriter strikes out this material; it leaves the corrected running text.}
\psentry{ps_134-7_062-v0011}{11}{typewriter deletion}{\xout{s}}{}{The typewriter strikes out this material; it leaves the corrected running text.}
\psentry{ps_134-7_062-v0012}{12}{typewriter deletion}{\xout{s}}{}{The typewriter strikes out this material; it leaves the corrected running text.}
\psentry{ps_134-7_062-v0013}{13}{typed insertion}{\inserted[typed]{class}}{class}{The typed insertion is incorporated at its indicated position; insertion brackets leave the running text.}
\psentry{ps_134-7_062-v0014}{14}{typewriter deletion}{\xout{\illegible{1}}}{}{The typewriter strikes out this material; it leaves the corrected running text.}
\psentry{ps_134-7_062-v0015}{15}{typewriter deletion}{\xout{\unsure{padic topolgy}}}{}{The typewriter strikes out this material; it leaves the corrected running text.}
\psentry{ps_134-7_062-v0016}{16}{typewriter deletion}{\xout{since}}{}{The typewriter strikes out this material; it leaves the corrected running text.}
\psentry{ps_134-7_062-v0017}{17}{typed insertion}{\inserted[typed]{Lecture Notes}}{Lecture Notes}{The typed insertion is incorporated at its indicated position; insertion brackets leave the running text.}
\psentry{ps_134-7_062-v0018}{18}{typed insertion}{\inserted[typed]{(filtering)}}{(filtering)}{The typed insertion is incorporated at its indicated position; insertion brackets leave the running text.}
\psentry{ps_134-7_062-v0019}{19}{deletion}{\deleted{homotopy}}{}{The author cancels this material by hand; it leaves the corrected running text.}
\psentry{ps_134-7_062-v0020}{20}{typewriter deletion}{\xout{was}}{}{The typewriter strikes out this material; it leaves the corrected running text.}
\psentry{ps_134-7_062-v0021}{21}{typewriter deletion}{\xout{they}}{}{The typewriter strikes out this material; it leaves the corrected running text.}
\psapparatusend
\fi
\pssource{ps_134-7_063}{134-7/63}{https://grothendieck.umontpellier.fr/web/viewer.html?file=../134-7.pdf\#page=64}
\psnumber{ps_134-7_063-v0001}{1}
\psnumber{ps_134-7_063-v0002}{2}
\psnumber{ps_134-7_063-v0003}{3}
\psnumber{ps_134-7_063-v0004}{4}
\psnumber{ps_134-7_063-v0005}{5}
\psnumber{ps_134-7_063-v0006}{6}
\psnumber{ps_134-7_063-v0007}{7}
\psnumber{ps_134-7_063-v0008}{8}
\psnumber{ps_134-7_063-v0009}{9}
\psnumber{ps_134-7_063-v0010}{10}
\psnumber{ps_134-7_063-v0011}{11}
\psnumber{ps_134-7_063-v0012}{12}
\psnumber{ps_134-7_063-v0013}{13}
\psnumber{ps_134-7_063-v0014}{14}
\psnumber{ps_134-7_063-v0015}{15}
\psnumber{ps_134-7_063-v0016}{16}
\psnumber{ps_134-7_063-v0017}{17}
\psnumber{ps_134-7_063-v0018}{18}
\psnumber{ps_134-7_063-v0019}{19}
\psnumber{ps_134-7_063-v0020}{20}
\psnumber{ps_134-7_063-v0021}{21}
\psnumber{ps_134-7_063-v0022}{22}
\psnumber{ps_134-7_063-v0023}{23}
\psnumber{ps_134-7_063-v0024}{24}
\psnumber{ps_134-7_063-v0025}{25}
\psnumber{ps_134-7_063-v0026}{26}
\psnumber{ps_134-7_063-v0027}{27}
\psnumber{ps_134-7_063-v0028}{28}
\psnumber{ps_134-7_063-v0029}{29}
\psnumber{ps_134-7_063-v0030}{30}
\psnumber{ps_134-7_063-v0031}{31}
%% BEGIN SOURCE ps_134-7_063
% Pursuing Stacks (cotes 134-2 to 134-8): archive cote 134-7, site page 63 (PDF page 64), typescript page 505.
% typescript. Draft compared: ps6.tex:3086-3138.
% AI-transcribed by Claude Opus 5.5 (claude-opus-5-5, adaptive thinking, effort high), Message Batches API, batch ps4 (msgbatch_012Aoc1NuskATPsyBeKSYK7J), 2026-11-12; reviewed against both supplied scan halves; see CHECK_DECISIONS.tsv.
\scanpage{134-7/63}
\mspage{505}

schematic homotopy types (over k=$\underline{Z}$) defined via\var{\inserted[typed]{semisimplicial}}{semisimplicial}{typed insertion}{ps_134-7_063-v0001} unipotent bundles. Definit\unsure{e}
-ly, an understanding of schematic homotopy types will have to include the
(by now classical) Bousfield-Kan ideas, \var{\inserted{and these}}{\hand{and these}}{insertion}{ps_134-7_063-v0002} \var{\deleted{which}}{}{deletion}{ps_134-7_063-v0003} are also relevant for my
reflections on\allowbreak{}``integration''\allowbreak{}and ``cointegration''\var{\inserted[typed]{operations}}{operations}{typed insertion}{ps_134-7_063-v0004} (in connection with the
notion of a derivator (section 69)\hand{)}, called in their book ``homotopy \var{\inserted{direct}}{\hand{direct}}{insertion}{ps_134-7_063-v0005} \var{\deleted{inverse}}{}{deletion}{ps_134-7_063-v0006}
limits'' and ``homotopy \var{\inserted{inverse}}{\hand{inverse}}{insertion}{ps_134-7_063-v0007} \var{\deleted{direct}}{}{deletion}{ps_134-7_063-v0008} limits'' (in the special case of the derivator
associated to ordinary homotopy types, if I got it right). It came as a
surprise that \var{\inserted{in their book,}}{\hand{in their book,}}{insertion}{ps_134-7_063-v0009} these operations are developped \var{\xout{\unsure{as}}}{}{typewriter deletion}{ps_134-7_063-v0010} mainly as technical tools
for developping \var{\inserted[typed]{their}}{their}{typed insertion}{ps_134-7_063-v0011} \var{\xout{the}}{}{typewriter deletion}{ps_134-7_063-v0012} theory of k-completions, whereas in my own reflections
they appeared from the start as ``the'' main operations in homotopical as
well as homological algebra. There had been quite a similar surprise when
Tim Porter had sent me a reprint (in \var{\over{\unsure{j}}{J}}{\hand{J}}{writing over}{ps_134-7_063-v0013}uly, just before I stopped \var{\inserted{with my}}{\hand{with my}}{insertion}{ps_134-7_063-v0014} \var{\deleted{the}}{}{deletion}{ps_134-7_063-v0015} note\unsure{s}
for a month or so) of Don Anderson's beautiful paper ``Fibrations and \var{\over{\unsure{g}}{G}}{\hand{G}}{writing over}{ps_134-7_063-v0016}eome-
tric \var{\over{\unsure{r}}{R}}{\hand{R}}{writing over}{ps_134-7_063-v0017}ealizations'' (Bulletin of the Amer. Math. Soc. September 1978),
where a very general and (as I feel) quite basic existence theorem for
\var{integration and cointegration operations\hand{,} of the type precisely I was af-
ter\hand{,} in the \xout{\illegible{1}} set-up of closed model categories of Quillen\hand{,} is barely}{integration and cointegration operations\hand{,} in the  set-up of closed model categories of Quillen\hand{,} of the type precisely I was af-
ter\hand{,} is barely}{authorial transposition}{ps_134-7_063-v0018}
alluded to at the end of the introduction, and \var{\inserted[typed]{comes}}{comes}{typed insertion}{ps_134-7_063-v0019} \var{\xout{\illegible{1} appears as a}}{}{typewriter deletion}{ps_134-7_063-v0020}
more or less as \var{\inserted{just}}{\hand{just}}{insertion}{ps_134-7_063-v0021} a by\var{\deleted{\unsure{e}}}{}{deletion}{ps_134-7_063-v0022}-product of \var{\inserted[typed]{work done in view of a}}{work done in view of a}{typed insertion}{ps_134-7_063-v0023} \var{\xout{\unsure{an exposition of}}}{}{typewriter deletion}{ps_134-7_063-v0024} a result on geometric
realizations which (to an outsider like me at any rate) looks highly te\hand{\unsure{c}}hni
-cal and not inspiring in\var{\deleted{\unsure{t}}}{}{deletion}{ps_134-7_063-v0025} the least !

It is becoming clear that I cannot put off much longer getting acquain
-ted with the main ideas and results of Bousfield-Kan's book, which defini-
tely looks like one of the few basic texts \var{\over{\unsure{i}}{o}}{\hand{o}}{writing over}{ps_134-7_063-v0026}n \var{\xout{homotopy theory}}{}{typewriter deletion}{ps_134-7_063-v0027} foundational
matters in homotopy theory. Still, before doing some \var{\inserted[typed]{basic}}{basic}{typed insertion}{ps_134-7_063-v0028} \var{\xout{\unsure{more}}}{}{typewriter deletion}{ps_134-7_063-v0029} reading, I
would like to write down the sporadic reflections on schematic homotopy
types I went into during the last \var{\inserted{weeks}}{\hand{weeks}}{insertion}{ps_134-7_063-v0030} \var{\deleted{month}}{}{deletion}{ps_134-7_063-v0031}, while they are still fresh in my
mind !

\margin{63}

% Visual collation: both supplied image halves inspected; see CHECK_DECISIONS_134_7_8.tsv.

%% END SOURCE ps_134-7_063
\ifnum\EditionMode=1
\psapparatusstart
\psentry{ps_134-7_063-v0001}{1}{typed insertion}{\inserted[typed]{semisimplicial}}{semisimplicial}{The typed insertion is incorporated at its indicated position; insertion brackets leave the running text.}
\psentry{ps_134-7_063-v0002}{2}{insertion}{\inserted{and these}}{\hand{and these}}{The handwritten insertion is incorporated at its indicated position; insertion brackets leave the running text.}
\psentry{ps_134-7_063-v0003}{3}{deletion}{\deleted{which}}{}{The author cancels this material by hand; it leaves the corrected running text.}
\psentry{ps_134-7_063-v0004}{4}{typed insertion}{\inserted[typed]{operations}}{operations}{The typed insertion is incorporated at its indicated position; insertion brackets leave the running text.}
\psentry{ps_134-7_063-v0005}{5}{insertion}{\inserted{direct}}{\hand{direct}}{The handwritten insertion is incorporated at its indicated position; insertion brackets leave the running text.}
\psentry{ps_134-7_063-v0006}{6}{deletion}{\deleted{inverse}}{}{The author cancels this material by hand; it leaves the corrected running text.}
\psentry{ps_134-7_063-v0007}{7}{insertion}{\inserted{inverse}}{\hand{inverse}}{The handwritten insertion is incorporated at its indicated position; insertion brackets leave the running text.}
\psentry{ps_134-7_063-v0008}{8}{deletion}{\deleted{direct}}{}{The author cancels this material by hand; it leaves the corrected running text.}
\psentry{ps_134-7_063-v0009}{9}{insertion}{\inserted{in their book,}}{\hand{in their book,}}{The handwritten insertion is incorporated at its indicated position; insertion brackets leave the running text.}
\psentry{ps_134-7_063-v0010}{10}{typewriter deletion}{\xout{\unsure{as}}}{}{The typewriter strikes out this material; it leaves the corrected running text.}
\psentry{ps_134-7_063-v0011}{11}{typed insertion}{\inserted[typed]{their}}{their}{The typed insertion is incorporated at its indicated position; insertion brackets leave the running text.}
\psentry{ps_134-7_063-v0012}{12}{typewriter deletion}{\xout{the}}{}{The typewriter strikes out this material; it leaves the corrected running text.}
\psentry{ps_134-7_063-v0013}{13}{writing over}{\over{\unsure{j}}{J}}{\hand{J}}{The superposed handwritten reading replaces the underlying reading.}
\psentry{ps_134-7_063-v0014}{14}{insertion}{\inserted{with my}}{\hand{with my}}{The handwritten insertion is incorporated at its indicated position; insertion brackets leave the running text.}
\psentry{ps_134-7_063-v0015}{15}{deletion}{\deleted{the}}{}{The author cancels this material by hand; it leaves the corrected running text.}
\psentry{ps_134-7_063-v0016}{16}{writing over}{\over{\unsure{g}}{G}}{\hand{G}}{The superposed handwritten reading replaces the underlying reading.}
\psentry{ps_134-7_063-v0017}{17}{writing over}{\over{\unsure{r}}{R}}{\hand{R}}{The superposed handwritten reading replaces the underlying reading.}
\psentry{ps_134-7_063-v0018}{18}{authorial transposition}{integration and cointegration operations\hand{,} of the type precisely I was af-
ter\hand{,} in the \xout{\illegible{1}} set-up of closed model categories of Quillen\hand{,} is barely}{integration and cointegration operations\hand{,} in the  set-up of closed model categories of Quillen\hand{,} of the type precisely I was af-
ter\hand{,} is barely}{A pen enclosure surrounds in the ... set-up of closed model categories of Quillen and connects it to an insertion caret immediately after operations. Keep typed order in diplomatic text and apply the relocation in reading mode. Existing pen punctuation is retained.}
\psentry{ps_134-7_063-v0019}{19}{typed insertion}{\inserted[typed]{comes}}{comes}{The typed insertion is incorporated at its indicated position; insertion brackets leave the running text.}
\psentry{ps_134-7_063-v0020}{20}{typewriter deletion}{\xout{\illegible{1} appears as a}}{}{The typewriter strikes out this material; it leaves the corrected running text.}
\psentry{ps_134-7_063-v0021}{21}{insertion}{\inserted{just}}{\hand{just}}{The handwritten insertion is incorporated at its indicated position; insertion brackets leave the running text.}
\psentry{ps_134-7_063-v0022}{22}{deletion}{\deleted{\unsure{e}}}{}{The author cancels this material by hand; it leaves the corrected running text.}
\psentry{ps_134-7_063-v0023}{23}{typed insertion}{\inserted[typed]{work done in view of a}}{work done in view of a}{The typed insertion is incorporated at its indicated position; insertion brackets leave the running text.}
\psentry{ps_134-7_063-v0024}{24}{typewriter deletion}{\xout{\unsure{an exposition of}}}{}{The typewriter strikes out this material; it leaves the corrected running text.}
\psentry{ps_134-7_063-v0025}{25}{deletion}{\deleted{\unsure{t}}}{}{The author cancels this material by hand; it leaves the corrected running text.}
\psentry{ps_134-7_063-v0026}{26}{writing over}{\over{\unsure{i}}{o}}{\hand{o}}{The superposed handwritten reading replaces the underlying reading.}
\psentry{ps_134-7_063-v0027}{27}{typewriter deletion}{\xout{homotopy theory}}{}{The typewriter strikes out this material; it leaves the corrected running text.}
\psentry{ps_134-7_063-v0028}{28}{typed insertion}{\inserted[typed]{basic}}{basic}{The typed insertion is incorporated at its indicated position; insertion brackets leave the running text.}
\psentry{ps_134-7_063-v0029}{29}{typewriter deletion}{\xout{\unsure{more}}}{}{The typewriter strikes out this material; it leaves the corrected running text.}
\psentry{ps_134-7_063-v0030}{30}{insertion}{\inserted{weeks}}{\hand{weeks}}{The handwritten insertion is incorporated at its indicated position; insertion brackets leave the running text.}
\psentry{ps_134-7_063-v0031}{31}{deletion}{\deleted{month}}{}{The author cancels this material by hand; it leaves the corrected running text.}
\psapparatusend
\fi
\pssource{ps_134-7_064}{134-7/64}{https://grothendieck.umontpellier.fr/web/viewer.html?file=../134-7.pdf\#page=65}
\psnumber{ps_134-7_064-v0001}{1}
\psnumber{ps_134-7_064-v0002}{2}
\psnumber{ps_134-7_064-v0003}{3}
\psnumber{ps_134-7_064-v0004}{4}
\psnumber{ps_134-7_064-v0005}{5}
\psnumber{ps_134-7_064-v0006}{6}
\psnumber{ps_134-7_064-v0007}{7}
\psnumber{ps_134-7_064-v0008}{8}
\psnumber{ps_134-7_064-v0009}{9}
\psnumber{ps_134-7_064-v0010}{10}
\psnumber{ps_134-7_064-v0011}{11}
\psnumber{ps_134-7_064-v0012}{12}
\psnumber{ps_134-7_064-v0013}{13}
\psnumber{ps_134-7_064-v0014}{14}
\psnumber{ps_134-7_064-v0015}{15}
\psnumber{ps_134-7_064-v0016}{16}
\psnumber{ps_134-7_064-v0017}{17}
\psnumber{ps_134-7_064-v0018}{18}
\psnumber{ps_134-7_064-v0019}{19}
\psnumber{ps_134-7_064-v0020}{20}
\psnumber{ps_134-7_064-v0021}{21}
\psnumber{ps_134-7_064-v0022}{22}
\psnumber{ps_134-7_064-v0023}{23}
\psnumber{ps_134-7_064-v0024}{24}
\psnumber{ps_134-7_064-v0025}{25}
\psnumber{ps_134-7_064-v0026}{26}
\psnumber{ps_134-7_064-v0027}{27}
%% BEGIN SOURCE ps_134-7_064
% Pursuing Stacks (cotes 134-2 to 134-8): archive cote 134-7, site page 64 (PDF page 65), typescript page 506.
% typescript. Draft compared: ps6.tex:3120-3190.
% AI-transcribed by Claude Opus 5.5 (claude-opus-5-5, adaptive thinking, effort high), Message Batches API, batch pilot (msgbatch_0189Nhf1xB615YWnztubtLfc), 2026-11-12; reviewed against both supplied scan halves; see CHECK_DECISIONS.tsv.
\scanpage{134-7/64}
\mspage{506}

28.9.
120.\quad Yesterday\var{\inserted{(prompted by the reflections from the day before, cf. section 119)}}{\hand{(prompted by the reflections from the day before, cf. section 119)}}{insertion}{ps_134-7_064-v0001} I pondered a little on the common features of the various
set-ups for ``commutative algebra'' (possibly too for corresponding notions
of ``schemes'') one gets when introducing extra operations on commutative
rings or algebras, such as divided power structure\var{\xout{s}}{}{typewriter deletion}{ps_134-7_064-v0002} (on a suitable ideal\hand{),}
or a \hand{$\lambda$}-structure with\var{\xout{n}}{}{typewriter deletion}{ps_134-7_064-v0003} operations \hand{$\lambda$}$^i$ paraphrasing exterior powers, or a
binomial structure with operations $x\hand{\mapsto}\binom{x}{n}$ paraphrasing binomial coeffi-%
cients, or an S-structure with operations $S^i$ of Adam's type paraphrasing
sums of i.th powers of roots of a polynomial (a weakened version of a \hand{$\lambda$}-%
structure). It seems that the unifying notion here is the notion of an
``analyzer'' (analyseur) of Lazare, ``containing'' the Lazare analyzer for
commutative rings (not necessarily with unit), so that the components
\margin{\textcircled{$n\geqslant -1$}} \hand{$\Omega$}$_n$ (n\hand{$\geqslant$}\var{\over{\unsure{1}}{-1}}{\hand{-1}}{writing over}{ps_134-7_064-v0004}) are commutative rings (not necessarily with unit). In case
the \var{\xout{structural}}{}{typewriter deletion}{ps_134-7_064-v0005}\allowbreak{}\var{\inserted[typed]{extra}}{extra}{typed insertion}{ps_134-7_064-v0006} operations we want to introduce on commutative rings are
to be defined on rings \uline{with units} (not just on a suitable ideal of such
a ring, as is the case for the divided power structure), and they all
can be defined in terms of operations involving just one argument, the
reasonable extra axiom on \var{\xout{the}}{}{typewriter deletion}{ps_134-7_064-v0007} the corresponding analyzer (as suggested by
the examples at hand) is that for n \hand{$\geqslant$} \var{\over{\unsure{2}}{1}}{\hand{1}}{writing over}{ps_134-7_064-v0008}, \hand{$\Omega$}$_n$ can be recovered in terms
of \hand{$\Omega$}$_0$ and \hand{$\Omega$}$_{-1}$ = k$_0$ (the latter act\var{\xout{\unsure{ion}}}{}{typewriter deletion}{ps_134-7_064-v0009}\var{\inserted[typed]{ing}}{ing}{typed insertion}{ps_134-7_064-v0010} as a ground ring for the theory\hand{)}
as the \var{\deleted{n}}{}{deletion}{ps_134-7_064-v0011}\var{\inserted{(n+1)}}{\hand{(n+1)}}{insertion}{ps_134-7_064-v0012}-fold tensor power of \hand{$\Omega$}$_0$ over k$_0$ . Thus the \var{\deleted{\illegible{1}}}{}{deletion}{ps_134-7_064-v0013}whole structure of
\margin{\textcircled{à la ligne}} the analyzor may be thought of as embodied in the system \var{\inserted{$\Omega = (k_0, \Omega_0)$,}}{\hand{$\Omega = (k_0, \Omega_0)$,}}{insertion}{ps_134-7_064-v0014} \var{\deleted{$(k_0, \Omega_0)$, where}}{}{deletion}{ps_134-7_064-v0015}
\var{\inserted{where}}{\hand{where}}{insertion}{ps_134-7_064-v0016} k$_0$ is a commutative ring (with unit, now), \hand{$\Omega$}$_0$ a commutative k$_0$-algebra
\margin{\textcircled{id}} with unit, endowed moreover with a composition operation \var{\deleted{$(F,G)\mapsto F\circ G$,}}{}{deletion}{ps_134-7_064-v0017}\var{\inserted{$(F,G)\longmapsto F\circ G$ ,}}{\hand{$(F,G)\longmapsto F\circ G$ ,}}{insertion}{ps_134-7_064-v0018}
satisfying a bunch of simple axioms I do'nt feel like writing down here.
The simplest case of all of course (correspo\var{\over[typed]{\unsure{n}}{n}}{n}{typewriter overwriting}{ps_134-7_064-v0019}ding to \var{\xout{the}}{}{typewriter deletion}{ps_134-7_064-v0020} \var{\xout{u}}{}{typewriter deletion}{ps_134-7_064-v0021}usual commutati%
ve algebra ``over k$_0$'' as a ground ring, with no extra structure on commu-%
tative algebras with unit over k$_0$) is \hand{$\Omega$}$_0$ = k\hand{[}T\hand{]}\hand{,} with the usual compositi%
on of polynomials, T acting as the two-sided unit for composition. In
the general case, \hand{$\Omega$}$_0$ and it's tensor powers \var{\inserted[typed]{$\hand{\Omega}_n$ over $k_0$}}{$\hand{\Omega}_n$ over $k_0$}{typed insertion}{ps_134-7_064-v0022} are going to play the part
played by polynomial rings in ordinary commutative algebra, \var{\deleted{and}}{}{deletion}{ps_134-7_064-v0023} \var{\over{t}{T}}{\hand{T}}{writing over}{ps_134-7_064-v0024}here
should be a ready generalization, in this spirit, of taking the symme-%
tric algebra of a k$_0$-module (which, for a module free and of finite type\hand{,}
will \var{\deleted{be}}{}{deletion}{ps_134-7_064-v0025}\var{\inserted{yield a $\Omega$-structure}}{\hand{yield a $\Omega$-structure}}{insertion}{ps_134-7_064-v0026} isomorphic to one of the \hand{$\Omega$}$_n$'s). An \hand{$\Omega$}-structure on a set \var{\xout{x}}{}{typewriter deletion}{ps_134-7_064-v0027} k
amounts to giving a structure of a commutative k$_0$-algebra with unit on k,
plus a map
\[ \hand{\Omega_0 \to} \mathrm{Maps}(k,k) \quad \hand{,} \quad \hand{F \longmapsto (x \longmapsto F(x))} \]

\margin{64}

% Visual collation: both supplied image halves inspected; see CHECK_DECISIONS_134_7_8.tsv.

%% END SOURCE ps_134-7_064
\ifnum\EditionMode=1
\psapparatusstart
\psentry{ps_134-7_064-v0001}{1}{insertion}{\inserted{(prompted by the reflections from the day before, cf. section 119)}}{\hand{(prompted by the reflections from the day before, cf. section 119)}}{The handwritten insertion is incorporated at its indicated position; insertion brackets leave the running text.}
\psentry{ps_134-7_064-v0002}{2}{typewriter deletion}{\xout{s}}{}{The typewriter strikes out this material; it leaves the corrected running text.}
\psentry{ps_134-7_064-v0003}{3}{typewriter deletion}{\xout{n}}{}{The typewriter strikes out this material; it leaves the corrected running text.}
\psentry{ps_134-7_064-v0004}{4}{writing over}{\over{\unsure{1}}{-1}}{\hand{-1}}{The superposed handwritten reading replaces the underlying reading.}
\psentry{ps_134-7_064-v0005}{5}{typewriter deletion}{\xout{structural}}{}{The typewriter strikes out this material; it leaves the corrected running text.}
\psentry{ps_134-7_064-v0006}{6}{typed insertion}{\inserted[typed]{extra}}{extra}{The typed insertion is incorporated at its indicated position; insertion brackets leave the running text.}
\psentry{ps_134-7_064-v0007}{7}{typewriter deletion}{\xout{the}}{}{The typewriter strikes out this material; it leaves the corrected running text.}
\psentry{ps_134-7_064-v0008}{8}{writing over}{\over{\unsure{2}}{1}}{\hand{1}}{The superposed handwritten reading replaces the underlying reading.}
\psentry{ps_134-7_064-v0009}{9}{typewriter deletion}{\xout{\unsure{ion}}}{}{The typewriter strikes out this material; it leaves the corrected running text.}
\psentry{ps_134-7_064-v0010}{10}{typed insertion}{\inserted[typed]{ing}}{ing}{The typed insertion is incorporated at its indicated position; insertion brackets leave the running text.}
\psentry{ps_134-7_064-v0011}{11}{deletion}{\deleted{n}}{}{The author cancels this material by hand; it leaves the corrected running text.}
\psentry{ps_134-7_064-v0012}{12}{insertion}{\inserted{(n+1)}}{\hand{(n+1)}}{The handwritten insertion is incorporated at its indicated position; insertion brackets leave the running text.}
\psentry{ps_134-7_064-v0013}{13}{deletion}{\deleted{\illegible{1}}}{}{The author cancels this material by hand; it leaves the corrected running text.}
\psentry{ps_134-7_064-v0014}{14}{insertion}{\inserted{$\Omega = (k_0, \Omega_0)$,}}{\hand{$\Omega = (k_0, \Omega_0)$,}}{The handwritten insertion is incorporated at its indicated position; insertion brackets leave the running text.}
\psentry{ps_134-7_064-v0015}{15}{deletion}{\deleted{$(k_0, \Omega_0)$, where}}{}{The author cancels this material by hand; it leaves the corrected running text.}
\psentry{ps_134-7_064-v0016}{16}{insertion}{\inserted{where}}{\hand{where}}{The handwritten insertion is incorporated at its indicated position; insertion brackets leave the running text.}
\psentry{ps_134-7_064-v0017}{17}{deletion}{\deleted{$(F,G)\mapsto F\circ G$,}}{}{The author cancels this material by hand; it leaves the corrected running text.}
\psentry{ps_134-7_064-v0018}{18}{insertion}{\inserted{$(F,G)\longmapsto F\circ G$ ,}}{\hand{$(F,G)\longmapsto F\circ G$ ,}}{The handwritten insertion is incorporated at its indicated position; insertion brackets leave the running text.}
\psentry{ps_134-7_064-v0019}{19}{typewriter overwriting}{\over[typed]{\unsure{n}}{n}}{n}{The superposed typed reading replaces the underlying reading.}
\psentry{ps_134-7_064-v0020}{20}{typewriter deletion}{\xout{the}}{}{The typewriter strikes out this material; it leaves the corrected running text.}
\psentry{ps_134-7_064-v0021}{21}{typewriter deletion}{\xout{u}}{}{The typewriter strikes out this material; it leaves the corrected running text.}
\psentry{ps_134-7_064-v0022}{22}{typed insertion}{\inserted[typed]{$\hand{\Omega}_n$ over $k_0$}}{$\hand{\Omega}_n$ over $k_0$}{The typed insertion is incorporated at its indicated position; insertion brackets leave the running text.}
\psentry{ps_134-7_064-v0023}{23}{deletion}{\deleted{and}}{}{The author cancels this material by hand; it leaves the corrected running text.}
\psentry{ps_134-7_064-v0024}{24}{writing over}{\over{t}{T}}{\hand{T}}{The superposed handwritten reading replaces the underlying reading.}
\psentry{ps_134-7_064-v0025}{25}{deletion}{\deleted{be}}{}{The author cancels this material by hand; it leaves the corrected running text.}
\psentry{ps_134-7_064-v0026}{26}{insertion}{\inserted{yield a $\Omega$-structure}}{\hand{yield a $\Omega$-structure}}{The handwritten insertion is incorporated at its indicated position; insertion brackets leave the running text.}
\psentry{ps_134-7_064-v0027}{27}{typewriter deletion}{\xout{x}}{}{The typewriter strikes out this material; it leaves the corrected running text.}
\psapparatusend
\fi
\pssource{ps_134-7_065}{134-7/65}{https://grothendieck.umontpellier.fr/web/viewer.html?file=../134-7.pdf\#page=66}
\psnumber{ps_134-7_065-v0001}{1}
\psnumber{ps_134-7_065-v0002}{2}
\psnumber{ps_134-7_065-v0003}{3}
\psnumber{ps_134-7_065-v0004}{4}
\psnumber{ps_134-7_065-v0005}{5}
\psnumber{ps_134-7_065-v0006}{6}
\psnumber{ps_134-7_065-v0007}{7}
\psnumber{ps_134-7_065-v0008}{8}
\psnumber{ps_134-7_065-v0009}{9}
%% BEGIN SOURCE ps_134-7_065
% Pursuing Stacks (cotes 134-2 to 134-8): archive cote 134-7, site page 65 (PDF page 66), typescript page 507.
% typescript. Draft compared: ps6.tex:3172-3233.
% AI-transcribed by Claude Opus 5.5 (claude-opus-5-5, adaptive thinking, effort high), Message Batches API, batch ps4 (msgbatch_012Aoc1NuskATPsyBeKSYK7J), 2026-11-12; reviewed against both supplied scan halves; see CHECK_DECISIONS.tsv.
\scanpage{134-7/65}
\mspage{507}
compatible with the structures of $k_o$-algebras as well as composition
operations,  and satisfying moreover two conditions for
\[ F(x+y) \ , \quad \text{resp. } F(xy) \]
in terms of two diagonal maps
% Hand-drawn arrows are marked blue; typed objects remain plain.
\[ \diagram{\begin{tikzcd} \hand{\Omega}_o \ar[color=blue,r,shift left,"\hand{\Delta_a}"] \ar[color=blue,r,shift right,"\hand{\Delta_m}"'] \& \hand{\Omega}_1 = \hand{\Omega}_o\hand{\otimes}_{k_o}\hand{\Omega}_o \end{tikzcd}} \]
(which may be described in terms of the composition structure on $\hand{\Omega}_o$ ,
as expressing the compositions  $Fo(G'+G")$  resp.  $Fo(G'G")$),  and moreo\var{\xout{x}}{}{typewriter deletion}{ps_134-7_065-v0001}-%
ver one trivial condition for $F(\hand{\lambda}))$  when $\hand{\lambda}$ in $k$ comes from $k_o$ , \var{\xout{\illegible{1}}}{}{typewriter deletion}{ps_134-7_065-v0002}
namely compatibility of the map $k_o \hand{\to} k$  with the operations of $\hand{\Omega}_o$ on
both $k_o$ and $k$. (NB T\var{\over[typed]{h}{h}}{h}{typewriter overwriting}{ps_134-7_065-v0003}e operation of $\hand{\Omega}_o$ on $k_o$ is defined by
\[ F(\hand{\lambda}) = Fo\hand{\lambda} \ , \]
where $k_o$ is identified to a subring of $\hand{\Omega}_o$ .)

I did'nt pursue much further these ponderings,  just one digression
among many in the main line of investigation !  I also read through
the preprint of Davi\var{\over{c}{d}}{\hand{d}}{writing over}{ps_134-7_065-v0004} W. Jones on Poly-T-complexes,  which Ronnie Brown
(\var{\inserted{acting as}}{\hand{acting as}}{insertion}{ps_134-7_065-v0005}David Jones' supervisor) had sent me a while ago.  There he develops
a notion of polyhedral cells,  with a view of using these instead of
simplices or cubes for doing combinatorial homotopy theory\unsure{.|,} \var{\xout{in a spirit}}{}{typewriter deletion}{ps_134-7_065-v0006}
As I had pondered a little along this direction (cf. sections 91\hand{,} topic
\hand{(}8),  and section 93),  I was hpping that some of the perplexities I
had been meeting would be solved in David Jones\var{\xout{\unsure{k}}}{}{typewriter deletion}{ps_134-7_065-v0007}\hand{'} notes  -  for instance
that there would be handy criteria for a category $M$ made up with such
polyhedral cells to be a weak test category,  namely that objects of $M^{\hand{\wedge}}$
may be used as models for homotopy types;  also\hand{,}  that the``standard''
\var{\xout{\unsure{simple}}}{}{typewriter deletion}{ps_134-7_065-v0008} chain complex constructed in section 93 is indeed an ``abelianiza-%
tor'' for $M$,  i.e. may be used for computing homology of objects of $M^{\hand{\wedge}}$.
David Jones' emphasis however is a rather different one  -  he seems
mainly interested in generalizing the theory of ``thin'' structures
of M.K. Dakin \var{\over{\unsure{t}}{f}}{\hand{f}}{writing over}{ps_134-7_065-v0009}rom the simplicial to the more general polyhedral set-up,
and prove a corresponding equivalence of categories.  Thus my perplexi-%
ties remain - they are admittedly rather marginal in the main line of
thought,  and I doubt I'll stop to try and solve them.

\margin{65}

% Visual collation: both supplied image halves inspected; see CHECK_DECISIONS_134_7_8.tsv.

%% END SOURCE ps_134-7_065
\ifnum\EditionMode=1
\psapparatusstart
\psentry{ps_134-7_065-v0001}{1}{typewriter deletion}{\xout{x}}{}{The typewriter strikes out this material; it leaves the corrected running text.}
\psentry{ps_134-7_065-v0002}{2}{typewriter deletion}{\xout{\illegible{1}}}{}{The typewriter strikes out this material; it leaves the corrected running text.}
\psentry{ps_134-7_065-v0003}{3}{typewriter overwriting}{\over[typed]{h}{h}}{h}{The superposed typed reading replaces the underlying reading.}
\psentry{ps_134-7_065-v0004}{4}{writing over}{\over{c}{d}}{\hand{d}}{The superposed handwritten reading replaces the underlying reading.}
\psentry{ps_134-7_065-v0005}{5}{insertion}{\inserted{acting as}}{\hand{acting as}}{The handwritten insertion is incorporated at its indicated position; insertion brackets leave the running text.}
\psentry{ps_134-7_065-v0006}{6}{typewriter deletion}{\xout{in a spirit}}{}{The typewriter strikes out this material; it leaves the corrected running text.}
\psentry{ps_134-7_065-v0007}{7}{typewriter deletion}{\xout{\unsure{k}}}{}{The typewriter strikes out this material; it leaves the corrected running text.}
\psentry{ps_134-7_065-v0008}{8}{typewriter deletion}{\xout{\unsure{simple}}}{}{The typewriter strikes out this material; it leaves the corrected running text.}
\psentry{ps_134-7_065-v0009}{9}{writing over}{\over{\unsure{t}}{f}}{\hand{f}}{The superposed handwritten reading replaces the underlying reading.}
\psapparatusend
\fi
\pssource{ps_134-7_066}{134-7/66}{https://grothendieck.umontpellier.fr/web/viewer.html?file=../134-7.pdf\#page=67}
\psnumber{ps_134-7_066-v0001}{1}
\psnumber{ps_134-7_066-v0002}{2}
\psnumber{ps_134-7_066-v0003}{3}
\psnumber{ps_134-7_066-v0004}{4}
\psnumber{ps_134-7_066-v0005}{5}
\psnumber{ps_134-7_066-v0006}{6}
\psnumber{ps_134-7_066-v0007}{7}
\psnumber{ps_134-7_066-v0008}{8}
\psnumber{ps_134-7_066-v0009}{9}
\psnumber{ps_134-7_066-v0010}{10}
\psnumber{ps_134-7_066-v0011}{11}
\psnumber{ps_134-7_066-v0012}{12}
\psnumber{ps_134-7_066-v0013}{13}
\psnumber{ps_134-7_066-v0014}{14}
\psnumber{ps_134-7_066-v0015}{15}
\psnumber{ps_134-7_066-v0016}{16}
\psnumber{ps_134-7_066-v0017}{17}
\psnumber{ps_134-7_066-v0018}{18}
\psnumber{ps_134-7_066-v0019}{19}
\psnumber{ps_134-7_066-v0020}{20}
\psnumber{ps_134-7_066-v0021}{21}
\psnumber{ps_134-7_066-v0022}{22}
\psnumber{ps_134-7_066-v0023}{23}
\psnumber{ps_134-7_066-v0024}{24}
\psnumber{ps_134-7_066-v0025}{25}
\psnumber{ps_134-7_066-v0026}{26}
\psnumber{ps_134-7_066-v0027}{27}
\psnumber{ps_134-7_066-v0028}{28}
\psnumber{ps_134-7_066-v0029}{29}
%% BEGIN SOURCE ps_134-7_066
% Pursuing Stacks (cotes 134-2 to 134-8): archive cote 134-7, site page 66 (PDF page 67), typescript page 508.
% typescript. Draft compared: ps6.tex:3218-3296.
% AI-transcribed by Claude Opus 5.5 (claude-opus-5-5, adaptive thinking, effort high), Message Batches API, batch ps4 (msgbatch_012Aoc1NuskATPsyBeKSYK7J), 2026-11-12; reviewed against both supplied scan halves; see CHECK_DECISIONS.tsv.
\scanpage{134-7/66}
\mspage{508}

121.    The tentative approach towards\var{\xout{a}}{}{typewriter deletion}{ps_134-7_066-v0001} defining and studying "schematic"
homotopy types I have been following lately relies heavily on a suitable
notion of "linearization" of such homotopy types.  One \var{\xout{may}}{}{typewriter deletion}{ps_134-7_066-v0002} can imagine
that \var{\deleted{very much}}{}{deletion}{ps_134-7_066-v0003}\var{\inserted{many}}{\hand{many}}{insertion}{ps_134-7_066-v0004} different appr\var{\xout{\unsure{w}}}{}{typewriter deletion}{ps_134-7_066-v0005}oaches (for instance via De Rham complexes
with devided powers,  or additive small categories with diagonal maps)
\margin{$\times$}may be devised for "schematic" homotopy types,  but in \uwave{any case} it seems
likely that a suitable notion of linearization will play an important
role\var{\over[typed]{,}{.}}{.}{typewriter overwriting}{ps_134-7_066-v0006}  \var{\xout{and}}{}{typewriter deletion}{ps_134-7_066-v0007}It may be worthwhile therefore to try and pin down the wished-%
for main features of such a theory,  with the hope maybe of getting an
axiomatic description for it,  with a corresponding unicity statement.
Before doing so,  the first thing to do seems to review some main formal
features of linearization for ordinary ("discrete") homotopy types.

Recall the definition

(1) \qquad Hotab $\overset{\mathrm{def}}{=}$ D.(Ab) = derived category of the category of
\var{\xout{chain}}{}{typewriter deletion}{ps_134-7_066-v0008} abelian chain complexes\hand{,} with
respect to quasi-isomorphisms\hand{,}

and the two canonical functors

% Hand-drawn arrows are marked blue; typed objects remain plain.
(2) \qquad \diagram{\begin{tikzcd} \mathrm{Hot} \ar[color=blue,r, shift left, "\hand{LH.}"] \& \mathrm{Hotab} \ar[color=blue,l, shift left, "\var{\deleted{\illegible{1}}}{}{deletion}{ps_134-7_066-v0009}\hand{K_\pi}"] \end{tikzcd}} \quad ,

Where \var{\deleted{\illegible{1}.}}{}{deletion}{ps_134-7_066-v0010}\var{\inserted{LH.}}{\hand{LH.}}{insertion}{ps_134-7_066-v0011} is the "abelianization functor",  and \var{\deleted{DP}}{}{deletion}{ps_134-7_066-v0012}\var{\inserted{$K_\pi$}}{\hand{$K_\pi$}}{insertion}{ps_134-7_066-v0013} is\var{\inserted{\deleted{$K_\pi$}}}{\hand{}}{insertion}{ps_134-7_066-v0014}\var{\inserted[typed]{defined via}}{defined via}{typed insertion}{ps_134-7_066-v0015} the Kan-Dold-Puppe
functor,  associating to a chain complex the \var{\xout{homotopy type}}{}{typewriter deletion}{ps_134-7_066-v0016} corresponding
semisimplicial abelian group .  The diagram (2) may be viewed\var{\inserted{(up to equivalence)}}{\hand{(up to equivalence)}}{insertion}{ps_134-7_066-v0017} as deduced
from the corresponding diagram

(2') \qquad \diagram{\begin{tikzcd} \hand{\Delta}^{\hand{\wedge}} \ar[color=blue,r, shift left, "W"] \& \hand{\Delta}^{\hand{\wedge}}_{ab} \ar[color=blue,l, shift left, "DP"] \end{tikzcd}} \quad ($\hand{\cong}$ Ch.(Ab))

by passing to the suitable localized categories Hot and Hotab.  In the
diagram (2')  \var{\xout{DP is right}}{}{typewriter deletion}{ps_134-7_066-v0018}\var{\inserted[typed]{W is left}}{W is left}{typed insertion}{ps_134-7_066-v0019} adjoint to DP ,  which is now just the forget-%
ful functor.   One main fact about (2) is

(3) \qquad \var{\over{K}{$K_\pi$}}{\hand{$K_\pi$}}{writing over}{ps_134-7_066-v0020} is right adjoint to LH. \quad ,

which is just a neater way \var{\over{to}{for}}{\hand{for}}{writing over}{ps_134-7_066-v0021} express\var{\inserted{ing}}{\hand{ing}}{insertion}{ps_134-7_066-v0022} the familiar fact that for given
abelian group $\hand{\pi}$ and  natural integer n,  the object K(\hand{$\pi$},n) in  H\var{\over{o}{\unsure{\v{o}}}}{\hand{\unsure{\v{o}}}}{writing over}{ps_134-7_066-v0023}t
(namely the image by \var{\deleted{K}}{}{deletion}{ps_134-7_066-v0024}\var{\inserted{$K_\pi$}}{\hand{$K_\pi$}}{insertion}{ps_134-7_066-v0025} of the chain complex\var{\inserted{$\pi[n]$}}{\hand{$\pi[n]$}}{insertion}{ps_134-7_066-v0026} reduced to \hand{$\pi$} in degree n)
represents the cohomology functor
\[ X \hand{\mapsto} H^n(X, \hand{\pi}) \hand{\simeq} \mathrm{Hom}_{\mathrm{Hotab}}(LH.(X), \hand{\pi}[n]) \quad . \]
The functor \var{\deleted{K}}{}{deletion}{ps_134-7_066-v0027}\var{\inserted{$K_\pi$}}{\hand{$K_\pi$}}{insertion}{ps_134-7_066-v0028} may be called the "\uline{Eilenberg-Maclane functor}",  as it's
values are immediate generalizations of the Eilenberg-Maclane objects
\var{\deleted{\unsure{K($\hand{\pi}$,n)}}}{}{deletion}{ps_134-7_066-v0029}.

\margin{66}

% Visual collation: both supplied image halves inspected; see CHECK_DECISIONS_134_7_8.tsv.

%% END SOURCE ps_134-7_066
\ifnum\EditionMode=1
\psapparatusstart
\psentry{ps_134-7_066-v0001}{1}{typewriter deletion}{\xout{a}}{}{The typewriter strikes out this material; it leaves the corrected running text.}
\psentry{ps_134-7_066-v0002}{2}{typewriter deletion}{\xout{may}}{}{The typewriter strikes out this material; it leaves the corrected running text.}
\psentry{ps_134-7_066-v0003}{3}{deletion}{\deleted{very much}}{}{The author cancels this material by hand; it leaves the corrected running text.}
\psentry{ps_134-7_066-v0004}{4}{insertion}{\inserted{many}}{\hand{many}}{The handwritten insertion is incorporated at its indicated position; insertion brackets leave the running text.}
\psentry{ps_134-7_066-v0005}{5}{typewriter deletion}{\xout{\unsure{w}}}{}{The typewriter strikes out this material; it leaves the corrected running text.}
\psentry{ps_134-7_066-v0006}{6}{typewriter overwriting}{\over[typed]{,}{.}}{.}{The superposed typed reading replaces the underlying reading.}
\psentry{ps_134-7_066-v0007}{7}{typewriter deletion}{\xout{and}}{}{The typewriter strikes out this material; it leaves the corrected running text.}
\psentry{ps_134-7_066-v0008}{8}{typewriter deletion}{\xout{chain}}{}{The typewriter strikes out this material; it leaves the corrected running text.}
\psentry{ps_134-7_066-v0009}{9}{deletion}{$\deleted{\illegible{1}}$}{}{The author cancels this material by hand; it leaves the corrected running text.}
\psentry{ps_134-7_066-v0010}{10}{deletion}{\deleted{\illegible{1}.}}{}{The author cancels this material by hand; it leaves the corrected running text.}
\psentry{ps_134-7_066-v0011}{11}{insertion}{\inserted{LH.}}{\hand{LH.}}{The handwritten insertion is incorporated at its indicated position; insertion brackets leave the running text.}
\psentry{ps_134-7_066-v0012}{12}{deletion}{\deleted{DP}}{}{The author cancels this material by hand; it leaves the corrected running text.}
\psentry{ps_134-7_066-v0013}{13}{insertion}{\inserted{$K_\pi$}}{\hand{$K_\pi$}}{The handwritten insertion is incorporated at its indicated position; insertion brackets leave the running text.}
\psentry{ps_134-7_066-v0014}{14}{insertion}{\inserted{\deleted{$K_\pi$}}}{\hand{}}{The handwritten insertion is incorporated at its indicated position; insertion brackets leave the running text.}
\psentry{ps_134-7_066-v0015}{15}{typed insertion}{\inserted[typed]{defined via}}{defined via}{The typed insertion is incorporated at its indicated position; insertion brackets leave the running text.}
\psentry{ps_134-7_066-v0016}{16}{typewriter deletion}{\xout{homotopy type}}{}{The typewriter strikes out this material; it leaves the corrected running text.}
\psentry{ps_134-7_066-v0017}{17}{insertion}{\inserted{(up to equivalence)}}{\hand{(up to equivalence)}}{The handwritten insertion is incorporated at its indicated position; insertion brackets leave the running text.}
\psentry{ps_134-7_066-v0018}{18}{typewriter deletion}{\xout{DP is right}}{}{The typewriter strikes out this material; it leaves the corrected running text.}
\psentry{ps_134-7_066-v0019}{19}{typed insertion}{\inserted[typed]{W is left}}{W is left}{The typed insertion is incorporated at its indicated position; insertion brackets leave the running text.}
\psentry{ps_134-7_066-v0020}{20}{writing over}{\over{K}{$K_\pi$}}{\hand{$K_\pi$}}{The superposed handwritten reading replaces the underlying reading.}
\psentry{ps_134-7_066-v0021}{21}{writing over}{\over{to}{for}}{\hand{for}}{The superposed handwritten reading replaces the underlying reading.}
\psentry{ps_134-7_066-v0022}{22}{insertion}{\inserted{ing}}{\hand{ing}}{The handwritten insertion is incorporated at its indicated position; insertion brackets leave the running text.}
\psentry{ps_134-7_066-v0023}{23}{writing over}{\over{o}{\unsure{\v{o}}}}{\hand{\unsure{\v{o}}}}{The superposed handwritten reading replaces the underlying reading.}
\psentry{ps_134-7_066-v0024}{24}{deletion}{\deleted{K}}{}{The author cancels this material by hand; it leaves the corrected running text.}
\psentry{ps_134-7_066-v0025}{25}{insertion}{\inserted{$K_\pi$}}{\hand{$K_\pi$}}{The handwritten insertion is incorporated at its indicated position; insertion brackets leave the running text.}
\psentry{ps_134-7_066-v0026}{26}{insertion}{\inserted{$\pi[n]$}}{\hand{$\pi[n]$}}{The handwritten insertion is incorporated at its indicated position; insertion brackets leave the running text.}
\psentry{ps_134-7_066-v0027}{27}{deletion}{\deleted{K}}{}{The author cancels this material by hand; it leaves the corrected running text.}
\psentry{ps_134-7_066-v0028}{28}{insertion}{\inserted{$K_\pi$}}{\hand{$K_\pi$}}{The handwritten insertion is incorporated at its indicated position; insertion brackets leave the running text.}
\psentry{ps_134-7_066-v0029}{29}{deletion}{\deleted{\unsure{K($\hand{\pi}$,n)}}}{}{The author cancels this material by hand; it leaves the corrected running text.}
\psapparatusend
\fi
\pssource{ps_134-7_067}{134-7/67}{https://grothendieck.umontpellier.fr/web/viewer.html?file=../134-7.pdf\#page=68}
\psnumber{ps_134-7_067-v0001}{1}
\psnumber{ps_134-7_067-v0002}{2}
\psnumber{ps_134-7_067-v0003}{3}
\psnumber{ps_134-7_067-v0004}{4}
\psnumber{ps_134-7_067-v0005}{5}
\psnumber{ps_134-7_067-v0006}{6}
\psnumber{ps_134-7_067-v0007}{7}
\psnumber{ps_134-7_067-v0008}{8}
\psnumber{ps_134-7_067-v0009}{9}
\psnumber{ps_134-7_067-v0010}{10}
\psnumber{ps_134-7_067-v0011}{11}
\psnumber{ps_134-7_067-v0012}{12}
\psnumber{ps_134-7_067-v0013}{13}
\psnumber{ps_134-7_067-v0014}{14}
%% BEGIN SOURCE ps_134-7_067
% Pursuing Stacks (cotes 134-2 to 134-8): archive cote 134-7, site page 67 (PDF page 68), typescript page 509.
% typescript. Draft compared: ps6.tex:3282-3365.
% AI-transcribed by Claude Opus 5.5 (claude-opus-5-5, adaptive thinking, effort high), Message Batches API, batch ps4 (msgbatch_012Aoc1NuskATPsyBeKSYK7J), 2026-11-12; reviewed against both supplied scan halves; see CHECK_DECISIONS.tsv.
\scanpage{134-7/67}
\mspage{509}
K(\hand{$\pi$},n). As any object in H\var{\over{o}{\unsure{$\theta$}}}{\hand{\unsure{$\theta$}}}{writing over}{ps_134-7_067-v0001}tab \var{\xout{can be written}}{}{typewriter deletion}{ps_134-7_067-v0002} is isomorphic to a
\margin{direct argument}product of objects $\hand{\pi}_n[n]$, it follows that \var{\inserted[typed]{in order}}{in order}{typed insertion}{ps_134-7_067-v0003} \var{\xout{the}}{}{typewriter deletion}{ps_134-7_067-v0004} to check the adjunction
formula between LH. and E it is enough to do so for objects in Hotab
of the type $\hand{\pi}[n]$, which is the ``familiar fact'' just recalled. The nota-%
tion $K_{\hand{\pi}}$ in (2) is meant to suggest the Eilenberg-Maclane K(\hand{$\pi$},n) objec\unsure{t}
generalized by the objects $K_{\hand{\pi}}(L.)$ , and also to recall that we recover
the homology invariants of \var{\xout{K(L.)}}{}{typewriter deletion}{ps_134-7_067-v0005} L. from $K_{\hand{\pi}}(L.)$ \var{\inserted[typed]{via the $\hand{\pi}_i$ invariants,}}{via the $\hand{\pi}_i$ invariants,}{typed insertion}{ps_134-7_067-v0006} by the formula
\[
(4) \qquad \hand{\pi}_i(K_{\hand{\pi}}(L.)) \hand{\simeq} H_i(L.) \quad ,
\]
which implies by the way that the functor $K_{\hand{\pi}}$ is ``conservative'', i.e.
a map in Hotab which is transformed into an isomorphism is an isomorphism\unsure{.}
This should not be confused with the stronger property of being fully
faithful, or equivalently of the left adjoint LH. being a lo\var{\over{\unsure{d}}{c}}{\hand{c}}{writing over}{ps_134-7_067-v0007}alization
functor, or equivalently still, the adjunction morphism
\[
(5) \qquad LH.(K_{\hand{\pi}}(L.)) \hand{\to} L.
\]
being an isomorphism in H\var{\over{o}{\u{o}}}{\hand{\u{o}}}{writing over}{ps_134-7_067-v0008}tab, which is definitely false !

The other adjunction morphism
\[
(6) \qquad X \hand{\to} K_{\hand{\pi}}(LH.(X)) \overset{\mathrm{def}}{=} X_{ab}
\]
is still more interesting, it's effect on the homotopy invariants $\hand{\pi}_i$ \var{\deleted{are}}{}{deletion}{ps_134-7_067-v0009}\var{\inserted{are}}{\hand{are}}{insertion}{ps_134-7_067-v0010}
the Hurewicz \var{\xout{maps}}{}{typewriter deletion}{ps_134-7_067-v0011} homomorphisms
\[
(7) \qquad \hand{\pi}_i(X) \hand{\to} H_i(X) \overset{\mathrm{def}}{=} \hand{\pi}_i(X_{ab}) \var{\deleted{\overset{\mathrm{def}}{=}}}{}{deletion}{ps_134-7_067-v0012}\hand{\simeq} H_i(LH.(X)) \quad ;
\]
\var{\deleted{thus}}{}{deletion}{ps_134-7_067-v0013} introducing the homotopy fiber of (6) (in the case of pointed
homotopy types) and denoting by $\hand{\gamma}_i(X)$ it's homotopy invariants, we get
the exact sequence of J.H.C. Whitehead (as recalled in a letter from R.
Brown I just got)
\[
\hand{(7')} \qquad \dots \hand{\to} \hand{\gamma}_i(X) \hand{\to} \hand{\pi}_i(X) \hand{\to} \hand{\widetilde{\phantom{H}}}\!\!H_i(X) \hand{\to} \hand{\gamma}_{i-1}(X) \hand{\to} \quad \dots \quad ,
\]
where $\hand{\widetilde{\phantom{H}}}\!\!H_i$ denotes the ``reduced'' homology group of a \uline{pointed} homotopy
type, equal to $H_i$ for $i\hand{\neq}0$ and to Coker $(H_o(pt) \hand{\to} H_o(X)) \hand{\simeq} H_o(X)/\underline{Z}$ for
i = 0.

The case of pointed homotopy types seems of importance for schema-%
tic homotopy types, and deserves some extra mention and care. We may
factor diagram (2) into
\margin{\textcircled{$\widetilde{DP}$}}
% Hand-drawn arrows are marked blue; typed objects remain plain.
\diagram{
\begin{tikzcd}
(8) \qquad \hand{\Delta^{\wedge}} \ar[color=blue,r, shift left, "\hand{\alpha}"] \& \hand{\Delta^{\wedge\bullet}} \ar[color=blue,l, shift left, "\hand{\beta}"] \ar[color=blue,r, shift left, "\hand{\widetilde{W}}"] \& \hand{\Delta^{\wedge}}_{ab} \ar[color=blue,l, shift left, "\hand{\widetilde{DP}}"]
\end{tikzcd} \quad ,
}
where $\hand{\Delta^{\wedge\bullet}}$ is the category of \uline{pointed} semisimplicial complexes, \hand{$\beta$} the
\var{\xout{the}}{}{typewriter deletion}{ps_134-7_067-v0014} forgetful functor from these to nonpointed complexes, and \hand{$\alpha$} it's
left adjoint, which may be interpreted as

\margin{67}

% Visual collation: both supplied image halves inspected; see CHECK_DECISIONS_134_7_8.tsv.

%% END SOURCE ps_134-7_067
\ifnum\EditionMode=1
\psapparatusstart
\psentry{ps_134-7_067-v0001}{1}{writing over}{\over{o}{\unsure{$\theta$}}}{\hand{\unsure{$\theta$}}}{The superposed handwritten reading replaces the underlying reading.}
\psentry{ps_134-7_067-v0002}{2}{typewriter deletion}{\xout{can be written}}{}{The typewriter strikes out this material; it leaves the corrected running text.}
\psentry{ps_134-7_067-v0003}{3}{typed insertion}{\inserted[typed]{in order}}{in order}{The typed insertion is incorporated at its indicated position; insertion brackets leave the running text.}
\psentry{ps_134-7_067-v0004}{4}{typewriter deletion}{\xout{the}}{}{The typewriter strikes out this material; it leaves the corrected running text.}
\psentry{ps_134-7_067-v0005}{5}{typewriter deletion}{\xout{K(L.)}}{}{The typewriter strikes out this material; it leaves the corrected running text.}
\psentry{ps_134-7_067-v0006}{6}{typed insertion}{\inserted[typed]{via the $\hand{\pi}_i$ invariants,}}{via the $\hand{\pi}_i$ invariants,}{The typed insertion is incorporated at its indicated position; insertion brackets leave the running text.}
\psentry{ps_134-7_067-v0007}{7}{writing over}{\over{\unsure{d}}{c}}{\hand{c}}{The superposed handwritten reading replaces the underlying reading.}
\psentry{ps_134-7_067-v0008}{8}{writing over}{\over{o}{\u{o}}}{\hand{\u{o}}}{The superposed handwritten reading replaces the underlying reading.}
\psentry{ps_134-7_067-v0009}{9}{deletion}{\deleted{are}}{}{The author cancels this material by hand; it leaves the corrected running text.}
\psentry{ps_134-7_067-v0010}{10}{insertion}{\inserted{are}}{\hand{are}}{The handwritten insertion is incorporated at its indicated position; insertion brackets leave the running text.}
\psentry{ps_134-7_067-v0011}{11}{typewriter deletion}{\xout{maps}}{}{The typewriter strikes out this material; it leaves the corrected running text.}
\psentry{ps_134-7_067-v0012}{12}{deletion}{$\deleted{\overset{\mathrm{def}}{=}}$}{}{The author cancels this material by hand; it leaves the corrected running text.}
\psentry{ps_134-7_067-v0013}{13}{deletion}{\deleted{thus}}{}{The author cancels this material by hand; it leaves the corrected running text.}
\psentry{ps_134-7_067-v0014}{14}{typewriter deletion}{\xout{the}}{}{The typewriter strikes out this material; it leaves the corrected running text.}
\psapparatusend
\fi
\pssource{ps_134-7_068}{134-7/68}{https://grothendieck.umontpellier.fr/web/viewer.html?file=../134-7.pdf\#page=69}
\psnumber{ps_134-7_068-v0001}{1}
\psnumber{ps_134-7_068-v0002}{2}
\psnumber{ps_134-7_068-v0003}{3}
\psnumber{ps_134-7_068-v0004}{4}
\psnumber{ps_134-7_068-v0005}{5}
\psnumber{ps_134-7_068-v0006}{6}
\psnumber{ps_134-7_068-v0007}{7}
\psnumber{ps_134-7_068-v0008}{8}
%% BEGIN SOURCE ps_134-7_068
% Pursuing Stacks (cotes 134-2 to 134-8): archive cote 134-7, site page 68 (PDF page 69), typescript page 510.
% typescript. Draft compared: ps6.tex:3348-3431.
% AI-transcribed by Claude Opus 5.5 (claude-opus-5-5, adaptive thinking, effort high), Message Batches API, batch ps4 (msgbatch_012Aoc1NuskATPsyBeKSYK7J), 2026-11-12; reviewed against both supplied scan halves; see CHECK_DECISIONS.tsv.
\scanpage{134-7/68}
\mspage{510}

(9) \qquad $\hand{\alpha}(X_{\hand{*}}) = X_* \hand{\amalg} e_*$ \qquad ,

where $e_*$ is the \var{\xout{unit}}{}{typewriter deletion}{ps_134-7_068-v0001} final objedt of $\hand{\Delta}^{\hand{\wedge}}$ , and the second member is
pointed by \unsure{it's|its} summand $e_*$ . The functor $\overset{\hand{\sim}}{\mathrm{DP}}$ comes from applying compo-%
nentwise the \var{\inserted[typed]{obvious}}{obvious}{typed insertion}{ps_134-7_068-v0002} functor from (Ab) to (Sets$^\hand{\bullet}$) (pointed sets)\var{\over{.}{,}}{\hand{,}}{writing over}{ps_134-7_068-v0003} \var{\inserted{$\overset{\sim}{W}$ is it's left adjoint.}}{\hand{$\overset{\sim}{W}$ is it's left adjoint.}}{insertion}{ps_134-7_068-v0004} Passing to
the suitable localized categories, we get from (8)

% Hand-drawn arrows are marked blue; typed objects remain plain.
(10) \diagram{\begin{tikzcd}
\mathrm{Hot} \ar[color=blue,r, shift left, "\hand{\alpha}"] \ar[color=blue,rr, bend left=40, "\hand{\mathrm{LH.}}"] \& \mathrm{Hot}^\hand{\bullet} \ar[color=blue,l, shift left, "\hand{\beta}"] \ar[color=blue,r, shift left, "\overset{\hand{\sim}}{\mathrm{LH.}}"] \& \mathrm{Hotab} \ar[color=blue,l, shift left, "\overset{\hand{\sim}}{K}_{\hand{\pi}}"] \ar[color=blue,ll, bend left=40, "\hand{K_\pi}"]
\end{tikzcd}} \qquad ,

factoring (2), where \hand{$\alpha$} is now defined by the formula similar to (9)

(9') \qquad $\hand{\alpha}(X) = X \hand{\amalg} e$ \quad ,

where e denotes the final object of H\unsure{o}t , and is used for defining the
pointed structure of the second member. The functor $\overset{\hand{\sim}}{W}$ in (8) can be
described (as is seen componentwise) as

(11) \qquad $\overset{\hand{\sim}}{W}\ (X_*) = W(\hand{\beta}(X_*))/ \mathrm{Im}\, W(e_*)$ \quad ,

where
\[ W(e_*) \hand{\longrightarrow} W(X_*) \qquad \text{i.e.} \qquad \underline{Z}^{(e\ )} \hand{\longrightarrow} \underline{Z}^{(X\ )} \]
is deduced from the pointing map $e_* \hand{\longrightarrow} X_*$ . Accordingly we get
an expression

\var{\xout{(12)}}{}{typewriter deletion}{ps_134-7_068-v0005} \qquad $\overset{\hand{\sim}}{\mathrm{LH}}_{\hand{\bullet}}(X) \hand{\simeq} \mathrm{LH}_{\hand{\bullet}}(X)/\mathrm{LH}_{\hand{\bullet}}(e)$ \qquad ,

more accurately an exact triangle

\var{\inserted{(12)}}{\hand{(12)}}{insertion}{ps_134-7_068-v0006} \var{\deleted{(12)}}{}{deletion}{ps_134-7_068-v0007} \diagram{\begin{tikzcd}[column sep=small]
\& \overset{\hand{\sim}}{\mathrm{LH}}.(X) \ar[color=blue,dl, "\hand{}"'] \& \\
\mathrm{LH}.(e)=Z\hand{[}0\hand{]} \ar[color=blue,rr] \& \& \mathrm{LH}.(X) \ar[color=blue,ul, "\hand{}"']
\end{tikzcd}} \qquad ,

where X is any object in Hot$^\hand{\bullet}$ and LH.(X) is short for LH.($\hand{\beta}$(X)). Of
course the functors $\overset{\hand{\sim}}{\mathrm{LH}}$ and $\overset{\hand{\sim}}{K}_{\hand{\pi}}$ are still adjoint\var{\inserted[typed]{(one hopes !)}}{(one hopes !)}{typed insertion}{ps_134-7_068-v0008}, hence for any pointed
homotopy type X an adjunction map in H\unsure{o}t$^\hand{\bullet}$

(6') \qquad $X \hand{\longrightarrow} \overset{\hand{\sim}}{K}\ (\overset{\hand{\sim}}{\mathrm{LH}}.(X))$ \quad ,

and (7) and (7') are deduced from (6') and it's homotopy fiber, rather
than from (6) where it does'nt really make sense because of lack of cano-%
nical base points for taking $\hand{\gamma}_i$'s and homotopy fibers.

\margin{68}

% Visual collation: both supplied image halves inspected; see CHECK_DECISIONS_134_7_8.tsv.

%% END SOURCE ps_134-7_068
\ifnum\EditionMode=1
\psapparatusstart
\psentry{ps_134-7_068-v0001}{1}{typewriter deletion}{\xout{unit}}{}{The typewriter strikes out this material; it leaves the corrected running text.}
\psentry{ps_134-7_068-v0002}{2}{typed insertion}{\inserted[typed]{obvious}}{obvious}{The typed insertion is incorporated at its indicated position; insertion brackets leave the running text.}
\psentry{ps_134-7_068-v0003}{3}{writing over}{\over{.}{,}}{\hand{,}}{The superposed handwritten reading replaces the underlying reading.}
\psentry{ps_134-7_068-v0004}{4}{insertion}{\inserted{$\overset{\sim}{W}$ is it's left adjoint.}}{\hand{$\overset{\sim}{W}$ is it's left adjoint.}}{The handwritten insertion is incorporated at its indicated position; insertion brackets leave the running text.}
\psentry{ps_134-7_068-v0005}{5}{typewriter deletion}{\xout{(12)}}{}{The typewriter strikes out this material; it leaves the corrected running text.}
\psentry{ps_134-7_068-v0006}{6}{insertion}{\inserted{(12)}}{\hand{(12)}}{The handwritten insertion is incorporated at its indicated position; insertion brackets leave the running text.}
\psentry{ps_134-7_068-v0007}{7}{deletion}{\deleted{(12)}}{}{The author cancels this material by hand; it leaves the corrected running text.}
\psentry{ps_134-7_068-v0008}{8}{typed insertion}{\inserted[typed]{(one hopes !)}}{(one hopes !)}{The typed insertion is incorporated at its indicated position; insertion brackets leave the running text.}
\psapparatusend
\fi
\pssource{ps_134-7_069}{134-7/69}{https://grothendieck.umontpellier.fr/web/viewer.html?file=../134-7.pdf\#page=70}
\psnumber{ps_134-7_069-v0001}{1}
\psnumber{ps_134-7_069-v0002}{2}
\psnumber{ps_134-7_069-v0003}{3}
\psnumber{ps_134-7_069-v0004}{4}
\psnumber{ps_134-7_069-v0005}{5}
\psnumber{ps_134-7_069-v0006}{6}
\psnumber{ps_134-7_069-v0007}{7}
\psnumber{ps_134-7_069-v0008}{8}
\psnumber{ps_134-7_069-v0009}{9}
\psnumber{ps_134-7_069-v0010}{10}
\psnumber{ps_134-7_069-v0011}{11}
\psnumber{ps_134-7_069-v0012}{12}
\psnumber{ps_134-7_069-v0013}{13}
\psnumber{ps_134-7_069-v0014}{14}
\psnumber{ps_134-7_069-v0015}{15}
\psnumber{ps_134-7_069-v0016}{16}
\psnumber{ps_134-7_069-v0017}{17}
\psnumber{ps_134-7_069-v0018}{18}
\psnumber{ps_134-7_069-v0019}{19}
\psnumber{ps_134-7_069-v0020}{20}
\psnumber{ps_134-7_069-v0021}{21}
\psnumber{ps_134-7_069-v0022}{22}
%% BEGIN SOURCE ps_134-7_069
% Pursuing Stacks (cotes 134-2 to 134-8): archive cote 134-7, site page 69 (PDF page 70), typescript page 511.
% typescript. Draft compared: ps6.tex:3411-3489.
% AI-transcribed by Claude Opus 5.5 (claude-opus-5-5, adaptive thinking, effort high), Message Batches API, batch ps4 (msgbatch_012Aoc1NuskATPsyBeKSYK7J), 2026-11-12; reviewed against both supplied scan halves; see CHECK_DECISIONS.tsv.
\scanpage{134-7/69}
\mspage{511}

2.10.

\var{\xout{122}}{}{typewriter deletion}{ps_134-7_069-v0001}
122. Since the last notes, I have been doing three day\var{\over{s}{\unsure{s'}}}{\hand{\unsure{s'}}}{writing over}{ps_134-7_069-v0002} \var{\inserted{(including today's)}}{\hand{(including today's)}}{insertion}{ps_134-7_069-v0003} scratchwork
on various questions around abelianization in general (for discrete homoto%
py types) and on Postnikov dévissage, in connection with the review on
some main formal properties of abelianization, started upon in \var{\inserted{the previous}}{\hand{the previous}}{insertion}{ps_134-7_069-v0004} section \var{\xout{12}}{}{typewriter deletion}{ps_134-7_069-v0005}
121. I did'nt get anything really new for me, rather it was just part of
the necessary rubbing against the things, in order to get a better fee\var{\xout{lin}}{}{typewriter deletion}{ps_134-7_069-v0006}\hand{-}%
ling of what they are like, or what \var{\xout{is likely}}{}{typewriter deletion}{ps_134-7_069-v0007} they are likely to be like
- what is likely to be true\hand{,} and what not. T\var{\over[typed]{h}{h}}{h}{typewriter overwriting}{ps_134-7_069-v0008}e most interesting maybe is
that I got an inkling of a fairly general version of a Kan-Dold-Puppe kind
of relationship, \var{\inserted{in terms of derived categories,}}{\hand{in terms of derived categories,}}{insertion}{ps_134-7_069-v0009} valid presumably for any \var{\inserted[typed]{local}}{local}{typed insertion}{ps_134-7_069-v0010} test category, and in particu-%
lar \var{\inserted[typed]{for}}{for}{typed insertion}{ps_134-7_069-v0011} \var{\xout{in}}{}{typewriter deletion}{ps_134-7_069-v0012} categories like \hand{$\Delta$}$_{/X}$ , with X in \hand{$\Delta$}$^{\hand{\wedge}}$. It would be untimely \var{\xout{thou}}{}{typewriter deletion}{ps_134-7_069-v0013}
though to build up still more the (already pretty high) \var{\inserted[typed]{tower}}{tower}{typed insertion}{ps_134-7_069-v0014} \var{\xout{\illegible{1}}}{}{typewriter deletion}{ps_134-7_069-v0015} of digres-%
sions, and for the time being I'll stick to what is relevant strictly
to my \var{\deleted{present}}{}{deletion}{ps_134-7_069-v0016} \var{\inserted{immediate}}{\hand{immediate}}{insertion}{ps_134-7_069-v0017} puppose - namely getting through with the wishful thinking
about schematization ! Thus I'll be content to work wi th the category
\hand{$\Delta$} in order to construct models for homotopy types and perform constructi-%
ons with them(such as abelianization), without getting involved at pre\var{\xout{s}}{}{typewriter deletion}{ps_134-7_069-v0018}-%
sent \var{\xout{\unsure{at}}}{}{typewriter deletion}{ps_134-7_069-v0019} in looking up \var{\xout{\unsure{on}}}{}{typewriter deletion}{ps_134-7_069-v0020} how much is going over(and how)to the case of
more general small categories A ...

It occured to me that the variant $\overset{\hand{\sim}}{W}$ or L$\overset{\hand{\sim}}{\mathrm{H}}$. \var{\deleted{which I introduced}}{}{deletion}{ps_134-7_069-v0021}
for the abelianization functors W and LH. , \var{\inserted{introduced in previous section}}{\hand{introduced in previous section}}{insertion}{ps_134-7_069-v0022} using a pointed structure
for the argument X in \hand{$\Delta$}$^{\hand{\wedge}}$ or Hot$^{\hand{\bullet}}$ , could beadvantageously defined with-%
out this extra structure. The construction for $\overset{\hand{\sim}}{W}$ which follows goes
through indeed in any topos whatever (not only \hand{$\Delta$}$^{\hand{\wedge}}$). Let X be an object
in \hand{$\Delta$}$^{\hand{\wedge}}$, then there is a canonical augmentation map
\[ \hand{\varepsilon} : W(X) = \underline{Z}^{(X)} \hand{\longrightarrow} \underline{Z}_{\hand{\Delta}} \]
towards the constant semisimplicial group $\underline{Z}_{\hand{\Delta}}$ , corresponding to the
constant map
\[ X \hand{\longrightarrow} \underline{Z}_{\hand{\Delta}} \]
with value 1 . We thus get a functorial exact sequence
(1) \[ 0 \hand{\longrightarrow} \overset{\hand{\sim}}{W}(X) \hand{\longrightarrow} W(X) \hand{\xrightarrow{\varepsilon}} \underline{Z} \hand{\longrightarrow} 0 \quad , \]
where $\overset{\hand{\sim}}{W}(X)$ denotes the kernel of the augmentation above, hence an exten-%
sion of $\underline{Z}$ by $\overset{\hand{\sim}}{W}(X)$ , or (what amounts essentially to the same) a torso
under the group object $\overset{\hand{\sim}}{W}(X)$, which we'll denote by
(2) \[ W(X)(1) = \hand{\varepsilon}^{-1}(1) \quad . \]

\margin{69}

% Visual collation: both supplied image halves inspected; see CHECK_DECISIONS_134_7_8.tsv.

%% END SOURCE ps_134-7_069
\ifnum\EditionMode=1
\psapparatusstart
\psentry{ps_134-7_069-v0001}{1}{typewriter deletion}{\xout{122}}{}{The typewriter strikes out this material; it leaves the corrected running text.}
\psentry{ps_134-7_069-v0002}{2}{writing over}{\over{s}{\unsure{s'}}}{\hand{\unsure{s'}}}{The superposed handwritten reading replaces the underlying reading.}
\psentry{ps_134-7_069-v0003}{3}{insertion}{\inserted{(including today's)}}{\hand{(including today's)}}{The handwritten insertion is incorporated at its indicated position; insertion brackets leave the running text.}
\psentry{ps_134-7_069-v0004}{4}{insertion}{\inserted{the previous}}{\hand{the previous}}{The handwritten insertion is incorporated at its indicated position; insertion brackets leave the running text.}
\psentry{ps_134-7_069-v0005}{5}{typewriter deletion}{\xout{12}}{}{The typewriter strikes out this material; it leaves the corrected running text.}
\psentry{ps_134-7_069-v0006}{6}{typewriter deletion}{\xout{lin}}{}{The typewriter strikes out this material; it leaves the corrected running text.}
\psentry{ps_134-7_069-v0007}{7}{typewriter deletion}{\xout{is likely}}{}{The typewriter strikes out this material; it leaves the corrected running text.}
\psentry{ps_134-7_069-v0008}{8}{typewriter overwriting}{\over[typed]{h}{h}}{h}{The superposed typed reading replaces the underlying reading.}
\psentry{ps_134-7_069-v0009}{9}{insertion}{\inserted{in terms of derived categories,}}{\hand{in terms of derived categories,}}{The handwritten insertion is incorporated at its indicated position; insertion brackets leave the running text.}
\psentry{ps_134-7_069-v0010}{10}{typed insertion}{\inserted[typed]{local}}{local}{The typed insertion is incorporated at its indicated position; insertion brackets leave the running text.}
\psentry{ps_134-7_069-v0011}{11}{typed insertion}{\inserted[typed]{for}}{for}{The typed insertion is incorporated at its indicated position; insertion brackets leave the running text.}
\psentry{ps_134-7_069-v0012}{12}{typewriter deletion}{\xout{in}}{}{The typewriter strikes out this material; it leaves the corrected running text.}
\psentry{ps_134-7_069-v0013}{13}{typewriter deletion}{\xout{thou}}{}{The typewriter strikes out this material; it leaves the corrected running text.}
\psentry{ps_134-7_069-v0014}{14}{typed insertion}{\inserted[typed]{tower}}{tower}{The typed insertion is incorporated at its indicated position; insertion brackets leave the running text.}
\psentry{ps_134-7_069-v0015}{15}{typewriter deletion}{\xout{\illegible{1}}}{}{The typewriter strikes out this material; it leaves the corrected running text.}
\psentry{ps_134-7_069-v0016}{16}{deletion}{\deleted{present}}{}{The author cancels this material by hand; it leaves the corrected running text.}
\psentry{ps_134-7_069-v0017}{17}{insertion}{\inserted{immediate}}{\hand{immediate}}{The handwritten insertion is incorporated at its indicated position; insertion brackets leave the running text.}
\psentry{ps_134-7_069-v0018}{18}{typewriter deletion}{\xout{s}}{}{The typewriter strikes out this material; it leaves the corrected running text.}
\psentry{ps_134-7_069-v0019}{19}{typewriter deletion}{\xout{\unsure{at}}}{}{The typewriter strikes out this material; it leaves the corrected running text.}
\psentry{ps_134-7_069-v0020}{20}{typewriter deletion}{\xout{\unsure{on}}}{}{The typewriter strikes out this material; it leaves the corrected running text.}
\psentry{ps_134-7_069-v0021}{21}{deletion}{\deleted{which I introduced}}{}{The author cancels this material by hand; it leaves the corrected running text.}
\psentry{ps_134-7_069-v0022}{22}{insertion}{\inserted{introduced in previous section}}{\hand{introduced in previous section}}{The handwritten insertion is incorporated at its indicated position; insertion brackets leave the running text.}
\psapparatusend
\fi
\pssource{ps_134-7_070}{134-7/70}{https://grothendieck.umontpellier.fr/web/viewer.html?file=../134-7.pdf\#page=71}
\psnumber{ps_134-7_070-v0001}{1}
\psnumber{ps_134-7_070-v0002}{2}
\psnumber{ps_134-7_070-v0003}{3}
\psnumber{ps_134-7_070-v0004}{4}
\psnumber{ps_134-7_070-v0005}{5}
\psnumber{ps_134-7_070-v0006}{6}
\psnumber{ps_134-7_070-v0007}{7}
\psnumber{ps_134-7_070-v0008}{8}
\psnumber{ps_134-7_070-v0009}{9}
\psnumber{ps_134-7_070-v0010}{10}
\psnumber{ps_134-7_070-v0011}{11}
%% BEGIN SOURCE ps_134-7_070
% Pursuing Stacks (cotes 134-2 to 134-8): archive cote 134-7, site page 70 (PDF page 71), typescript page 512.
% typescript. Draft compared: ps6.tex:3468-3545.
% AI-transcribed by Claude Opus 5.5 (claude-opus-5-5, adaptive thinking, effort high), Message Batches API, batch ps4 (msgbatch_012Aoc1NuskATPsyBeKSYK7J), 2026-11-12; reviewed against both supplied scan halves; see CHECK_DECISIONS.tsv.
\scanpage{134-7/70}
\mspage{512}

Moreover the canonical map
\[ X \hand{\longrightarrow} W(X) = \underline{Z}^{(X)} \]
factors (by construction) through
\[ (3) \qquad X \hand{\longrightarrow} W(X)(1) \quad ( \hand{\hookrightarrow} W(X) ) , \]
and this map is universal for all maps of X into torsors under abelian
group objects of $\hand{\Delta}^{\hand{\wedge}}$ . It is an immediate consequence of Whiteheads
theorem tha\var{\over{r}{t}}{\hand{t}}{writing over}{ps_134-7_070-v0001} a weak equivalence $X \hand{\to} Y$ in $\hand{\Delta}^{\hand{\wedge}}$ induces \var{\xout{\unsure{w}}}{}{typewriter deletion}{ps_134-7_070-v0002} weak equivalen-%
ce\hand{s} in $\hand{\Delta}^{\hand{\wedge}}$
\[ (4) \qquad \overset{\hand{\sim}}{W}(X) \hand{\longrightarrow} \overset{\hand{\sim}}{W}(Y) \quad , \qquad W(X)(1) \hand{\longrightarrow} W(Y)(1) \]
(and also $W(X)(n) \hand{\to} W(Y)(n)$ for any $n \hand{\in} \underline{Z}$ ), and accordingly that
the map between the normalized chain complexes corresponding to $\overset{\hand{\sim}}{W}(X)$ and
$\overset{\hand{\sim}}{W}(Y)$, namely by definition
\[ (4') \qquad L\overset{\hand{\sim}}{H}.(X) \hand{\longrightarrow} L\overset{\hand{\sim}}{H}.(Y) \]
is a quasi-isomorphism. Thus we get \var{\xout{functors}}{}{typewriter deletion}{ps_134-7_070-v0003} a functor
\[ \overset{\hand{\sim}}{\mathrm{LH}}. : \mathrm{Hot} \hand{\longrightarrow} \mathrm{Hotab} \quad , \]
whose composition\var{\inserted{$L\widetilde{H}.\circ\beta$ ,}}{\hand{$L\widetilde{H}.\circ\beta$ ,}}{insertion}{ps_134-7_070-v0004} with the forgetful functor
\[ \hand{\beta} : \mathrm{Hot}^{\hand{\bullet}} \hand{\longrightarrow} \mathrm{Hot} \]
is \var{\xout{the funct}}{}{typewriter deletion}{ps_134-7_070-v0005} canonically isomorphic to the functor $L\overset{\hand{\sim}}{H}.$ of p.510 (formula
9). More specifically, when X is in $\hand{\Delta}^{\hand{\wedge}\hand{\bullet}}$ , then the map
\[ \underline{Z} = \underline{Z}^{(e)} \hand{\longrightarrow} \underline{Z}^{(X)} \quad \text{deduced from } e \hand{\to} X \]
defines a splitting of the extension 1, i.e. an isomorphism
\[ (5) \qquad \underline{Z}^{(X)} = W(X) \hand{\simeq} \overset{\hand{\sim}}{W}(X) \hand{\oplus} \underline{Z} \quad , \]
and accordingly $\overset{\hand{\sim}}{W}(X)$ may be interpreted, at will, as a quotient group
of W(X) (which we did on p.510)\hand{,} or as a subobject of W(X),\var{\deleted{(}}{}{deletion}{ps_134-7_070-v0006}which is the
better choice, because it makes a good sense without using any pointed
structure. Again, without using the pointed structure, we\var{\xout{r}}{}{typewriter deletion}{ps_134-7_070-v0007} get a canoni-%
cal exact trian\var{\xout{\unsure{g}}}{}{typewriter deletion}{ps_134-7_070-v0008}gle in D(Ab) interpreting (1)
% Hand-drawn arrows are marked blue; typed objects remain plain.
\diagram{
\begin{tikzcd}[column sep=small]
 \& \underline{Z} \arrow[color=blue,dl] \& \\
L\overset{\hand{\sim}}{H}.(X) \arrow[color=blue,rr] \& \& LH.(X) \arrow[color=blue,ul]
\end{tikzcd}
}
(6)

replacing ((12) p. 510), which \var{\deleted{should be replaced}}{}{deletion}{ps_134-7_070-v0009} (in the case considered
there, namely X pointed\hand{)} \var{\inserted{should be replaced}}{\hand{should be replaced}}{insertion}{ps_134-7_070-v0010} by the more precise relationship
\[ (\var{\xout{4}}{}{typewriter deletion}{ps_134-7_070-v0011}\, 5') \qquad LH.(X) \hand{\simeq} L\overset{\hand{\sim}}{H}.(X) \hand{\oplus} \underline{Z} \quad , \]
an isomorphism functorial for X in Hot\hand{$^{\bullet}$}.

\margin{70}

% Visual collation: both supplied image halves inspected; see CHECK_DECISIONS_134_7_8.tsv.

%% END SOURCE ps_134-7_070
\ifnum\EditionMode=1
\psapparatusstart
\psentry{ps_134-7_070-v0001}{1}{writing over}{\over{r}{t}}{\hand{t}}{The superposed handwritten reading replaces the underlying reading.}
\psentry{ps_134-7_070-v0002}{2}{typewriter deletion}{\xout{\unsure{w}}}{}{The typewriter strikes out this material; it leaves the corrected running text.}
\psentry{ps_134-7_070-v0003}{3}{typewriter deletion}{\xout{functors}}{}{The typewriter strikes out this material; it leaves the corrected running text.}
\psentry{ps_134-7_070-v0004}{4}{insertion}{\inserted{$L\widetilde{H}.\circ\beta$ ,}}{\hand{$L\widetilde{H}.\circ\beta$ ,}}{The handwritten insertion is incorporated at its indicated position; insertion brackets leave the running text.}
\psentry{ps_134-7_070-v0005}{5}{typewriter deletion}{\xout{the funct}}{}{The typewriter strikes out this material; it leaves the corrected running text.}
\psentry{ps_134-7_070-v0006}{6}{deletion}{\deleted{(}}{}{The author cancels this material by hand; it leaves the corrected running text.}
\psentry{ps_134-7_070-v0007}{7}{typewriter deletion}{\xout{r}}{}{The typewriter strikes out this material; it leaves the corrected running text.}
\psentry{ps_134-7_070-v0008}{8}{typewriter deletion}{\xout{\unsure{g}}}{}{The typewriter strikes out this material; it leaves the corrected running text.}
\psentry{ps_134-7_070-v0009}{9}{deletion}{\deleted{should be replaced}}{}{The author cancels this material by hand; it leaves the corrected running text.}
\psentry{ps_134-7_070-v0010}{10}{insertion}{\inserted{should be replaced}}{\hand{should be replaced}}{The handwritten insertion is incorporated at its indicated position; insertion brackets leave the running text.}
\psentry{ps_134-7_070-v0011}{11}{typewriter deletion}{$\xout{4}$}{}{The typewriter strikes out this material; it leaves the corrected running text.}
\psapparatusend
\fi
\pssource{ps_134-7_071}{134-7/71}{https://grothendieck.umontpellier.fr/web/viewer.html?file=../134-7.pdf\#page=72}
\psnumber{ps_134-7_071-v0001}{1}
\psnumber{ps_134-7_071-v0002}{2}
\psnumber{ps_134-7_071-v0003}{3}
\psnumber{ps_134-7_071-v0004}{4}
\psnumber{ps_134-7_071-v0005}{5}
\psnumber{ps_134-7_071-v0006}{6}
\psnumber{ps_134-7_071-v0007}{7}
\psnumber{ps_134-7_071-v0008}{8}
\psnumber{ps_134-7_071-v0009}{9}
\psnumber{ps_134-7_071-v0010}{10}
\psnumber{ps_134-7_071-v0011}{11}
\psnumber{ps_134-7_071-v0012}{12}
\psnumber{ps_134-7_071-v0013}{13}
\psnumber{ps_134-7_071-v0014}{14}
%% BEGIN SOURCE ps_134-7_071
% Pursuing Stacks (cotes 134-2 to 134-8): archive cote 134-7, site page 71 (PDF page 72), typescript page 513.
% typescript. Draft compared: ps6.tex:3526-3605.
% AI-transcribed by Claude Opus 5.5 (claude-opus-5-5, adaptive thinking, effort high), Message Batches API, batch ps4 (msgbatch_012Aoc1NuskATPsyBeKSYK7J), 2026-11-12; reviewed against both supplied scan halves; see CHECK_DECISIONS.tsv.
\scanpage{134-7/71}
\mspage{513}

Truth to tell, these generalities are more interesting still \var{\xout{in}}{}{typewriter deletion}{ps_134-7_071-v0001} for
an argument X in a category like $\hand{\Delta}^{\hand{\wedge}}{}_{/Y} \hand{\simeq} (\hand{\Delta}_{/Y})^{\hand{\wedge}}$ , where Y is an arbitrary
object in $\hand{\Delta}^{\hand{\wedge}}$, rather than in $\hand{\Delta}^{\hand{\wedge}}$ itself. In the latter case, X always
admits a pointed structure i.e. a section over the final object (provided
X is non-empty), hence there always exists a splitting for (1), and
hence \var{\inserted{one}}{\hand{one}}{insertion}{ps_134-7_071-v0002} for (6); whereas for an object X over Y, there does not always
exist a section over Y, and accordingly it may well happen that the exten%
% NOTE: the right edge of the scan cuts the end of this line and of some lines below
-sion similar to (1) (but taken ``over Y'') does not split, i.e. that the
torsor $W_{/Y}(X)(1)$ is not trivial. T\var{\over[typed]{h}{h}}{h}{typewriter overwriting}{ps_134-7_071-v0003}e class of this torsor is an element
\var{\xout{in}}{}{typewriter deletion}{ps_134-7_071-v0004}
\[ (7) \qquad c(X/Y)\hand{\in} H^1(Y, \overset{\hand{\sim}}{W}_{/Y}(X)) \quad , \]
a very interesting invariant indeed, giving rise to the Postnikov inva-%
riants, when Y is one of the $X_n$'s occuring in the Postnikov dévissage
of X . It was while trying to understand the precise relationship between
a cohomology group as in (7), with ``continu\var{\xout{\unsure{o}}}{}{typewriter deletion}{ps_134-7_071-v0005}ous'' coefficients (by which
I mean to suggest that $\overset{\hand{\sim}}{W}_{/Y}(X)$ is viewed intuitively as a fiberspace over
the ``space'' X, whose fibers are topological \var{\inserted[typed]{abelian}}{abelian}{typed insertion}{ps_134-7_071-v0006} groups which are by no means
discrete, but \var{\xout{\illegible{5}}}{}{typewriter deletion}{ps_134-7_071-v0007}\var{\inserted[typed]{got a bunch of non vanishing $\hand{\pi}_i$'s !),}}{got a bunch of non vanishing $\hand{\pi}_i$'s !),}{typed insertion}{ps_134-7_071-v0008} and a cohomology group
with ``discrete'' coefficients, more accurately with coefficients in a
complex of chains in the category of abelian sheaves over X i.e. over $\hand{\Delta}_{/X}$,
whose homology sheaves are definitely of a ``diserete'' nature\var{\xout{,}}{}{typewriter deletion}{ps_134-7_071-v0009} (for instanc%
they are locally constant if X$\hand{\to}$Y is a Kan fibration, and their fibers
are $\underline{Z}$-modules of finite type provided we make a mild finite-type assumpti%
on on the fibers of X$\hand{\to}$Y), that I got involved in a more general reflec-%
tion on a Kan\hand{-}Dold-Puppe type of relationship, \var{\xout{which}}{}{typewriter deletion}{ps_134-7_071-v0010} I hope to come back
to \var{\inserted[typed]{this}}{this}{typed insertion}{ps_134-7_071-v0011} when getting back to the general review o\var{\over{n}{f}}{\hand{f}}{writing over}{ps_134-7_071-v0012} linearization begun in part
V of these notes, and abruptly interrupted after section 109, when
getting \var{\xout{involved}}{}{typewriter deletion}{ps_134-7_071-v0013} caught unwittingly by the enticing mystery of schematiza-%
tion !

Another afterthought to the previous notes, is that I am going to
denote by K , $\overset{\hand{\sim}}{K}$ , the functors
\[ (7) \qquad K\colon \mathrm{Hotab} \hand{\longrightarrow} \mathrm{Hot} \quad , \quad \overset{\hand{\sim}}{K} : \mathrm{Hotab} \hand{\to} \mathrm{Hot}^{\hand{\bullet}} \hand{,} \quad \hand{K=\beta\circ\widetilde{K} ,} \]
denoted by $K_{\hand{\pi}}$ , $\overset{\hand{\sim}}{K}_{\hand{\pi}}$ in section 121. This gives a formula such as
\[ (8) \qquad K(\hand{\pi}\hand{[}n\hand{]}) = K(\hand{\pi},n) \quad , \]
where $\hand{\pi}\hand{[}n\hand{]}$ denotes the chain complex in (Ab) which is $\hand{\pi}$ placed in degree
n . I was intending to complement the \var{\xout{latter}}{}{typewriter deletion}{ps_134-7_071-v0014} former notation $K_{\hand{\pi}}$ by a
similar notation $K_H$ (where H stands for ``homology''\hand{,} as $\hand{\pi}$ was standing

\margin{71}

% Visual collation: both supplied image halves inspected; see CHECK_DECISIONS_134_7_8.tsv.

%% END SOURCE ps_134-7_071
\ifnum\EditionMode=1
\psapparatusstart
\psentry{ps_134-7_071-v0001}{1}{typewriter deletion}{\xout{in}}{}{The typewriter strikes out this material; it leaves the corrected running text.}
\psentry{ps_134-7_071-v0002}{2}{insertion}{\inserted{one}}{\hand{one}}{The handwritten insertion is incorporated at its indicated position; insertion brackets leave the running text.}
\psentry{ps_134-7_071-v0003}{3}{typewriter overwriting}{\over[typed]{h}{h}}{h}{The superposed typed reading replaces the underlying reading.}
\psentry{ps_134-7_071-v0004}{4}{typewriter deletion}{\xout{in}}{}{The typewriter strikes out this material; it leaves the corrected running text.}
\psentry{ps_134-7_071-v0005}{5}{typewriter deletion}{\xout{\unsure{o}}}{}{The typewriter strikes out this material; it leaves the corrected running text.}
\psentry{ps_134-7_071-v0006}{6}{typed insertion}{\inserted[typed]{abelian}}{abelian}{The typed insertion is incorporated at its indicated position; insertion brackets leave the running text.}
\psentry{ps_134-7_071-v0007}{7}{typewriter deletion}{\xout{\illegible{5}}}{}{The typewriter strikes out this material; it leaves the corrected running text.}
\psentry{ps_134-7_071-v0008}{8}{typed insertion}{\inserted[typed]{got a bunch of non vanishing $\hand{\pi}_i$'s !),}}{got a bunch of non vanishing $\hand{\pi}_i$'s !),}{The typed insertion is incorporated at its indicated position; insertion brackets leave the running text.}
\psentry{ps_134-7_071-v0009}{9}{typewriter deletion}{\xout{,}}{}{The typewriter strikes out this material; it leaves the corrected running text.}
\psentry{ps_134-7_071-v0010}{10}{typewriter deletion}{\xout{which}}{}{The typewriter strikes out this material; it leaves the corrected running text.}
\psentry{ps_134-7_071-v0011}{11}{typed insertion}{\inserted[typed]{this}}{this}{The typed insertion is incorporated at its indicated position; insertion brackets leave the running text.}
\psentry{ps_134-7_071-v0012}{12}{writing over}{\over{n}{f}}{\hand{f}}{The superposed handwritten reading replaces the underlying reading.}
\psentry{ps_134-7_071-v0013}{13}{typewriter deletion}{\xout{involved}}{}{The typewriter strikes out this material; it leaves the corrected running text.}
\psentry{ps_134-7_071-v0014}{14}{typewriter deletion}{\xout{latter}}{}{The typewriter strikes out this material; it leaves the corrected running text.}
\psapparatusend
\fi
\pssource{ps_134-7_072}{134-7/72}{https://grothendieck.umontpellier.fr/web/viewer.html?file=../134-7.pdf\#page=73}
\psnumber{ps_134-7_072-v0001}{1}
\psnumber{ps_134-7_072-v0002}{2}
\psnumber{ps_134-7_072-v0003}{3}
\psnumber{ps_134-7_072-v0004}{4}
\psnumber{ps_134-7_072-v0005}{5}
\psnumber{ps_134-7_072-v0006}{6}
\psnumber{ps_134-7_072-v0007}{7}
\psnumber{ps_134-7_072-v0008}{8}
\psnumber{ps_134-7_072-v0009}{9}
\psnumber{ps_134-7_072-v0010}{10}
\psnumber{ps_134-7_072-v0011}{11}
\psnumber{ps_134-7_072-v0012}{12}
\psnumber{ps_134-7_072-v0013}{13}
\psnumber{ps_134-7_072-v0014}{14}
\psnumber{ps_134-7_072-v0015}{15}
\psnumber{ps_134-7_072-v0016}{16}
\psnumber{ps_134-7_072-v0017}{17}
\psnumber{ps_134-7_072-v0018}{18}
\psnumber{ps_134-7_072-v0019}{19}
\psnumber{ps_134-7_072-v0020}{20}
\psnumber{ps_134-7_072-v0021}{21}
\psnumber{ps_134-7_072-v0022}{22}
%% BEGIN SOURCE ps_134-7_072
% Pursuing Stacks (cotes 134-2 to 134-8): archive cote 134-7, site page 72 (PDF page 73), typescript page 514.
% typescript. Draft compared: ps6.tex:3584-3660.
% AI-transcribed by Claude Opus 5.5 (claude-opus-5-5, adaptive thinking, effort high), Message Batches API, batch ps4 (msgbatch_012Aoc1NuskATPsyBeKSYK7J), 2026-11-12; reviewed against both supplied scan halves; see CHECK_DECISIONS.tsv.
\scanpage{134-7/72}
\mspage{514}

for ``homotopy'', and K means ``complex'' ), but I finally found the notation
% NOTE: the last letters of "notation" run off the right edge of the scan
S (suggestive of ``\unsure{spheres|spherex}'') more congenial, giving rise to

\margin{\textcircled{H}}(8') \qquad $S(\var{\over{\pi}{H}}{\hand{H}}{writing over}{ps_134-7_072-v0001}[n]) = S(\var{\over{\pi}{H}}{\hand{H}}{writing over}{ps_134-7_072-v0002},n)$ = sphere-like homotopy type whose $\mathrm{L}\overset{\hand{\sim}}{\mathrm{H}}.$
is isomorphic to $\var{\over{\pi}{H}}{\hand{H}}{writing over}{ps_134-7_072-v0003}[n]$.

But I am anticipating somewhat on \var{\xout{what follows}}{}{typewriter deletion}{ps_134-7_072-v0004} some of the wishful thin-%
king still ahead, involving the \var{\inserted{would-be}}{\hand{would-be}}{insertion}{ps_134-7_072-v0005} description (preferably in the schematic
set-up, but more elementarily maybe in the discrete one of ordinary homo-%
topy types) of a functor
\[ \var{\xout{\unsure{(9)}}}{}{typewriter deletion}{ps_134-7_072-v0006} \qquad S: \ \mathrm{Hotab} = \mathrm{D.(Ab)} \hand{\longrightarrow} \mathrm{Hot} \quad , \]
whose most important formal property should be
\[ \var{\xout{\unsure{(10)|(9)}}}{}{typewriter deletion}{ps_134-7_072-v0007} \qquad \mathrm{LH.}(S(L.)) \hand{\simeq} L. \quad , \]
compare with formula (8') for $L. = \hand{\pi}[n]$. However such a formula cannot
hold for any L. , i.e. an arbitrary \var{\xout{L. in D.(Ab)}}{}{typewriter deletion}{ps_134-7_072-v0008} chain complex in (Ab)
is not quasi-isomorphic to an object LH.(X) with X in (Hot), \var{\inserted{as}}{\hand{as}}{insertion}{ps_134-7_072-v0009} a necessary
condition for this is that $\mathrm{H}_o(L.)$ should be free. Thus we better \var{\xout{assume}}{}{typewriter deletion}{ps_134-7_072-v0010}
restrict to the 0-connected case, and take L. subject to \var{\xout{H}}{}{typewriter deletion}{ps_134-7_072-v0011}
\[ \mathrm{H}_o(L.) = 0 \quad \text{i.e. } L. \var{\inserted[typed]{in}}{in}{typed insertion}{ps_134-7_072-v0012}\ \mathrm{D}_{\hand{\geqslant} 1}(\mathrm{Ab}) \overset{\hand{def}}{=}\var{\xout{x}}{}{typewriter deletion}{ps_134-7_072-v0013} \mathrm{Hotab(0\text{-}conn)}\hand{,} \]
and \var{\xout{define S as a \unsure{functor}}}{}{typewriter deletion}{ps_134-7_072-v0014} \var{\inserted[typed]{describe S, or better still}}{describe S, or better still}{typed insertion}{ps_134-7_072-v0015} $\hand{\widetilde{\phantom{S}}}\!S$\hand{,} as a functor
\[ (9) \qquad \hand{\widetilde{\phantom{S}}}\!S: \ \mathrm{H\unsure{o}tab(0\text{-}conn)} \hand{\longrightarrow} \mathrm{Hot}^{\hand{\bullet}}\mathrm{(0\text{-}conn)} \hand{\hookrightarrow \mathrm{Hot}^{\bullet}}\hand{,} \]
subject to the condition
\[ (10) \qquad \mathrm{L}\overset{\hand{\sim}}{\mathrm{H}}.(\hand{\widetilde{\phantom{S}}}\!S(L.)) \hand{\simeq} L. \quad , \]
giving rise to
\[ (10') \qquad \overset{\hand{\sim}}{\mathrm{H}}_i(\hand{\widetilde{\phantom{S}}}\!S(L.)) \hand{\simeq} \mathrm{H}_i(L.) \quad , \]
\margin{\{}which should be compared to the \var{\inserted{familiar}}{\hand{familiar}}{insertion}{ps_134-7_072-v0016} formula
\[ (11) \qquad \hand{\pi}_i(\hand{\widetilde{\phantom{K}}}\!K(L.)) \hand{\simeq} \mathrm{H}_i(L.) \quad . \]

Maybe formula (10) is not quite enough for \var{\over{c}{ch}}{\hand{ch}}{writing over}{ps_134-7_072-v0017}aracterizing the functor
$\hand{\widetilde{\phantom{S}}}\!S$ up to unique isomorphism, and I confess I did'nt try to construct
in some canonical way a functor (9) satisfying (10) - and I am not \var{\inserted{quite}}{\hand{quite}}{insertion}{ps_134-7_072-v0018} sure
even whether such a functor exists. One main evidence for existence of
such a functor would be the existence and unicity (up to unique isomorp\var{\xout{h}}{}{typewriter deletion}{ps_134-7_072-v0019}
\margin{\textcircled{H}}phism) of the homotopy types $S(\var{\over{\pi}{H}}{\hand{H}}{writing over}{ps_134-7_072-v0020},n)$ in (8') - and I have a vague
remembrance of having looked through a paper, fifteen o\var{\over{f}{r}}{\hand{r}}{writing over}{ps_134-7_072-v0021} twenty years
ago, where such spaces were indeed introduced and studied. \var{\xout{Here}}{}{typewriter deletion}{ps_134-7_072-v0022} Thus
a functor $\hand{\widetilde{\phantom{S}}}\!S$ is perhaps nowadays more or less standard knowledge in homo-%

\margin{72}

% Visual collation: both supplied image halves inspected; see CHECK_DECISIONS_134_7_8.tsv.

%% END SOURCE ps_134-7_072
\ifnum\EditionMode=1
\psapparatusstart
\psentry{ps_134-7_072-v0001}{1}{writing over}{$\over{\pi}{H}$}{$\hand{H}$}{The superposed handwritten reading replaces the underlying reading.}
\psentry{ps_134-7_072-v0002}{2}{writing over}{$\over{\pi}{H}$}{$\hand{H}$}{The superposed handwritten reading replaces the underlying reading.}
\psentry{ps_134-7_072-v0003}{3}{writing over}{$\over{\pi}{H}$}{$\hand{H}$}{The superposed handwritten reading replaces the underlying reading.}
\psentry{ps_134-7_072-v0004}{4}{typewriter deletion}{\xout{what follows}}{}{The typewriter strikes out this material; it leaves the corrected running text.}
\psentry{ps_134-7_072-v0005}{5}{insertion}{\inserted{would-be}}{\hand{would-be}}{The handwritten insertion is incorporated at its indicated position; insertion brackets leave the running text.}
\psentry{ps_134-7_072-v0006}{6}{typewriter deletion}{$\xout{\unsure{(9)}}$}{}{The typewriter strikes out this material; it leaves the corrected running text.}
\psentry{ps_134-7_072-v0007}{7}{typewriter deletion}{$\xout{\unsure{(10)|(9)}}$}{}{The typewriter strikes out this material; it leaves the corrected running text.}
\psentry{ps_134-7_072-v0008}{8}{typewriter deletion}{\xout{L. in D.(Ab)}}{}{The typewriter strikes out this material; it leaves the corrected running text.}
\psentry{ps_134-7_072-v0009}{9}{insertion}{\inserted{as}}{\hand{as}}{The handwritten insertion is incorporated at its indicated position; insertion brackets leave the running text.}
\psentry{ps_134-7_072-v0010}{10}{typewriter deletion}{\xout{assume}}{}{The typewriter strikes out this material; it leaves the corrected running text.}
\psentry{ps_134-7_072-v0011}{11}{typewriter deletion}{\xout{H}}{}{The typewriter strikes out this material; it leaves the corrected running text.}
\psentry{ps_134-7_072-v0012}{12}{typed insertion}{$\inserted[typed]{in}$}{$in$}{The typed insertion is incorporated at its indicated position; insertion brackets leave the running text.}
\psentry{ps_134-7_072-v0013}{13}{typewriter deletion}{$\xout{x}$}{}{The typewriter strikes out this material; it leaves the corrected running text.}
\psentry{ps_134-7_072-v0014}{14}{typewriter deletion}{\xout{define S as a \unsure{functor}}}{}{The typewriter strikes out this material; it leaves the corrected running text.}
\psentry{ps_134-7_072-v0015}{15}{typed insertion}{\inserted[typed]{describe S, or better still}}{describe S, or better still}{The typed insertion is incorporated at its indicated position; insertion brackets leave the running text.}
\psentry{ps_134-7_072-v0016}{16}{insertion}{\inserted{familiar}}{\hand{familiar}}{The handwritten insertion is incorporated at its indicated position; insertion brackets leave the running text.}
\psentry{ps_134-7_072-v0017}{17}{writing over}{\over{c}{ch}}{\hand{ch}}{The superposed handwritten reading replaces the underlying reading.}
\psentry{ps_134-7_072-v0018}{18}{insertion}{\inserted{quite}}{\hand{quite}}{The handwritten insertion is incorporated at its indicated position; insertion brackets leave the running text.}
\psentry{ps_134-7_072-v0019}{19}{typewriter deletion}{\xout{h}}{}{The typewriter strikes out this material; it leaves the corrected running text.}
\psentry{ps_134-7_072-v0020}{20}{writing over}{$\over{\pi}{H}$}{$\hand{H}$}{The superposed handwritten reading replaces the underlying reading.}
\psentry{ps_134-7_072-v0021}{21}{writing over}{\over{f}{r}}{\hand{r}}{The superposed handwritten reading replaces the underlying reading.}
\psentry{ps_134-7_072-v0022}{22}{typewriter deletion}{\xout{Here}}{}{The typewriter strikes out this material; it leaves the corrected running text.}
\psapparatusend
\fi
\pssource{ps_134-7_073}{134-7/73}{https://grothendieck.umontpellier.fr/web/viewer.html?file=../134-7.pdf\#page=74}
\psnumber{ps_134-7_073-v0001}{1}
\psnumber{ps_134-7_073-v0002}{2}
\psnumber{ps_134-7_073-v0003}{3}
\psnumber{ps_134-7_073-v0004}{4}
\psnumber{ps_134-7_073-v0005}{5}
\psnumber{ps_134-7_073-v0006}{6}
\psnumber{ps_134-7_073-v0007}{7}
\psnumber{ps_134-7_073-v0008}{8}
\psnumber{ps_134-7_073-v0009}{9}
\psnumber{ps_134-7_073-v0010}{10}
\psnumber{ps_134-7_073-v0011}{11}
\psnumber{ps_134-7_073-v0012}{12}
\psnumber{ps_134-7_073-v0013}{13}
\psnumber{ps_134-7_073-v0014}{14}
\psnumber{ps_134-7_073-v0015}{15}
\psnumber{ps_134-7_073-v0016}{16}
\psnumber{ps_134-7_073-v0017}{17}
%% BEGIN SOURCE ps_134-7_073
% Pursuing Stacks (cotes 134-2 to 134-8): archive cote 134-7, site page 73 (PDF page 74), typescript page 515.
% typescript. Draft compared: ps6.tex:3648-3720.
% AI-transcribed by Claude Opus 5.5 (claude-opus-5-5, adaptive thinking, effort high), Message Batches API, batch ps4 (msgbatch_012Aoc1NuskATPsyBeKSYK7J), 2026-11-12; reviewed against both supplied scan halves; see CHECK_DECISIONS.tsv.
\scanpage{134-7/73}
\mspage{515}

topy theory.  \var{\xout{Of}}{}{typewriter deletion}{ps_134-7_073-v0001} I was led to postulate the existence of a functor $\overset{\hand{\sim}}{S}$  in
the context of \uline{shhematic} homotopy types,  by reasons of symmetry in the
LH.  and L\hand{$\pi$}. formalism,  which I'll try to get through a little later.
It should be clear from the putset though that in the schematic set-up,
even the mere existence of (schematic) homotopy types $S(\underline{Z},n)$  (correspon-%
ding to ordinary spheres),  and \var{\xout{\illegible{1}}}{}{typewriter deletion}{ps_134-7_073-v0002} even for $n=2$ or $3$\var{\inserted{only}}{\hand{only}}{insertion}{ps_134-7_073-v0003}, is very far from
trivial,  and as a matter of fact is'nt even known (in the set-up of
semisimplicial unipotent bundles).  This shows at the same time that if
such a functor $\overset{\hand{\sim}}{S}$ can actually be constructed in the schematic case,  it is
likely to be a lot more interesting still than in the discrete case,  as
it will presumably give information on homotopy groups of spheres (via the
desciption of the \hand{$\pi$}$_i$'s  of a semisimplicial unipotent bundle via the Lie
functor  (cf. section 118)\hand{)}.

T\var{\over[typed]{h}{h}}{h}{typewriter overwriting}{ps_134-7_073-v0004}ere  is a fourth tentative functor in between the categories
Hotab  and  Hot ,  or more accurately;  from a suitable subcategory
H\var{\over{o}{\unsure{$\breve{o}$}}}{\hand{\unsure{$\breve{o}$}}}{writing over}{ps_134-7_073-v0005}t$^\hand{\bullet}$(0)  of Hot$^\hand{\bullet}$  (corresponding to 0-connected pointed homotopy types
whose \hand{$\pi$}$_1$ is abelian and acts trivially on the higher \hand{$\pi$}$_i$'s) to
Hotab (0-conn) = H\var{\over{o}{\unsure{$\hat{o}$}}}{\hand{\unsure{$\hat{o}$}}}{writing over}{ps_134-7_073-v0006}tab(0),  strongly suggested by the schematic case of
section 118,  namely a functor

(*) \qquad $\mathrm{L}\hand{\pi}. : \mathrm{Hot}^\hand{\bullet}(0) \hand{\longrightarrow} \mathrm{Hotab}^\hand{\bullet}(0) \quad ???$  ,

whose main functorial property should be

(\var{\over{**}{12}}{\hand{12}}{writing over}{ps_134-7_073-v0007}) \qquad $H_i(\mathrm{L}\hand{\pi}.(X)) \hand{\simeq} \hand{\pi}_i(X)$ \quad \var{\deleted{(?)}}{}{deletion}{ps_134-7_073-v0008} .

To play really safe,  one might hope such a functor to be defined at any
rate for 1-connected pointed homotopy types.  Another equally important
property, suggested by the schematic set-up,\var{\inserted[typed]{as wehh ad by (11)}}{as wehh ad by (11)}{typed insertion}{ps_134-7_073-v0009}\var{\inserted{and (**),}}{\hand{and (**),}}{insertion}{ps_134-7_073-v0010} is the formula

(\var{\over{***}{13}}{\hand{13}}{writing over}{ps_134-7_073-v0011}) \qquad $\mathrm{L}\hand{\pi}.(\overset{\hand{\sim}}{K}(L.)) \hand{\simeq} L.\hand{\unsure{;}}$ \quad \var{\deleted{(?)}}{}{deletion}{ps_134-7_073-v0012} ,

similar to formula (10)  (with the pair  ($\overset{\hand{\sim}}{S}$, $\overset{\hand{\sim}}{\mathrm{LH}}$. )  replaced by the pair
($\overset{\hand{\sim}}{K}$ , L\hand{$\pi$}. )).  Of course,  for X of the type $\overset{\hand{\sim}}{K}$(L.),  (\var{\over{**}{12}}{\hand{12}}{writing over}{ps_134-7_073-v0013}) follows from
(\var{\over{***}{13}}{\hand{13}}{writing over}{ps_134-7_073-v0014}) .

3.10.

\var{\xout{Applying}}{}{typewriter deletion}{ps_134-7_073-v0015} We'll see though that a functor  L\hand{$\pi$}.  \var{\xout{(*)}}{}{typewriter deletion}{ps_134-7_073-v0016} as in (*),  satis-%
fying the properties above,  \uline{does not exist}.  Indeed,  applying L\hand{$\pi$}. to
the adjunction morphism
\[ X \hand{\longrightarrow} \overset{\hand{\sim}}{K}\ \overset{\hand{\sim}}{\mathrm{LH}}.(X) \]
(cf. p.509, (6) ),\var{\inserted{and using (13),}}{\hand{and using (13),}}{insertion}{ps_134-7_073-v0017} we should get a Hurewicz  homomorphism

\margin{73}

% Visual collation: both supplied image halves inspected; see CHECK_DECISIONS_134_7_8.tsv.

%% END SOURCE ps_134-7_073
\ifnum\EditionMode=1
\psapparatusstart
\psentry{ps_134-7_073-v0001}{1}{typewriter deletion}{\xout{Of}}{}{The typewriter strikes out this material; it leaves the corrected running text.}
\psentry{ps_134-7_073-v0002}{2}{typewriter deletion}{\xout{\illegible{1}}}{}{The typewriter strikes out this material; it leaves the corrected running text.}
\psentry{ps_134-7_073-v0003}{3}{insertion}{\inserted{only}}{\hand{only}}{The handwritten insertion is incorporated at its indicated position; insertion brackets leave the running text.}
\psentry{ps_134-7_073-v0004}{4}{typewriter overwriting}{\over[typed]{h}{h}}{h}{The superposed typed reading replaces the underlying reading.}
\psentry{ps_134-7_073-v0005}{5}{writing over}{\over{o}{\unsure{$\breve{o}$}}}{\hand{\unsure{$\breve{o}$}}}{The superposed handwritten reading replaces the underlying reading.}
\psentry{ps_134-7_073-v0006}{6}{writing over}{\over{o}{\unsure{$\hat{o}$}}}{\hand{\unsure{$\hat{o}$}}}{The superposed handwritten reading replaces the underlying reading.}
\psentry{ps_134-7_073-v0007}{7}{writing over}{\over{**}{12}}{\hand{12}}{The superposed handwritten reading replaces the underlying reading.}
\psentry{ps_134-7_073-v0008}{8}{deletion}{\deleted{(?)}}{}{The author cancels this material by hand; it leaves the corrected running text.}
\psentry{ps_134-7_073-v0009}{9}{typed insertion}{\inserted[typed]{as wehh ad by (11)}}{as wehh ad by (11)}{The typed insertion is incorporated at its indicated position; insertion brackets leave the running text.}
\psentry{ps_134-7_073-v0010}{10}{insertion}{\inserted{and (**),}}{\hand{and (**),}}{The handwritten insertion is incorporated at its indicated position; insertion brackets leave the running text.}
\psentry{ps_134-7_073-v0011}{11}{writing over}{\over{***}{13}}{\hand{13}}{The superposed handwritten reading replaces the underlying reading.}
\psentry{ps_134-7_073-v0012}{12}{deletion}{\deleted{(?)}}{}{The author cancels this material by hand; it leaves the corrected running text.}
\psentry{ps_134-7_073-v0013}{13}{writing over}{\over{**}{12}}{\hand{12}}{The superposed handwritten reading replaces the underlying reading.}
\psentry{ps_134-7_073-v0014}{14}{writing over}{\over{***}{13}}{\hand{13}}{The superposed handwritten reading replaces the underlying reading.}
\psentry{ps_134-7_073-v0015}{15}{typewriter deletion}{\xout{Applying}}{}{The typewriter strikes out this material; it leaves the corrected running text.}
\psentry{ps_134-7_073-v0016}{16}{typewriter deletion}{\xout{(*)}}{}{The typewriter strikes out this material; it leaves the corrected running text.}
\psentry{ps_134-7_073-v0017}{17}{insertion}{\inserted{and using (13),}}{\hand{and using (13),}}{The handwritten insertion is incorporated at its indicated position; insertion brackets leave the running text.}
\psapparatusend
\fi
\pssource{ps_134-7_074}{134-7/74}{https://grothendieck.umontpellier.fr/web/viewer.html?file=../134-7.pdf\#page=75}
\psnumber{ps_134-7_074-v0001}{1}
\psnumber{ps_134-7_074-v0002}{2}
\psnumber{ps_134-7_074-v0003}{3}
\psnumber{ps_134-7_074-v0004}{4}
\psnumber{ps_134-7_074-v0005}{5}
\psnumber{ps_134-7_074-v0006}{6}
\psnumber{ps_134-7_074-v0007}{7}
\psnumber{ps_134-7_074-v0008}{8}
\psnumber{ps_134-7_074-v0009}{9}
\psnumber{ps_134-7_074-v0010}{10}
\psnumber{ps_134-7_074-v0011}{11}
\psnumber{ps_134-7_074-v0012}{12}
\psnumber{ps_134-7_074-v0013}{13}
\psnumber{ps_134-7_074-v0014}{14}
\psnumber{ps_134-7_074-v0015}{15}
\psnumber{ps_134-7_074-v0016}{16}
%% BEGIN SOURCE ps_134-7_074
% Pursuing Stacks (cotes 134-2 to 134-8): archive cote 134-7, site page 74 (PDF page 75), typescript page 516.
% typescript. Draft compared: ps6.tex:3698-3774.
% AI-transcribed by Claude Opus 5.5 (claude-opus-5-5, adaptive thinking, effort high), Message Batches API, batch ps4 (msgbatch_012Aoc1NuskATPsyBeKSYK7J), 2026-11-12; reviewed against both supplied scan halves; see CHECK_DECISIONS.tsv.
\scanpage{134-7/74}
\mspage{516}

(14) \qquad $L\hand{\pi}.(X) \hand{\longrightarrow} L\overset{\hand{\sim}}{H}.(X)$

more precise than the separate homomorphisms
\[ \hand{\pi}_i(X) \hand{\to} \overset{\hand{\sim}}{H}_i(X) \ , \]
and applying this to an object $\overset{\hand{\sim}}{K}(L.)$ and applying (13) again, we should
get a functorial homomorphism\var{\xout{s}}{}{typewriter deletion}{ps_134-7_074-v0001}

(15) \qquad $L. \hand{\to} L\overset{\hand{\sim}}{H}.(\overset{\hand{\sim}}{K}(L\hand{.}))$

in opposite direction from the adjunction morphism

(1\var{\over{\unsure{5}}{6}}{\hand{6}}{writing over}{ps_134-7_074-v0002}*) \qquad $L\overset{\hand{\sim}}{H}.(\overset{\hand{\sim}}{K}(L\hand{.})) \hand{\longrightarrow} L.$

(p. 509, (5)), the composition of the two being the identity in L. .
In other words, (15) should be a canonical splitting of the natural
adjunction morphism (16). Take for instance $L. =\hand{\pi}\hand{[}n\hand{]}$, then the first
member of (16) is the Eilenberg-Maclane homology $L\overset{\hand{\sim}}{H}.(\hand{\pi},n)$, whose first
non vanishing homology group is $\hand{\pi}$ , and (16) is just the canonical
augmentation
\[ L\overset{\hand{\sim}}{H}.(\hand{\pi},n)\hand{\to}\hand{\pi}\hand{[}n\hand{]} \ , \]
and taking the $\mathrm{Ext}^i_{\underline{Z}}$ of both members with $M\hand{[}n\hand{]}$ (M any object in (Ab))
yields the transposed canonical homomorphism

(17) \qquad $\mathrm{Ext}^i_{\hand{Z}}(\hand{\pi},M) \hand{\longrightarrow} H^{\var{\over{\illegible{1}}{n+i}}{\hand{n+i}}{writing over}{ps_134-7_074-v0003}}(\hand{\pi},n; M)$ \quad ,

which is an isomorphism for $\var{\over{\illegible{1}}{i=0}}{\hand{i=0}}{writing over}{ps_134-7_074-v0004}$ (whereas for $i\hand{<} \var{\over{\illegible{1}}{0}}{\hand{0}}{writing over}{ps_134-7_074-v0005}$ both members are
zero ). For i = 1 , we get an exact sequence (``universal coefficients'')

(18) \qquad $0 \hand{\longrightarrow} \mathrm{Ext}^1_{\hand{Z}}(\hand{\pi},M) \hand{\longrightarrow} H^{\var{\over[typed]{\illegible{1}}{n+1}}{n+1}{typewriter overwriting}{ps_134-7_074-v0006}}(\hand{\pi},n;M) \hand{\longrightarrow} \mathrm{Hom}_{\hand{Z}}(\var{\over{\illegible{1}}{H_{n+1}}}{\hand{H_{n+1}}}{writing over}{ps_134-7_074-v0007}(\hand{\pi},n)\hand{, M)} \hand{\to 0}$ .

A canonical splitting (15) would \var{\over{\unsure{t}}{y}}{\hand{y}}{writing over}{ps_134-7_074-v0008}ield\var{\deleted{s}}{}{deletion}{ps_134-7_074-v0009} a canonical splitting of this
exact sequence, which definitely looks somehow as ``against nature'' !
Surely, it should'nt be hard to find, for any given inte\var{\xout{r}}{}{typewriter deletion}{ps_134-7_074-v0010}ger $n \hand{\geq} 1$,
\var{\xout{modules}}{}{typewriter deletion}{ps_134-7_074-v0011}\var{\inserted{suitable}}{\hand{suitable}}{insertion}{ps_134-7_074-v0012} abelian groups $\hand{\pi}$ and M\hand{,} such that there does not ex\var{\over{\unsure{ast}}{ist}}{\hand{ist}}{writing over}{ps_134-7_074-v0013} a split-%
ting of (18) stab\var{\xout{m}}{}{typewriter deletion}{ps_134-7_074-v0014}le under the action of the group
\[ G \ = \ \mathrm{Aut}(\hand{\pi}) \ \mathrm{x} \ \mathrm{Aut}(M) \]
(acting by ``transport de structure'' on the three terms of (18)); presu-%
mably even\hand{,} looking at Aut($\hand{\pi}$) should be enough. As I am not at all fami-%
liar with Eilenberg-Maclane cohomology, except a little in the case n=1
i.e. for usual group cohomology, I did'nt try to construct an example
for any n, only one for n=1 i.e. in a non simply connected case, which
again is a little less convincing as if it was \var{\inserted{one}}{\hand{one}}{insertion}{ps_134-7_074-v0015} for n=2...\var{\xout{Thus \illegible{1}}}{}{typewriter deletion}{ps_134-7_074-v0016}

Thus let's take n=1 , a vector space of finite dimension over
the prime field $\underline{F}_2$ , and $M = \underline{F}_2$ , in this case standard calculations
\margin{74}

% Visual collation: both supplied image halves inspected; see CHECK_DECISIONS_134_7_8.tsv.

%% END SOURCE ps_134-7_074
\ifnum\EditionMode=1
\psapparatusstart
\psentry{ps_134-7_074-v0001}{1}{typewriter deletion}{\xout{s}}{}{The typewriter strikes out this material; it leaves the corrected running text.}
\psentry{ps_134-7_074-v0002}{2}{writing over}{\over{\unsure{5}}{6}}{\hand{6}}{The superposed handwritten reading replaces the underlying reading.}
\psentry{ps_134-7_074-v0003}{3}{writing over}{$\over{\illegible{1}}{n+i}$}{$\hand{n+i}$}{The superposed handwritten reading replaces the underlying reading.}
\psentry{ps_134-7_074-v0004}{4}{writing over}{$\over{\illegible{1}}{i=0}$}{$\hand{i=0}$}{The superposed handwritten reading replaces the underlying reading.}
\psentry{ps_134-7_074-v0005}{5}{writing over}{$\over{\illegible{1}}{0}$}{$\hand{0}$}{The superposed handwritten reading replaces the underlying reading.}
\psentry{ps_134-7_074-v0006}{6}{typewriter overwriting}{$\over[typed]{\illegible{1}}{n+1}$}{$n+1$}{The superposed typed reading replaces the underlying reading.}
\psentry{ps_134-7_074-v0007}{7}{writing over}{$\over{\illegible{1}}{H_{n+1}}$}{$\hand{H_{n+1}}$}{The superposed handwritten reading replaces the underlying reading.}
\psentry{ps_134-7_074-v0008}{8}{writing over}{\over{\unsure{t}}{y}}{\hand{y}}{The superposed handwritten reading replaces the underlying reading.}
\psentry{ps_134-7_074-v0009}{9}{deletion}{\deleted{s}}{}{The author cancels this material by hand; it leaves the corrected running text.}
\psentry{ps_134-7_074-v0010}{10}{typewriter deletion}{\xout{r}}{}{The typewriter strikes out this material; it leaves the corrected running text.}
\psentry{ps_134-7_074-v0011}{11}{typewriter deletion}{\xout{modules}}{}{The typewriter strikes out this material; it leaves the corrected running text.}
\psentry{ps_134-7_074-v0012}{12}{insertion}{\inserted{suitable}}{\hand{suitable}}{The handwritten insertion is incorporated at its indicated position; insertion brackets leave the running text.}
\psentry{ps_134-7_074-v0013}{13}{writing over}{\over{\unsure{ast}}{ist}}{\hand{ist}}{The superposed handwritten reading replaces the underlying reading.}
\psentry{ps_134-7_074-v0014}{14}{typewriter deletion}{\xout{m}}{}{The typewriter strikes out this material; it leaves the corrected running text.}
\psentry{ps_134-7_074-v0015}{15}{insertion}{\inserted{one}}{\hand{one}}{The handwritten insertion is incorporated at its indicated position; insertion brackets leave the running text.}
\psentry{ps_134-7_074-v0016}{16}{typewriter deletion}{\xout{Thus \illegible{1}}}{}{The typewriter strikes out this material; it leaves the corrected running text.}
\psapparatusend
\fi
\pssource{ps_134-7_075}{134-7/75}{https://grothendieck.umontpellier.fr/web/viewer.html?file=../134-7.pdf\#page=76}
\psnumber{ps_134-7_075-v0001}{1}
\psnumber{ps_134-7_075-v0002}{2}
\psnumber{ps_134-7_075-v0003}{3}
\psnumber{ps_134-7_075-v0004}{4}
\psnumber{ps_134-7_075-v0005}{5}
\psnumber{ps_134-7_075-v0006}{6}
\psnumber{ps_134-7_075-v0007}{7}
\psnumber{ps_134-7_075-v0008}{8}
%% BEGIN SOURCE ps_134-7_075
% Pursuing Stacks (cotes 134-2 to 134-8): archive cote 134-7, site page 75 (PDF page 76), typescript page 517.
% typescript. Draft compared: ps6.tex:3759-3814.
% AI-transcribed by Claude Opus 5.5 (claude-opus-5-5, adaptive thinking, effort high), Message Batches API, batch ps4 (msgbatch_012Aoc1NuskATPsyBeKSYK7J), 2026-11-12; reviewed against both supplied scan halves; see CHECK_DECISIONS.tsv.
\scanpage{134-7/75}
\mspage{517}

show that the exact sequence (18) can be identified with the familiar
exact sequence
\[ (18') \qquad 0 \hand{\longrightarrow} V' \hand{\longrightarrow} \mathrm{Sym}^2(V') \hand{\longrightarrow} \hand{\bigwedge^{2}} V' \hand{\longrightarrow} 0 \quad , \]
where the first arrows associates to every linear form on V it's square
i.e. the same form (!) but viewed as being a quadratic form on V; and
the second arrow associates to every quadratic form on V the associated
bilinear form, which is alternate because of char. 2 . It is well
known I guess that for dim V $\hand{\geqslant}$ 2, there does not exist a splitting of
(18') which is stable under Aut(V)\hand{$\simeq$}Aut(V').

\uline{Remark}. If we admit that a similar example can be found for \var{\inserted{non-splitting of}}{\hand{non-splitting of}}{insertion}{ps_134-7_075-v0001} (18), with
n $\hand{\geqslant}$2 (which should'nt be hard I guess for someone familiar with Eilen-%
berg-Maclane homology), this shows that there does \uline{not} exist (as was
contemplated at the beginning, cf. section 111) an equivalence of catego-%
ries between shhematic 1-connected homotopy types (defined via the
model category $\hand{\mathcal{M}}_1(\underline{Z})$ of semisimplicial unipotent bundles \var{\inserted[typed]{over $\underline{Z}$}}{over $\underline{Z}$}{typed insertion}{ps_134-7_075-v0002} satisfying
$X_{\bullet}$$=X_1=e$ ) and 1-connected pointed usual homotopy types, as this would
imply the existence of a functor L\hand{$\Pi$}$_{\bullet}$ , \var{\xout{\unsure{(**)}}}{}{typewriter deletion}{ps_134-7_075-v0003} hence of a canonical splitting
of (18). The same argument will show that the once hoped-for comparison
theorem between usual Eilenberg-Maclane homology (or cohomology), and
the schematic one, is false, \var{\xout{as the}}{}{typewriter deletion}{ps_134-7_075-v0004} \var{\inserted{because}}{\hand{because}}{insertion}{ps_134-7_075-v0005} in the schematic set-up L\hand{$\Pi$}$_{\bullet}$ does
exist, and the direct construction of (15) (a \var{\inserted[typed]{functorial}}{functorial}{typed insertion}{ps_134-7_075-v0006} splitting of (16)) is
then anyhow a tautology. Thus after all we do'nt have to rely on
delicate results of Breen's on \var{\xout{\unsure{the Ext$^{\hand{i}}$ \illegible{1} of $\underline{G}$}}}{}{typewriter deletion}{ps_134-7_075-v0007} Ext$^{i}$($\underline{G}_a$,$\underline{G}_a$) over the
prime fields $\underline{F}_p$ \var{\xout{in order to use}}{}{typewriter deletion}{ps_134-7_075-v0008} (as suggested by Illusie), in order
to get this"negative" result, causing unreasonable expectations to
crash ...

\margin{75}

% Visual collation: both supplied image halves inspected; see CHECK_DECISIONS_134_7_8.tsv.

%% END SOURCE ps_134-7_075
\ifnum\EditionMode=1
\psapparatusstart
\psentry{ps_134-7_075-v0001}{1}{insertion}{\inserted{non-splitting of}}{\hand{non-splitting of}}{The handwritten insertion is incorporated at its indicated position; insertion brackets leave the running text.}
\psentry{ps_134-7_075-v0002}{2}{typed insertion}{\inserted[typed]{over $\underline{Z}$}}{over $\underline{Z}$}{The typed insertion is incorporated at its indicated position; insertion brackets leave the running text.}
\psentry{ps_134-7_075-v0003}{3}{typewriter deletion}{\xout{\unsure{(**)}}}{}{The typewriter strikes out this material; it leaves the corrected running text.}
\psentry{ps_134-7_075-v0004}{4}{typewriter deletion}{\xout{as the}}{}{The typewriter strikes out this material; it leaves the corrected running text.}
\psentry{ps_134-7_075-v0005}{5}{insertion}{\inserted{because}}{\hand{because}}{The handwritten insertion is incorporated at its indicated position; insertion brackets leave the running text.}
\psentry{ps_134-7_075-v0006}{6}{typed insertion}{\inserted[typed]{functorial}}{functorial}{The typed insertion is incorporated at its indicated position; insertion brackets leave the running text.}
\psentry{ps_134-7_075-v0007}{7}{typewriter deletion}{\xout{\unsure{the Ext$^{\hand{i}}$ \illegible{1} of $\underline{G}$}}}{}{The typewriter strikes out this material; it leaves the corrected running text.}
\psentry{ps_134-7_075-v0008}{8}{typewriter deletion}{\xout{in order to use}}{}{The typewriter strikes out this material; it leaves the corrected running text.}
\psapparatusend
\fi
\pssource{ps_134-7_076}{134-7/76}{https://grothendieck.umontpellier.fr/web/viewer.html?file=../134-7.pdf\#page=77}
\psnumber{ps_134-7_076-v0001}{1}
\psnumber{ps_134-7_076-v0002}{2}
\psnumber{ps_134-7_076-v0003}{3}
\psnumber{ps_134-7_076-v0004}{4}
\psnumber{ps_134-7_076-v0005}{5}
\psnumber{ps_134-7_076-v0006}{6}
\psnumber{ps_134-7_076-v0007}{7}
\psnumber{ps_134-7_076-v0008}{8}
\psnumber{ps_134-7_076-v0009}{9}
\psnumber{ps_134-7_076-v0010}{10}
\psnumber{ps_134-7_076-v0011}{11}
\psnumber{ps_134-7_076-v0012}{12}
\psnumber{ps_134-7_076-v0013}{13}
%% BEGIN SOURCE ps_134-7_076
% Pursuing Stacks (cotes 134-2 to 134-8): archive cote 134-7, site page 76 (PDF page 77), typescript page 518.
% typescript. Draft compared: ps6.tex:3795-3857.
% AI-transcribed by Claude Opus 5.5 (claude-opus-5-5, adaptive thinking, effort high), Message Batches API, batch ps4 (msgbatch_012Aoc1NuskATPsyBeKSYK7J), 2026-11-12; reviewed against both supplied scan halves; see CHECK_DECISIONS.tsv.
\scanpage{134-7/76}
\mspage{518}

4.10.

123.\quad For the last few days  - since I resumed the daily mathematical
ponderings (interrupted more or less for nearly one month),  the wind in
my sails has been rather slack I feel.  It has been nearly six weeks
from now that I started pondering on schematization and on schematic
homotopy types - in the process I got rid of some misplaced expectations,
very well.  Still,  the unpleasant feeling remains of not having really
any hold yet on those would-be schematic homotopy types,  due mainly to
my uncapacity so far of performing (in the model category of semisimpli-%
cial unipotent bundles say) the basic homotopy operations of taking
homotopy fibers and cofibers of maps.  In lack of this,  I am not \var{\deleted{even}}{}{deletion}{ps_134-7_076-v0001}
(``morally'') sure\var{\inserted[typed]{yet}}{yet}{typed insertion}{ps_134-7_076-v0002} that there does exist indeed a substantial \var{\deleted{the}}{}{deletion}{ps_134-7_076-v0003} reality
of the kind I have been trying to foreshadow.  I am not even wholly clear
yet of how to define the notion of ``weak equivalence'' in the model catego%
% NOTE: the right edge of this line and the next is cut off by the scan
ry $\hand{\mathcal{M}}_\bullet$(k) (or in the smaller one $\hand{\mathcal{M}}_1$(k), to play safe, as even the defini%
tion of $\hand{\mathcal{M}}_\bullet$(k) is'nt too clear yet) -  there are three ways to define
weak equivalences,  using LH. , or L\hand{$\pi$}. , or the sections functor
from $\hand{\mathcal{M}}_1$(k) to ordinary semisimplicial complexes,  and it is'nt clear
yet whether these are indeed equivalent.  We may of course call ``weak
equivalence'' in $\hand{\mathcal{M}}_1$(k)  a map which \var{\deleted{is transformed by}}{}{deletion}{ps_134-7_076-v0004}\var{\inserted{becomes}}{\hand{becomes}}{insertion}{ps_134-7_076-v0005} an isomorphism by
any of these three functors.  But I would'nt say I am wholly confident
yet that there exists \var{\xout{say}}{}{typewriter deletion}{ps_134-7_076-v0006} a localization Hot$_1$(k) of $\hand{\mathcal{M}}_1$(k)\hand{,} wi th respect
to this notion of weak equivalence or some finer one,  in such a way that
in H\var{\over[typed]{\unsure{e}}{o}}{o}{typewriter overwriting}{ps_134-7_076-v0007}t$_1$(k) one may introduce the two types of ``fibration'' and ``cofibrati-%
on'' sequences with the usual properties,  and the sections functor
\[ \mathrm{Hot}_1(k) \hand{\longrightarrow} \mathrm{Hot}_1 \]
should respect this structure,  and moreover give rise to functorial
isomorphisms
\[ H_i(L\hand{\pi}.(X_*)) \hand{\simeq} \hand{\pi}_i(X_*(k)) \qquad ; \]
nor even do I feel wholly confident that a theory of this kind can be
developped,  \var{\deleted{using}}{}{deletion}{ps_134-7_076-v0008} possibly\var{\inserted{with}}{\hand{with}}{insertion}{ps_134-7_076-v0009} different kinds of models for k-homotopy
types from the ones I have been using.  At any rate,  it seems to me \var{\deleted{that}}{}{deletion}{ps_134-7_076-v0010}
that the two kind of conditions or features I have just been stating,
are indeed the crucial ones,  plus (I should add)  existence in H\var{\over[typed]{\unsure{e}}{o}}{o}{typewriter overwriting}{ps_134-7_076-v0011}t$_1$(k)
of an object \var{\deleted{\illegible{1}}}{}{deletion}{ps_134-7_076-v0012}
\[ S^2_{\hand{k}} = S(2,\var{\over{\illegible{1}}{k}}{\hand{k}}{writing over}{ps_134-7_076-v0013}) \quad , \]
standing for the  2-sphere,  with
\[ \hand{\widetilde{\phantom{LH}}}\!\!\!\!\mathrm{LH}.(S^2_{\hand{k}}) \hand{\simeq} \underline{Z}\hand{[}2\hand{]} , \]

\margin{76}

% Visual collation: both supplied image halves inspected; see CHECK_DECISIONS_134_7_8.tsv.

%% END SOURCE ps_134-7_076
\ifnum\EditionMode=1
\psapparatusstart
\psentry{ps_134-7_076-v0001}{1}{deletion}{\deleted{even}}{}{The author cancels this material by hand; it leaves the corrected running text.}
\psentry{ps_134-7_076-v0002}{2}{typed insertion}{\inserted[typed]{yet}}{yet}{The typed insertion is incorporated at its indicated position; insertion brackets leave the running text.}
\psentry{ps_134-7_076-v0003}{3}{deletion}{\deleted{the}}{}{The author cancels this material by hand; it leaves the corrected running text.}
\psentry{ps_134-7_076-v0004}{4}{deletion}{\deleted{is transformed by}}{}{The author cancels this material by hand; it leaves the corrected running text.}
\psentry{ps_134-7_076-v0005}{5}{insertion}{\inserted{becomes}}{\hand{becomes}}{The handwritten insertion is incorporated at its indicated position; insertion brackets leave the running text.}
\psentry{ps_134-7_076-v0006}{6}{typewriter deletion}{\xout{say}}{}{The typewriter strikes out this material; it leaves the corrected running text.}
\psentry{ps_134-7_076-v0007}{7}{typewriter overwriting}{\over[typed]{\unsure{e}}{o}}{o}{The superposed typed reading replaces the underlying reading.}
\psentry{ps_134-7_076-v0008}{8}{deletion}{\deleted{using}}{}{The author cancels this material by hand; it leaves the corrected running text.}
\psentry{ps_134-7_076-v0009}{9}{insertion}{\inserted{with}}{\hand{with}}{The handwritten insertion is incorporated at its indicated position; insertion brackets leave the running text.}
\psentry{ps_134-7_076-v0010}{10}{deletion}{\deleted{that}}{}{The author cancels this material by hand; it leaves the corrected running text.}
\psentry{ps_134-7_076-v0011}{11}{typewriter overwriting}{\over[typed]{\unsure{e}}{o}}{o}{The superposed typed reading replaces the underlying reading.}
\psentry{ps_134-7_076-v0012}{12}{deletion}{\deleted{\illegible{1}}}{}{The author cancels this material by hand; it leaves the corrected running text.}
\psentry{ps_134-7_076-v0013}{13}{writing over}{$\over{\illegible{1}}{k}$}{$\hand{k}$}{The superposed handwritten reading replaces the underlying reading.}
\psapparatusend
\fi
\pssource{ps_134-7_077}{134-7/77}{https://grothendieck.umontpellier.fr/web/viewer.html?file=../134-7.pdf\#page=78}
\psnumber{ps_134-7_077-v0001}{1}
\psnumber{ps_134-7_077-v0002}{2}
\psnumber{ps_134-7_077-v0003}{3}
\psnumber{ps_134-7_077-v0004}{4}
\psnumber{ps_134-7_077-v0005}{5}
\psnumber{ps_134-7_077-v0006}{6}
\psnumber{ps_134-7_077-v0007}{7}
\psnumber{ps_134-7_077-v0008}{8}
\psnumber{ps_134-7_077-v0009}{9}
\psnumber{ps_134-7_077-v0010}{10}
\psnumber{ps_134-7_077-v0011}{11}
\psnumber{ps_134-7_077-v0012}{12}
\psnumber{ps_134-7_077-v0013}{13}
\psnumber{ps_134-7_077-v0014}{14}
%% BEGIN SOURCE ps_134-7_077
% Pursuing Stacks (cotes 134-2 to 134-8): archive cote 134-7, site page 77 (PDF page 78), typescript page 519.
% typescript. Draft compared: ps6.tex:3842-3898.
% AI-transcribed by Claude Opus 5.5 (claude-opus-5-5, adaptive thinking, effort high), Message Batches API, batch ps4 (msgbatch_012Aoc1NuskATPsyBeKSYK7J), 2026-11-12; reviewed against both supplied scan halves; see CHECK_DECISIONS.tsv.
\scanpage{134-7/77}
\mspage{519}

\var{\deleted{and}}{}{deletion}{ps_134-7_077-v0001} whose image in Hot$_1$ (in case k=$\underline{Z}$ say) should be the (homotopy type
of\var{\deleted{\unsure{f}}}{}{deletion}{ps_134-7_077-v0002} the) ordinary 2-sphere. Taking suspensions, we'll get from this
the n-spheres $S^n_k$\var{\inserted{$= S(n,k),$}}{\hand{$= S(n,k),$}}{insertion}{ps_134-7_077-v0003} over k, for any n$\hand{\geqslant}$2. (NB In case the relevant
homotopy theoretic structures can be introduced in \var{\xout{Hot$_{\unsure{0}}$(}}{}{typewriter deletion}{ps_134-7_077-v0004} a suitable lar-%
ger Hot$_{\hand{\bullet}}$(k), containing the objects $\overset{\hand{\sim}}{K}$(L.) for L. \var{\xout{in}}{}{typewriter deletion}{ps_134-7_077-v0005} a 0-connected
chain complex of k-modules, Hot$_{\hand{\bullet}}$(k) contains an object S(1,k)=$\overset{\hand{\sim}}{K}$(1,k),
and $S^2_k$ may be simply described as it's suspension.) One main motivation
for trying to push through a theory of schematic homotopy types \var{\deleted{\illegible{1}}}{}{deletion}{ps_134-7_077-v0006}\var{\inserted{may}}{\hand{may}}{insertion}{ps_134-7_077-v0007} be
\var{\deleted{through}}{}{deletion}{ps_134-7_077-v0008} the hope that this may provide new insights into the homotopy
groups of spheres. A first\var{\xout{\illegible{1}}}{}{typewriter deletion}{ps_134-7_077-v0009} interesting consequence would be
that for any given sphere $S^n$ (in the set-up now of ordinary homotpy theo-%
ry), the series of all it's homotopy groups $\hand{\pi}_i(S^n) = \mathrm{Hom}_{\mathrm{Hot}}(S^i,S^n)$
may be viewed as being the homology modules of an object of D.(Ab),
namely $L\hand{\pi}.(S^n_{Z})$ , canonically associated to $S^n$ . As was seen by an
example\var{\deleted{,}}{}{deletion}{ps_134-7_077-v0010}\var{\inserted[typed]{in the previous section,}}{in the previous section,}{typed insertion}{ps_134-7_077-v0011} the existence of such a \uline{canonical} representation is by now
means a trivial or inocuous fact, it is indeed \var{\xout{false for}}{}{typewriter deletion}{ps_134-7_077-v0012} definitely
false for arbitrary homotopy types - here then it would turn out as a
rather special feature of the full subcategory of Hot$^\bullet$ made up with
the spheres $S^n$ (n$\hand{\geqslant}$2). Thus, for any $\hand{\alpha\in\pi}_i(S^n)$ i.e. for any map
\[\hand{\alpha} :\ S^i \hand{\to} S^n\]
in Hot$^\bullet$ , there should be an induced map
\[L\hand{\pi}.(\hand{\alpha}) : L\hand{\pi}.(S^i) \hand{\to} L\hand{\pi}.(S^{\hand{n}}) \qquad ;\]
with associativity condition for a composition
\[S^i \hand{\xrightarrow{\alpha}} S^n \hand{\xrightarrow{\beta}} S^m \qquad .\]
If we write
\[{}^{\hand{n}}\hand{\Pi}. \overset{\mathrm{def}}{=} L\ .(S^n) \quad ,\]
this defines on the set of chain complexes (more accurately, abelian
homotopy types) ${}^n\hand{\Pi}.$ ( n$\hand{\geqslant}$ 2) quite a specific structure, which merits
to be investigated a priori (under the assumption, of course, of the
existence of a reasonable theory of s\var{\over{\unsure{che}}{che}}{\hand{che}}{writing over}{ps_134-7_077-v0013}matic homotopy types, satisfying
the criteria above). Maybe after all it will\hand{\unsure{,}} turn out that the homotopy
groups of spheres \uline{cannot} be fitted into such an uncompassing structure
- so much the better, as this will show that the kind of theory I star-%
ted trying to dig out does'nt exist, which will clear up the situation
a great deal ! But if there does exist such a \var{\inserted[typed]{Canonical}}{Canonical}{typed insertion}{ps_134-7_077-v0014} structure, it surely
should'nt be ignored, and it should be pleasant work to try and pin
down exactly what extra information on homotopy groups of spheres is

\margin{77}

% Visual collation: both supplied image halves inspected; see CHECK_DECISIONS_134_7_8.tsv.

%% END SOURCE ps_134-7_077
\ifnum\EditionMode=1
\psapparatusstart
\psentry{ps_134-7_077-v0001}{1}{deletion}{\deleted{and}}{}{The author cancels this material by hand; it leaves the corrected running text.}
\psentry{ps_134-7_077-v0002}{2}{deletion}{\deleted{\unsure{f}}}{}{The author cancels this material by hand; it leaves the corrected running text.}
\psentry{ps_134-7_077-v0003}{3}{insertion}{\inserted{$= S(n,k),$}}{\hand{$= S(n,k),$}}{The handwritten insertion is incorporated at its indicated position; insertion brackets leave the running text.}
\psentry{ps_134-7_077-v0004}{4}{typewriter deletion}{\xout{Hot$_{\unsure{0}}$(}}{}{The typewriter strikes out this material; it leaves the corrected running text.}
\psentry{ps_134-7_077-v0005}{5}{typewriter deletion}{\xout{in}}{}{The typewriter strikes out this material; it leaves the corrected running text.}
\psentry{ps_134-7_077-v0006}{6}{deletion}{\deleted{\illegible{1}}}{}{The author cancels this material by hand; it leaves the corrected running text.}
\psentry{ps_134-7_077-v0007}{7}{insertion}{\inserted{may}}{\hand{may}}{The handwritten insertion is incorporated at its indicated position; insertion brackets leave the running text.}
\psentry{ps_134-7_077-v0008}{8}{deletion}{\deleted{through}}{}{The author cancels this material by hand; it leaves the corrected running text.}
\psentry{ps_134-7_077-v0009}{9}{typewriter deletion}{\xout{\illegible{1}}}{}{The typewriter strikes out this material; it leaves the corrected running text.}
\psentry{ps_134-7_077-v0010}{10}{deletion}{\deleted{,}}{}{The author cancels this material by hand; it leaves the corrected running text.}
\psentry{ps_134-7_077-v0011}{11}{typed insertion}{\inserted[typed]{in the previous section,}}{in the previous section,}{The typed insertion is incorporated at its indicated position; insertion brackets leave the running text.}
\psentry{ps_134-7_077-v0012}{12}{typewriter deletion}{\xout{false for}}{}{The typewriter strikes out this material; it leaves the corrected running text.}
\psentry{ps_134-7_077-v0013}{13}{writing over}{\over{\unsure{che}}{che}}{\hand{che}}{The superposed handwritten reading replaces the underlying reading.}
\psentry{ps_134-7_077-v0014}{14}{typed insertion}{\inserted[typed]{Canonical}}{Canonical}{The typed insertion is incorporated at its indicated position; insertion brackets leave the running text.}
\psapparatusend
\fi
\pssource{ps_134-7_078}{134-7/78}{https://grothendieck.umontpellier.fr/web/viewer.html?file=../134-7.pdf\#page=79}
\psnumber{ps_134-7_078-v0001}{1}
\psnumber{ps_134-7_078-v0002}{2}
\psnumber{ps_134-7_078-v0003}{3}
\psnumber{ps_134-7_078-v0004}{4}
\psnumber{ps_134-7_078-v0005}{5}
\psnumber{ps_134-7_078-v0006}{6}
\psnumber{ps_134-7_078-v0007}{7}
\psnumber{ps_134-7_078-v0008}{8}
\psnumber{ps_134-7_078-v0009}{9}
\psnumber{ps_134-7_078-v0010}{10}
\psnumber{ps_134-7_078-v0011}{11}
\psnumber{ps_134-7_078-v0012}{12}
\psnumber{ps_134-7_078-v0013}{13}
\psnumber{ps_134-7_078-v0014}{14}
\psnumber{ps_134-7_078-v0015}{15}
\psnumber{ps_134-7_078-v0016}{16}
\psnumber{ps_134-7_078-v0017}{17}
%% BEGIN SOURCE ps_134-7_078
% Pursuing Stacks (cotes 134-2 to 134-8): archive cote 134-7, site page 78 (PDF page 79), typescript page 520.
% typescript. Draft compared: ps6.tex:3886-3944.
% AI-transcribed by Claude Opus 5.5 (claude-opus-5-5, adaptive thinking, effort high), Message Batches API, batch ps4 (msgbatch_012Aoc1NuskATPsyBeKSYK7J), 2026-11-12; reviewed against both supplied scan halves; see CHECK_DECISIONS.tsv.
\scanpage{134-7/78}
\mspage{520}

invol\var{\over[typed]{\illegible{1}}{v}}{v}{typewriter overwriting}{ps_134-7_078-v0001}ed in such extra structure.  It may be noted for instance that,
when taking i=n  hence $\hand{\pi}_i(S^n)=\underline{Z}$ ,  we get an operation of the multiplica-%
tive monoid  $\underline{Z}^{\mathrm{mult}}$ \hand{on} ${}^n\hand{\Pi}_\bullet$,  whose action on the homology groups of ${}^n\hand{\Pi}$,
i.e. on the homotopy groups  $\hand{\pi}_j(S^n_{\var{\xout{\illegible{1}}}{}{typewriter deletion}{ps_134-7_078-v0002}})$\hand{,} is surely an important extra struc-%
ture on these,  which I hope has been studied by the homotopy people
extensively.  Surely it has been known too for a long time (at any rate
since Artin-Mazur's foundational work on profinite homotopy types) that
when $\hand{\pi}_j(S^n)$ is finite (i.e. practically in all cases except j=n) this
ac\hand{t}ion comes from a continuous action of  the multiplicative \var{\xout{groupoid}}{}{typewriter deletion}{ps_134-7_078-v0003}
monoid  \var{\xout{$\underline{Z}^\wedge$}}{}{typewriter deletion}{ps_134-7_078-v0004}
\[ M \ = \ \underline{Z}^{\wedge\mathrm{mult}} \ , \]
where  $\underline{Z}^\wedge$ is the completion of $\underline{Z}$ with respect to the topology defined by
subgroups of finite index:
\[ \underline{Z}^\wedge = \lim \underline{Z}/n\underline{Z} \hand{\simeq} \hand{\prod}_p \underline{Z}_p \qquad , \]
where in the last member the product is taken over all primes,  and $\underline{Z}_p$
denotes p-adic inte\var{\xout{r}}{}{typewriter deletion}{ps_134-7_078-v0005}gers.  Taking homotopy types over the ring k = $\underline{Z}^\wedge$,
we should get that this monoid  M operates on the object
\[ {}^n\hand{\Pi}_\bullet \overset{L}{\hand{\otimes}}_{\underline{Z}}\underline{Z}^\wedge \qquad \text{in} \quad D_\bullet(\underline{\mathbf{Ab}}_{\underline{Z}^\wedge}) \qquad , \]
as this may be interpreted as  $L\hand{\Pi}_\bullet(S^n_k)$\hand{,} and  $k^{\mathrm{mult}}$  operates on $S^n_k$ =S(n,k\hand{)}
for any ring k\var{\over{.}{,}}{\hand{,}}{writing over}{ps_134-7_078-v0006} \hand{(}if indeed  S(n,$\hand{\pi}$)  depends functorially on the variable
k-module $\hand{\pi}$\hand{)}.  \var{\deleted{(}}{}{deletion}{ps_134-7_078-v0007}At any rate,  it is surely well\unsure{\hand{-}} known that \var{\xout{$\underline{Z}^\wedge$ operat}}{}{typewriter deletion}{ps_134-7_078-v0008}
M = $\underline{Z}^{\wedge\mathrm{mult}}$  operates on the profinite completions of all spheres.  For
the odd-dimensional sp\var{\over[typed]{\unsure{hh}}{he}}{he}{typewriter overwriting}{ps_134-7_078-v0009}res,  and taking the restriction of this operati-%
on to the \var{\xout{maximal}}{}{typewriter deletion}{ps_134-7_078-v0010}\var{\inserted{largest}}{\hand{largest}}{insertion}{ps_134-7_078-v0011} subgroup
\[ M^\times = Z^{\wedge\times} \hand{\simeq} \hand{\prod}_p \underline{Z}_p^\times \]
of M,  this group may be viewed as the \var{\inserted[typed]{profinite}}{profinite}{typed insertion}{ps_134-7_078-v0012} Galois group of the maximal
cyclotomic extension of the prime field $\underline{Q}$  (deduced by adjoining\unsure{\hand{,}} all
roots of unity),  and it\hand{'}s action o\var{\over[typed]{f}{n}}{n}{typewriter overwriting}{ps_134-7_078-v0013}  the profinite completion $S^{n\hand{\wedge}}$ of $S^n$
may be viewed as the canonical Galois action,  when $S^{n=\hand{(}2m\var{\xout{+}}{}{typewriter deletion}{ps_134-7_078-v0014}-1\hand{)}}$ is inter-%
preted as the homotopy type of affine \var{\xout{m-space}}{}{typewriter deletion}{ps_134-7_078-v0015}\var{\inserted{(complex)}}{\hand{(complex)}}{insertion}{ps_134-7_078-v0016} m-space minus the origin
(which makes a sense as a scheme over \var{\xout{\illegible{1}}}{}{typewriter deletion}{ps_134-7_078-v0017} the prime field $\underline{Q}$).  This was
the interpretation in my mind,  when stating that the action of $\underline{Z}^{\mathrm{mult}}$
and of it's completion $\underline{Z}^{\wedge\mathrm{mult}}$  on the homotopy groups of spheres was
an important extra structure on these.

\margin{78}

% Visual collation: both supplied image halves inspected; see CHECK_DECISIONS_134_7_8.tsv.

%% END SOURCE ps_134-7_078
\ifnum\EditionMode=1
\psapparatusstart
\psentry{ps_134-7_078-v0001}{1}{typewriter overwriting}{\over[typed]{\illegible{1}}{v}}{v}{The superposed typed reading replaces the underlying reading.}
\psentry{ps_134-7_078-v0002}{2}{typewriter deletion}{$\xout{\illegible{1}}$}{}{The typewriter strikes out this material; it leaves the corrected running text.}
\psentry{ps_134-7_078-v0003}{3}{typewriter deletion}{\xout{groupoid}}{}{The typewriter strikes out this material; it leaves the corrected running text.}
\psentry{ps_134-7_078-v0004}{4}{typewriter deletion}{\xout{$\underline{Z}^\wedge$}}{}{The typewriter strikes out this material; it leaves the corrected running text.}
\psentry{ps_134-7_078-v0005}{5}{typewriter deletion}{\xout{r}}{}{The typewriter strikes out this material; it leaves the corrected running text.}
\psentry{ps_134-7_078-v0006}{6}{writing over}{\over{.}{,}}{\hand{,}}{The superposed handwritten reading replaces the underlying reading.}
\psentry{ps_134-7_078-v0007}{7}{deletion}{\deleted{(}}{}{The author cancels this material by hand; it leaves the corrected running text.}
\psentry{ps_134-7_078-v0008}{8}{typewriter deletion}{\xout{$\underline{Z}^\wedge$ operat}}{}{The typewriter strikes out this material; it leaves the corrected running text.}
\psentry{ps_134-7_078-v0009}{9}{typewriter overwriting}{\over[typed]{\unsure{hh}}{he}}{he}{The superposed typed reading replaces the underlying reading.}
\psentry{ps_134-7_078-v0010}{10}{typewriter deletion}{\xout{maximal}}{}{The typewriter strikes out this material; it leaves the corrected running text.}
\psentry{ps_134-7_078-v0011}{11}{insertion}{\inserted{largest}}{\hand{largest}}{The handwritten insertion is incorporated at its indicated position; insertion brackets leave the running text.}
\psentry{ps_134-7_078-v0012}{12}{typed insertion}{\inserted[typed]{profinite}}{profinite}{The typed insertion is incorporated at its indicated position; insertion brackets leave the running text.}
\psentry{ps_134-7_078-v0013}{13}{typewriter overwriting}{\over[typed]{f}{n}}{n}{The superposed typed reading replaces the underlying reading.}
\psentry{ps_134-7_078-v0014}{14}{typewriter deletion}{$\xout{+}$}{}{The typewriter strikes out this material; it leaves the corrected running text.}
\psentry{ps_134-7_078-v0015}{15}{typewriter deletion}{\xout{m-space}}{}{The typewriter strikes out this material; it leaves the corrected running text.}
\psentry{ps_134-7_078-v0016}{16}{insertion}{\inserted{(complex)}}{\hand{(complex)}}{The handwritten insertion is incorporated at its indicated position; insertion brackets leave the running text.}
\psentry{ps_134-7_078-v0017}{17}{typewriter deletion}{\xout{\illegible{1}}}{}{The typewriter strikes out this material; it leaves the corrected running text.}
\psapparatusend
\fi
\pssource{ps_134-7_079}{134-7/79}{https://grothendieck.umontpellier.fr/web/viewer.html?file=../134-7.pdf\#page=80}
\psnumber{ps_134-7_079-v0001}{1}
\psnumber{ps_134-7_079-v0002}{2}
\psnumber{ps_134-7_079-v0003}{3}
\psnumber{ps_134-7_079-v0004}{4}
\psnumber{ps_134-7_079-v0005}{5}
\psnumber{ps_134-7_079-v0006}{6}
\psnumber{ps_134-7_079-v0007}{7}
\psnumber{ps_134-7_079-v0008}{8}
\psnumber{ps_134-7_079-v0009}{9}
\psnumber{ps_134-7_079-v0010}{10}
\psnumber{ps_134-7_079-v0011}{11}
\psnumber{ps_134-7_079-v0012}{12}
\psnumber{ps_134-7_079-v0013}{13}
\psnumber{ps_134-7_079-v0014}{14}
\psnumber{ps_134-7_079-v0015}{15}
%% BEGIN SOURCE ps_134-7_079
% Pursuing Stacks (cotes 134-2 to 134-8): archive cote 134-7, site page 79 (PDF page 80), typescript page 521.
% typescript. Draft compared: ps6.tex:3929-3995.
% AI-transcribed by Claude Opus 5.5 (claude-opus-5-5, adaptive thinking, effort high), Message Batches API, batch ps4 (msgbatch_012Aoc1NuskATPsyBeKSYK7J), 2026-11-12; reviewed against both supplied scan halves; see CHECK_DECISIONS.tsv.
\scanpage{134-7/79}
\mspage{521}

124.   I guess I'll skip giving a\var{\xout{n}}{}{typewriter deletion}{ps_134-7_079-v0001} more \var{\inserted{(rather!)}}{\hand{(rather!)}}{insertion}{ps_134-7_079-v0002} or less complete axiomatic
description of $\mathrm{Hot}^{\bullet}_{o}$ or of $\mathrm{Hot}^{\bullet}_{o}(k)$,  in terms of the pair of adjoint
functors $\overset{\hand{\sim}}{\mathrm{LH}}.$ and $\overset{\hand{\sim}}{K}$ and Postnikov dévissage,  although I did go through
this exercise lately - it does'nt really shed new light on the approach
we are following here towards ``schematic'' homotopy types over arbitrary
ground rings.  In contrast with the two other approaches I have been thin%
\hand{-}king of before (namely De Rham complexes with devided powers,  and small
categories with diagon\var{\xout{\unsure{r}}}{}{typewriter deletion}{ps_134-7_079-v0003}al maps,  cf. sections 94 and 109 respectively),
the characteristic feature of this approach seems to be that it takes in-%
to account the existence of the canonical functor  \var{\xout{$\mathrm{Hot}_o$\illegible{1}D.(Ab)}}{}{typewriter deletion}{ps_134-7_079-v0004}
\var{\xout{D.(Ab)\illegible{1}Hot$_o$}}{}{typewriter deletion}{ps_134-7_079-v0005}  $\overset{\hand{\sim}}{K} : \mathrm{D.(Ab)} \hand{\longrightarrow} (\mathrm{Hot})$,  paraphrased by a functor
\[ (1) \qquad \overset{\hand{\sim}}{K} : \mathrm{D.(Ab}_k)(\text{0-conn}) \hand{\longrightarrow} \mathrm{Hot}_o(k) \]
compatible with ring extension $k \hand{\to} k'$ ;  while all three approaches have
in common \var{\xout{the existence}}{}{typewriter deletion}{ps_134-7_079-v0006} that they yield a paraphrase
\[ (2) \qquad \overset{\hand{\sim}}{\mathrm{LH}}. :\ \mathrm{Hot}_o(k) \hand{\longrightarrow} \mathrm{D.(Ab}_k)(\text{0-conn.}) \]
of the \var{\xout{tot}}{}{typewriter deletion}{ps_134-7_079-v0007} (total) homology functor,  in a way again compatible with
extension of ground rings.  In the two earlier approaches this \var{\inserted{latter}}{\hand{latter}}{insertion}{ps_134-7_079-v0008} functor
comes out in a wholly tautological way,\var{\xout{whereas I suspect}}{}{typewriter deletion}{ps_134-7_079-v0009}  whereas in the
present approach via semisimplicial schemes,  it is'nt quite so tautologi%
cal indeed.  \var{\inserted[typed]{If we want to define a functor $\overset{\hand{\sim}}{K}$ in the context \inserted{say of}}}{If we want to define a functor $\overset{\hand{\sim}}{K}$ in the context \hand{say of}}{typed insertion}{ps_134-7_079-v0010} \var{\xout{As a matter of fact in the two earlier approaches, it is}}{}{typewriter deletion}{ps_134-7_079-v0011}
\var{\inserted[typed]{De Rham algebras with devided \xout{induced by} powers.}}{De Rham algebras with devided  powers.}{typed insertion}{ps_134-7_079-v0012} \hand{\unsure{k}}
\var{\xout{is easy to define a functor K too as the composition}}{}{typewriter deletion}{ps_134-7_079-v0013}  one may think
of associating\var{\inserted[typed]{, first to a chain}}{, first to a chain}{typed insertion}{ps_134-7_079-v0014} \var{\xout{D.(Ab$_k$)\illegible{1}\unsure{Hom}}}{}{typewriter deletion}{ps_134-7_079-v0015} complex in $\mathrm{Ab}_k$  the corresponding semi-%
simplicial k-module,  view this as a semisimplicial set and take it's
De Rham complex with devided powers,  and coefficients in k (as we want
an object over k).  But this is visibly silly,  except possibly for $k=\underline{Z}$,
as the result does'nt depend on the k-module structure of the chain com-%
plex we start with.  It is doubtful anyhow that in this context,  the
tautological functor (2) admits a right adjoint (which we then would call
$\overset{\hand{\sim}}{K}$  of course). - Another noteworthy difference between the \unsure{présent|present} app-%
roach towards defining ``schematic'' homotopy types,  and the two earlier
ones (``earlier'' in these notes, at any rate !)  is that for $k=\underline{Z}$ say,
the category of ``homotopy types over $\underline{Z}$ '' we get here maps into the cate-%
gory of ordinary homotopy types,  whereas it was the opposite with the
two other approaches.

What however I still would like to do,  is a little more daydreaming
about the expected formal properties of the basic functors in between
the two categories appearing in (1) and (2)  - namely,  essentially,
between ``schematic homotopy types''  and ``abelian'' ones.  It will be con-%

\margin{79}

% Visual collation: both supplied image halves inspected; see CHECK_DECISIONS_134_7_8.tsv.

%% END SOURCE ps_134-7_079
\ifnum\EditionMode=1
\psapparatusstart
\psentry{ps_134-7_079-v0001}{1}{typewriter deletion}{\xout{n}}{}{The typewriter strikes out this material; it leaves the corrected running text.}
\psentry{ps_134-7_079-v0002}{2}{insertion}{\inserted{(rather!)}}{\hand{(rather!)}}{The handwritten insertion is incorporated at its indicated position; insertion brackets leave the running text.}
\psentry{ps_134-7_079-v0003}{3}{typewriter deletion}{\xout{\unsure{r}}}{}{The typewriter strikes out this material; it leaves the corrected running text.}
\psentry{ps_134-7_079-v0004}{4}{typewriter deletion}{\xout{$\mathrm{Hot}_o$\illegible{1}D.(Ab)}}{}{The typewriter strikes out this material; it leaves the corrected running text.}
\psentry{ps_134-7_079-v0005}{5}{typewriter deletion}{\xout{D.(Ab)\illegible{1}Hot$_o$}}{}{The typewriter strikes out this material; it leaves the corrected running text.}
\psentry{ps_134-7_079-v0006}{6}{typewriter deletion}{\xout{the existence}}{}{The typewriter strikes out this material; it leaves the corrected running text.}
\psentry{ps_134-7_079-v0007}{7}{typewriter deletion}{\xout{tot}}{}{The typewriter strikes out this material; it leaves the corrected running text.}
\psentry{ps_134-7_079-v0008}{8}{insertion}{\inserted{latter}}{\hand{latter}}{The handwritten insertion is incorporated at its indicated position; insertion brackets leave the running text.}
\psentry{ps_134-7_079-v0009}{9}{typewriter deletion}{\xout{whereas I suspect}}{}{The typewriter strikes out this material; it leaves the corrected running text.}
\psentry{ps_134-7_079-v0010}{10}{typed insertion}{\inserted[typed]{If we want to define a functor $\overset{\hand{\sim}}{K}$ in the context \inserted{say of}}}{If we want to define a functor $\overset{\hand{\sim}}{K}$ in the context \hand{say of}}{The typed insertion is incorporated at its indicated position; insertion brackets leave the running text.}
\psentry{ps_134-7_079-v0011}{11}{typewriter deletion}{\xout{As a matter of fact in the two earlier approaches, it is}}{}{The typewriter strikes out this material; it leaves the corrected running text.}
\psentry{ps_134-7_079-v0012}{12}{typed insertion}{\inserted[typed]{De Rham algebras with devided \xout{induced by} powers.}}{De Rham algebras with devided  powers.}{The typed insertion is incorporated at its indicated position; insertion brackets leave the running text.}
\psentry{ps_134-7_079-v0013}{13}{typewriter deletion}{\xout{is easy to define a functor K too as the composition}}{}{The typewriter strikes out this material; it leaves the corrected running text.}
\psentry{ps_134-7_079-v0014}{14}{typed insertion}{\inserted[typed]{, first to a chain}}{, first to a chain}{The typed insertion is incorporated at its indicated position; insertion brackets leave the running text.}
\psentry{ps_134-7_079-v0015}{15}{typewriter deletion}{\xout{D.(Ab$_k$)\illegible{1}\unsure{Hom}}}{}{The typewriter strikes out this material; it leaves the corrected running text.}
\psapparatusend
\fi
\pssource{ps_134-7_080}{134-7/80}{https://grothendieck.umontpellier.fr/web/viewer.html?file=../134-7.pdf\#page=81}
\psnumber{ps_134-7_080-v0001}{1}
\psnumber{ps_134-7_080-v0002}{2}
\psnumber{ps_134-7_080-v0003}{3}
\psnumber{ps_134-7_080-v0004}{4}
\psnumber{ps_134-7_080-v0005}{5}
\psnumber{ps_134-7_080-v0006}{6}
\psnumber{ps_134-7_080-v0007}{7}
\psnumber{ps_134-7_080-v0008}{8}
\psnumber{ps_134-7_080-v0009}{9}
\psnumber{ps_134-7_080-v0010}{10}
\psnumber{ps_134-7_080-v0011}{11}
\psnumber{ps_134-7_080-v0012}{12}
\psnumber{ps_134-7_080-v0013}{13}
\psnumber{ps_134-7_080-v0014}{14}
\psnumber{ps_134-7_080-v0015}{15}
\psnumber{ps_134-7_080-v0016}{16}
\psnumber{ps_134-7_080-v0017}{17}
\psnumber{ps_134-7_080-v0018}{18}
\psnumber{ps_134-7_080-v0019}{19}
%% BEGIN SOURCE ps_134-7_080
% Pursuing Stacks (cotes 134-2 to 134-8): archive cote 134-7, site page 80 (PDF page 81), typescript page 522.
% typescript. Draft compared: ps6.tex:3979-4037.
% AI-transcribed by Claude Opus 5.5 (claude-opus-5-5, adaptive thinking, effort high), Message Batches API, batch ps4 (msgbatch_012Aoc1NuskATPsyBeKSYK7J), 2026-11-12; reviewed against both supplied scan halves; see CHECK_DECISIONS.tsv.
\scanpage{134-7/80}
\mspage{522}
venient to denote the category of the latter by
\[ (3) \qquad Hotab_0(k) \]
rather than $D.(Ab)(0\text{-conn})$,  the subscript 0 denoting the 0-connectedness
condition.  The reason for this change in notation is that,  according
to the choices made in the basic definitions,  the ``k-linear'' algebraic
interpretation of this category may vary still.  Strictly speaking,  in
the approach we have been following in terms of semisimplicial unipotent
bundles,  the category (3) cannot really be described as the subcategory
of 0-connected objects of $D.(Ab_k)$,  we have seen rather that in order to
define a functor $\overset{\hand{\sim}}{LH}.$ ,  we have to work with \var{\xout{semisimplicial}}{}{typewriter deletion}{ps_134-7_080-v0001} chain comple-%
xes in $Pro(Ab_k)$  rather than \var{\xout{in}}{}{typewriter deletion}{ps_134-7_080-v0002} in $(Ab_k)$; \var{\xout{(}}{}{typewriter deletion}{ps_134-7_080-v0003}cf. section 11\var{\over{\illegible{1}}{5}}{\hand{5}}{writing over}{ps_134-7_080-v0004}, where this
point is made rather forcefully. (Recall that taking projective limits
instead would'nt help,  because then the linearization functor $\overset{\hand{\sim}}{LH}.$
would no longer commute to ground ring extension ! ).  The trouble then
is that in order to define
\[ \overset{\hand{\sim}}{K} : \ Hotab_0(k) \hand{\longrightarrow} Hot_0(k) \quad , \]
we \var{\deleted{are}}{}{deletion}{ps_134-7_080-v0005}\var{\inserted{are}}{\hand{are}}{insertion}{ps_134-7_080-v0006} obliged to enlarge accordingly the category of models used for
describing  $Hot_0(k)$,  namely take semisimplicial \uline{pro}-unipotent bundles ,
rather than just semisimplicial unipotent bundles.  This sudden invasion
of the picture by proobjects may appear forbidding  -  but maybe it will
appear less so,  or at any rate kind of natural and unescapable,  if we re%
\hand{-}call that when it will come to working with homotopy types of general
topoi,  these are anyhow ``prohomotopy types'' rather,  and they can be des-%
cribed only \var{\xout{in terms}}{}{typewriter deletion}{ps_134-7_080-v0007} by working systematically in terms of proob\var{\over{\unsure{g}}{\unsure{j}}}{\hand{\unsure{j}}}{writing over}{ps_134-7_080-v0008}ects of
various categories.  We'll see \var{\deleted{also}}{}{deletion}{ps_134-7_080-v0009}\var{\inserted{however}}{\hand{however}}{insertion}{ps_134-7_080-v0010} that there may be a way out of the
``pro''-mess,  by using a slightly different,  more or less dual approach
towards the idea of ``unipotent bundles''.  Of course,  one may also think
of using the ``pointed linearization'' $L_{pt}$ for $LH.$ ,  instead of $L$ ,
\var{\xout{in defining \unsure{LH LH (X}}}{}{typewriter deletion}{ps_134-7_080-v0011} but as was pointed out in section 117,  this
will lead to ``formal'' homology and cohomology invariants rather than to
``schematic'' ones,  and the relation between these and corresponding inva-%
riants for ordinary homotopy types will be looser still;  at any rate,
it is suited for \var{\xout{defining}}{}{typewriter deletion}{ps_134-7_080-v0012} describing Postnikov dévissage in the set-up
of ``formal schematic'' homotopy types,\var{\inserted{only}}{\hand{only}}{insertion}{ps_134-7_080-v0013} not for schematic homotopy types.
\var{\deleted{But}}{}{deletion}{ps_134-7_080-v0014}\var{\inserted{Still,}}{\hand{Still,}}{insertion}{ps_134-7_080-v0015} the theory of the former may have an interest in it's own right,  even
though it's relation to ordinary homotopy types is'nt so clear at present.
Thus we may \var{\inserted[typed]{state}}{state}{typed insertion}{ps_134-7_080-v0016} that there are \var{\inserted{around}}{\hand{around}}{insertion}{ps_134-7_080-v0017} at present \var{\deleted{around}}{}{deletion}{ps_134-7_080-v0018} three or four different can-%
didates for possibly fitting the daydreaming I want to write down.  \var{\deleted{\unsure{With}}}{}{deletion}{ps_134-7_080-v0019}
% NOTE: the right edge of the page is cut at "if we re" and at the struck word after "write down."
\margin{80}

% Visual collation: both supplied image halves inspected; see CHECK_DECISIONS_134_7_8.tsv.

%% END SOURCE ps_134-7_080
\ifnum\EditionMode=1
\psapparatusstart
\psentry{ps_134-7_080-v0001}{1}{typewriter deletion}{\xout{semisimplicial}}{}{The typewriter strikes out this material; it leaves the corrected running text.}
\psentry{ps_134-7_080-v0002}{2}{typewriter deletion}{\xout{in}}{}{The typewriter strikes out this material; it leaves the corrected running text.}
\psentry{ps_134-7_080-v0003}{3}{typewriter deletion}{\xout{(}}{}{The typewriter strikes out this material; it leaves the corrected running text.}
\psentry{ps_134-7_080-v0004}{4}{writing over}{\over{\illegible{1}}{5}}{\hand{5}}{The superposed handwritten reading replaces the underlying reading.}
\psentry{ps_134-7_080-v0005}{5}{deletion}{\deleted{are}}{}{The author cancels this material by hand; it leaves the corrected running text.}
\psentry{ps_134-7_080-v0006}{6}{insertion}{\inserted{are}}{\hand{are}}{The handwritten insertion is incorporated at its indicated position; insertion brackets leave the running text.}
\psentry{ps_134-7_080-v0007}{7}{typewriter deletion}{\xout{in terms}}{}{The typewriter strikes out this material; it leaves the corrected running text.}
\psentry{ps_134-7_080-v0008}{8}{writing over}{\over{\unsure{g}}{\unsure{j}}}{\hand{\unsure{j}}}{The superposed handwritten reading replaces the underlying reading.}
\psentry{ps_134-7_080-v0009}{9}{deletion}{\deleted{also}}{}{The author cancels this material by hand; it leaves the corrected running text.}
\psentry{ps_134-7_080-v0010}{10}{insertion}{\inserted{however}}{\hand{however}}{The handwritten insertion is incorporated at its indicated position; insertion brackets leave the running text.}
\psentry{ps_134-7_080-v0011}{11}{typewriter deletion}{\xout{in defining \unsure{LH LH (X}}}{}{The typewriter strikes out this material; it leaves the corrected running text.}
\psentry{ps_134-7_080-v0012}{12}{typewriter deletion}{\xout{defining}}{}{The typewriter strikes out this material; it leaves the corrected running text.}
\psentry{ps_134-7_080-v0013}{13}{insertion}{\inserted{only}}{\hand{only}}{The handwritten insertion is incorporated at its indicated position; insertion brackets leave the running text.}
\psentry{ps_134-7_080-v0014}{14}{deletion}{\deleted{But}}{}{The author cancels this material by hand; it leaves the corrected running text.}
\psentry{ps_134-7_080-v0015}{15}{insertion}{\inserted{Still,}}{\hand{Still,}}{The handwritten insertion is incorporated at its indicated position; insertion brackets leave the running text.}
\psentry{ps_134-7_080-v0016}{16}{typed insertion}{\inserted[typed]{state}}{state}{The typed insertion is incorporated at its indicated position; insertion brackets leave the running text.}
\psentry{ps_134-7_080-v0017}{17}{insertion}{\inserted{around}}{\hand{around}}{The handwritten insertion is incorporated at its indicated position; insertion brackets leave the running text.}
\psentry{ps_134-7_080-v0018}{18}{deletion}{\deleted{around}}{}{The author cancels this material by hand; it leaves the corrected running text.}
\psentry{ps_134-7_080-v0019}{19}{deletion}{\deleted{\unsure{With}}}{}{The author cancels this material by hand; it leaves the corrected running text.}
\psapparatusend
\fi
\pssource{ps_134-7_081}{134-7/81}{https://grothendieck.umontpellier.fr/web/viewer.html?file=../134-7.pdf\#page=82}
\psnumber{ps_134-7_081-v0001}{1}
\psnumber{ps_134-7_081-v0002}{2}
\psnumber{ps_134-7_081-v0003}{3}
\psnumber{ps_134-7_081-v0004}{4}
\psnumber{ps_134-7_081-v0005}{5}
\psnumber{ps_134-7_081-v0006}{6}
\psnumber{ps_134-7_081-v0007}{7}
\psnumber{ps_134-7_081-v0008}{8}
\psnumber{ps_134-7_081-v0009}{9}
\psnumber{ps_134-7_081-v0010}{10}
\psnumber{ps_134-7_081-v0011}{11}
\psnumber{ps_134-7_081-v0012}{12}
\psnumber{ps_134-7_081-v0013}{13}
%% BEGIN SOURCE ps_134-7_081
% Pursuing Stacks (cotes 134-2 to 134-8): archive cote 134-7, site page 81 (PDF page 82), typescript page 523.
% typescript. Draft compared: ps6.tex:4025-4111.
% AI-transcribed by Claude Opus 5.5 (claude-opus-5-5, adaptive thinking, effort high), Message Batches API, batch ps4 (msgbatch_012Aoc1NuskATPsyBeKSYK7J), 2026-11-12; reviewed against both supplied scan halves; see CHECK_DECISIONS.tsv.
\scanpage{134-7/81}
\mspage{523}
With this ih mind, we should'nt be too specific \var{\inserted{(for the time being)}}{\hand{(for the time being)}}{insertion}{ps_134-7_081-v0001} about the exact meaning
to be given to the sy\var{\xout{\illegible{1}}}{}{typewriter deletion}{ps_134-7_081-v0002}mbol\hand{s} Hotab$_0$(k) and Hot$_0$(k). It may be safer \var{\xout{too}}{}{typewriter deletion}{ps_134-7_081-v0003}
to replace these by Hotab$_1$(k) and Hot$_1$(k) (where the subscript 1 means
``1-connected''), all the more so as we hav'nt been able yet, in the
context of unipotent bundles, to give even a tentative precise definiti-%
on of Hot$_0$(k) i.e. of \hand{$\mathcal{M}$}$_0$(k) (mimicking the condition of abelian \var{\xout{f\illegible{1}}}{}{typewriter deletion}{ps_134-7_081-v0004} \hand{$\pi$}$_1$
with trivial action on the higher \hand{$\pi$}$_i$\hand{'}s), except by expressing things
via the associated ordinary homotopy type (passing ov er to X$_*$(k)), which
looks kind of stupid indeed ! But nothing is ``safe'' here anyhow, and
the subscript 0 looks definitely more natural here than subscript 1 ,
so we may as well keep it !

5.10.

What we are mainly interested in\hand{,} for expressing the relationships
between the non additive category Hot$_0$(k) and it's additive counterpart
Hotab$_0$(k), is an array of four functors \hand{\unsure{,}}in between these, two of
which being $L\overset{\hand{\sim}}{H}.$ (\var{\inserted{(total)}}{\hand{(total)}}{insertion}{ps_134-7_081-v0005}\var{\xout{t}}{}{typewriter deletion}{ps_134-7_081-v0006}``homology'',or ``\var{\deleted{abeliani}}{}{deletion}{ps_134-7_081-v0007}\allowbreak{}\var{\inserted{lineari}}{\hand{lineari}}{insertion}{ps_134-7_081-v0008}zation'') and $\overset{\hand{\sim}}{K}$ (Eilenberg-Mac-%
lane functor) in (1) and (2), the two remaining one's being L\hand{$\pi$}. (``total
homotopy'') and $\overset{\hand{\sim}}{S}$ (the ``spherical functor'' , compare p; 514 (9)), fit\var{\xout{tin}}{}{typewriter deletion}{ps_134-7_081-v0009}%
ting in\var{\inserted{to}}{\hand{to}}{insertion}{ps_134-7_081-v0010} the following diagram

\diagram{
% Hand-drawn arrows are marked blue; typed objects remain plain.
\begin{tikzcd}[column sep=huge]
\text{(4)}\qquad \text{Hotab}_0(k) \ar[color=blue,r, shift left, "\hand{\widetilde{K}}"] \ar[color=blue,r, shift right, "\hand{\widetilde{S}}"'] \& \text{Hot}_0(k) \ar[color=blue,l, bend right=40, "\hand{L\widetilde{H}.}"'] \ar[color=blue,l, bend left=40, "\hand{L\pi.}"]
\end{tikzcd}
}

I would now like to discuss the main formal properties to be expected
from this set of functors.

A) \uline{Adjunction} \uline{properties:}
\[
(5)\quad \hand{\left\{\rule{0pt}{2.5em}\right.}
\begin{array}{l}
\text{a) (LH.,}\hand{\widetilde{K}}\text{) is a pair of adjoint functors ,}\\
\text{b) (}\hand{\widetilde{S}}\text{, L\ .) is a pair of adjoint functors .}
\end{array}
\]

Thus we get adjunction homomorphisms (for L. in Hotab$_0$(k) and X
in Hot$_0$(k) )
\[
(6\ \text{a})\quad \hand{\left\{\rule{0pt}{2.5em}\right.}
\begin{array}{ll}
\text{X} \hand{\longrightarrow} \overset{\hand{\sim}}{\text{K}}(\text{L}\overset{\hand{\sim}}{\text{H}}.(\text{X})) & (\overset{\text{def}}{=} \overset{\hand{\sim}}{\text{W}}(\text{X}) \quad (\text{``pointed }\var{\deleted{abeliani}}{}{deletion}{ps_134-7_081-v0011}\var{\inserted{lineari}}{\hand{lineari}}{insertion}{ps_134-7_081-v0012}\text{zation''}))\\
\text{L}\overset{\hand{\sim}}{\text{H}}.(\overset{\hand{\sim}}{\text{K}}(\text{L.})) \hand{\longrightarrow} \text{L.} & ,
\end{array}
\]

and

\var{\xout{$S(L\illegible{1}(X))\illegible{1}X$}}{}{typewriter deletion}{ps_134-7_081-v0013}

\margin{81}

% Visual collation: both supplied image halves inspected; see CHECK_DECISIONS_134_7_8.tsv.

%% END SOURCE ps_134-7_081
\ifnum\EditionMode=1
\psapparatusstart
\psentry{ps_134-7_081-v0001}{1}{insertion}{\inserted{(for the time being)}}{\hand{(for the time being)}}{The handwritten insertion is incorporated at its indicated position; insertion brackets leave the running text.}
\psentry{ps_134-7_081-v0002}{2}{typewriter deletion}{\xout{\illegible{1}}}{}{The typewriter strikes out this material; it leaves the corrected running text.}
\psentry{ps_134-7_081-v0003}{3}{typewriter deletion}{\xout{too}}{}{The typewriter strikes out this material; it leaves the corrected running text.}
\psentry{ps_134-7_081-v0004}{4}{typewriter deletion}{\xout{f\illegible{1}}}{}{The typewriter strikes out this material; it leaves the corrected running text.}
\psentry{ps_134-7_081-v0005}{5}{insertion}{\inserted{(total)}}{\hand{(total)}}{The handwritten insertion is incorporated at its indicated position; insertion brackets leave the running text.}
\psentry{ps_134-7_081-v0006}{6}{typewriter deletion}{\xout{t}}{}{The typewriter strikes out this material; it leaves the corrected running text.}
\psentry{ps_134-7_081-v0007}{7}{deletion}{\deleted{abeliani}}{}{The author cancels this material by hand; it leaves the corrected running text.}
\psentry{ps_134-7_081-v0008}{8}{insertion}{\inserted{lineari}}{\hand{lineari}}{The handwritten insertion is incorporated at its indicated position; insertion brackets leave the running text.}
\psentry{ps_134-7_081-v0009}{9}{typewriter deletion}{\xout{tin}}{}{The typewriter strikes out this material; it leaves the corrected running text.}
\psentry{ps_134-7_081-v0010}{10}{insertion}{\inserted{to}}{\hand{to}}{The handwritten insertion is incorporated at its indicated position; insertion brackets leave the running text.}
\psentry{ps_134-7_081-v0011}{11}{deletion}{$\deleted{abeliani}$}{}{The author cancels this material by hand; it leaves the corrected running text.}
\psentry{ps_134-7_081-v0012}{12}{insertion}{$\inserted{lineari}$}{$\hand{lineari}$}{The handwritten insertion is incorporated at its indicated position; insertion brackets leave the running text.}
\psentry{ps_134-7_081-v0013}{13}{typewriter deletion}{\xout{$S(L\illegible{1}(X))\illegible{1}X$}}{}{The typewriter strikes out this material; it leaves the corrected running text.}
\psapparatusend
\fi
\pssource{ps_134-7_082}{134-7/82}{https://grothendieck.umontpellier.fr/web/viewer.html?file=../134-7.pdf\#page=83}
\psnumber{ps_134-7_082-v0001}{1}
\psnumber{ps_134-7_082-v0002}{2}
\psnumber{ps_134-7_082-v0003}{3}
\psnumber{ps_134-7_082-v0004}{4}
\psnumber{ps_134-7_082-v0005}{5}
\psnumber{ps_134-7_082-v0006}{6}
\psnumber{ps_134-7_082-v0007}{7}
\psnumber{ps_134-7_082-v0008}{8}
\psnumber{ps_134-7_082-v0009}{9}
\psnumber{ps_134-7_082-v0010}{10}
\psnumber{ps_134-7_082-v0011}{11}
\psnumber{ps_134-7_082-v0012}{12}
\psnumber{ps_134-7_082-v0013}{13}
\psnumber{ps_134-7_082-v0014}{14}
\psnumber{ps_134-7_082-v0015}{15}
\psnumber{ps_134-7_082-v0016}{16}
\psnumber{ps_134-7_082-v0017}{17}
%% BEGIN SOURCE ps_134-7_082
% Pursuing Stacks (cotes 134-2 to 134-8): archive cote 134-7, site page 82 (PDF page 83), typescript page 524.
% typescript. Draft compared: ps6.tex:4083-4170.
% AI-transcribed by Claude Opus 5.5 (claude-opus-5-5, adaptive thinking, effort high), Message Batches API, batch ps4 (msgbatch_012Aoc1NuskATPsyBeKSYK7J), 2026-11-12; reviewed against both supplied scan halves; see CHECK_DECISIONS.tsv.
\scanpage{134-7/82}
\mspage{524}
\[
(6\ b)\qquad \hand{\left\{\rule{0pt}{2.5em}\right.}\begin{array}{l}
\overset{\hand{\sim}}{S}(L\hand{\pi}.(X)) \hand{\longrightarrow} \ X \\
L. \hand{\to} L\hand{\pi}.(\overset{\hand{\sim}}{S}(L.)) \qquad .
\end{array}
\]

T\var{\over[typed]{h}{h}}{h}{typewriter overwriting}{ps_134-7_082-v0001}e adjunction formulas for the two pairs in (5) are
\[
(7)\qquad \hand{\left\{\rule{0pt}{2.5em}\right.}\begin{array}{ll}
a) & Hom(L\overset{\hand{\sim}}{H}.(X),L.) \hand{\simeq} Hom(X, \overset{\hand{\sim}}{K}(L.)) \\
b) & Hom(L., L\hand{\pi}.(X)) \hand{\simeq} Hom(\overset{\hand{\sim}}{S}(L),X) \qquad .
\end{array}
\]

T\var{\over[typed]{h}{h}}{h}{typewriter overwriting}{ps_134-7_082-v0002}e first hand side of a) should be viewed as \uline{cohomology} of X with
coefficients\var{\xout{x}}{}{typewriter deletion}{ps_134-7_082-v0003} in L. (more specifically, \var{\inserted{as}}{\hand{as}}{insertion}{ps_134-7_082-v0004} $\overset{\hand{\sim}}{H}{}^{0}(X, L.)$, and in case $L. =\hand{\pi}\hand{[}n\hand{]}$\hand{,}
\var{\xout{hence}}{}{typewriter deletion}{ps_134-7_082-v0005} posing
\[
(8\ a)\qquad K(\hand{\pi}\hand{[}n\hand{]}) \overset{\text{def}}{=} K(\hand{\pi},n) \qquad ,
\]
we get that the Eilenberg-Maclane object $K(\hand{\pi},n)$ represents the cohomology
functor
\[
X \hand{\longmapsto} \overset{\hand{\sim}}{H}{}^{0}(X,\hand{\pi}\hand{[}n\hand{]}) = H^{n}(X,\hand{\pi}) \quad .
\]

Symetrically, the first hand side of (7 b) should be viewed as
a ``mixed homotopy \var{\deleted{group}}{}{deletion}{ps_134-7_082-v0006}\var{\inserted{module}}{\hand{module}}{insertion}{ps_134-7_082-v0007}'' of X (relative to the ``co-coefficient'' L.),
I am tempted to denote it \var{\xout{by}}{}{typewriter deletion}{ps_134-7_082-v0008} as
\margin{\textcircled{H}}
\[
\hand{\pi}_{0}(X,L.)\ ,\quad resp.\ \ \hand{\pi}_{n}(X,\hand{H})\quad if\ \ L. =\hand{H}\hand{[}n\hand{]}\ ;
\]
thus, posing
\[
(8\ b)\qquad \overset{\hand{\sim}}{S}(\hand{H}\hand{[}n\hand{]}) \overset{\text{def}}{=} S(\hand{H},n) \qquad \text{(sph\var{\inserted{e}}{\hand{e}}{insertion}{ps_134-7_082-v0009}re-like object for $\hand{H}$,n) ,}
\]
we get that the $\hand{H}$-sphere $S(\hand{H},n)$ \var{\inserted{over k}}{\hand{over k}}{insertion}{ps_134-7_082-v0010} represents the (covariant) functor \var{\inserted{on $Hot_{0}(k)$}}{\hand{on $Hot_{0}(k)$}}{insertion}{ps_134-7_082-v0011}
\[
X \hand{\longmapsto} \hand{\pi}_{n}(X,\hand{H}) \quad .
\]
In case $\hand{H}=$\var{\deleted{\illegible{1}}}{}{deletion}{ps_134-7_082-v0012}\hand{k}\hand{,} we get
\[
\hand{\pi}_{n}(X,k) \ = \ Hom(k\hand{[}n\hand{]},L\hand{\pi}.(X)) = H_{n}(L\hand{\pi}.(X)) \overset{\text{def}}{=} \hand{\pi}_{n}(X) \quad ,
\]
and we get that the \var{\xout{functor}}{}{typewriter deletion}{ps_134-7_082-v0013} ``usual'' homotopy \var{\deleted{group}}{}{deletion}{ps_134-7_082-v0014}\var{\inserted{module}}{\hand{module}}{insertion}{ps_134-7_082-v0015} functor
\[
X \hand{\longrightarrow} \hand{\pi}_{n}(X)
\]
is represented by the ``usual'' n-sphere (over k however !) S(n,k), as
it should. ( This is of course the main justification why we expect
$(\overset{\hand{\sim}}{S},L\hand{\pi}.)$ to be a pair of adjoint functors, whereas adjunction for the
pair $(L\overset{\hand{\sim}}{H}.,\overset{\hand{\sim}}{K})$ is already fairly familiar from the set-up of ordinary
homotopy types.)

The adjunction formulae (7) show that the objects $\overset{\hand{\sim}}{K}(L.)$
\margin{\textcircled{$\widetilde{W}(X)$}}
in $H\unsure{o}t_{0}(k)$ a\unsure{r}e group objects, whereas the objects $\overset{\hand{\sim}}{S}(L.)$ are co-group
objects. This gives some \var{\inserted{inner}}{\hand{inner}}{insertion}{ps_134-7_082-v0016} justification for calling \var{\inserted{$(\widetilde{W}(X)=$}}{\hand{$(\widetilde{W}(X)=$}}{insertion}{ps_134-7_082-v0017} $\overset{\hand{\sim}}{K}(L\overset{\hand{\sim}}{H}.(X))$ the
\margin{82}

% Visual collation: both supplied image halves inspected; see CHECK_DECISIONS_134_7_8.tsv.

%% END SOURCE ps_134-7_082
\ifnum\EditionMode=1
\psapparatusstart
\psentry{ps_134-7_082-v0001}{1}{typewriter overwriting}{\over[typed]{h}{h}}{h}{The superposed typed reading replaces the underlying reading.}
\psentry{ps_134-7_082-v0002}{2}{typewriter overwriting}{\over[typed]{h}{h}}{h}{The superposed typed reading replaces the underlying reading.}
\psentry{ps_134-7_082-v0003}{3}{typewriter deletion}{\xout{x}}{}{The typewriter strikes out this material; it leaves the corrected running text.}
\psentry{ps_134-7_082-v0004}{4}{insertion}{\inserted{as}}{\hand{as}}{The handwritten insertion is incorporated at its indicated position; insertion brackets leave the running text.}
\psentry{ps_134-7_082-v0005}{5}{typewriter deletion}{\xout{hence}}{}{The typewriter strikes out this material; it leaves the corrected running text.}
\psentry{ps_134-7_082-v0006}{6}{deletion}{\deleted{group}}{}{The author cancels this material by hand; it leaves the corrected running text.}
\psentry{ps_134-7_082-v0007}{7}{insertion}{\inserted{module}}{\hand{module}}{The handwritten insertion is incorporated at its indicated position; insertion brackets leave the running text.}
\psentry{ps_134-7_082-v0008}{8}{typewriter deletion}{\xout{by}}{}{The typewriter strikes out this material; it leaves the corrected running text.}
\psentry{ps_134-7_082-v0009}{9}{insertion}{\inserted{e}}{\hand{e}}{The handwritten insertion is incorporated at its indicated position; insertion brackets leave the running text.}
\psentry{ps_134-7_082-v0010}{10}{insertion}{\inserted{over k}}{\hand{over k}}{The handwritten insertion is incorporated at its indicated position; insertion brackets leave the running text.}
\psentry{ps_134-7_082-v0011}{11}{insertion}{\inserted{on $Hot_{0}(k)$}}{\hand{on $Hot_{0}(k)$}}{The handwritten insertion is incorporated at its indicated position; insertion brackets leave the running text.}
\psentry{ps_134-7_082-v0012}{12}{deletion}{\deleted{\illegible{1}}}{}{The author cancels this material by hand; it leaves the corrected running text.}
\psentry{ps_134-7_082-v0013}{13}{typewriter deletion}{\xout{functor}}{}{The typewriter strikes out this material; it leaves the corrected running text.}
\psentry{ps_134-7_082-v0014}{14}{deletion}{\deleted{group}}{}{The author cancels this material by hand; it leaves the corrected running text.}
\psentry{ps_134-7_082-v0015}{15}{insertion}{\inserted{module}}{\hand{module}}{The handwritten insertion is incorporated at its indicated position; insertion brackets leave the running text.}
\psentry{ps_134-7_082-v0016}{16}{insertion}{\inserted{inner}}{\hand{inner}}{The handwritten insertion is incorporated at its indicated position; insertion brackets leave the running text.}
\psentry{ps_134-7_082-v0017}{17}{insertion}{\inserted{$(\widetilde{W}(X)=$}}{\hand{$(\widetilde{W}(X)=$}}{The handwritten insertion is incorporated at its indicated position; insertion brackets leave the running text.}
\psapparatusend
\fi
\pssource{ps_134-7_083}{134-7/83}{https://grothendieck.umontpellier.fr/web/viewer.html?file=../134-7.pdf\#page=84}
\psnumber{ps_134-7_083-v0001}{1}
\psnumber{ps_134-7_083-v0002}{2}
\psnumber{ps_134-7_083-v0003}{3}
\psnumber{ps_134-7_083-v0004}{4}
\psnumber{ps_134-7_083-v0005}{5}
\psnumber{ps_134-7_083-v0006}{6}
\psnumber{ps_134-7_083-v0007}{7}
\psnumber{ps_134-7_083-v0008}{8}
\psnumber{ps_134-7_083-v0009}{9}
\psnumber{ps_134-7_083-v0010}{10}
\psnumber{ps_134-7_083-v0011}{11}
\psnumber{ps_134-7_083-v0012}{12}
\psnumber{ps_134-7_083-v0013}{13}
\psnumber{ps_134-7_083-v0014}{14}
\psnumber{ps_134-7_083-v0015}{15}
\psnumber{ps_134-7_083-v0016}{16}
\psnumber{ps_134-7_083-v0017}{17}
\psnumber{ps_134-7_083-v0018}{18}
\psnumber{ps_134-7_083-v0019}{19}
\psnumber{ps_134-7_083-v0020}{20}
\psnumber{ps_134-7_083-v0021}{21}
\psnumber{ps_134-7_083-v0022}{22}
\psnumber{ps_134-7_083-v0023}{23}
%% BEGIN SOURCE ps_134-7_083
% Pursuing Stacks (cotes 134-2 to 134-8): archive cote 134-7, site page 83 (PDF page 84), typescript page 525.
% typescript. Draft compared: ps6.tex:4154-4237.
% AI-transcribed by Claude Opus 5.5 (claude-opus-5-5, adaptive thinking, effort high), Message Batches API, batch ps4 (msgbatch_012Aoc1NuskATPsyBeKSYK7J), 2026-11-12; reviewed against both supplied scan halves; see CHECK_DECISIONS.tsv.
\scanpage{134-7/83}
\mspage{525}

(pointed) ``\var{\deleted{abelian}}{}{deletion}{ps_134-7_083-v0001}\var{\inserted{linear}}{\hand{linear}}{insertion}{ps_134-7_083-v0002}ization'' of the schematic homotopy type X, \var{\xout{whereas}}{}{typewriter deletion}{ps_134-7_083-v0003}
\var{\xout{and to view the adjunction morphism}}{}{typewriter deletion}{ps_134-7_083-v0004} which maps into \var{\deleted{it's abelianization}}{}{deletion}{ps_134-7_083-v0005}\var{\inserted{the former}}{\hand{the former}}{insertion}{ps_134-7_083-v0006}
\var{\xout{Dually,}}{}{typewriter deletion}{ps_134-7_083-v0007} by the first adjunction morphism in (6 a). Dually, we may call
$\overset{\hand{\sim}}{S}(L\hand{\pi}.(X))$ the (pointed) ``co-\var{\deleted{abelian}}{}{deletion}{ps_134-7_083-v0008}\var{\inserted{linear}}{\hand{linear}}{insertion}{ps_134-7_083-v0009}ization'' of X, it maps into X by the
first adjunction morphism in (6 b).

The \var{\xout{first member in}}{}{typewriter deletion}{ps_134-7_083-v0010}\var{\inserted[typed]{source}}{source}{typed insertion}{ps_134-7_083-v0011} $L\overset{\hand{\sim}}{H}.(\overset{\hand{\sim}}{K}(L.))$ in the second adjunction map (6 a)
may be viewed as Eilenberg-Maclane type \var{\inserted[typed]{(total)}}{(total)}{typed insertion}{ps_134-7_083-v0012} homology, corresponding to the
chain complex L. \hand{,} (reducing to the \var{\xout{familiar}}{}{typewriter deletion}{ps_134-7_083-v0013} more familiar $L\overset{\hand{\sim}}{H}.(\hand{\pi},n)$
when $L. =\hand{\pi}\hand{[}n\hand{]}$\unsure{\hand{)}} (but in the schematic sense of course ! ), whereas
the target $L\hand{\pi}.(\overset{\hand{\sim}}{S}(L.))$ in the second adjunction map (6 b) may be viewed
as ``total homotopy'' of a sphere-type space, reducing \var{\xout{(}}{}{typewriter deletion}{ps_134-7_083-v0014}in case $L. =\hand{\pi}\hand{[}n\hand{]}$\var{\xout{)}}{}{typewriter deletion}{ps_134-7_083-v0015}
to the total homotopy of the standard n-sphere over k, $S(n,\hand{\pi})$. In case
$k=\underline{Z}$ , we expect of course \var{\xout{this to}}{}{typewriter deletion}{ps_134-7_083-v0016} the homology groups of this chain
complex to be the ordinary homotopy groups of the usual sphere $S^n$ -
cf. \hand{G)} below.

B) \uline{Inversion formulae}: two functorial isomorphisms
\[
(9)\qquad \hand{\Bigg\{}\begin{array}{ll}
a) & L\hand{\pi}.(\overset{\hand{\sim}}{K}(L.)) \hand{\xrightarrow{\sim}} L. \quad , \\
b) & L\overset{\hand{\sim}}{H}.(\overset{\hand{\sim}}{S}(L.)) \hand{\xrightarrow{\sim}} L. \quad .
\end{array}
\]

Maybe we should have begun with these, as they kind of fix the mea-%
ning of the two functors $\overset{\hand{\sim}}{K},\overset{\hand{\sim}}{S}$ in terms of $L\overset{\hand{\sim}}{H}.$ , $L\hand{\pi}.$ , which may be
viewed as embodying respectively the two main sets of invariants of a\var{\xout{n}}{}{typewriter deletion}{ps_134-7_083-v0017}
homotopy type, namely homology and homotopy.

C) \uline{Hurewicz map} :
\[
(10)\qquad L\hand{\pi}.(X) \hand{\longrightarrow} L\overset{\hand{\sim}}{H}.(X) \quad .
\]

We get such a map by applying $L\hand{\pi}.$ to the \var{\deleted{abelian}}{}{deletion}{ps_134-7_083-v0018}\var{\inserted{lineari}}{\hand{lineari}}{insertion}{ps_134-7_083-v0019}zation map in
(6 a), and using (9 a); symetrically, we may apply $L\overset{\hand{\sim}}{H}.$ to the
\var{\xout{abelianization}}{}{typewriter deletion}{ps_134-7_083-v0020}\var{\inserted[typed]{colineari}}{colineari}{typed insertion}{ps_134-7_083-v0021}zation map (6 b), and use (9 b). We get a priori two maps
(10), and the statement is that these two maps are the same. Moreover
we have the all-important \uline{Hurewicz} \uline{theorem}: The first non-vanishing
homology object\hand{s} for $L\hand{\pi}.(X)$ and $L\overset{\hand{\sim}}{H}.(X)$ occur in the same dimension, n
say, and (10) induces an isomorphism for $H_n$ , \hand{an epimorphism for $H_{n+1}$.}

D) \uline{Exactness properties}: they can be stated shortly by saying that
in each one of the two pairs of adjoint functors (5), the left adjoint
one \var{\xout{commutes \unsure{with} cof}}{}{typewriter deletion}{ps_134-7_083-v0022} respects cofibration sequences, whereas the right
adjoint respects fibration\var{\xout{s}}{}{typewriter deletion}{ps_134-7_083-v0023} sequences. This may be detailed as four
exactness statements, one for each one of the four basic functors.

\margin{83}

% Visual collation: both supplied image halves inspected; see CHECK_DECISIONS_134_7_8.tsv.

%% END SOURCE ps_134-7_083
\ifnum\EditionMode=1
\psapparatusstart
\psentry{ps_134-7_083-v0001}{1}{deletion}{\deleted{abelian}}{}{The author cancels this material by hand; it leaves the corrected running text.}
\psentry{ps_134-7_083-v0002}{2}{insertion}{\inserted{linear}}{\hand{linear}}{The handwritten insertion is incorporated at its indicated position; insertion brackets leave the running text.}
\psentry{ps_134-7_083-v0003}{3}{typewriter deletion}{\xout{whereas}}{}{The typewriter strikes out this material; it leaves the corrected running text.}
\psentry{ps_134-7_083-v0004}{4}{typewriter deletion}{\xout{and to view the adjunction morphism}}{}{The typewriter strikes out this material; it leaves the corrected running text.}
\psentry{ps_134-7_083-v0005}{5}{deletion}{\deleted{it's abelianization}}{}{The author cancels this material by hand; it leaves the corrected running text.}
\psentry{ps_134-7_083-v0006}{6}{insertion}{\inserted{the former}}{\hand{the former}}{The handwritten insertion is incorporated at its indicated position; insertion brackets leave the running text.}
\psentry{ps_134-7_083-v0007}{7}{typewriter deletion}{\xout{Dually,}}{}{The typewriter strikes out this material; it leaves the corrected running text.}
\psentry{ps_134-7_083-v0008}{8}{deletion}{\deleted{abelian}}{}{The author cancels this material by hand; it leaves the corrected running text.}
\psentry{ps_134-7_083-v0009}{9}{insertion}{\inserted{linear}}{\hand{linear}}{The handwritten insertion is incorporated at its indicated position; insertion brackets leave the running text.}
\psentry{ps_134-7_083-v0010}{10}{typewriter deletion}{\xout{first member in}}{}{The typewriter strikes out this material; it leaves the corrected running text.}
\psentry{ps_134-7_083-v0011}{11}{typed insertion}{\inserted[typed]{source}}{source}{The typed insertion is incorporated at its indicated position; insertion brackets leave the running text.}
\psentry{ps_134-7_083-v0012}{12}{typed insertion}{\inserted[typed]{(total)}}{(total)}{The typed insertion is incorporated at its indicated position; insertion brackets leave the running text.}
\psentry{ps_134-7_083-v0013}{13}{typewriter deletion}{\xout{familiar}}{}{The typewriter strikes out this material; it leaves the corrected running text.}
\psentry{ps_134-7_083-v0014}{14}{typewriter deletion}{\xout{(}}{}{The typewriter strikes out this material; it leaves the corrected running text.}
\psentry{ps_134-7_083-v0015}{15}{typewriter deletion}{\xout{)}}{}{The typewriter strikes out this material; it leaves the corrected running text.}
\psentry{ps_134-7_083-v0016}{16}{typewriter deletion}{\xout{this to}}{}{The typewriter strikes out this material; it leaves the corrected running text.}
\psentry{ps_134-7_083-v0017}{17}{typewriter deletion}{\xout{n}}{}{The typewriter strikes out this material; it leaves the corrected running text.}
\psentry{ps_134-7_083-v0018}{18}{deletion}{\deleted{abelian}}{}{The author cancels this material by hand; it leaves the corrected running text.}
\psentry{ps_134-7_083-v0019}{19}{insertion}{\inserted{lineari}}{\hand{lineari}}{The handwritten insertion is incorporated at its indicated position; insertion brackets leave the running text.}
\psentry{ps_134-7_083-v0020}{20}{typewriter deletion}{\xout{abelianization}}{}{The typewriter strikes out this material; it leaves the corrected running text.}
\psentry{ps_134-7_083-v0021}{21}{typed insertion}{\inserted[typed]{colineari}}{colineari}{The typed insertion is incorporated at its indicated position; insertion brackets leave the running text.}
\psentry{ps_134-7_083-v0022}{22}{typewriter deletion}{\xout{commutes \unsure{with} cof}}{}{The typewriter strikes out this material; it leaves the corrected running text.}
\psentry{ps_134-7_083-v0023}{23}{typewriter deletion}{\xout{s}}{}{The typewriter strikes out this material; it leaves the corrected running text.}
\psapparatusend
\fi
\pssource{ps_134-7_084}{134-7/84}{https://grothendieck.umontpellier.fr/web/viewer.html?file=../134-7.pdf\#page=85}
\psnumber{ps_134-7_084-v0001}{1}
\psnumber{ps_134-7_084-v0002}{2}
\psnumber{ps_134-7_084-v0003}{3}
\psnumber{ps_134-7_084-v0004}{4}
\psnumber{ps_134-7_084-v0005}{5}
\psnumber{ps_134-7_084-v0006}{6}
\psnumber{ps_134-7_084-v0007}{7}
\psnumber{ps_134-7_084-v0008}{8}
\psnumber{ps_134-7_084-v0009}{9}
\psnumber{ps_134-7_084-v0010}{10}
\psnumber{ps_134-7_084-v0011}{11}
\psnumber{ps_134-7_084-v0012}{12}
\psnumber{ps_134-7_084-v0013}{13}
\psnumber{ps_134-7_084-v0014}{14}
\psnumber{ps_134-7_084-v0015}{15}
\psnumber{ps_134-7_084-v0016}{16}
\psnumber{ps_134-7_084-v0017}{17}
%% BEGIN SOURCE ps_134-7_084
% Pursuing Stacks (cotes 134-2 to 134-8): archive cote 134-7, site page 84 (PDF page 85), typescript page 526.
% typescript. Draft compared: ps6.tex:4220-4292.
% AI-transcribed by Claude Opus 5.5 (claude-opus-5-5, adaptive thinking, effort high), Message Batches API, batch ps4 (msgbatch_012Aoc1NuskATPsyBeKSYK7J), 2026-11-12; reviewed against both supplied scan halves; see CHECK_DECISIONS.tsv.
\scanpage{134-7/84}
\mspage{526}

Thus $\mathrm{L}\overset{\hand{\sim}}{\mathrm{H}}_{\bullet}$ respects cofibration sequences, which means essentially
that for such a sequence in H\unsure{o}t$_{\circ}$(k)
\[ Y \hand{\longrightarrow} X \hand{\longrightarrow} Z \quad , \ Z \text{ the homotopy cofiber of } Y \hand{\rightarrow} X \ , \]
we get a corresponding exact triangle in Hotab$_{\circ}$(k)
\[ \text{(11 b)} \qquad
% Hand-drawn arrows are marked blue; typed objects remain plain.
\diagram{\begin{tikzcd}[column sep=small]
 \& \mathrm{L}\overset{\hand{\sim}}{\mathrm{H}}_{\bullet}(\var{\over{\unsure{X}}{Z}}{\hand{Z}}{writing over}{ps_134-7_084-v0001}) \ar[color=blue,dl, "\hand{1}"'] \& \\
 \mathrm{L}\overset{\hand{\sim}}{\mathrm{H}}_{\bullet}(\var{\over{Z}{Y}}{\hand{Y}}{writing over}{ps_134-7_084-v0002}) \ar[color=blue,rr] \& \& \mathrm{L}\overset{\hand{\sim}}{\mathrm{H}}_{\bullet}(\var{\over{\unsure{Y}}{X}}{\hand{X}}{writing over}{ps_134-7_084-v0003}) \ar[color=blue,ul]
\end{tikzcd}} \qquad , \]
which is the most complete and elegant way, I guess, for expressing
behaviour of homology with respect to \var{\xout{susp}}{}{typewriter deletion}{ps_134-7_084-v0004} homotopy cofibers and suspen-%
sions. DUally, $\mathrm{L}\hand{\Pi}_{\bullet}$ respects fibration sequences, which means that
for such a sequence in Hot$_{\circ}$(k)
\[ Z \hand{\rightarrow} X \hand{\rightarrow} Y \quad , \ Z \text{ the \var{\inserted{(connected)}}{\hand{(connected)}}{insertion}{ps_134-7_084-v0005} homotopy fiber of } X \hand{\rightarrow} Y \ , \]
we get a corresponding exact triangle in Hotab$_{\bullet}$(k)
\[ \text{(11 a)} \qquad
% Hand-drawn arrows are marked blue; typed objects remain plain.
\diagram{\begin{tikzcd}[column sep=small]
 \& \mathrm{L}\hand{\Pi}_{\bullet}(\var{\over{Z}{Y}}{\hand{Y}}{writing over}{ps_134-7_084-v0006}) \ar[color=blue,dl, "\hand{1}"'] \& \\
 \mathrm{L}\hand{\Pi}_{\bullet}(Z) \ar[color=blue,rr] \& \& \mathrm{L}\hand{\Pi}_{\bullet}(X) \ar[color=blue,ul]
\end{tikzcd}} \qquad , \]
which is a more complete and elegant way of stating the exact homotopy
sequence of a fibration, in terms of total homotopy.

\var{\inserted[typed]{Formula}}{Formula}{typed insertion}{ps_134-7_084-v0007} (11 b) implies that $\mathrm{L}\overset{\hand{\sim}}{\mathrm{H}}_{\bullet}$ commutes with suspension functors $\hand{\Sigma}$ ,
the latter in Hotab$_{\bullet}$(k) (visualized as a derived chain complex category)
is just the shift functor
\[ \hand{\Sigma}_{ab} : \ L_{\bullet} \hand{\longmapsto} L_{\bullet}[1] \qquad ( L[1]_{n} = L_{n-1}) \quad \var{\deleted{\unsure{\ast}}}{}{deletion}{ps_134-7_084-v0008} , \]
which is a \uline{fully faithful} functor . Dually, formula (11 a) implies
that $\mathrm{L}\hand{\Pi}_{\bullet}$ commutes to the loop-space functors $\hand{\Omega^{\circ}}$; the latter in
Hotab$_{\bullet}$(k) cannot be described as just a shift in opposite direction
\var{\xout{in}}{}{typewriter deletion}{ps_134-7_084-v0009}
\[ L_{\bullet} \hand{\longmapsto} L_{\bullet}[-1] \ , \]
as this will get us out of \var{\inserted[typed]{0-connected}}{0-connected}{typed insertion}{ps_134-7_084-v0010} chain complexes\unsure{\var{\over{.}{,}}{\hand{,}}{writing over}{ps_134-7_084-v0011}} we have to truncate moreover
at dimension \var{\over{0}{1}}{\hand{1}}{writing over}{ps_134-7_084-v0012} (afterwards, or at dimension \var{\xout{\illegible{1}}}{}{typewriter deletion}{ps_134-7_084-v0013} 2 beforehand). This func%
% NOTE: the right edge of the page is cut at the end of this line; a hyphen after "func" may be hidden
tor is not fully faithful therefore - we loose something when passing
from $L_{\bullet}$ to $\hand{\Omega^{\circ}}_{ab}(L_{\bullet})$.

As for the two functors $\overset{\hand{\sim}}{K},\overset{\hand{\sim}}{S}$ \var{\xout{from}}{}{typewriter deletion}{ps_134-7_084-v0014} in the opposite direction, from
Hotab$_{\bullet}$(k) (i.e. essentially chain complexes) to Hot$_{\bullet}$(k), that the
first one respects fibration sequences is surely quite familiar a fact
(and kind of tautological) in the set-up of ordinary homotopy types. \var{\inserted{That}}{\hand{That}}{insertion}{ps_134-7_084-v0015}
\var{\over{T}{t}}{\hand{t}}{writing over}{ps_134-7_084-v0016}\var{\over[typed]{h}{h}}{h}{typewriter overwriting}{ps_134-7_084-v0017}e second, less familiar functor $\overset{\hand{\sim}}{S}$ should respect cofibration sequen-%

\margin{84}

% Visual collation: both supplied image halves inspected; see CHECK_DECISIONS_134_7_8.tsv.

%% END SOURCE ps_134-7_084
\ifnum\EditionMode=1
\psapparatusstart
\psentry{ps_134-7_084-v0001}{1}{writing over}{$\over{\unsure{X}}{Z}$}{$\hand{Z}$}{The superposed handwritten reading replaces the underlying reading.}
\psentry{ps_134-7_084-v0002}{2}{writing over}{$\over{Z}{Y}$}{$\hand{Y}$}{The superposed handwritten reading replaces the underlying reading.}
\psentry{ps_134-7_084-v0003}{3}{writing over}{$\over{\unsure{Y}}{X}$}{$\hand{X}$}{The superposed handwritten reading replaces the underlying reading.}
\psentry{ps_134-7_084-v0004}{4}{typewriter deletion}{\xout{susp}}{}{The typewriter strikes out this material; it leaves the corrected running text.}
\psentry{ps_134-7_084-v0005}{5}{insertion}{\inserted{(connected)}}{\hand{(connected)}}{The handwritten insertion is incorporated at its indicated position; insertion brackets leave the running text.}
\psentry{ps_134-7_084-v0006}{6}{writing over}{$\over{Z}{Y}$}{$\hand{Y}$}{The superposed handwritten reading replaces the underlying reading.}
\psentry{ps_134-7_084-v0007}{7}{typed insertion}{\inserted[typed]{Formula}}{Formula}{The typed insertion is incorporated at its indicated position; insertion brackets leave the running text.}
\psentry{ps_134-7_084-v0008}{8}{deletion}{$\deleted{\unsure{\ast}}$}{}{The author cancels this material by hand; it leaves the corrected running text.}
\psentry{ps_134-7_084-v0009}{9}{typewriter deletion}{\xout{in}}{}{The typewriter strikes out this material; it leaves the corrected running text.}
\psentry{ps_134-7_084-v0010}{10}{typed insertion}{\inserted[typed]{0-connected}}{0-connected}{The typed insertion is incorporated at its indicated position; insertion brackets leave the running text.}
\psentry{ps_134-7_084-v0011}{11}{writing over}{\over{.}{,}}{\hand{,}}{The superposed handwritten reading replaces the underlying reading.}
\psentry{ps_134-7_084-v0012}{12}{writing over}{\over{0}{1}}{\hand{1}}{The superposed handwritten reading replaces the underlying reading.}
\psentry{ps_134-7_084-v0013}{13}{typewriter deletion}{\xout{\illegible{1}}}{}{The typewriter strikes out this material; it leaves the corrected running text.}
\psentry{ps_134-7_084-v0014}{14}{typewriter deletion}{\xout{from}}{}{The typewriter strikes out this material; it leaves the corrected running text.}
\psentry{ps_134-7_084-v0015}{15}{insertion}{\inserted{That}}{\hand{That}}{The handwritten insertion is incorporated at its indicated position; insertion brackets leave the running text.}
\psentry{ps_134-7_084-v0016}{16}{writing over}{\over{T}{t}}{\hand{t}}{The superposed handwritten reading replaces the underlying reading.}
\psentry{ps_134-7_084-v0017}{17}{typewriter overwriting}{\over[typed]{h}{h}}{h}{The superposed typed reading replaces the underlying reading.}
\psapparatusend
\fi
\pssource{ps_134-7_085}{134-7/85}{https://grothendieck.umontpellier.fr/web/viewer.html?file=../134-7.pdf\#page=86}
\psnumber{ps_134-7_085-v0001}{1}
\psnumber{ps_134-7_085-v0002}{2}
\psnumber{ps_134-7_085-v0003}{3}
\psnumber{ps_134-7_085-v0004}{4}
\psnumber{ps_134-7_085-v0005}{5}
\psnumber{ps_134-7_085-v0006}{6}
\psnumber{ps_134-7_085-v0007}{7}
\psnumber{ps_134-7_085-v0008}{8}
\psnumber{ps_134-7_085-v0009}{9}
\psnumber{ps_134-7_085-v0010}{10}
\psnumber{ps_134-7_085-v0011}{11}
\psnumber{ps_134-7_085-v0012}{12}
\psnumber{ps_134-7_085-v0013}{13}
\psnumber{ps_134-7_085-v0014}{14}
\psnumber{ps_134-7_085-v0015}{15}
\psnumber{ps_134-7_085-v0016}{16}
\psnumber{ps_134-7_085-v0017}{17}
\psnumber{ps_134-7_085-v0018}{18}
\psnumber{ps_134-7_085-v0019}{19}
\psnumber{ps_134-7_085-v0020}{20}
\psnumber{ps_134-7_085-v0021}{21}
\psnumber{ps_134-7_085-v0022}{22}
\psnumber{ps_134-7_085-v0023}{23}
\psnumber{ps_134-7_085-v0024}{24}
%% BEGIN SOURCE ps_134-7_085
% Pursuing Stacks (cotes 134-2 to 134-8): archive cote 134-7, site page 85 (PDF page 86), typescript page 527.
% typescript. Draft compared: ps6.tex:4277-4338.
% AI-transcribed by Claude Opus 5.5 (claude-opus-5-5, adaptive thinking, effort high), Message Batches API, batch ps4 (msgbatch_012Aoc1NuskATPsyBeKSYK7J), 2026-11-12; reviewed against both supplied scan halves; see CHECK_DECISIONS.tsv.
\scanpage{134-7/85}
\mspage{527}

ces should be useful in order to give a more or less explicit constructi-%
on of $\overset{\hand{\sim}}{S}$(L.), for given L., in terms of the ``spheres over k'' \var{\xout{of}}{}{typewriter deletion}{ps_134-7_085-v0001} $S(\underline{K},n)$.

\uline{Remarks}. 1)The superscript 0 for the loop functor $\hand{\Omega}^{0}$ or $\hand{\Omega}^{0}_{ab}$ should
not be confused with the subscript n in the iterated loop-space functor
$\hand{\Omega}_n$ , it is \var{\xout{put in}}{}{typewriter deletion}{ps_134-7_085-v0002} added here to suggest that we are taking the \uline{neutral}
\uline{connected component} of the ``true'' full loop-space, \var{\deleted{as}}{}{deletion}{ps_134-7_085-v0003} \var{\inserted{this being}}{\hand{this being}}{insertion}{ps_134-7_085-v0004} imposed by the
restriction of working throughout with 0-connected objects\hand{;} \var{\deleted{--}}{}{deletion}{ps_134-7_085-v0005} likewise,
the ``homotopy fiber'' operation in the present context should be viewed
as meaning ``connected component at the marked point of the full homotopy
fiber''. These necessary readjustments of the usual notions is \var{\inserted{being}}{\hand{being}}{insertion}{ps_134-7_085-v0006} felt as an
unwelcome feature (of which I have become aware only at this very moment,
I confess, through the writing of the notes). Intuitively, the restric-%
tion to 0-connected \var{\inserted[typed]{(pointed)}}{(pointed)}{typed insertion}{ps_134-7_085-v0007} homotopy types does'nt seem so imperative, technical
\hand{-}ly however, when working with \var{\inserted[typed]{semisimplicial}}{semisimplicial}{typed insertion}{ps_134-7_085-v0008} schematic models, we had felt like intro-%
ducing the condition $X_0$=e , which may be viewed as a strong form of a
0-connectedness condition. Truth to tell, it is'nt so clear that this
condition is \var{\xout{of}}{}{typewriter deletion}{ps_134-7_085-v0009} going to be of great utility - anyhow, it is'nt enough
to ensure what we're really after at present (namely abelian $\hand{\pi}_1$ and tri-%
vial action on the higher $\hand{\pi}_i$'s ), a condition which is anyhow indepen\var{\xout{\illegible{1}}}{}{typewriter deletion}{ps_134-7_085-v0010}
dant of any 0-connectedness type assumption, and \var{\inserted[typed]{is}}{is}{typed insertion}{ps_134-7_085-v0011} moreover (it would
seem) stable under the basic fiber and cofiber operations. To sum up,
it may well turn out that we better replace the categories of 0-connected
schematic homotopy types Hotab$_0$(k) and Hot$_0$(k), by slightly larger ones,
so as to get rid of the 0-connectedness restriction. This however is at
present a relatively minor point, and therefore we'll leave the \var{\xout{\unsure{notati-}}}{}{typewriter deletion}{ps_134-7_085-v0012} \var{\inserted{notations}}{\hand{notations}}{insertion}{ps_134-7_085-v0013}
ons as they are.

2) Behaviour of $L\overset{\hand{\sim}}{H}.$ with respect to \uline{fibration} sequences \var{\xout{is a}}{}{typewriter deletion}{ps_134-7_085-v0014}
(instead of cofibration sequences) \var{\xout{are}}{}{typewriter deletion}{ps_134-7_085-v0015} is \var{\inserted{a}}{\hand{a}}{insertion}{ps_134-7_085-v0016} relatively delicate \var{\inserted{matter}}{\hand{matter}}{insertion}{ps_134-7_085-v0017}, it is
governed by the Leray spectral sequence, whose initial term is the
homology of the base Y with ``coefficients'' in the homology of the \var{\xout{\unsure{(varia-}}}{}{typewriter deletion}{ps_134-7_085-v0018}
\var{\xout{ble)}}{}{typewriter deletion}{ps_134-7_085-v0019} fiber. I wonder if there is anything similar for the behaviour
of L$\hand{\pi}$. with respect to cofibration sequences ?

3) In the display of the main expected properties of the four
basic functors, there is a striking \uline{symmetry}, which we tried to stress
\var{\inserted[typed]{by}}{by}{typed insertion}{ps_134-7_085-v0020} \var{\xout{through}}{}{typewriter deletion}{ps_134-7_085-v0021} the way of presenting the main formulae. O\var{\over[typed]{n}{n}}{n}{typewriter overwriting}{ps_134-7_085-v0022}e way to express
this symmetry is to say that things look \var{\xout{\unsure{more or}}}{}{typewriter deletion}{ps_134-7_085-v0023} as if there was an auto-%
duality in the pair of categories (H\var{\over[typed]{o}{o}}{o}{typewriter overwriting}{ps_134-7_085-v0024}tab$_0$(k), Hot$_0$(k)), namely a pair
of contravariant involutive functors ($D_{ab}$,D) , each of which interchan\unsure{g}

\margin{85}

% Visual collation: both supplied image halves inspected; see CHECK_DECISIONS_134_7_8.tsv.

%% END SOURCE ps_134-7_085
\ifnum\EditionMode=1
\psapparatusstart
\psentry{ps_134-7_085-v0001}{1}{typewriter deletion}{\xout{of}}{}{The typewriter strikes out this material; it leaves the corrected running text.}
\psentry{ps_134-7_085-v0002}{2}{typewriter deletion}{\xout{put in}}{}{The typewriter strikes out this material; it leaves the corrected running text.}
\psentry{ps_134-7_085-v0003}{3}{deletion}{\deleted{as}}{}{The author cancels this material by hand; it leaves the corrected running text.}
\psentry{ps_134-7_085-v0004}{4}{insertion}{\inserted{this being}}{\hand{this being}}{The handwritten insertion is incorporated at its indicated position; insertion brackets leave the running text.}
\psentry{ps_134-7_085-v0005}{5}{deletion}{\deleted{--}}{}{The author cancels this material by hand; it leaves the corrected running text.}
\psentry{ps_134-7_085-v0006}{6}{insertion}{\inserted{being}}{\hand{being}}{The handwritten insertion is incorporated at its indicated position; insertion brackets leave the running text.}
\psentry{ps_134-7_085-v0007}{7}{typed insertion}{\inserted[typed]{(pointed)}}{(pointed)}{The typed insertion is incorporated at its indicated position; insertion brackets leave the running text.}
\psentry{ps_134-7_085-v0008}{8}{typed insertion}{\inserted[typed]{semisimplicial}}{semisimplicial}{The typed insertion is incorporated at its indicated position; insertion brackets leave the running text.}
\psentry{ps_134-7_085-v0009}{9}{typewriter deletion}{\xout{of}}{}{The typewriter strikes out this material; it leaves the corrected running text.}
\psentry{ps_134-7_085-v0010}{10}{typewriter deletion}{\xout{\illegible{1}}}{}{The typewriter strikes out this material; it leaves the corrected running text.}
\psentry{ps_134-7_085-v0011}{11}{typed insertion}{\inserted[typed]{is}}{is}{The typed insertion is incorporated at its indicated position; insertion brackets leave the running text.}
\psentry{ps_134-7_085-v0012}{12}{typewriter deletion}{\xout{\unsure{notati-}}}{}{The typewriter strikes out this material; it leaves the corrected running text.}
\psentry{ps_134-7_085-v0013}{13}{insertion}{\inserted{notations}}{\hand{notations}}{The handwritten insertion is incorporated at its indicated position; insertion brackets leave the running text.}
\psentry{ps_134-7_085-v0014}{14}{typewriter deletion}{\xout{is a}}{}{The typewriter strikes out this material; it leaves the corrected running text.}
\psentry{ps_134-7_085-v0015}{15}{typewriter deletion}{\xout{are}}{}{The typewriter strikes out this material; it leaves the corrected running text.}
\psentry{ps_134-7_085-v0016}{16}{insertion}{\inserted{a}}{\hand{a}}{The handwritten insertion is incorporated at its indicated position; insertion brackets leave the running text.}
\psentry{ps_134-7_085-v0017}{17}{insertion}{\inserted{matter}}{\hand{matter}}{The handwritten insertion is incorporated at its indicated position; insertion brackets leave the running text.}
\psentry{ps_134-7_085-v0018}{18}{typewriter deletion}{\xout{\unsure{(varia-}}}{}{The typewriter strikes out this material; it leaves the corrected running text.}
\psentry{ps_134-7_085-v0019}{19}{typewriter deletion}{\xout{ble)}}{}{The typewriter strikes out this material; it leaves the corrected running text.}
\psentry{ps_134-7_085-v0020}{20}{typed insertion}{\inserted[typed]{by}}{by}{The typed insertion is incorporated at its indicated position; insertion brackets leave the running text.}
\psentry{ps_134-7_085-v0021}{21}{typewriter deletion}{\xout{through}}{}{The typewriter strikes out this material; it leaves the corrected running text.}
\psentry{ps_134-7_085-v0022}{22}{typewriter overwriting}{\over[typed]{n}{n}}{n}{The superposed typed reading replaces the underlying reading.}
\psentry{ps_134-7_085-v0023}{23}{typewriter deletion}{\xout{\unsure{more or}}}{}{The typewriter strikes out this material; it leaves the corrected running text.}
\psentry{ps_134-7_085-v0024}{24}{typewriter overwriting}{\over[typed]{o}{o}}{o}{The superposed typed reading replaces the underlying reading.}
\psapparatusend
\fi
\pssource{ps_134-7_086}{134-7/86}{https://grothendieck.umontpellier.fr/web/viewer.html?file=../134-7.pdf\#page=87}
\psnumber{ps_134-7_086-v0001}{1}
\psnumber{ps_134-7_086-v0002}{2}
\psnumber{ps_134-7_086-v0003}{3}
\psnumber{ps_134-7_086-v0004}{4}
\psnumber{ps_134-7_086-v0005}{5}
\psnumber{ps_134-7_086-v0006}{6}
\psnumber{ps_134-7_086-v0007}{7}
\psnumber{ps_134-7_086-v0008}{8}
\psnumber{ps_134-7_086-v0009}{9}
\psnumber{ps_134-7_086-v0010}{10}
\psnumber{ps_134-7_086-v0011}{11}
\psnumber{ps_134-7_086-v0012}{12}
\psnumber{ps_134-7_086-v0013}{13}
\psnumber{ps_134-7_086-v0014}{14}
\psnumber{ps_134-7_086-v0015}{15}
%% BEGIN SOURCE ps_134-7_086
% Pursuing Stacks (cotes 134-2 to 134-8): archive cote 134-7, site page 86 (PDF page 87), typescript page 528.
% typescript. Draft compared: ps6.tex:4326-4385.
% AI-transcribed by Claude Opus 5.5 (claude-opus-5-5, adaptive thinking, effort high), Message Batches API, batch ps4 (msgbatch_012Aoc1NuskATPsyBeKSYK7J), 2026-11-12; reviewed against both supplied scan halves; see CHECK_DECISIONS.tsv.
\scanpage{134-7/86}
\mspage{528}

\var{\inserted{ges}}{\hand{ges}}{insertion}{ps_134-7_086-v0001}fibration and cofibration sequences i.e. transforms fibers in\var{\inserted{to}}{\hand{to}}{insertion}{ps_134-7_086-v0002} cofibers
and vice-versa, and the pair of functors interchanging $L\overset{\hand{\sim}}{H}.$ and $L\hand{\pi}.$
on the one hand, $\overset{\hand{\sim}}{K}$ and $\overset{\hand{\sim}}{S}$ on the other. This heuristic formulation is
compatible with all we have stated so far, except one thing - namely
\margin{\textcircled{$\Sigma_{ab}$}}the suspension functor $\hand{\Sigma}_{ab}$ in Hotab$_{\bullet}$(k) is fully faithful, whereas the
(supposedly ``dual'') loop-space functor $\hand{\Omega}^{\unsure{o|0}}_{ab}$ is not. Thus it will be more
accurate to say that, if we view the formulaire developped so far as
being \var{\xout{an}}{}{typewriter deletion}{ps_134-7_086-v0003} the description of a certain structure type, whose basic
ingredients are two categories $\hand{\mathcal{H}}_{ab}$ and $\hand{\mathcal{H}}$ , \var{\inserted[typed]{endowed with fiber and cofiber operations,}}{endowed with fiber and cofiber operations,}{typed insertion}{ps_134-7_086-v0004} tied by four functors
as above satisfying a bunch of properties, then the axioms are autodual
in an obvious sense; namely if they are satisfied for a system\var{\xout{\unsure{s}}}{}{typewriter deletion}{ps_134-7_086-v0005} \var{\xout{they}}{}{typewriter deletion}{ps_134-7_086-v0006}
($\hand{\mathcal{H}}_{ab}$, $\hand{\mathcal{H}}$, $L\overset{\hand{\sim}}{H}.$,$\overset{\hand{\sim}}{K}$,$\overset{\hand{\sim}}{S}$,$L\hand{\pi}.$) , they are equally satisfied by the system
($\hand{\mathcal{H}}_{ab}{}^{op}$,$\hand{\mathcal{H}}^{op}$, $L\hand{\pi}.$,$\overset{\hand{\sim}}{S}$,$\overset{\hand{\sim}}{K}$,$L\overset{\hand{\sim}}{H}.$). (Of course, among the properties of the
structure type, we are \uline{not} going to include that the\var{\xout{ an}}{}{typewriter deletion}{ps_134-7_086-v0007} \var{\xout{loop-space}}{}{typewriter deletion}{ps_134-7_086-v0008} \var{\inserted[typed]{suspension}}{suspension}{typed insertion}{ps_134-7_086-v0009} functor
in $\hand{\mathcal{H}}_{ab}$ is fully faithful !). As already stated before, the main rea-%
son for introducing the fourth functor $\overset{\hand{\sim}}{S}$ in the picture, was because it
was felt that this was lacking in \var{\xout{the picture}}{}{typewriter deletion}{ps_134-7_086-v0010} order to round it up. Thus
for any kind of notion or statement in this set-up, suggested by some
kind or other of geometric insight, it becomes automatic to look at
\margin{$\ulcorner$}it's dual and see whether it makes sense.\hand{$\rceil$}

As for a formal autoduality,
we already noticed there is none, even in Hotab$_{\bullet}$(k) just by itself.
Thinking of this category as D.(Ab$_{k}$), namely as a full subcategory of
the derived category D(Ab$_{k}$), we cannot help\var{\xout{\illegible{1}}}{}{typewriter deletion}{ps_134-7_086-v0011} though but thinking of
the standard ``dualizing'' functor
\[ L. \hand{\longmapsto} RHOM(L.,k) \;:\; D^{-}(Ab_{k}) \hand{\longrightarrow} D^{+}(Ab_{k}) \quad , \]
inducing a perfect duality \var{\xout{between}}{}{typewriter deletion}{ps_134-7_086-v0012} within the category of ``perfect''
objects in D(Ab$_{k}$), namely objects which are isomorphic to those which
may be described by complexes in Ab$_{k}$ having a bounded span of degrees,
and all of whose components are projective of finite type. But this aut\unsure{o}%
% NOTE: the end of this line is cut at the right edge of the scan
duality of course transforms \uline{chain} complexes into \uline{cochain} complexes
- we have to shift these in order to get chain complexes again. This
suggests that maybe we'll hit upon \var{\xout{actual}}{}{typewriter deletion}{ps_134-7_086-v0013} an actual autoduality,
provided we go over from Hot$_{\bullet}$(k) to a suitable``stabilized''category,
the suspension category say (deduced from the initial one by \var{\xout{making}}{}{typewriter deletion}{ps_134-7_086-v0014}
introducing formally a quasi-inverse for the suspension functor). The
homology functor $L\overset{\hand{\sim}}{H}.$ extends to a functor from the suspension category
to the corresponding category for Hotab$_{\bullet}$(k), say D$^{\unsure{-}}$(Ab$_{k}$). T\var{\over[typed]{h}{h}}{h}{typewriter overwriting}{ps_134-7_086-v0015}e $L\hand{\pi}.$

\margin{86}

% Visual collation: both supplied image halves inspected; see CHECK_DECISIONS_134_7_8.tsv.

%% END SOURCE ps_134-7_086
\ifnum\EditionMode=1
\psapparatusstart
\psentry{ps_134-7_086-v0001}{1}{insertion}{\inserted{ges}}{\hand{ges}}{The handwritten insertion is incorporated at its indicated position; insertion brackets leave the running text.}
\psentry{ps_134-7_086-v0002}{2}{insertion}{\inserted{to}}{\hand{to}}{The handwritten insertion is incorporated at its indicated position; insertion brackets leave the running text.}
\psentry{ps_134-7_086-v0003}{3}{typewriter deletion}{\xout{an}}{}{The typewriter strikes out this material; it leaves the corrected running text.}
\psentry{ps_134-7_086-v0004}{4}{typed insertion}{\inserted[typed]{endowed with fiber and cofiber operations,}}{endowed with fiber and cofiber operations,}{The typed insertion is incorporated at its indicated position; insertion brackets leave the running text.}
\psentry{ps_134-7_086-v0005}{5}{typewriter deletion}{\xout{\unsure{s}}}{}{The typewriter strikes out this material; it leaves the corrected running text.}
\psentry{ps_134-7_086-v0006}{6}{typewriter deletion}{\xout{they}}{}{The typewriter strikes out this material; it leaves the corrected running text.}
\psentry{ps_134-7_086-v0007}{7}{typewriter deletion}{\xout{ an}}{}{The typewriter strikes out this material; it leaves the corrected running text.}
\psentry{ps_134-7_086-v0008}{8}{typewriter deletion}{\xout{loop-space}}{}{The typewriter strikes out this material; it leaves the corrected running text.}
\psentry{ps_134-7_086-v0009}{9}{typed insertion}{\inserted[typed]{suspension}}{suspension}{The typed insertion is incorporated at its indicated position; insertion brackets leave the running text.}
\psentry{ps_134-7_086-v0010}{10}{typewriter deletion}{\xout{the picture}}{}{The typewriter strikes out this material; it leaves the corrected running text.}
\psentry{ps_134-7_086-v0011}{11}{typewriter deletion}{\xout{\illegible{1}}}{}{The typewriter strikes out this material; it leaves the corrected running text.}
\psentry{ps_134-7_086-v0012}{12}{typewriter deletion}{\xout{between}}{}{The typewriter strikes out this material; it leaves the corrected running text.}
\psentry{ps_134-7_086-v0013}{13}{typewriter deletion}{\xout{actual}}{}{The typewriter strikes out this material; it leaves the corrected running text.}
\psentry{ps_134-7_086-v0014}{14}{typewriter deletion}{\xout{making}}{}{The typewriter strikes out this material; it leaves the corrected running text.}
\psentry{ps_134-7_086-v0015}{15}{typewriter overwriting}{\over[typed]{h}{h}}{h}{The superposed typed reading replaces the underlying reading.}
\psapparatusend
\fi
\pssource{ps_134-7_087}{134-7/87}{https://grothendieck.umontpellier.fr/web/viewer.html?file=../134-7.pdf\#page=88}
\psnumber{ps_134-7_087-v0001}{1}
\psnumber{ps_134-7_087-v0002}{2}
\psnumber{ps_134-7_087-v0003}{3}
\psnumber{ps_134-7_087-v0004}{4}
\psnumber{ps_134-7_087-v0005}{5}
\psnumber{ps_134-7_087-v0006}{6}
\psnumber{ps_134-7_087-v0007}{7}
\psnumber{ps_134-7_087-v0008}{8}
\psnumber{ps_134-7_087-v0009}{9}
\psnumber{ps_134-7_087-v0010}{10}
\psnumber{ps_134-7_087-v0011}{11}
\psnumber{ps_134-7_087-v0012}{12}
\psnumber{ps_134-7_087-v0013}{13}
\psnumber{ps_134-7_087-v0014}{14}
\psnumber{ps_134-7_087-v0015}{15}
\psnumber{ps_134-7_087-v0016}{16}
%% BEGIN SOURCE ps_134-7_087
% Pursuing Stacks (cotes 134-2 to 134-8): archive cote 134-7, site page 87 (PDF page 88), typescript page 529.
% typescript. Draft compared: ps6.tex:4373-4449.
% AI-transcribed by Claude Opus 5.5 (claude-opus-5-5, adaptive thinking, effort high), Message Batches API, batch ps4 (msgbatch_012Aoc1NuskATPsyBeKSYK7J), 2026-11-12; reviewed against both supplied scan halves; see CHECK_DECISIONS.tsv.
\scanpage{134-7/87}
\mspage{529}

\var{\inserted{functor}}{\hand{functor}}{insertion}{ps_134-7_087-v0001} though is lost on our way - I am affraid I am confusing the kind of dua-%
lity I am after here, with a rather different type of duality, kin to
Poincaré duality, and \var{\inserted[typed]{discovered}}{discovered}{typed insertion}{ps_134-7_087-v0002} \var{\xout{due}}{}{typewriter deletion}{ps_134-7_087-v0003} I believe by J.H.C. Whitehead, in the context
of spaces having the homotopy type of a finite complex. For such a
complex, the idea \var{\inserted[typed]{(as far as I remember)}}{(as far as I remember)}{typed insertion}{ps_134-7_087-v0004} is to embed X into a large-dimensional sphere, and
to take the complement $X\hand{^{\vee}}$ of an open tubular neighbourhood of X. Up to
suspension, the homotopy type of X does not depend on the choices made,
and \var{\inserted[typed]{(shifting back by $n$)}}{(shifting back by $n$)}{typed insertion}{ps_134-7_087-v0005} we thus get a canonical object in the suspension category, depending
contravariantly on X. T\var{\over[typed]{h}{h}}{h}{typewriter overwriting}{ps_134-7_087-v0006}e functor $X\hand{\mapsto}X\hand{^{\vee}}$ (if I got it right) is an auto-%
duality of the relevant full subcategory of the suspension category,
and the (reduced) homology functor $L\overset{\hand{\sim}}{H}.$ \var{\xout{commutes to}}{}{typewriter deletion}{ps_134-7_087-v0007} maps this subcategory
into perfect complexes, and commutes to the autodualities. The whole
story seems tailored towards a study of duality relations for the functor
$L\overset{\hand{\sim}}{H}.$ exclusively - without any reference to the homotopy invariants $\hand{\pi}_{i}$.
I do'nt know if one can devise a similar story for $L\hand{\pi}.$ , by stabilizing
with respect to the loop space functor (or is this just nonsense ?). At
any rate, there does'nt seem any autoduality in view, exchanging the
homology and homotopy invariants ...

E) \uline{Conservativity properties:} The functors $L\overset{\hand{\sim}}{H}.$ and $L\hand{\pi}.$ are both
conservative, i.e. a map in $\mathrm{Hot}_{\bullet}(k)$ \var{\inserted[typed]{which}}{which}{typed insertion}{ps_134-7_087-v0008} \var{\xout{mapped}}{}{typewriter deletion}{ps_134-7_087-v0009} by either of these functors
becomes an isomorphism, is an isomorphism. (Here of course the 0-connec%
\hand{-}tedness assumption for the \var{\xout{objects we are}}{}{typewriter deletion}{ps_134-7_087-v0010} homotopy types we are working
with is essential, as far as the functor $L\hand{\pi}.$ is concerned at any rate.)

F) \uline{Base change properties}: All four functors, and the adjunction
and inversion maps (7) and (9), are compatible with ring extension
$k\hand{\to} k'$ , it being understood that such a ring homomorphism defines
functors
\[ (12)\qquad \mathrm{Hotab}_{\bullet}(k)\hand{\longrightarrow}\mathrm{Hotab}_{\bullet}(k') \ , \quad \mathrm{Hot}_{o}(k)\hand{\longrightarrow}\mathrm{Hot}_{\bullet}(k') \]
compatible with the fibration and cofibration structures.

In opposite direction, there should be too \var{\inserted{a}}{\hand{a}}{insertion}{ps_134-7_087-v0011} ``restriction'' functor\var{\xout{s}}{}{typewriter deletion}{ps_134-7_087-v0012}
(less important though I feel), I did'nt try to find out what \var{\xout{their}}{}{typewriter deletion}{ps_134-7_087-v0013} \var{\inserted{it's}}{\hand{it's}}{insertion}{ps_134-7_087-v0014}
formal properties should be with respect to the formalism developped her\unsure{e}

G) \uline{Comparison with ordinary homotopy types }: We got functors
\[ (13)\qquad \mathrm{Hotab}_{\bullet}(k)\hand{\longrightarrow} \mathrm{Hotab} = D.(Ab) \ , \quad \mathrm{Hot}_{o}(k)\hand{\longrightarrow} \mathrm{Hot}_{\bullet} \quad , \]
(the first being interpreted as ``forget k'' functor, the second \var{\deleted{\illegible{1}}}{}{deletion}{ps_134-7_087-v0015} \var{\inserted{as a}}{\hand{as a}}{insertion}{ps_134-7_087-v0016}
sections functor on semisimplicial schematic models). This pair of

\margin{87}

% Visual collation: both supplied image halves inspected; see CHECK_DECISIONS_134_7_8.tsv.

%% END SOURCE ps_134-7_087
\ifnum\EditionMode=1
\psapparatusstart
\psentry{ps_134-7_087-v0001}{1}{insertion}{\inserted{functor}}{\hand{functor}}{The handwritten insertion is incorporated at its indicated position; insertion brackets leave the running text.}
\psentry{ps_134-7_087-v0002}{2}{typed insertion}{\inserted[typed]{discovered}}{discovered}{The typed insertion is incorporated at its indicated position; insertion brackets leave the running text.}
\psentry{ps_134-7_087-v0003}{3}{typewriter deletion}{\xout{due}}{}{The typewriter strikes out this material; it leaves the corrected running text.}
\psentry{ps_134-7_087-v0004}{4}{typed insertion}{\inserted[typed]{(as far as I remember)}}{(as far as I remember)}{The typed insertion is incorporated at its indicated position; insertion brackets leave the running text.}
\psentry{ps_134-7_087-v0005}{5}{typed insertion}{\inserted[typed]{(shifting back by $n$)}}{(shifting back by $n$)}{The typed insertion is incorporated at its indicated position; insertion brackets leave the running text.}
\psentry{ps_134-7_087-v0006}{6}{typewriter overwriting}{\over[typed]{h}{h}}{h}{The superposed typed reading replaces the underlying reading.}
\psentry{ps_134-7_087-v0007}{7}{typewriter deletion}{\xout{commutes to}}{}{The typewriter strikes out this material; it leaves the corrected running text.}
\psentry{ps_134-7_087-v0008}{8}{typed insertion}{\inserted[typed]{which}}{which}{The typed insertion is incorporated at its indicated position; insertion brackets leave the running text.}
\psentry{ps_134-7_087-v0009}{9}{typewriter deletion}{\xout{mapped}}{}{The typewriter strikes out this material; it leaves the corrected running text.}
\psentry{ps_134-7_087-v0010}{10}{typewriter deletion}{\xout{objects we are}}{}{The typewriter strikes out this material; it leaves the corrected running text.}
\psentry{ps_134-7_087-v0011}{11}{insertion}{\inserted{a}}{\hand{a}}{The handwritten insertion is incorporated at its indicated position; insertion brackets leave the running text.}
\psentry{ps_134-7_087-v0012}{12}{typewriter deletion}{\xout{s}}{}{The typewriter strikes out this material; it leaves the corrected running text.}
\psentry{ps_134-7_087-v0013}{13}{typewriter deletion}{\xout{their}}{}{The typewriter strikes out this material; it leaves the corrected running text.}
\psentry{ps_134-7_087-v0014}{14}{insertion}{\inserted{it's}}{\hand{it's}}{The handwritten insertion is incorporated at its indicated position; insertion brackets leave the running text.}
\psentry{ps_134-7_087-v0015}{15}{deletion}{\deleted{\illegible{1}}}{}{The author cancels this material by hand; it leaves the corrected running text.}
\psentry{ps_134-7_087-v0016}{16}{insertion}{\inserted{as a}}{\hand{as a}}{The handwritten insertion is incorporated at its indicated position; insertion brackets leave the running text.}
\psapparatusend
\fi
\pssource{ps_134-7_088}{134-7/88}{https://grothendieck.umontpellier.fr/web/viewer.html?file=../134-7.pdf\#page=89}
\psnumber{ps_134-7_088-v0001}{1}
\psnumber{ps_134-7_088-v0002}{2}
\psnumber{ps_134-7_088-v0003}{3}
\psnumber{ps_134-7_088-v0004}{4}
\psnumber{ps_134-7_088-v0005}{5}
\psnumber{ps_134-7_088-v0006}{6}
\psnumber{ps_134-7_088-v0007}{7}
\psnumber{ps_134-7_088-v0008}{8}
\psnumber{ps_134-7_088-v0009}{9}
\psnumber{ps_134-7_088-v0010}{10}
\psnumber{ps_134-7_088-v0011}{11}
\psnumber{ps_134-7_088-v0012}{12}
\psnumber{ps_134-7_088-v0013}{13}
\psnumber{ps_134-7_088-v0014}{14}
%% BEGIN SOURCE ps_134-7_088
% Pursuing Stacks (cotes 134-2 to 134-8): archive cote 134-7, site page 88 (PDF page 89), typescript page 530.
% typescript. Draft compared: ps6.tex:4433-4501.
% AI-transcribed by Claude Opus 5.5 (claude-opus-5-5, adaptive thinking, effort high), Message Batches API, batch ps4 (msgbatch_012Aoc1NuskATPsyBeKSYK7J), 2026-11-12; reviewed against both supplied scan halves; see CHECK_DECISIONS.tsv.
\scanpage{134-7/88}
\mspage{530}
functors is compatible with the functors $\overset{\hand{\sim}}{K}$ , but \var{\inserted[typed]{(even if k=Z)}}{(even if k=Z)}{typed insertion}{ps_134-7_088-v0001} definitely \uline{not} with \var{\xout{its}}{}{typewriter deletion}{ps_134-7_088-v0002} \var{\inserted[typed]{their}}{their}{typed insertion}{ps_134-7_088-v0003}
left adjoints $L\overset{\hand{\sim}}{H}.$ namely with homology. Despite this fact, we hope \var{\inserted{if k=Z}}{\hand{if k=Z}}{insertion}{ps_134-7_088-v0004} that
\var{\inserted[typed]{the functors}}{the functors}{typed insertion}{ps_134-7_088-v0005} (13) are compatible with the functors $\overset{\hand{\sim}}{S}$ (assuming that $\overset{\hand{\sim}}{S}$ can actually
be constructed in the discrete set-up too)\hand{,} \var{\deleted{\hand{if k=Z},}}{}{deletion}{ps_134-7_088-v0006} so that spheres are
transformed into spheres. Compatibility with L\hand{$\Pi$}. does'nt have a meaning
strictly speaking, as this functor is not defined on ordinary homotopy
types, however the functors \hand{$\Pi$}$_i$ are defined in both contexts, and the
sections functor should commute to these (for arbitrary k). Thus the
only serious uncompatibility trouble concerns the homology functor LH. ,
\margin{\textcircled{H}} and should not arise however for sphere-like objects S(\hand{$H$},n). \var{\xout{-}}{}{typewriter deletion}{ps_134-7_088-v0007} I for-%
got to state \var{\inserted{from outset}}{\hand{from outset}}{insertion}{ps_134-7_088-v0008} that the two functors (13) are expected to \var{\xout{\illegible{1}}}{}{typewriter deletion}{ps_134-7_088-v0009} respect
fibration and cofibration sequence structures, of course.

13.10.

125. During the last eight days, I have been \var{\xout{involved}}{}{typewriter deletion}{ps_134-7_088-v0010} busy with a
number of things, which left little leisure for mathematical ponderings.
I got a letter from H.J. Baues, who had seen my notes, which induced
him to send me his preprint ``On the homotopy classification problem''
(Chapters I to V + chapter Ext). I spent two evenings looking through
parts of these, where Baues carries as far as possible (namely quite far
indeed) the homotopy formalism\var{\xout{x}}{}{typewriter deletion}{ps_134-7_088-v0011} \var{\xout{with}}{}{typewriter deletion}{ps_134-7_088-v0012} in the context of his so-called
``cofibration (or fibration) categories'', using as his leading thread
his ideas on ``obstruction theory''. He makes the point that he tried by
his basic notion to pinpoi\unsure{\hand{/}}nt the weakest axiomatic set-up, sufficient
however \hand{f}or develop\var{\deleted{p}}{}{deletion}{ps_134-7_088-v0013}ing all the major familiar (and even some not at all
familiar !) notions, operations and statements of usual homotopy theory.
In his letter, he suggested that maybe in any ``universe'' where homotopy
constructions make a sense, one or the other of his two mutually dual
axiomatic set-ups should be around. Such suggestion was of course quite
interesting for my present reflections, as I do have the hope indeed
that there exists a ``universe'' of schematic homotopy types, which may
be described in terms of the models (namely semisimplicial unipotent
bundles) I have been using so far, or at any rate by closely related
kind of models. More specifically, I do hope for the two kind of opera-%
tions to make sense, namely ``integration'' (including homotopy cofibers\var{\xout{)}}{}{typewriter deletion}{ps_134-7_088-v0014}
and suspensions) and ``cointegration'' (including homotopy fibers and
loop objects), which in Baues set-up should correspond on the category
of models to a structure of a co\hand{f}ibration category and a fibration cate-%
gory respectively. As the kind of models I have been working with
\margin{88}

% Visual collation: both supplied image halves inspected; see CHECK_DECISIONS_134_7_8.tsv.

%% END SOURCE ps_134-7_088
\ifnum\EditionMode=1
\psapparatusstart
\psentry{ps_134-7_088-v0001}{1}{typed insertion}{\inserted[typed]{(even if k=Z)}}{(even if k=Z)}{The typed insertion is incorporated at its indicated position; insertion brackets leave the running text.}
\psentry{ps_134-7_088-v0002}{2}{typewriter deletion}{\xout{its}}{}{The typewriter strikes out this material; it leaves the corrected running text.}
\psentry{ps_134-7_088-v0003}{3}{typed insertion}{\inserted[typed]{their}}{their}{The typed insertion is incorporated at its indicated position; insertion brackets leave the running text.}
\psentry{ps_134-7_088-v0004}{4}{insertion}{\inserted{if k=Z}}{\hand{if k=Z}}{The handwritten insertion is incorporated at its indicated position; insertion brackets leave the running text.}
\psentry{ps_134-7_088-v0005}{5}{typed insertion}{\inserted[typed]{the functors}}{the functors}{The typed insertion is incorporated at its indicated position; insertion brackets leave the running text.}
\psentry{ps_134-7_088-v0006}{6}{deletion}{\deleted{\hand{if k=Z},}}{}{The author cancels this material by hand; it leaves the corrected running text.}
\psentry{ps_134-7_088-v0007}{7}{typewriter deletion}{\xout{-}}{}{The typewriter strikes out this material; it leaves the corrected running text.}
\psentry{ps_134-7_088-v0008}{8}{insertion}{\inserted{from outset}}{\hand{from outset}}{The handwritten insertion is incorporated at its indicated position; insertion brackets leave the running text.}
\psentry{ps_134-7_088-v0009}{9}{typewriter deletion}{\xout{\illegible{1}}}{}{The typewriter strikes out this material; it leaves the corrected running text.}
\psentry{ps_134-7_088-v0010}{10}{typewriter deletion}{\xout{involved}}{}{The typewriter strikes out this material; it leaves the corrected running text.}
\psentry{ps_134-7_088-v0011}{11}{typewriter deletion}{\xout{x}}{}{The typewriter strikes out this material; it leaves the corrected running text.}
\psentry{ps_134-7_088-v0012}{12}{typewriter deletion}{\xout{with}}{}{The typewriter strikes out this material; it leaves the corrected running text.}
\psentry{ps_134-7_088-v0013}{13}{deletion}{\deleted{p}}{}{The author cancels this material by hand; it leaves the corrected running text.}
\psentry{ps_134-7_088-v0014}{14}{typewriter deletion}{\xout{)}}{}{The typewriter strikes out this material; it leaves the corrected running text.}
\psapparatusend
\fi
\pssource{ps_134-7_089}{134-7/89}{https://grothendieck.umontpellier.fr/web/viewer.html?file=../134-7.pdf\#page=90}
\psnumber{ps_134-7_089-v0001}{1}
\psnumber{ps_134-7_089-v0002}{2}
\psnumber{ps_134-7_089-v0003}{3}
\psnumber{ps_134-7_089-v0004}{4}
\psnumber{ps_134-7_089-v0005}{5}
\psnumber{ps_134-7_089-v0006}{6}
\psnumber{ps_134-7_089-v0007}{7}
\psnumber{ps_134-7_089-v0008}{8}
\psnumber{ps_134-7_089-v0009}{9}
\psnumber{ps_134-7_089-v0010}{10}
\psnumber{ps_134-7_089-v0011}{11}
\psnumber{ps_134-7_089-v0012}{12}
\psnumber{ps_134-7_089-v0013}{13}
\psnumber{ps_134-7_089-v0014}{14}
\psnumber{ps_134-7_089-v0015}{15}
\psnumber{ps_134-7_089-v0016}{16}
%% BEGIN SOURCE ps_134-7_089
% Pursuing Stacks (cotes 134-2 to 134-8): archive cote 134-7, site page 89 (PDF page 90), typescript page 531.
% typescript. Draft compared: ps6.tex:4489-4547.
% AI-transcribed by Claude Opus 5.5 (claude-opus-5-5, adaptive thinking, effort high), Message Batches API, batch ps4 (msgbatch_012Aoc1NuskATPsyBeKSYK7J), 2026-11-12; reviewed against both supplied scan halves; see CHECK_DECISIONS.tsv.
\scanpage{134-7/89}
\mspage{531}

do'nt allow \var{\xout{for}}{}{typewriter deletion}{ps_134-7_089-v0001} an  amalgamated sum (or ``push-out'') construction,  except
in the trivial case when one of the two arrows to be amalgamated is an
isomorphism,  it is clear that we cannot hope to get with these models
a cofibration category.  There may be however a fibration category struc-%
ture,  and I started playing around a little\var{\xout{x}}{}{typewriter deletion}{ps_134-7_089-v0002} with a possible notion of
Kan fibration -  without coming to a definite conclusion yet,  however.
I feel I should'nt dwell much longer still,  for the time being,  on get-%
ting off the ground a theory of schematic homotopy types,  maybe I'll
come back upon it later,  after next chapter,  when (hopefully) I'll be
a little more``in the game'' of Kan type conditions,  closed model category
structures  and Baues' ``halved'' variants of these.

At any rate,  whether or not a \var{\inserted[typed]{homotopy}}{homotopy}{typed insertion}{ps_134-7_089-v0003} cofiber construction  can be carried
through for schematic homotopy types,  it seems rather clear to me that
Baues' suggestion or expectation\hand{,} about the set-up of cofibration and
fibration categories he ended up with,  is not quite justified.  Already
by the time \var{\inserted{(in the late \unsure{Sixties}) when}}{\hand{(in the late \unsure{Sixties}) when}}{insertion}{ps_134-7_089-v0004} I first heard from Quillen about his approach towards axioma-%
tic homotopy theory,  it struck me that his approach (and the same applie\unsure{d}
to Baues one\unsure{\hand{'}}s), \var{\xout{is}}{}{typewriter deletion}{ps_134-7_089-v0005}  when applied say to semisimplicial objects in some
category A as ``models'',  is pretty strongly relying on the existence of
``enough'' projectives (or dually,  of enough injectives,  if working with
cosemisimplicial objects) in A.  When A is a topos say,  then the
category $\underline{\mathrm{Hom}}(\hand{\Delta}^{\circ},A)$ of semisimplicial objects of A does'nt carry \var{\xout{\illegible{1}}}{}{typewriter deletion}{ps_134-7_089-v0006}\var{\inserted{it seems}}{\hand{it seems}}{insertion}{ps_134-7_089-v0007}
(except in very special cases,  say of a totally disconnected topos) a
\var{\inserted[typed]{reasonable}}{reasonable}{typed insertion}{ps_134-7_089-v0008}structure of a \var{\xout{closed}}{}{typewriter deletion}{ps_134-7_089-v0009}\var{\inserted[typed]{Quillen}}{Quillen}{typed insertion}{ps_134-7_089-v0010} model category,  nor even (I would think) of a
cofibration or a fibration category in the sense of Baues;  however I am
pretty sure that the derived category \var{\xout{\unsure{(in th D(A)}}}{}{typewriter deletion}{ps_134-7_089-v0011}  (with respect to the
\margin{$\times$} notion of  weak equivalence introduced by Illusie \hand{[\quad]})  has important
geometric meaning and is indeed a ``universe for homotopy theory''  - and
the first steps in developing such theory have been taken by Illusie
already in Chapter I of his thesis.  Quite similarly,  if A is an abelian
category,  Verdier's derived category  D(A) , deduced from the category
C(A) of complexes in A by localizing with respect to the set W of quasi-%
isomorphisms,  \var{\xout{it}}{}{typewriter deletion}{ps_134-7_089-v0012} does allow for a homotopy formalism (as does any
``triangulated category''in the sense of Verdier),  however it does'nt seem
that this formalism may be deduced from a structure of a cofibration cate
\hand{-}gory or a fibration category on C(A).   T\var{\over[typed]{h}{h}}{h}{typewriter overwriting}{ps_134-7_089-v0013}e same I guess holds for the
\margin{\textcircled{\shortstack{$C^{-}(A)$\\$D^{-}(A)$}}} subcategories $C^{\var{\over{\unsure{+}}{-}}{\hand{-}}{writing over}{ps_134-7_089-v0014}}(A)$ and $C^{+}(A)$,  giving rise to $D^{\var{\over{\unsure{+}}{-}}{\hand{-}}{writing over}{ps_134-7_089-v0015}}(A)$ and $D^{+}(A)$,  which
are more important still than D(A) in\var{\xout{t}}{}{typewriter deletion}{ps_134-7_089-v0016} the everyday cohomology formalism
of ``spaces'' of all kinds (namely,  essentially,  of ringed topoi).  When

\margin{89}

% Visual collation: both supplied image halves inspected; see CHECK_DECISIONS_134_7_8.tsv.

%% END SOURCE ps_134-7_089
\ifnum\EditionMode=1
\psapparatusstart
\psentry{ps_134-7_089-v0001}{1}{typewriter deletion}{\xout{for}}{}{The typewriter strikes out this material; it leaves the corrected running text.}
\psentry{ps_134-7_089-v0002}{2}{typewriter deletion}{\xout{x}}{}{The typewriter strikes out this material; it leaves the corrected running text.}
\psentry{ps_134-7_089-v0003}{3}{typed insertion}{\inserted[typed]{homotopy}}{homotopy}{The typed insertion is incorporated at its indicated position; insertion brackets leave the running text.}
\psentry{ps_134-7_089-v0004}{4}{insertion}{\inserted{(in the late \unsure{Sixties}) when}}{\hand{(in the late \unsure{Sixties}) when}}{The handwritten insertion is incorporated at its indicated position; insertion brackets leave the running text.}
\psentry{ps_134-7_089-v0005}{5}{typewriter deletion}{\xout{is}}{}{The typewriter strikes out this material; it leaves the corrected running text.}
\psentry{ps_134-7_089-v0006}{6}{typewriter deletion}{\xout{\illegible{1}}}{}{The typewriter strikes out this material; it leaves the corrected running text.}
\psentry{ps_134-7_089-v0007}{7}{insertion}{\inserted{it seems}}{\hand{it seems}}{The handwritten insertion is incorporated at its indicated position; insertion brackets leave the running text.}
\psentry{ps_134-7_089-v0008}{8}{typed insertion}{\inserted[typed]{reasonable}}{reasonable}{The typed insertion is incorporated at its indicated position; insertion brackets leave the running text.}
\psentry{ps_134-7_089-v0009}{9}{typewriter deletion}{\xout{closed}}{}{The typewriter strikes out this material; it leaves the corrected running text.}
\psentry{ps_134-7_089-v0010}{10}{typed insertion}{\inserted[typed]{Quillen}}{Quillen}{The typed insertion is incorporated at its indicated position; insertion brackets leave the running text.}
\psentry{ps_134-7_089-v0011}{11}{typewriter deletion}{\xout{\unsure{(in th D(A)}}}{}{The typewriter strikes out this material; it leaves the corrected running text.}
\psentry{ps_134-7_089-v0012}{12}{typewriter deletion}{\xout{it}}{}{The typewriter strikes out this material; it leaves the corrected running text.}
\psentry{ps_134-7_089-v0013}{13}{typewriter overwriting}{\over[typed]{h}{h}}{h}{The superposed typed reading replaces the underlying reading.}
\psentry{ps_134-7_089-v0014}{14}{writing over}{$\over{\unsure{+}}{-}$}{$\hand{-}$}{The superposed handwritten reading replaces the underlying reading.}
\psentry{ps_134-7_089-v0015}{15}{writing over}{$\over{\unsure{+}}{-}$}{$\hand{-}$}{The superposed handwritten reading replaces the underlying reading.}
\psentry{ps_134-7_089-v0016}{16}{typewriter deletion}{\xout{t}}{}{The typewriter strikes out this material; it leaves the corrected running text.}
\psapparatusend
\fi
\pssource{ps_134-7_090}{134-7/90}{https://grothendieck.umontpellier.fr/web/viewer.html?file=../134-7.pdf\#page=91}
\psnumber{ps_134-7_090-v0001}{1}
\psnumber{ps_134-7_090-v0002}{2}
\psnumber{ps_134-7_090-v0003}{3}
\psnumber{ps_134-7_090-v0004}{4}
\psnumber{ps_134-7_090-v0005}{5}
\psnumber{ps_134-7_090-v0006}{6}
\psnumber{ps_134-7_090-v0007}{7}
\psnumber{ps_134-7_090-v0008}{8}
\psnumber{ps_134-7_090-v0009}{9}
\psnumber{ps_134-7_090-v0010}{10}
\psnumber{ps_134-7_090-v0011}{11}
\psnumber{ps_134-7_090-v0012}{12}
\psnumber{ps_134-7_090-v0013}{13}
\psnumber{ps_134-7_090-v0014}{14}
\psnumber{ps_134-7_090-v0015}{15}
\psnumber{ps_134-7_090-v0016}{16}
\psnumber{ps_134-7_090-v0017}{17}
\psnumber{ps_134-7_090-v0018}{18}
\psnumber{ps_134-7_090-v0019}{19}
\psnumber{ps_134-7_090-v0020}{20}
\psnumber{ps_134-7_090-v0021}{21}
\psnumber{ps_134-7_090-v0022}{22}
\psnumber{ps_134-7_090-v0023}{23}
\psnumber{ps_134-7_090-v0024}{24}
\psnumber{ps_134-7_090-v0025}{25}
\psnumber{ps_134-7_090-v0026}{26}
%% BEGIN SOURCE ps_134-7_090
% Pursuing Stacks (cotes 134-2 to 134-8): archive cote 134-7, site page 90 (PDF page 91), typescript page 532.
% typescript. Draft compared: ps6.tex:4533-4590.
% AI-transcribed by Claude Opus 5.5 (claude-opus-5-5, adaptive thinking, effort high), Message Batches API, batch ps4 (msgbatch_012Aoc1NuskATPsyBeKSYK7J), 2026-11-12; reviewed against both supplied scan halves; see CHECK_DECISIONS.tsv.
\scanpage{134-7/90}
\mspage{532}

A admits enough projectives or enough injectives respectively,  these
categories\var{\inserted[typed]{(as pointed out by Quillen)}}{(as pointed out by Quillen)}{typed insertion}{ps_134-7_090-v0001} are associated to closed model structures on $C^{-}(A)$ and $C^{+}(A)$
respectively,  \var{\deleted{and}}{}{deletion}{ps_134-7_090-v0002} \var{\inserted{but (it seems)}}{\hand{but (it seems)}}{insertion}{ps_134-7_090-v0003} not otherwise.

It strikes me that nobody apparently so far has tried to develop
homotopy theory,  starting as basic data with a category of models \var{\xout{a}}{}{typewriter deletion}{ps_134-7_090-v0004} $M$
together with a set $W$ of weak equivalences in $M$ ,  satisfying suitable
assumptions,  without giving moreover a notion of ``fibration'' or of
``cofibrations'' in M.  Various examples,   indluding the all-important
case (from my point of view) of (Cat),  suggest that the choice of the
notion of fibration or cofibration is'nt really so imperative,  that it
is to a certain extent arbitrary,  different choices (compatible with the
basic data W) leading to \var{essentially\hand{/}``the same''\hand{/}homotopy}{essentially ``the same'' homotopy}{authorial spacing}{ps_134-7_090-v0005} theory.  This \var{\xout{\illegible{1}}}{}{typewriter deletion}{ps_134-7_090-v0006}
is due to the fact that they lead to the same notions of ``integration'' and
``cointegration'' of homotopy types,  which depend\var{\xout{s}}{}{typewriter deletion}{ps_134-7_090-v0007} indeed only on W and
which,  in my eyes,  \var{\deleted{\xout{is}}}{}{deletion}{ps_134-7_090-v0008}\var{\inserted[typed]{are}}{are}{typed insertion}{ps_134-7_090-v0009} \uline{the} two main operations of homotopy theory
(compare section 69).  They seem to me the key for defining such a thing
as a \var{\inserted[typed]{specific}}{specific}{typed insertion}{ps_134-7_090-v0010} ``homotopy theory'' ,  independently of any particular choice of
a category of models ( + extra structure\var{\inserted[typed]{on it}}{on it}{typed insertion}{ps_134-7_090-v0011}, and notably\hand{,}weak equivalences)
used for describing it.  The precise technical notion achieving this is
of course the notion of a ``derivator'' - and I do hope that it should'nt
be too awfully hard to construct,  for instance,  the ``derivator of
schematic homotopy types'',  and maybe even \var{\over{\unsure{ch}}{ch}}{\hand{ch}}{writing over}{ps_134-7_090-v0012}aracterize it axiomatically
(as well as the derivator of usual homotopy types),  without having to
\var{\inserted[typed]{make any mention of models}}{make any mention of models}{typed insertion}{ps_134-7_090-v0013}\var{\deleted{\xout{\illegible{1}}}}{}{deletion}{ps_134-7_090-v0014} for such characterization.  One may say that,  after \var{\xout{the}}{}{typewriter deletion}{ps_134-7_090-v0015} \var{\inserted{the}}{\hand{the}}{insertion}{ps_134-7_090-v0016}\var{\inserted[typed]{one}}{one}{typed insertion}{ps_134-7_090-v0017} major
step in \var{\inserted[typed]{the}}{the}{typed insertion}{ps_134-7_090-v0018} foundations of homological \var{algebra,\hand{|}consisting}{algebra, consisting}{authorial spacing}{ps_134-7_090-v0019} in introducing \var{\xout{and}}{}{typewriter deletion}{ps_134-7_090-v0020}
the derived category of an abelian category (and systematically working
with \var{\xout{this}}{}{typewriter deletion}{ps_134-7_090-v0021} derived categories for stating the main \var{\xout{theorems}}{}{typewriter deletion}{ps_134-7_090-v0022} facts about
cohomo\hand{lo}gy of all kind of ``spaces'' namely topoi,  such as usual topologica\unsure{l}
spaces,  schemes and the like ...),  the \var{\xout{foundational}}{}{typewriter deletion}{ps_134-7_090-v0023} work on foundations
more or less stopped short,  while the next step to take was to come to
a grasp of the full structure involved in derived categories,  namely the
structure of a derivator.  And it turns out that \var{\xout{both}}{}{typewriter deletion}{ps_134-7_090-v0024}\var{\inserted[typed]{as well}}{as well}{typed insertion}{ps_134-7_090-v0025} the step of taking
derived categories (namely localizations of suitable model categories),
as the next namely introducing the corresponding structure of a derivator\hand{,}
make good sense also in non abelian set-ups,  namely for doing \var{\xout{homo}}{}{typewriter deletion}{ps_134-7_090-v0026}
``homotopy theory'',  thus viewed as a non-comm\hand{u}tative  version of homology
theory.

\margin{90}

% Visual collation: both supplied image halves inspected; see CHECK_DECISIONS_134_7_8.tsv.

%% END SOURCE ps_134-7_090
\ifnum\EditionMode=1
\psapparatusstart
\psentry{ps_134-7_090-v0001}{1}{typed insertion}{\inserted[typed]{(as pointed out by Quillen)}}{(as pointed out by Quillen)}{The typed insertion is incorporated at its indicated position; insertion brackets leave the running text.}
\psentry{ps_134-7_090-v0002}{2}{deletion}{\deleted{and}}{}{The author cancels this material by hand; it leaves the corrected running text.}
\psentry{ps_134-7_090-v0003}{3}{insertion}{\inserted{but (it seems)}}{\hand{but (it seems)}}{The handwritten insertion is incorporated at its indicated position; insertion brackets leave the running text.}
\psentry{ps_134-7_090-v0004}{4}{typewriter deletion}{\xout{a}}{}{The typewriter strikes out this material; it leaves the corrected running text.}
\psentry{ps_134-7_090-v0005}{5}{authorial spacing}{essentially\hand{/}``the same''\hand{/}homotopy}{essentially ``the same'' homotopy}{The two oblique pen strokes mark the crowded word/quotation boundaries. The diplomatic checked source keeps the strokes; the exact master variant supplies ordinary separation in reading mode.}
\psentry{ps_134-7_090-v0006}{6}{typewriter deletion}{\xout{\illegible{1}}}{}{The typewriter strikes out this material; it leaves the corrected running text.}
\psentry{ps_134-7_090-v0007}{7}{typewriter deletion}{\xout{s}}{}{The typewriter strikes out this material; it leaves the corrected running text.}
\psentry{ps_134-7_090-v0008}{8}{deletion}{\deleted{\xout{is}}}{}{The author cancels this material by hand; it leaves the corrected running text.}
\psentry{ps_134-7_090-v0009}{9}{typed insertion}{\inserted[typed]{are}}{are}{The typed insertion is incorporated at its indicated position; insertion brackets leave the running text.}
\psentry{ps_134-7_090-v0010}{10}{typed insertion}{\inserted[typed]{specific}}{specific}{The typed insertion is incorporated at its indicated position; insertion brackets leave the running text.}
\psentry{ps_134-7_090-v0011}{11}{typed insertion}{\inserted[typed]{on it}}{on it}{The typed insertion is incorporated at its indicated position; insertion brackets leave the running text.}
\psentry{ps_134-7_090-v0012}{12}{writing over}{\over{\unsure{ch}}{ch}}{\hand{ch}}{The superposed handwritten reading replaces the underlying reading.}
\psentry{ps_134-7_090-v0013}{13}{typed insertion}{\inserted[typed]{make any mention of models}}{make any mention of models}{The typed insertion is incorporated at its indicated position; insertion brackets leave the running text.}
\psentry{ps_134-7_090-v0014}{14}{deletion}{\deleted{\xout{\illegible{1}}}}{}{The author cancels this material by hand; it leaves the corrected running text.}
\psentry{ps_134-7_090-v0015}{15}{typewriter deletion}{\xout{the}}{}{The typewriter strikes out this material; it leaves the corrected running text.}
\psentry{ps_134-7_090-v0016}{16}{insertion}{\inserted{the}}{\hand{the}}{The handwritten insertion is incorporated at its indicated position; insertion brackets leave the running text.}
\psentry{ps_134-7_090-v0017}{17}{typed insertion}{\inserted[typed]{one}}{one}{The typed insertion is incorporated at its indicated position; insertion brackets leave the running text.}
\psentry{ps_134-7_090-v0018}{18}{typed insertion}{\inserted[typed]{the}}{the}{The typed insertion is incorporated at its indicated position; insertion brackets leave the running text.}
\psentry{ps_134-7_090-v0019}{19}{authorial spacing}{algebra,\hand{|}consisting}{algebra, consisting}{The pen vertical stroke separates the crowded comma and following word; retain it diplomatically and apply a space in the reading variant.}
\psentry{ps_134-7_090-v0020}{20}{typewriter deletion}{\xout{and}}{}{The typewriter strikes out this material; it leaves the corrected running text.}
\psentry{ps_134-7_090-v0021}{21}{typewriter deletion}{\xout{this}}{}{The typewriter strikes out this material; it leaves the corrected running text.}
\psentry{ps_134-7_090-v0022}{22}{typewriter deletion}{\xout{theorems}}{}{The typewriter strikes out this material; it leaves the corrected running text.}
\psentry{ps_134-7_090-v0023}{23}{typewriter deletion}{\xout{foundational}}{}{The typewriter strikes out this material; it leaves the corrected running text.}
\psentry{ps_134-7_090-v0024}{24}{typewriter deletion}{\xout{both}}{}{The typewriter strikes out this material; it leaves the corrected running text.}
\psentry{ps_134-7_090-v0025}{25}{typed insertion}{\inserted[typed]{as well}}{as well}{The typed insertion is incorporated at its indicated position; insertion brackets leave the running text.}
\psentry{ps_134-7_090-v0026}{26}{typewriter deletion}{\xout{homo}}{}{The typewriter strikes out this material; it leaves the corrected running text.}
\psapparatusend
\fi
\pssource{ps_134-7_091}{134-7/91}{https://grothendieck.umontpellier.fr/web/viewer.html?file=../134-7.pdf\#page=92}
\psnumber{ps_134-7_091-v0001}{1}
\psnumber{ps_134-7_091-v0002}{2}
\psnumber{ps_134-7_091-v0003}{3}
\psnumber{ps_134-7_091-v0004}{4}
\psnumber{ps_134-7_091-v0005}{5}
\psnumber{ps_134-7_091-v0006}{6}
\psnumber{ps_134-7_091-v0007}{7}
\psnumber{ps_134-7_091-v0008}{8}
\psnumber{ps_134-7_091-v0009}{9}
\psnumber{ps_134-7_091-v0010}{10}
\psnumber{ps_134-7_091-v0011}{11}
\psnumber{ps_134-7_091-v0012}{12}
\psnumber{ps_134-7_091-v0013}{13}
\psnumber{ps_134-7_091-v0014}{14}
\psnumber{ps_134-7_091-v0015}{15}
\psnumber{ps_134-7_091-v0016}{16}
\psnumber{ps_134-7_091-v0017}{17}
\psnumber{ps_134-7_091-v0018}{18}
%% BEGIN SOURCE ps_134-7_091
% Pursuing Stacks (cotes 134-2 to 134-8): archive cote 134-7, site page 91 (PDF page 92), typescript page 533.
% typescript. Draft compared: ps6.tex:4574-4645.
% AI-transcribed by Claude Opus 5.5 (claude-opus-5-5, adaptive thinking, effort high), Message Batches API, batch ps4 (msgbatch_012Aoc1NuskATPsyBeKSYK7J), 2026-11-12; reviewed against both supplied scan halves; see CHECK_DECISIONS.tsv.
\scanpage{134-7/91}
\var{\xout{126}}{}{typewriter deletion}{ps_134-7_091-v0001} \mspage{533}

126.      Let's come back though to schematic homotopy types.  Last thing
we looked at \var{\xout{were}}{}{typewriter deletion}{ps_134-7_091-v0002}\var{\inserted[typed]{was}}{was}{typed insertion}{ps_134-7_091-v0003} the four basic functors between "schematic homotopy
types"  and "linear schematic homotopy types" over a given ground ring k,
making up the two basic categories
\[ \mathrm{Ho\var{\xout{m}}{}{typewriter deletion}{ps_134-7_091-v0004}tab}_o(k) \quad \text{and} \quad \mathrm{Hot}_o(k) \quad . \]
\var{\inserted{We have been}}{\hand{We have been}}{insertion}{ps_134-7_091-v0005} dwelling somewhat on the remarkable forma\var{\over{\unsure{lly}}{l}}{\hand{l}}{writing over}{ps_134-7_091-v0006} symmetry to be expected
for these relations.  It \var{\xout{was}}{}{typewriter deletion}{ps_134-7_091-v0007} is tempting then to try and dualize any
kind of basic notion or construction which makes sense in terms of the
basic data,  namely the four functors and the adjunction and inversion
maps relating them.  Maybe the very first thing which forces attention
is the P\var{\over[typed]{o}{o}}{o}{typewriter overwriting}{ps_134-7_091-v0008}stnikov \var{\xout{fibration}}{}{typewriter deletion}{ps_134-7_091-v0009} dévissage of an object of $\mathrm{Hot}_o(k)$,  which
had been (together with abelianization) the very starting point for our
approach towards schematic homotopy types.   One main ingredient of this
dévissage is the "tower" of the Cartan-Serre type functors
\[ (1) \qquad X \hand{\longmapsto} X_n \]
and maps
\[ (2) \qquad X \hand{\longrightarrow} X_n \]
(with n a nat\var{\xout{y}}{}{typewriter deletion}{ps_134-7_091-v0010}ural integer),  where $X_n$ is deduced from X by "killing it's
$\hand{\pi}_i$'s  for  $i > n$" .  We may \var{\xout{say}}{}{typewriter deletion}{ps_134-7_091-v0011} call an object Y of $\mathrm{H\var{\over[typed]{\unsure{g}}{o}}{o}{typewriter overwriting}{ps_134-7_091-v0012}t}_o(k)$
\uline{n-}\uuline{co}\uline{-connected}  (a notion in a way \var{\xout{dual}}{}{typewriter deletion}{ps_134-7_091-v0013} symmetric to n-connectedness)
if  $\hand{\pi}_i(X) = 0$  for  $i > n$ ,  and denote by
\[ (3) \qquad \mathrm{Hot}_o(k)_{\hand{\leqslant} n} \]
the full subcategory of $\mathrm{Hot}_o(k)$ made up with these objects,  which is
therefore the inverse image by the functor
\[ L\hand{\pi}. : \mathrm{Hot}_o(k) \hand{\longrightarrow} \mathrm{H\var{\over[typed]{\unsure{g}}{o}}{o}{typewriter overwriting}{ps_134-7_091-v0014}tab}_o(k) \]
of the carresponding full subcategory
\[ (\var{\over{5}{4}}{\hand{4}}{writing over}{ps_134-7_091-v0015}) \qquad \mathrm{Hotab}_o(k)_{\hand{\leqslant} n} \]
of $\mathrm{Hotab}_o(k)$.  The most natural way for defining $X_n$ in terms of X. \hand{(}in
an axiomatic set up with the four functors as basic data,\hand{)}  together with
the canonical "fibration" (2),  is by \var{\xout{defining}}{}{typewriter deletion}{ps_134-7_091-v0016}\var{\inserted[typed]{describing}}{describing}{typed insertion}{ps_134-7_091-v0017} (2) as the "universal"
map (in $\mathrm{H\var{\over[typed]{\unsure{g}}{o}}{o}{typewriter overwriting}{ps_134-7_091-v0018}t}_o(k)$)  of X into an  n-co-connected object.  In other words,
we are surmising  that the inclusion functor
\[ \mathrm{Hot}_o(k)_{\hand{\leqslant} n} \hand{\hookrightarrow} \mathrm{Hot}_o(k) \]
admits a left adjoint,  and the latter is denoted by  $X \hand{\longmapsto} X_n$ .   This

\margin{91}

% Visual collation: both supplied image halves inspected; see CHECK_DECISIONS_134_7_8.tsv.

%% END SOURCE ps_134-7_091
\ifnum\EditionMode=1
\psapparatusstart
\psentry{ps_134-7_091-v0001}{1}{typewriter deletion}{\xout{126}}{}{The typewriter strikes out this material; it leaves the corrected running text.}
\psentry{ps_134-7_091-v0002}{2}{typewriter deletion}{\xout{were}}{}{The typewriter strikes out this material; it leaves the corrected running text.}
\psentry{ps_134-7_091-v0003}{3}{typed insertion}{\inserted[typed]{was}}{was}{The typed insertion is incorporated at its indicated position; insertion brackets leave the running text.}
\psentry{ps_134-7_091-v0004}{4}{typewriter deletion}{$\xout{m}$}{}{The typewriter strikes out this material; it leaves the corrected running text.}
\psentry{ps_134-7_091-v0005}{5}{insertion}{\inserted{We have been}}{\hand{We have been}}{The handwritten insertion is incorporated at its indicated position; insertion brackets leave the running text.}
\psentry{ps_134-7_091-v0006}{6}{writing over}{\over{\unsure{lly}}{l}}{\hand{l}}{The superposed handwritten reading replaces the underlying reading.}
\psentry{ps_134-7_091-v0007}{7}{typewriter deletion}{\xout{was}}{}{The typewriter strikes out this material; it leaves the corrected running text.}
\psentry{ps_134-7_091-v0008}{8}{typewriter overwriting}{\over[typed]{o}{o}}{o}{The superposed typed reading replaces the underlying reading.}
\psentry{ps_134-7_091-v0009}{9}{typewriter deletion}{\xout{fibration}}{}{The typewriter strikes out this material; it leaves the corrected running text.}
\psentry{ps_134-7_091-v0010}{10}{typewriter deletion}{\xout{y}}{}{The typewriter strikes out this material; it leaves the corrected running text.}
\psentry{ps_134-7_091-v0011}{11}{typewriter deletion}{\xout{say}}{}{The typewriter strikes out this material; it leaves the corrected running text.}
\psentry{ps_134-7_091-v0012}{12}{typewriter overwriting}{$\over[typed]{\unsure{g}}{o}$}{$o$}{The superposed typed reading replaces the underlying reading.}
\psentry{ps_134-7_091-v0013}{13}{typewriter deletion}{\xout{dual}}{}{The typewriter strikes out this material; it leaves the corrected running text.}
\psentry{ps_134-7_091-v0014}{14}{typewriter overwriting}{$\over[typed]{\unsure{g}}{o}$}{$o$}{The superposed typed reading replaces the underlying reading.}
\psentry{ps_134-7_091-v0015}{15}{writing over}{$\over{5}{4}$}{$\hand{4}$}{The superposed handwritten reading replaces the underlying reading.}
\psentry{ps_134-7_091-v0016}{16}{typewriter deletion}{\xout{defining}}{}{The typewriter strikes out this material; it leaves the corrected running text.}
\psentry{ps_134-7_091-v0017}{17}{typed insertion}{\inserted[typed]{describing}}{describing}{The typed insertion is incorporated at its indicated position; insertion brackets leave the running text.}
\psentry{ps_134-7_091-v0018}{18}{typewriter overwriting}{$\over[typed]{\unsure{g}}{o}$}{$o$}{The superposed typed reading replaces the underlying reading.}
\psapparatusend
\fi
\pssource{ps_134-7_092}{134-7/92}{https://grothendieck.umontpellier.fr/web/viewer.html?file=../134-7.pdf\#page=93}
\psnumber{ps_134-7_092-v0001}{1}
\psnumber{ps_134-7_092-v0002}{2}
\psnumber{ps_134-7_092-v0003}{3}
\psnumber{ps_134-7_092-v0004}{4}
\psnumber{ps_134-7_092-v0005}{5}
\psnumber{ps_134-7_092-v0006}{6}
\psnumber{ps_134-7_092-v0007}{7}
\psnumber{ps_134-7_092-v0008}{8}
\psnumber{ps_134-7_092-v0009}{9}
\psnumber{ps_134-7_092-v0010}{10}
\psnumber{ps_134-7_092-v0011}{11}
\psnumber{ps_134-7_092-v0012}{12}
\psnumber{ps_134-7_092-v0013}{13}
\psnumber{ps_134-7_092-v0014}{14}
\psnumber{ps_134-7_092-v0015}{15}
\psnumber{ps_134-7_092-v0016}{16}
%% BEGIN SOURCE ps_134-7_092
% Pursuing Stacks (cotes 134-2 to 134-8): archive cote 134-7, site page 92 (PDF page 93), typescript page 534.
% typescript. Draft compared: ps6.tex:4629-4690.
% AI-transcribed by Claude Opus 5.5 (claude-opus-5-5, adaptive thinking, effort high), Message Batches API, batch ps4 (msgbatch_012Aoc1NuskATPsyBeKSYK7J), 2026-11-12; reviewed against both supplied scan halves; see CHECK_DECISIONS.tsv.
\scanpage{134-7/92}
\mspage{534}

description gives rise at once to the familiar ``tower'' structure for
variable n
\[ (\var{\over{\unsure{4}}{5}}{\hand{5}}{writing over}{ps_134-7_092-v0001}) \qquad X_{n+1} \hand{\to} X_n \hand{\to} \dots \hand{\to} X_1 \hand{\to} X_o \quad ( \unsure{=} e ) \quad . \]
Intuitively, the maps in (\var{\over{\unsure{4}}{5}}{\hand{5}}{writing over}{ps_134-7_092-v0002}) are viewed as being (surjective) fibrations
between (connected) ``spaces'', X being viewed as a kind of inverse limit
of the $X_n$'s . Dually, we would expect to get a sequence of inclusions
\[ \hand{(*)} \qquad {}_oX \hand{\hookrightarrow} {}_1X \hand{\hookrightarrow} \dots \hand{\hookrightarrow} {}_nX \hand{\hookrightarrow} {}_{(n+1)}X \hand{\hookrightarrow} \cdots , \]
with X appearing as \var{\xout{the}}{}{typewriter deletion}{ps_134-7_092-v0003} a kind of direct limit. In the set up of ordina-%
ry homotopy types, modelized by semisimplicial sets, one \var{\deleted{is}}{}{deletion}{ps_134-7_092-v0004}\var{\inserted{will}}{\hand{will}}{insertion}{ps_134-7_092-v0005} think\var{\deleted{ing}}{}{deletion}{ps_134-7_092-v0006}
at once of the filtration by skeleta - which however is'nt quite the
right thing surely, because if we dualize the familiar Cartan-Serre
requirement on (2) (namely that it induces an isomorphism for $\hand{\pi}_i$ for
$i \hand{\leqslant} n$), we see that the ``inclusion''
\[ {}_nX \hand{\hookrightarrow} X \]
should induce an isomorphism\var{\xout{s}}{}{typewriter deletion}{ps_134-7_092-v0007} for $H_i$ for $i \hand{\leqslant} n$ , which is'nt quite
true for the skeletal filtration (it is OK for $i \hand{\leqslant} n-1$ only); the condi-%
tion
\[ H_i(X_n) = 0 \quad \text{for} \quad i > n \]
is OK though for this filtration. So the next idea would be to modify
a little the n-skeleton ${}_nX$ , to straighten things out .\var{\xout{\illegible{1}the analogy}}{}{typewriter deletion}{ps_134-7_092-v0008}
I played around some along these lines, and after some initial optimism
came to the feeling that there does not exist (in the discrete nor in the
schematic set-up) such an increasing \var{\inserted[typed]{canonical}}{canonical}{typed insertion}{ps_134-7_092-v0009} filtration of a homotopy type. I
did'nt make any formal statement and proof for this (in the set-up of
ordinary homotopy types say), however in the process of playing around
it became soon clear that in various ways, \uline{there is} \var{\deleted{an}}{}{deletion}{ps_134-7_092-v0010}\var{\inserted{\uline{some}}}{\hand{\uline{some}}}{insertion}{ps_134-7_092-v0011} \uline{essential dissy-}%
\uline{metry} between the seemingly ``dual'' situations, when trying to get the
two types of ``filtrations'' of the object X. O\var{\over{\unsure{n}}{N}}{\hand{N}}{writing over}{ps_134-7_092-v0012}e dissymetry occurs alrea-%
dy in the very definition of the subcategory (3), and of the correspon-%
ding ``dual'' \var{\xout{definition}}{}{typewriter deletion}{ps_134-7_092-v0013} subcategory
\[ (\var{\over{5}{6}}{\hand{6}}{writing over}{ps_134-7_092-v0014}) \qquad {}_{(\hand{\leqslant} n)}Hot_o(k) \]
of objects satisfying
\[ H_i(Y) = 0 \quad \text{for} \quad i > n \ . \]
Namely, both subcategories (3) and (\var{\over{5}{6}}{\hand{6}}{writing over}{ps_134-7_092-v0015}) are defined in terms of the \uline{same}
subcategory of $Hotab_o(k)$, namely (\var{\over{\unsure{3}}{4}}{\hand{4}}{writing over}{ps_134-7_092-v0016}), as the inverse image of the

\margin{92}

% Visual collation: both supplied image halves inspected; see CHECK_DECISIONS_134_7_8.tsv.

%% END SOURCE ps_134-7_092
\ifnum\EditionMode=1
\psapparatusstart
\psentry{ps_134-7_092-v0001}{1}{writing over}{$\over{\unsure{4}}{5}$}{$\hand{5}$}{The superposed handwritten reading replaces the underlying reading.}
\psentry{ps_134-7_092-v0002}{2}{writing over}{\over{\unsure{4}}{5}}{\hand{5}}{The superposed handwritten reading replaces the underlying reading.}
\psentry{ps_134-7_092-v0003}{3}{typewriter deletion}{\xout{the}}{}{The typewriter strikes out this material; it leaves the corrected running text.}
\psentry{ps_134-7_092-v0004}{4}{deletion}{\deleted{is}}{}{The author cancels this material by hand; it leaves the corrected running text.}
\psentry{ps_134-7_092-v0005}{5}{insertion}{\inserted{will}}{\hand{will}}{The handwritten insertion is incorporated at its indicated position; insertion brackets leave the running text.}
\psentry{ps_134-7_092-v0006}{6}{deletion}{\deleted{ing}}{}{The author cancels this material by hand; it leaves the corrected running text.}
\psentry{ps_134-7_092-v0007}{7}{typewriter deletion}{\xout{s}}{}{The typewriter strikes out this material; it leaves the corrected running text.}
\psentry{ps_134-7_092-v0008}{8}{typewriter deletion}{\xout{\illegible{1}the analogy}}{}{The typewriter strikes out this material; it leaves the corrected running text.}
\psentry{ps_134-7_092-v0009}{9}{typed insertion}{\inserted[typed]{canonical}}{canonical}{The typed insertion is incorporated at its indicated position; insertion brackets leave the running text.}
\psentry{ps_134-7_092-v0010}{10}{deletion}{\deleted{an}}{}{The author cancels this material by hand; it leaves the corrected running text.}
\psentry{ps_134-7_092-v0011}{11}{insertion}{\inserted{\uline{some}}}{\hand{\uline{some}}}{The handwritten insertion is incorporated at its indicated position; insertion brackets leave the running text.}
\psentry{ps_134-7_092-v0012}{12}{writing over}{\over{\unsure{n}}{N}}{\hand{N}}{The superposed handwritten reading replaces the underlying reading.}
\psentry{ps_134-7_092-v0013}{13}{typewriter deletion}{\xout{definition}}{}{The typewriter strikes out this material; it leaves the corrected running text.}
\psentry{ps_134-7_092-v0014}{14}{writing over}{$\over{5}{6}$}{$\hand{6}$}{The superposed handwritten reading replaces the underlying reading.}
\psentry{ps_134-7_092-v0015}{15}{writing over}{\over{5}{6}}{\hand{6}}{The superposed handwritten reading replaces the underlying reading.}
\psentry{ps_134-7_092-v0016}{16}{writing over}{\over{\unsure{3}}{4}}{\hand{4}}{The superposed handwritten reading replaces the underlying reading.}
\psapparatusend
\fi
\pssource{ps_134-7_093}{134-7/93}{https://grothendieck.umontpellier.fr/web/viewer.html?file=../134-7.pdf\#page=94}
\psnumber{ps_134-7_093-v0001}{1}
\psnumber{ps_134-7_093-v0002}{2}
\psnumber{ps_134-7_093-v0003}{3}
\psnumber{ps_134-7_093-v0004}{4}
\psnumber{ps_134-7_093-v0005}{5}
\psnumber{ps_134-7_093-v0006}{6}
\psnumber{ps_134-7_093-v0007}{7}
\psnumber{ps_134-7_093-v0008}{8}
\psnumber{ps_134-7_093-v0009}{9}
\psnumber{ps_134-7_093-v0010}{10}
\psnumber{ps_134-7_093-v0011}{11}
\psnumber{ps_134-7_093-v0012}{12}
\psnumber{ps_134-7_093-v0013}{13}
\psnumber{ps_134-7_093-v0014}{14}
\psnumber{ps_134-7_093-v0015}{15}
\psnumber{ps_134-7_093-v0016}{16}
\psnumber{ps_134-7_093-v0017}{17}
\psnumber{ps_134-7_093-v0018}{18}
\psnumber{ps_134-7_093-v0019}{19}
\psnumber{ps_134-7_093-v0020}{20}
\psnumber{ps_134-7_093-v0021}{21}
\psnumber{ps_134-7_093-v0022}{22}
\psnumber{ps_134-7_093-v0023}{23}
\psnumber{ps_134-7_093-v0024}{24}
\psnumber{ps_134-7_093-v0025}{25}
%% BEGIN SOURCE ps_134-7_093
% Pursuing Stacks (cotes 134-2 to 134-8): archive cote 134-7, site page 93 (PDF page 94), typescript page 535.
% typescript. Draft compared: ps6.tex:4674-4736.
% AI-transcribed by Claude Opus 5.5 (claude-opus-5-5, adaptive thinking, effort high), Message Batches API, batch ps4 (msgbatch_012Aoc1NuskATPsyBeKSYK7J), 2026-11-12; reviewed against both supplied scan halves; see CHECK_DECISIONS.tsv.
\scanpage{134-7/93}
\mspage{535}

latter by either L\hand{$\pi$}. , or LH. . Now the point is that the properties
of this subcategory, with respect to the inner structure of Hotab$_o$(k)
(involving the``left'' and the ``right'' homotopy operations, notably the
suspension and the loop functors\var{\over{x}{\unsure{,}}}{\hand{\unsure{,}}}{writing over}{ps_134-7_093-v0001} respectively) are by no means autosym-%
metric. The main dissymetry, it would seem, turns out in this, that
(\var{\over{5}{4}}{\hand{4}}{writing over}{ps_134-7_093-v0002}) is stable under the loop functor, and by no means under the suspen-%
sion functor. This \var{\xout{implies that}}{}{typewriter deletion}{ps_134-7_093-v0003}\var{\inserted[typed]{is the reason why,}}{is the reason why,}{typed insertion}{ps_134-7_093-v0004} even if we assume that (in analogy
to what happens for the categories (3)) the inclusion functors from the
categories (5) into Hot$_o$(k) do admit the relevant (namely right) adjoint
(which I greatly doubt anyhow ...), and thus give rise for any object X
to an increasing filtration (*) , the Postnikov-type relationships\var{\xout{betw}}{}{typewriter deletion}{ps_134-7_093-v0005}
between X$_n$ and X$_{n+1}$ cannot be quite dualized to a similar relationship
between $_n$X and $_{(n+1)}$X . We may think \var{\inserted{next,}}{\hand{next,}}{insertion}{ps_134-7_093-v0006} of course\hand{,} of defining an increa-%
sing sequence of subcategories \var{\xout{of}}{}{typewriter deletion}{ps_134-7_093-v0007} (5) of H\var{\over[typed]{o}{o}}{o}{typewriter overwriting}{ps_134-7_093-v0008}t$_o$(k) in terms of LH. and a
correspondance sequence of subcategories of Hotab$_o$(k), different from
the categories (\var{\over{5}{4}}{\hand{4}}{writing over}{ps_134-7_093-v0009}), and stable under suspensions. \var{\xout{Here the natural}}{}{typewriter deletion}{ps_134-7_093-v0010}
\var{\xout{choice would be to}}{}{typewriter deletion}{ps_134-7_093-v0011} But there \unsure{doesn't|does'nt} seem to be anything reasonable \var{\xout{aro}}{}{typewriter deletion}{ps_134-7_093-v0012}
around along these lines.

To sum up, it does'nt seem one should overemphasize the somewhat
startling symmetry which appeared in section 124 between ``homotopy''
\var{\xout{invariants}}{}{typewriter deletion}{ps_134-7_093-v0013} embodied in L\hand{$\pi$}. , and``homology''\var{\xout{invariants}}{}{typewriter deletion}{ps_134-7_093-v0014} embodied in LH. -
in some essential respects, it would seem that the corresponding two
functors do have non mutually symmetric properties. I guess I have to
apologize for having taken that long for coming to a conclusion which,
presumably, must be felt as a kind of self-evidence \var{\over{\illegible{1}}{by}}{\hand{by}}{writing over}{ps_134-7_093-v0015} \var{\inserted[typed]{all}}{all}{typed insertion}{ps_134-7_093-v0016} homotopy
people ! \margin{$\lfloor$\\$\ulcorner$}

127. Before leaving (for the time being) the topic of schematic
homotopy types and schematization, I would like still to add a few com-%
\var{\xout{ments}}{}{typewriter deletion}{ps_134-7_093-v0017} ments, about various possibilities for working with different kind
of models for defining schematic homotopy types. My point here is not
in \var{\xout{changing}}{}{typewriter deletion}{ps_134-7_093-v0018}\var{\inserted[typed]{replacing}}{replacing}{typed insertion}{ps_134-7_093-v0019} the \var{\inserted[typed]{basic}}{basic}{typed insertion}{ps_134-7_093-v0020} test-category we are working with, here \hand{$\Delta$} , by some
other (say the category \hand{\unsure{$\bigcirc$}} of \var{\inserted[typed]{standard}}{standard}{typed insertion}{ps_134-7_093-v0021} hemispheres) - \var{\xout{we'll keep}}{}{typewriter deletion}{ps_134-7_093-v0022} this choice I
feel should be more or less irrelevant, so we may as well keep \hand{$\Delta$} .
Thus we are going to work with semisimplicial objects, and the \var{\deleted{whole}}{}{deletion}{ps_134-7_093-v0023} \var{\inserted{main}}{\hand{main}}{insertion}{ps_134-7_093-v0024}
question then is (for a given ground ring k) to say \var{\xout{wh}}{}{typewriter deletion}{ps_134-7_093-v0025} precisely what
kind of objects we are allowing (or imposing !) as components for our
complexes. At any rate they should be ``objects over k'', and the most

\margin{93}

% Visual collation: both supplied image halves inspected; see CHECK_DECISIONS_134_7_8.tsv.

%% END SOURCE ps_134-7_093
\ifnum\EditionMode=1
\psapparatusstart
\psentry{ps_134-7_093-v0001}{1}{writing over}{\over{x}{\unsure{,}}}{\hand{\unsure{,}}}{The superposed handwritten reading replaces the underlying reading.}
\psentry{ps_134-7_093-v0002}{2}{writing over}{\over{5}{4}}{\hand{4}}{The superposed handwritten reading replaces the underlying reading.}
\psentry{ps_134-7_093-v0003}{3}{typewriter deletion}{\xout{implies that}}{}{The typewriter strikes out this material; it leaves the corrected running text.}
\psentry{ps_134-7_093-v0004}{4}{typed insertion}{\inserted[typed]{is the reason why,}}{is the reason why,}{The typed insertion is incorporated at its indicated position; insertion brackets leave the running text.}
\psentry{ps_134-7_093-v0005}{5}{typewriter deletion}{\xout{betw}}{}{The typewriter strikes out this material; it leaves the corrected running text.}
\psentry{ps_134-7_093-v0006}{6}{insertion}{\inserted{next,}}{\hand{next,}}{The handwritten insertion is incorporated at its indicated position; insertion brackets leave the running text.}
\psentry{ps_134-7_093-v0007}{7}{typewriter deletion}{\xout{of}}{}{The typewriter strikes out this material; it leaves the corrected running text.}
\psentry{ps_134-7_093-v0008}{8}{typewriter overwriting}{\over[typed]{o}{o}}{o}{The superposed typed reading replaces the underlying reading.}
\psentry{ps_134-7_093-v0009}{9}{writing over}{\over{5}{4}}{\hand{4}}{The superposed handwritten reading replaces the underlying reading.}
\psentry{ps_134-7_093-v0010}{10}{typewriter deletion}{\xout{Here the natural}}{}{The typewriter strikes out this material; it leaves the corrected running text.}
\psentry{ps_134-7_093-v0011}{11}{typewriter deletion}{\xout{choice would be to}}{}{The typewriter strikes out this material; it leaves the corrected running text.}
\psentry{ps_134-7_093-v0012}{12}{typewriter deletion}{\xout{aro}}{}{The typewriter strikes out this material; it leaves the corrected running text.}
\psentry{ps_134-7_093-v0013}{13}{typewriter deletion}{\xout{invariants}}{}{The typewriter strikes out this material; it leaves the corrected running text.}
\psentry{ps_134-7_093-v0014}{14}{typewriter deletion}{\xout{invariants}}{}{The typewriter strikes out this material; it leaves the corrected running text.}
\psentry{ps_134-7_093-v0015}{15}{writing over}{\over{\illegible{1}}{by}}{\hand{by}}{The superposed handwritten reading replaces the underlying reading.}
\psentry{ps_134-7_093-v0016}{16}{typed insertion}{\inserted[typed]{all}}{all}{The typed insertion is incorporated at its indicated position; insertion brackets leave the running text.}
\psentry{ps_134-7_093-v0017}{17}{typewriter deletion}{\xout{ments}}{}{The typewriter strikes out this material; it leaves the corrected running text.}
\psentry{ps_134-7_093-v0018}{18}{typewriter deletion}{\xout{changing}}{}{The typewriter strikes out this material; it leaves the corrected running text.}
\psentry{ps_134-7_093-v0019}{19}{typed insertion}{\inserted[typed]{replacing}}{replacing}{The typed insertion is incorporated at its indicated position; insertion brackets leave the running text.}
\psentry{ps_134-7_093-v0020}{20}{typed insertion}{\inserted[typed]{basic}}{basic}{The typed insertion is incorporated at its indicated position; insertion brackets leave the running text.}
\psentry{ps_134-7_093-v0021}{21}{typed insertion}{\inserted[typed]{standard}}{standard}{The typed insertion is incorporated at its indicated position; insertion brackets leave the running text.}
\psentry{ps_134-7_093-v0022}{22}{typewriter deletion}{\xout{we'll keep}}{}{The typewriter strikes out this material; it leaves the corrected running text.}
\psentry{ps_134-7_093-v0023}{23}{deletion}{\deleted{whole}}{}{The author cancels this material by hand; it leaves the corrected running text.}
\psentry{ps_134-7_093-v0024}{24}{insertion}{\inserted{main}}{\hand{main}}{The handwritten insertion is incorporated at its indicated position; insertion brackets leave the running text.}
\psentry{ps_134-7_093-v0025}{25}{typewriter deletion}{\xout{wh}}{}{The typewriter strikes out this material; it leaves the corrected running text.}
\psapparatusend
\fi
\pssource{ps_134-7_094}{134-7/94}{https://grothendieck.umontpellier.fr/web/viewer.html?file=../134-7.pdf\#page=95}
\psnumber{ps_134-7_094-v0001}{1}
\psnumber{ps_134-7_094-v0002}{2}
\psnumber{ps_134-7_094-v0003}{3}
\psnumber{ps_134-7_094-v0004}{4}
\psnumber{ps_134-7_094-v0005}{5}
\psnumber{ps_134-7_094-v0006}{6}
\psnumber{ps_134-7_094-v0007}{7}
\psnumber{ps_134-7_094-v0008}{8}
\psnumber{ps_134-7_094-v0009}{9}
\psnumber{ps_134-7_094-v0010}{10}
\psnumber{ps_134-7_094-v0011}{11}
\psnumber{ps_134-7_094-v0012}{12}
\psnumber{ps_134-7_094-v0013}{13}
\psnumber{ps_134-7_094-v0014}{14}
\psnumber{ps_134-7_094-v0015}{15}
\psnumber{ps_134-7_094-v0016}{16}
\psnumber{ps_134-7_094-v0017}{17}
\psnumber{ps_134-7_094-v0018}{18}
\psnumber{ps_134-7_094-v0019}{19}
\psnumber{ps_134-7_094-v0020}{20}
\psnumber{ps_134-7_094-v0021}{21}
%% BEGIN SOURCE ps_134-7_094
% Pursuing Stacks (cotes 134-2 to 134-8): archive cote 134-7, site page 94 (PDF page 95), typescript page 536.
% typescript. Draft compared: ps6.tex:4724-4796.
% AI-transcribed by Claude Opus 5.5 (claude-opus-5-5, adaptive thinking, effort high), Message Batches API, batch ps4 (msgbatch_012Aoc1NuskATPsyBeKSYK7J), 2026-11-12; reviewed against both supplied scan halves; see CHECK_DECISIONS.tsv.
\scanpage{134-7/94}
\mspage{536}

encompassing choice for such objects seems to be sheaves on the category
of all \var{\xout{pre}}{}{typewriter deletion}{ps_134-7_094-v0001}schemes  over k (or equivalently,  of all affine schemes over
k),  for a suitable topology such as the fpqc topology  (compare section
111,  pages 446-447\var{\over{\unsure{r}}{)}}{\hand{)}}{writing over}{ps_134-7_094-v0002}.  Apart from the technico-logical nuisance of this
not being a $\underline{U}$-site,  where $\underline{U}$ is our basic universe  (cf. p. 492\var{\over{\unsure{r}}{)}}{\hand{)}}{writing over}{ps_134-7_094-v0003},
which we'll ignore here (as it is'nt really too serious a difficulty),
\var{\inserted{when}}{\hand{when}}{insertion}{ps_134-7_094-v0004} working with semisimplicial sheaves,  the embryo of foundational work
of Illusie's on the derived category of a topos via semisimplicial
sheaves\unsure{x}\hand{[} \var{\xout{\unsure{x}}}{}{typewriter deletion}{ps_134-7_094-v0005} chapter 1\hand{]} becomes available.  Thus for a map of semi-\margin{x}
simplicial sheaves over k ,
\[
\text{(1)}\qquad X_{*} \hand{\longrightarrow} Y_{*} \quad ,
\]
we know already what it means to be a quasi-isomorphism\var{\inserted[typed]{(= weak equivalence)}}{(= weak equivalence)}{typed insertion}{ps_134-7_094-v0006}. It seems like-%
ly that,  even when drastically restricting the sheaves allowed as com-%
ponents for our semisimplicial models,  the notion of quasi-isomorphism
relevant for these models should be the same as Illusie's\var{\xout{\unsure{,}}}{}{typewriter deletion}{ps_134-7_094-v0007}  for the
encompassing topos.  This is one point we have been neglecting so far.
I am not going here to recall Illusie's definition (in terms of the
homotopy sheaves  $\underline{\hand{\pi}}_{i}(X_{*}\ , s)$  where s is a section of $X_{*}$ ),  or trans-
late it into cohomological terms.   One point I want to make here,  is
that it is by no means automatic  that \var{\xout{the}}{}{typewriter deletion}{ps_134-7_094-v0008} if (1) is a \var{\xout{quasi-isomorphism}}{}{typewriter deletion}{ps_134-7_094-v0009}
weak equivalence,  the same holds for the corresponding map of semisim-%
plicial sets
\[
\text{(2)}\qquad X_{*}(k) \hand{\longrightarrow} Y_{*}(k) \quad .
\]
Thus \var{\xout{t}}{}{typewriter deletion}{ps_134-7_094-v0010}there is no \var{\xout{\unsure{tau}}}{}{typewriter deletion}{ps_134-7_094-v0011} more or less tautological ``sections'' functor from
Illusie's derived category\var{\xout{\unsure{,}}}{}{typewriter deletion}{ps_134-7_094-v0012} to the category (Hot) of usual homotopy
types. \hand{(}There surely \uline{is} \var{\xout{such}}{}{typewriter deletion}{ps_134-7_094-v0013} a \var{\inserted[typed]{canonical}}{canonical}{typed insertion}{ps_134-7_094-v0014} functor though,  in the case of a general
topos \var{\inserted[typed]{T}}{T}{typed insertion}{ps_134-7_094-v0015} and Illusie\hand{'}s corresponding d\var{\xout{\unsure{r}}}{}{typewriter deletion}{ps_134-7_094-v0016}erived category,  namely what we would
like to call the ``\uline{cointegration over} T ''  \uline{functor},  which should come
out of the formalism of stacks we hav'nt \var{\xout{yet}}{}{typewriter deletion}{ps_134-7_094-v0017} begun to  develop yet.\hand{)}
Presumably,\var{\inserted{however}}{\hand{however}}{insertion}{ps_134-7_094-v0018}  when working with \var{\xout{models}}{}{typewriter deletion}{ps_134-7_094-v0019} semisimplicial models whose compo-%
nents are restricted  to be unipotent bundles or something pretty close
to these (see below for examples),  whenever (1) is a weak equivalence,
the same \var{\inserted{will}}{\hand{will}}{insertion}{ps_134-7_094-v0020} hold\var{\deleted{s}}{}{deletion}{ps_134-7_094-v0021} for (2) .  As the choice of objects (let's call them
simply the ``bundles'') we are allowing,  should clearly be stable under
ring extension,  it will then follow that more generally,  for any
algebra k' over k,  the corresponding map
\[
\text{(2')}\qquad X_{*}(k') \hand{\longrightarrow} Y_{*}(k')
\]

\margin{94}

% Visual collation: both supplied image halves inspected; see CHECK_DECISIONS_134_7_8.tsv.

%% END SOURCE ps_134-7_094
\ifnum\EditionMode=1
\psapparatusstart
\psentry{ps_134-7_094-v0001}{1}{typewriter deletion}{\xout{pre}}{}{The typewriter strikes out this material; it leaves the corrected running text.}
\psentry{ps_134-7_094-v0002}{2}{writing over}{\over{\unsure{r}}{)}}{\hand{)}}{The superposed handwritten reading replaces the underlying reading.}
\psentry{ps_134-7_094-v0003}{3}{writing over}{\over{\unsure{r}}{)}}{\hand{)}}{The superposed handwritten reading replaces the underlying reading.}
\psentry{ps_134-7_094-v0004}{4}{insertion}{\inserted{when}}{\hand{when}}{The handwritten insertion is incorporated at its indicated position; insertion brackets leave the running text.}
\psentry{ps_134-7_094-v0005}{5}{typewriter deletion}{\xout{\unsure{x}}}{}{The typewriter strikes out this material; it leaves the corrected running text.}
\psentry{ps_134-7_094-v0006}{6}{typed insertion}{\inserted[typed]{(= weak equivalence)}}{(= weak equivalence)}{The typed insertion is incorporated at its indicated position; insertion brackets leave the running text.}
\psentry{ps_134-7_094-v0007}{7}{typewriter deletion}{\xout{\unsure{,}}}{}{The typewriter strikes out this material; it leaves the corrected running text.}
\psentry{ps_134-7_094-v0008}{8}{typewriter deletion}{\xout{the}}{}{The typewriter strikes out this material; it leaves the corrected running text.}
\psentry{ps_134-7_094-v0009}{9}{typewriter deletion}{\xout{quasi-isomorphism}}{}{The typewriter strikes out this material; it leaves the corrected running text.}
\psentry{ps_134-7_094-v0010}{10}{typewriter deletion}{\xout{t}}{}{The typewriter strikes out this material; it leaves the corrected running text.}
\psentry{ps_134-7_094-v0011}{11}{typewriter deletion}{\xout{\unsure{tau}}}{}{The typewriter strikes out this material; it leaves the corrected running text.}
\psentry{ps_134-7_094-v0012}{12}{typewriter deletion}{\xout{\unsure{,}}}{}{The typewriter strikes out this material; it leaves the corrected running text.}
\psentry{ps_134-7_094-v0013}{13}{typewriter deletion}{\xout{such}}{}{The typewriter strikes out this material; it leaves the corrected running text.}
\psentry{ps_134-7_094-v0014}{14}{typed insertion}{\inserted[typed]{canonical}}{canonical}{The typed insertion is incorporated at its indicated position; insertion brackets leave the running text.}
\psentry{ps_134-7_094-v0015}{15}{typed insertion}{\inserted[typed]{T}}{T}{The typed insertion is incorporated at its indicated position; insertion brackets leave the running text.}
\psentry{ps_134-7_094-v0016}{16}{typewriter deletion}{\xout{\unsure{r}}}{}{The typewriter strikes out this material; it leaves the corrected running text.}
\psentry{ps_134-7_094-v0017}{17}{typewriter deletion}{\xout{yet}}{}{The typewriter strikes out this material; it leaves the corrected running text.}
\psentry{ps_134-7_094-v0018}{18}{insertion}{\inserted{however}}{\hand{however}}{The handwritten insertion is incorporated at its indicated position; insertion brackets leave the running text.}
\psentry{ps_134-7_094-v0019}{19}{typewriter deletion}{\xout{models}}{}{The typewriter strikes out this material; it leaves the corrected running text.}
\psentry{ps_134-7_094-v0020}{20}{insertion}{\inserted{will}}{\hand{will}}{The handwritten insertion is incorporated at its indicated position; insertion brackets leave the running text.}
\psentry{ps_134-7_094-v0021}{21}{deletion}{\deleted{s}}{}{The author cancels this material by hand; it leaves the corrected running text.}
\psapparatusend
\fi

\end{document}
